% Options for packages loaded elsewhere
\PassOptionsToPackage{unicode}{hyperref}
\PassOptionsToPackage{hyphens}{url}
%
\documentclass[
]{book}
\usepackage{amsmath,amssymb}
\usepackage{iftex}
\ifPDFTeX
  \usepackage[T1]{fontenc}
  \usepackage[utf8]{inputenc}
  \usepackage{textcomp} % provide euro and other symbols
\else % if luatex or xetex
  \usepackage{unicode-math} % this also loads fontspec
  \defaultfontfeatures{Scale=MatchLowercase}
  \defaultfontfeatures[\rmfamily]{Ligatures=TeX,Scale=1}
\fi
\usepackage{lmodern}
\ifPDFTeX\else
  % xetex/luatex font selection
\fi
% Use upquote if available, for straight quotes in verbatim environments
\IfFileExists{upquote.sty}{\usepackage{upquote}}{}
\IfFileExists{microtype.sty}{% use microtype if available
  \usepackage[]{microtype}
  \UseMicrotypeSet[protrusion]{basicmath} % disable protrusion for tt fonts
}{}
\makeatletter
\@ifundefined{KOMAClassName}{% if non-KOMA class
  \IfFileExists{parskip.sty}{%
    \usepackage{parskip}
  }{% else
    \setlength{\parindent}{0pt}
    \setlength{\parskip}{6pt plus 2pt minus 1pt}}
}{% if KOMA class
  \KOMAoptions{parskip=half}}
\makeatother
\usepackage{xcolor}
\usepackage{color}
\usepackage{fancyvrb}
\newcommand{\VerbBar}{|}
\newcommand{\VERB}{\Verb[commandchars=\\\{\}]}
\DefineVerbatimEnvironment{Highlighting}{Verbatim}{commandchars=\\\{\}}
% Add ',fontsize=\small' for more characters per line
\usepackage{framed}
\definecolor{shadecolor}{RGB}{248,248,248}
\newenvironment{Shaded}{\begin{snugshade}}{\end{snugshade}}
\newcommand{\AlertTok}[1]{\textcolor[rgb]{0.94,0.16,0.16}{#1}}
\newcommand{\AnnotationTok}[1]{\textcolor[rgb]{0.56,0.35,0.01}{\textbf{\textit{#1}}}}
\newcommand{\AttributeTok}[1]{\textcolor[rgb]{0.13,0.29,0.53}{#1}}
\newcommand{\BaseNTok}[1]{\textcolor[rgb]{0.00,0.00,0.81}{#1}}
\newcommand{\BuiltInTok}[1]{#1}
\newcommand{\CharTok}[1]{\textcolor[rgb]{0.31,0.60,0.02}{#1}}
\newcommand{\CommentTok}[1]{\textcolor[rgb]{0.56,0.35,0.01}{\textit{#1}}}
\newcommand{\CommentVarTok}[1]{\textcolor[rgb]{0.56,0.35,0.01}{\textbf{\textit{#1}}}}
\newcommand{\ConstantTok}[1]{\textcolor[rgb]{0.56,0.35,0.01}{#1}}
\newcommand{\ControlFlowTok}[1]{\textcolor[rgb]{0.13,0.29,0.53}{\textbf{#1}}}
\newcommand{\DataTypeTok}[1]{\textcolor[rgb]{0.13,0.29,0.53}{#1}}
\newcommand{\DecValTok}[1]{\textcolor[rgb]{0.00,0.00,0.81}{#1}}
\newcommand{\DocumentationTok}[1]{\textcolor[rgb]{0.56,0.35,0.01}{\textbf{\textit{#1}}}}
\newcommand{\ErrorTok}[1]{\textcolor[rgb]{0.64,0.00,0.00}{\textbf{#1}}}
\newcommand{\ExtensionTok}[1]{#1}
\newcommand{\FloatTok}[1]{\textcolor[rgb]{0.00,0.00,0.81}{#1}}
\newcommand{\FunctionTok}[1]{\textcolor[rgb]{0.13,0.29,0.53}{\textbf{#1}}}
\newcommand{\ImportTok}[1]{#1}
\newcommand{\InformationTok}[1]{\textcolor[rgb]{0.56,0.35,0.01}{\textbf{\textit{#1}}}}
\newcommand{\KeywordTok}[1]{\textcolor[rgb]{0.13,0.29,0.53}{\textbf{#1}}}
\newcommand{\NormalTok}[1]{#1}
\newcommand{\OperatorTok}[1]{\textcolor[rgb]{0.81,0.36,0.00}{\textbf{#1}}}
\newcommand{\OtherTok}[1]{\textcolor[rgb]{0.56,0.35,0.01}{#1}}
\newcommand{\PreprocessorTok}[1]{\textcolor[rgb]{0.56,0.35,0.01}{\textit{#1}}}
\newcommand{\RegionMarkerTok}[1]{#1}
\newcommand{\SpecialCharTok}[1]{\textcolor[rgb]{0.81,0.36,0.00}{\textbf{#1}}}
\newcommand{\SpecialStringTok}[1]{\textcolor[rgb]{0.31,0.60,0.02}{#1}}
\newcommand{\StringTok}[1]{\textcolor[rgb]{0.31,0.60,0.02}{#1}}
\newcommand{\VariableTok}[1]{\textcolor[rgb]{0.00,0.00,0.00}{#1}}
\newcommand{\VerbatimStringTok}[1]{\textcolor[rgb]{0.31,0.60,0.02}{#1}}
\newcommand{\WarningTok}[1]{\textcolor[rgb]{0.56,0.35,0.01}{\textbf{\textit{#1}}}}
\usepackage{longtable,booktabs,array}
\usepackage{calc} % for calculating minipage widths
% Correct order of tables after \paragraph or \subparagraph
\usepackage{etoolbox}
\makeatletter
\patchcmd\longtable{\par}{\if@noskipsec\mbox{}\fi\par}{}{}
\makeatother
% Allow footnotes in longtable head/foot
\IfFileExists{footnotehyper.sty}{\usepackage{footnotehyper}}{\usepackage{footnote}}
\makesavenoteenv{longtable}
\usepackage{graphicx}
\makeatletter
\def\maxwidth{\ifdim\Gin@nat@width>\linewidth\linewidth\else\Gin@nat@width\fi}
\def\maxheight{\ifdim\Gin@nat@height>\textheight\textheight\else\Gin@nat@height\fi}
\makeatother
% Scale images if necessary, so that they will not overflow the page
% margins by default, and it is still possible to overwrite the defaults
% using explicit options in \includegraphics[width, height, ...]{}
\setkeys{Gin}{width=\maxwidth,height=\maxheight,keepaspectratio}
% Set default figure placement to htbp
\makeatletter
\def\fps@figure{htbp}
\makeatother
\setlength{\emergencystretch}{3em} % prevent overfull lines
\providecommand{\tightlist}{%
  \setlength{\itemsep}{0pt}\setlength{\parskip}{0pt}}
\setcounter{secnumdepth}{5}
% 页面设置
\usepackage{geometry}
\geometry{a4paper,left=2.5cm,right=2.5cm,top=2.5cm,bottom=2.5cm}

% 中文支持
\usepackage{ctex}

% 消除空白页
\let\cleardoublepage\clearpage
\ifLuaTeX
  \usepackage{selnolig}  % disable illegal ligatures
\fi
\usepackage[]{natbib}
\bibliographystyle{plainnat}
\usepackage{bookmark}
\IfFileExists{xurl.sty}{\usepackage{xurl}}{} % add URL line breaks if available
\urlstyle{same}
\hypersetup{
  pdftitle={统计学：从概率论到深度学习},
  pdfauthor={穿行冷风},
  hidelinks,
  pdfcreator={LaTeX via pandoc}}

\title{统计学：从概率论到深度学习}
\author{穿行冷风}
\date{2026-04-30}

\begin{document}
\maketitle

{
\setcounter{tocdepth}{1}
\tableofcontents
}
\chapter*{前言}\label{ux524dux8a00}
\addcontentsline{toc}{chapter}{前言}

这是一本讲统计的书。如果将来某天，这本书能出版，估计和放在大学图书馆的大部分''教科书''一样，平平无奇，无人在意。但是对我而言，这将会是在我平平无奇的大学生活中最美好的一笔回忆。

大学读的经济统计，选择这个专业谈不上心之所向，但也还没那么糟。学的越久，会越觉得''经济''和''统计''二字的含义的丰富与深刻。我喜欢这种深刻，并且希望能把我学习中收获的苦与乐，浅浅的写下来。

统计学，一方面有着数学的外壳和概率论的基础，看起来充满科学和理性。另一方面，连接着机器学习和大数据，感觉是引领潮流的领航人。但是，统计学的内核其实是哲学，是人生，是世界观，也是方法论，是沟通所见与所知的桥梁。``不确定''一词，诉说了人生的苦楚。我们终其一生，就是在不确定的世界中，拼命争取一丝确定性的解脱，哪怕只是暂时的幻觉。``观测''一词，道尽了认识的真谛。我们永远无法知道世界的真相是什么，我们穷尽一生，只能用有限的观测去逼近真实的世界，确定性的世界。

是决定，因果，宿命？还是概率，运气，玄学？总不能成功时归功于因果，失败时归功于运气吧。这些宏大的命题不是我这短短的前言能说明白的，我只希望能在后面的内容中，透过统计的视角，去思考，并尝试做出一些回答。

\section{不确定性：世界的底色}\label{ux4e0dux786eux5b9aux6027ux4e16ux754cux7684ux5e95ux8272}

我们生活在一个不确定的世界中。明天的天气、股市的涨跌、下一次掷硬币的结果------这些事件的共同特征是：在发生之前，我们无法确切知道结果。统计学将这种不确定性分为两类：

\begin{itemize}
\tightlist
\item
  \textbf{偶然性（Aleatory Uncertainty）}：由系统内在的随机性产生，不可消除。例如，即使知道硬币是公平的，仍无法预测下一次抛掷的结果。概率论为描述这类不确定性提供了语言。
\item
  \textbf{认知性（Epistemic Uncertainty）}：源于知识的不足，可以通过收集更多数据来减少。例如，一枚硬币是否公平？我们不确定，但随着抛掷次数的增加，我们对硬币性质的认知会越来越精确。统计推断的目标正是量化并减少这类不确定性。
\end{itemize}

这两类不确定性的区分，是理解统计学一切方法的逻辑起点。

\section{信息熵与贝叶斯更新}\label{ux4fe1ux606fux71b5ux4e0eux8d1dux53f6ux65afux66f4ux65b0}

1948年，香农（Claude Shannon）提出了一个根本性问题：\textbf{如何度量不确定性本身？} 他的答案是信息熵。对于一个取 \(K\) 个可能值的离散随机变量 \(X\)，其熵定义为：

\[H(X) = -\sum_{k=1}^{K} p_k \log_2 p_k\]

当分布完全均匀时熵最大------我们对 \(X\) 的取值最不确定；当 \(X\) 以概率 1 取某个值时熵为零------我们完全确定。换言之，\textbf{熵就是''惊讶程度''的期望值}。这个视角贯穿全书：最大似然估计（Ch06）本质上是在最小化经验分布与模型分布之间的信息损失；决策树的分裂准则（Ch18）使用熵来量化节点的''纯度''。

如果说信息熵提供了\textbf{状态}的度量，贝叶斯更新则提供了\textbf{变化}的规则：

\[P(\theta \mid \text{data}) = \frac{P(\text{data} \mid \theta) \cdot P(\theta)}{P(\text{data})}\]

它将学习过程归纳为先验 → 似然 → 后验三个步骤。这个过程是递归的：今天的后验就是明天的先验。每当你看到新数据，就用贝叶斯规则更新一次信念------这正是人类学习的理想化模型。

\begin{quote}
我们的一生又何尝不是一场旷日持久的贝叶斯更新。
\end{quote}

本书从古典概率（Ch01）出发，经历频率学派的统计推断（Ch05-07），最终回归到贝叶斯的完整框架（Ch08、Ch09）。频率学派和贝叶斯学派的分歧，本质上是关于''概率是什么''这一元问题的不同回答。

\section{观测、模型与真实}\label{ux89c2ux6d4bux6a21ux578bux4e0eux771fux5b9e}

统计学的核心张力存在于三者之间：

\begin{verbatim}
真实世界（不可直接观测）
    ↓ 产生
观测数据（有限、含噪声）
    ↓ 拟合
统计模型（近似、有偏）
    ↓ 推断
对真实世界的认知
\end{verbatim}

格利文科定理（Ch05）告诉我们：当样本量趋于无穷时，经验分布函数会一致收敛到真实分布。这是统计学最深刻的保证------足够多的观测可以逼近真实。但现实中的样本总是有限的，因此统计学的大部分工作是在''数据有限''的约束下做出尽可能好的推断。这引出了贯穿全书的三个核心权衡：

\begin{itemize}
\tightlist
\item
  \textbf{偏差-方差权衡}：简单模型偏差大但方差小，复杂模型偏差小但方差大（Ch10、Ch17）。
\item
  \textbf{信息-计算权衡}：精确的贝叶斯后验往往无法计算，需要 MCMC（Ch09）或变分推断来近似。
\item
  \textbf{可解释性-预测力权衡}：线性回归可解释但表达能力有限，深度网络预测力强但难以解释。
\end{itemize}

\section{统计思维的核心循环}\label{ux7edfux8ba1ux601dux7ef4ux7684ux6838ux5fc3ux5faaux73af}

一切统计工作，无论简单还是复杂，都遵循同一个基本循环：

\begin{verbatim}
提出问题 → 收集数据 → 建立模型 → 统计推断 → 做出决策
    ↑                                          |
    └──────────── 新问题/新数据 ←──────────────┘
\end{verbatim}

\subsection{提出问题}\label{ux63d0ux51faux95eeux9898}

一个好的统计问题始于对不确定性的识别。你需要问自己：我想知道什么？（估计一个参数？检验一个假说？预测未来？）不确定性的来源是什么？（抽样误差？测量误差？模型误设？）我能承受多大的错误风险？这一步没有数学公式，但它决定了后续一切方法的选择------Ch07 中你会看到，问题的提法直接决定了拒绝域的形式。

\subsection{收集数据}\label{ux6536ux96c6ux6570ux636e}

数据是统计的血肉。随机抽样（Ch05）保证了样本对总体的代表性；随机分配（Ch22）是因果推断的黄金标准；数据的类型（连续/离散/类别/截断）决定了可用的模型族（Ch10、Ch11、Ch12）。统计学中有一条无法绕过的铁律：\textbf{垃圾进，垃圾出}------再精巧的模型也无法挽救糟糕的数据。

\subsection{建立模型}\label{ux5efaux7acbux6a21ux578b}

模型是统计的骨架，在''足够真实''和''足够简单''之间取得平衡。概率模型（Ch01-04）用随机变量和分布描述数据生成过程；参数模型（Ch06）用少量参数概括分布形态，高效但有偏；非参数模型（Ch13）让数据自己说话，灵活但需要更多样本。所有模型都是错的，但有些模型是有用的------你将在 Ch10 的残差分析和 Ch17 的偏差-方差分解中反复品味这句话。

\subsection{统计推断}\label{ux7edfux8ba1ux63a8ux65ad}

推断是统计的心脏------从样本到总体、从已知到未知的跳跃。估计（Ch06）回答''真实参数最可能的值是多少''；检验（Ch07）回答''观测到的效应是真实信号还是随机噪声''；预测（Ch10、Ch17-20）回答''给定历史，未来会怎样''。频率学派和贝叶斯学派在''如何做推断''上存在深刻分歧，本书平等地呈现了两者（Ch05-07 为频率路径，Ch08、Ch09 为贝叶斯路径）。

\subsection{做出决策}\label{ux505aux51faux51b3ux7b56}

统计分析的终点不是 p 值，而是决策。Neyman-Pearson 框架（Ch07）本质上是一个决策理论：控制第 I 类错误的同时最小化第 II 类错误。分类器（Ch18）最小化期望误分类代价。因果推断（Ch22）更直接服务于''是否干预、如何干预''的决策问题。

\section{关于本书}\label{ux5173ux4e8eux672cux4e66}

统计学的内容很庞杂，我希望能把这本书越写越厚，然后越写越薄。也就是说，这会是一本一直''正在进行时''的书，直到我认为足够好才会停止。

本书的组织遵循从基础到前沿、从理论到计算的认知路径，分为五个部分：

\begin{verbatim}
Part I │ 概率论基础（4章）    │ 不确定性的数学语言
Part II│ 经典统计推断（5章）  │ 频率学派与贝叶斯学派，从理论到计算
Part III│ 现代统计方法（5章） │ 回归 → GLM → 类别数据 → 非参数 → 时间序列
Part IV │ 统计计算（2章）     │ 优化方法与矩阵计算
Part V │ 机器学习（6章）     │ 从经典分类到深度学习，以及多元与因果
\end{verbatim}

\textbf{各部分的逻辑关系}：

\begin{itemize}
\tightlist
\item
  Part I 构建了概率论的基础设施，所有后续章节都依赖它。
\item
  Part II 在两个哲学范式下建立了统计推断的核心逻辑：频率学派（Ch05-07）和贝叶斯学派（Ch08-09）。Ch08 建立理论框架，Ch09 紧接其后提供计算实现------二者形成完整的贝叶斯路径。
\item
  Part III 将推断方法应用到各种实际场景：连续响应（Ch10）、离散响应（Ch11）、类别数据（Ch12，紧接 GLM 以复用 logit 框架）、无分布假设（Ch13）、时间依赖（Ch14）。
\item
  Part IV 为前面所有方法提供计算支撑：优化算法（Ch15）、矩阵分解与数值方法（Ch16）。
\item
  Part V 从统计学自然延伸到机器学习：基础与评估（Ch17）、分类器（Ch18）、集成方法（Ch19）、深度学习（Ch20）、多元分析（Ch21）、因果推断（Ch22）。
\end{itemize}

\textbf{不同读者的阅读路径}：

\begin{itemize}
\tightlist
\item
  从头系统学习：按章节顺序阅读。
\item
  快速上手回归与 GLM：Part I → Ch10 → Ch11。
\item
  专注于机器学习：Part I → Ch06-07 → Part V。
\item
  对贝叶斯方法感兴趣：Part I → Ch08 → Ch09（理论与计算紧密衔接）。
\item
  关注因果推断：Part I → Ch10 → Ch22。
\end{itemize}

\mainmatter

\chapter{概率初步与随机变量的分布}\label{ux6982ux7387ux521dux6b65ux4e0eux968fux673aux53d8ux91cfux7684ux5206ux5e03}

\begin{quote}
你是一个观测者。你面对世界，唯一确定的事情是：你不可能提前知道一切。概率论，就是我们为``不知道''发明的语言。
\end{quote}

\section{观测者与可能性空间}\label{ux89c2ux6d4bux8005ux4e0eux53efux80fdux6027ux7a7aux95f4}

一个观测者站在时间之前。他不知道明天会不会下雨，不知道下一次掷硬币的结果，不知道自己还能活多少年。他唯一拥有的，是\textbf{对可能性的列举}------所有可能发生的结果的集合。

这就是\textbf{样本空间} \(\Omega\)：观测者能想象到的、互不重叠的全部可能结果的集合。每一个具体的结果 \(\omega\) 是一个\textbf{样本点}。

样本空间的大小，度量了观测者\textbf{无知的程度}。当 \(\Omega\) 只有一个元素时，观测者全知------他知道必然发生什么。当 \(\Omega\) 有无穷多个元素时，观测者的无知是彻底的。

但无知不等于无能为力。观测者可以做两件事：

\begin{enumerate}
\def\labelenumi{\arabic{enumi}.}
\tightlist
\item
  \textbf{划定事件}：在 \(\Omega\) 中圈出一个子集 \(A\)，说``我关心这个子集是否出现''------这就是\textbf{随机事件}。
\item
  \textbf{分配信念}：对每个事件赋予一个 0 到 1 之间的数值，表达``我认为它发生的可能性''------这就是\textbf{概率}。
\end{enumerate}

\section{随机事件的运算}\label{ux968fux673aux4e8bux4ef6ux7684ux8fd0ux7b97}

事件本质上是 \(\Omega\) 的子集。因此，事件之间的运算就是集合运算------但每一个运算背后都有直观的含义。

\begin{itemize}
\tightlist
\item
  \textbf{并} \(A \cup B\)：\(A\) 或 \(B\) 至少有一个发生。这是在\textbf{扩大}可能性------只要落入任一区域就算。
\item
  \textbf{交} \(A \cap B\)：\(A\) 和 \(B\) 同时发生。这是在\textbf{缩小}可能性------必须同时满足两个条件。
\item
  \textbf{补} \(\bar{A}\)：\(A\) 不发生。这是观测者视野的翻转------关注的是 \(A\) 之外的一切。
\item
  \textbf{差} \(A - B\)：\(A\) 发生而 \(B\) 不发生。这是精确排除------``我要 \(A\)，但不要 \(B\) 附带进来''。
\end{itemize}

运算遵循对偶律：事件并的对立等于对立的交；事件交的对立等于对立的并。这是逻辑中德摩根律的概率表达。

\textbf{事件域} \(\mathcal{F}\) 是 \(\Omega\) 中那些``可以被赋予概率''的子集的全体。并非所有子集都有资格------我们只对``封闭于可列次运算''的子集类分配概率。\((\Omega, \mathcal{F})\) 构成了概率空间的骨架。

\section{概率是什么？}\label{ux6982ux7387ux662fux4ec0ux4e48}

这个问题争论了三百年。四种回答，分属两个阵营，追问同一个元问题：

\begin{quote}
\textbf{概率是世界的属性，还是认知的属性？}
\end{quote}

\subsection{阵营一：概率是世界的属性------客观概率}\label{ux9635ux8425ux4e00ux6982ux7387ux662fux4e16ux754cux7684ux5c5eux6027ux5ba2ux89c2ux6982ux7387}

\textbf{古典概率}（拉普拉斯，1812）是最早的系统回答。当 \(\Omega\) 中的每个样本点''理由均等''时：

\[P(A) = \frac{A \text{ 所含样本点数}}{\Omega \text{ 中样本点总数}}\]

拉普拉斯称之为''不充分理由原理''------如果你没有理由认为某些结果比另一些更可能，就认为它们等可能。古典概率的美，在于它将无知转化为等权。但它只适用于有限且对称的样本空间，无法处理有偏硬币，也无法回答''明天股市崩盘的概率''。

\textbf{频率概率}（冯·米塞斯，1919）继承了客观主义的路线，但换了一种操作化的进路：概率不是枚举出来的，而是在大量重复试验中，事件出现频率的\textbf{稳定值}。抛硬币十万次，正面频率稳定在 0.5 附近------这个稳定值就是 \(1/2\)。

频率方法的问题是：它假设世界可以无限重复。但很多事件不可重复------``明天股市崩盘的概率''没有频率解释。这个问题要到 Ch03 的大数定律才能给出数学担保：至少在可重复的场景中，频率收敛到概率是一个可以被证明的定理，而不只是经验观察。

古典概率和频率概率共享同一个信念：\textbf{概率客观存在于随机机制中，与观测者无关。}

\subsection{阵营二：概率是认知的属性------主观概率}\label{ux9635ux8425ux4e8cux6982ux7387ux662fux8ba4ux77e5ux7684ux5c5eux6027ux4e3bux89c2ux6982ux7387}

\textbf{主观概率}（拉姆齐 1926，德·菲内蒂 1937）给出完全不同的回答：概率是观测者根据经验和信息对事件发生的\textbf{个人信念}的数值表达。你的概率和我的概率可以不同------因为我们掌握的信息不同。硬币的''公平性''不是硬币的物理属性，而是观测者没有理由偏向任何一方的认知状态。

主观概率通向贝叶斯统计（Ch08）：概率是观测者认知状态的度量。新数据流入，概率更新。这个世界观下，不存在绝对的''真实概率''------只存在不断被数据修正的后验信念。

\begin{quote}
两派的共识比你想象的更多。频率学派在大样本下大量使用贝叶斯式的思维（MLE 的渐近正态性，Ch06）；贝叶斯学派在后验收敛的意义上同样依赖大数定律（Ch03）。Ch06 末尾将通过 KL 散度展示：最大似然估计和贝叶斯后验模态，在信息论上是同一种操作的两个视角。
\end{quote}

\subsection{超越争论：概率的公理化定义}\label{ux8d85ux8d8aux4e89ux8bbaux6982ux7387ux7684ux516cux7406ux5316ux5b9aux4e49}

科尔莫戈罗夫做了一个所有争论方都能接受的智慧之举------\textbf{绕过''概率是什么''的本体论问题，直接规定''概率应该如何运算''}：

\begin{itemize}
\tightlist
\item
  \textbf{非负性}：\(P(A) \geq 0\)。概率不能是负的------你不能比''绝不可能''更不可能。
\item
  \textbf{正则性}：\(P(\Omega) = 1\)。必然事件的概率为 1------可能性空间必须有一个度量上限。
\item
  \textbf{可列可加性}：若 \(A_1, A_2, \ldots\) 互不相容，则 \(P(\cup A_i) = \sum P(A_i)\)。互斥事件的概率可加------这是概率运算的核心规则。
\end{itemize}

公理化定义把概率变成了纯粹的数学对象。无论你认为概率是客观频率还是主观信念，三条公理都成立。\textbf{它提供的不是答案，而是一个中立的语法平台------在这个平台上，两派的对话可以继续。}

\begin{quote}
这四种回答之间的关系贯穿全书：Ch05-07 走频率路径（抽样分布→MLE→假设检验），Ch08 走贝叶斯路径（先验→后验→共轭），Ch03-04 为两者提供共同的渐近地基。它们不是敌人，而是从不同方向攀登同一座山。
\end{quote}

从公理出发，可以推导出概率的运算性质：有限可加性、补事件概率、单调性、加法公式、连续性。这些性质构成了概率运算的语法规则。

\section{信息熵：不确定性的数学度量}\label{ux4fe1ux606fux71b5ux4e0dux786eux5b9aux6027ux7684ux6570ux5b66ux5ea6ux91cf}

前面讨论了概率的四种定义，但它们都没有回答一个更根本的问题：\textbf{不确定性本身如何度量？}

1948 年，香农（Claude Shannon）在《通信的数学理论》中给出了答案。对于一个取 \(K\) 个可能值的离散随机变量 \(X\)，其分布为 \(P(X = k) = p_k\)，香农定义了\textbf{信息熵}：

\[H(X) = -\sum_{k=1}^{K} p_k \log_2 p_k\]

这个公式的含义出人意料地朴素：\(-\log_2 p_k\) 是事件 \(\{X = k\}\) 发生时所携带的\textbf{信息量}（单位为比特）------概率越小的事件，一旦发生，信息量越大。而熵，是信息量的\textbf{期望值}。换言之：

\begin{quote}
\textbf{熵 = 你在观测之前，预期会接收到的''惊讶程度''的平均值。}
\end{quote}

\subsection{熵的取值区间}\label{ux71b5ux7684ux53d6ux503cux533aux95f4}

当 \(X\) 完全均匀分布（\(p_k = 1/K\)）时，熵取到最大值 \(\log_2 K\)。此时你对 \(X\) 的取值最不确定------每一种可能性都同等合理，观测会给你最大的平均惊讶。

当 \(X\) 以概率 1 取某个值（退化分布）时，熵为零。你完全确定 \(X\) 会取什么------观测不会带来任何惊讶。

\textbf{一切其他分布的熵，都落在} \([0, \log_2 K]\) 之间。 也就是说，熵给出了一个从''完全确定''到''完全无知''的连续标尺。

\subsection{熵的深层含义}\label{ux71b5ux7684ux6df1ux5c42ux542bux4e49}

熵不仅仅是一个数字，它蕴含了几个贯穿全书的哲学命题：

\begin{enumerate}
\def\labelenumi{\arabic{enumi}.}
\item
  \textbf{无知的可量化性}：样本空间的大小给出了''无知的绝对边界''（\(\log_2 K\)），但熵告诉你在这个边界之内，你的无知处于什么水平。先验均匀分布 \(\Leftrightarrow\) 最大熵 \(\Leftrightarrow\) 最诚实的无知。
\item
  \textbf{信息就是熵的减少}：观测前熵为 \(H(X)\)，观测到 \(Y=y\) 后熵变为 \(H(X \mid Y=y)\)。两者之差，就是 \(Y=y\) 带来的\textbf{信息增益}。这正是决策树（Ch18）选择分裂变量的核心准则。
\item
  \textbf{最大熵原理}：在只知道部分约束（如均值和方差）的情况下，选择熵最大的分布是最''不武断''的选择。正态分布正是给定均值和方差下熵最大的连续分布------Ch08 的贝叶斯先验选择将回到这个原理。
\item
  \textbf{熵与似然的等价性}：最小化模型分布与经验分布之间的交叉熵，等价于最大化模型似然（Ch06）。统计学中最核心的估计方法和机器学习中最核心的损失函数，共享同一个信息论根基。
\end{enumerate}

\begin{quote}
前言中提到了''信息熵与贝叶斯更新''作为全书的底层哲学。后文讲完分布之后，我们将回到信息论------引入 KL 散度和交叉熵------届时上述第 4 点的等价性将得到完整的阐释。这条线索将贯穿全书：MLE 的信息论解释（Ch06）→ 决策树的信息增益（Ch18）→ 交叉熵损失（Ch17、Ch20）。
\end{quote}

\section{条件概率：新信息流入导致的可能性空间塌缩}\label{ux6761ux4ef6ux6982ux7387ux65b0ux4fe1ux606fux6d41ux5165ux5bfcux81f4ux7684ux53efux80fdux6027ux7a7aux95f4ux584cux7f29}

这是全书最重要的概念之一。不是因为它难，而是因为它道出了学习的本质。

\subsection{从无条件到有条件}\label{ux4eceux65e0ux6761ux4ef6ux5230ux6709ux6761ux4ef6}

你原本面对样本空间 \(\Omega\)，对事件 \(A\) 的概率评估是 \(P(A)\)。现在你得知：事件 \(B\) 已经发生了。

这一刻发生了什么？你的可能性空间\textbf{塌缩}了------从整个 \(\Omega\) 缩小到 \(B\)。\(B\) 之外的一切结果都变得不可能。在塌缩后的新空间里，\(A\) 必须落在 \(B\) 之内才会有可能发生。

\[P(A|B) = \frac{P(A \cap B)}{P(B)}\]

分子是 \(A\) 在 \(B\) 中的那部分，分母 \(P(B)\) 是归一化因子------确保塌缩后的概率仍然加起来等于 1。

这就是\textbf{条件概率}。它描述的，是观测者接收到新信息后的认知更新。

\begin{quote}
每当你听到一个消息，你的大脑就在做条件概率计算------只是它不写公式。
\end{quote}

\subsection{乘法公式：信息的链式流入}\label{ux4e58ux6cd5ux516cux5f0fux4fe1ux606fux7684ux94feux5f0fux6d41ux5165}

\[P(A_1 A_2 \cdots A_n) = P(A_1) P(A_2|A_1) P(A_3|A_1 A_2) \cdots P(A_n|A_1 A_2 \cdots A_{n-1})\]

事件依次发生，每次发生都改变后续事件的条件概率。这是时间顺序的数学表达------先发生什么，影响后发生什么。

\subsection{全概率公式：穷举所有塌缩路径}\label{ux5168ux6982ux7387ux516cux5f0fux7a77ux4e3eux6240ux6709ux584cux7f29ux8defux5f84}

\[P(A) = \sum_{i=1}^n P(B_i) P(A|B_i)\]

将 \(\Omega\) 划分成互不相容的 \(B_1, \ldots, B_n\)，\(A\) 的总概率等于：先塌缩到每个 \(B_i\) 的概率，乘以在 \(B_i\) 中 \(A\) 发生的条件概率，再求和。

这是在说：如果你不知道 \(A\) 的总概率，可以先问``如果世界是 \(B_1\) 会怎样''，再问``如果世界是 \(B_2\) 会怎样''，最后加权平均。先分解，再合成------这是分析复杂问题的通用策略。

\subsection{贝叶斯公式：逆着信息流回推}\label{ux8d1dux53f6ux65afux516cux5f0fux9006ux7740ux4fe1ux606fux6d41ux56deux63a8}

\[P(B_i|A) = \frac{P(B_i) P(A|B_i)}{\sum_{j=1}^n P(B_j) P(A|B_j)}\]

全概率公式是从原因推结果（\(B_i \to A\)），贝叶斯公式是从结果反推原因（\(A \to B_i\)）。

观测者看到了结果 \(A\)。他问：这个结果是由哪个原因导致的？在 \(A\) 发生之前，他对各原因的信念是 \(P(B_i)\)（先验概率）；在 \(A\) 发生之后，他更新为 \(P(B_i|A)\)（后验概率）。

\textbf{这个过程是递归的}：今天的后验就是明天的先验。每看到一条新数据，就用贝叶斯公式更新一次认知。人类学习的理想化模型，不过如此。

\begin{quote}
这是统计学习的核心引擎：Ch08 将其升华为完整的贝叶斯推断体系，Ch09 处理贝叶斯后验的计算实现。
\end{quote}

\begin{quote}
贝叶斯公式是信息流动的逆运算。前向流动是世界→数据（似然），反向流动是数据→世界（后验）。这个双向性将在 Ch08 得到完整展开。
\end{quote}

\subsection{独立性：信息通道的断裂}\label{ux72ecux7acbux6027ux4fe1ux606fux901aux9053ux7684ux65adux88c2}

\[P(AB) = P(A)P(B)\]

如果知道 \(B\) 是否发生，不改变你对 \(A\) 的概率评估------即 \(P(A|B) = P(A)\)------则称 \(A\) 与 \(B\) 独立。

\textbf{独立性的本质}：\(A\) 和 \(B\) 之间不存在信息通道。观测到一个，不会带来另一个的任何信息。

\begin{quote}
独立性假设是现代统计学中使用最频繁、也最容易滥用的假设。如果说``所有模型都是错的''，那么``错误地假设独立''是最常见的那类错误。Ch02 将展示相关性如何打破独立性，Ch13 将提供不依赖独立假设的推断方法。
\end{quote}

\section{随机变量}\label{ux968fux673aux53d8ux91cf}

样本点 \(\omega\) 可能是``正面''或``反面''，可能是``晴天''\,``阴天''``雨天''------这些不是数字，不便运算。\textbf{随机变量} \(X(\omega)\) 做的事，就是把每一个样本点映射到一个实数：

\[X: \Omega \to \mathbb{R}\]

映射之后，随机事件 \(\{\omega: X(\omega) \leq a\}\) 就成了实数轴上的一个区间。由此可以定义\textbf{分布函数}：

\[F(x) = P(X \leq x)\]

分布函数有三条基本性质：单调不减（值越大，累积概率越大）、右连续、在 \(-\infty\) 处为 0、在 \(+\infty\) 处为 1。

分布函数\textbf{完全}刻画了随机变量的统计规律。知道分布函数，就知道 \(X\) 落在任何区间的概率。反过来，两个随机变量同分布，当且仅当它们的分布函数处处相等。

随机变量分为两大类：

\begin{itemize}
\tightlist
\item
  \textbf{离散型}：\(X\) 只能取有限个或可列个值。用\textbf{分布列} \(p(x_i) = P(X = x_i)\) 描述。
\item
  \textbf{连续型}：\(X\) 可以取某个区间内的任意值。用\textbf{概率密度函数} \(p(x)\) 描述，满足 \(F(x) = \int_{-\infty}^x p(t)dt\)。
\end{itemize}

密度函数 \(p(x)\) 在单点 \(x\) 处的值\textbf{不是}概率（单点的概率为 0），而是概率的``瞬时密度''------\(p(x)dx\) 才是 \(X\) 落在 \([x, x+dx]\) 内的近似概率。

\section{期望与方差}\label{ux671fux671bux4e0eux65b9ux5dee}

如果分布函数是随机变量的完整肖像，那么期望和方差就是这张肖像上最重要的两个特征。

\subsection{数学期望：消除随机性的手段}\label{ux6570ux5b66ux671fux671bux6d88ux9664ux968fux673aux6027ux7684ux624bux6bb5}

\[E(X) = \begin{cases} \sum_i x_i p(x_i) & \text{离散} \\ \int_{-\infty}^{\infty} x p(x) dx & \text{连续} \end{cases}\]

\textbf{期望不是``最可能的值''，而是``长期平均''}。如果你能无限次重复随机试验，所有观测值的平均会收敛到期望------这是大数定律（Ch03）的承诺。

\begin{quote}
期望是概率分布的\textbf{重心}。如果把概率质量 \(p(x)\) 看作物理质量沿实数轴的分布，那么 \(E(X)\) 就是质心。它是一阶矩 \(\mu_1 = E(X)\)。
\end{quote}

一个重要性质：对任意函数 \(g\)，

\[E[g(X)] = \sum_i g(x_i) p(x_i) \quad \text{或} \quad \int_{-\infty}^{\infty} g(x) p(x) dx\]

这个公式被称为``\textbf{无意识统计学家的法则}''------你不需要先求出 \(Y = g(X)\) 的分布，就可以直接计算 \(E(Y)\)。

\subsection{方差：对确定性掌控程度的度量}\label{ux65b9ux5deeux5bf9ux786eux5b9aux6027ux638cux63a7ux7a0bux5ea6ux7684ux5ea6ux91cf}

\[Var(X) = E[(X - E(X))^2] = E(X^2) - [E(X)]^2\]

期望告诉你``中心在哪''，方差告诉你``离中心有多远''。

\begin{quote}
方差是概率分布的\textbf{转动惯量}。如果 \(E(X)\) 是质心，那么 \(Var(X)\) 就是概率质量绕质心的转动惯量。方差越大，分布越``散''。
\end{quote}

标准差 \(\sigma(X) = \sqrt{Var(X)}\) 与 \(X\) 有相同的量纲，更适合直观理解。``3σ 原则''（正态分布下，约 99.7\% 的概率落在 \(\mu \pm 3\sigma\) 内）是对方差最常用的直观解释。

\subsection{切比雪夫不等式：方差对偏离的控制力}\label{ux5207ux6bd4ux96eaux592bux4e0dux7b49ux5f0fux65b9ux5deeux5bf9ux504fux79bbux7684ux63a7ux5236ux529b}

\[P(|X - E(X)| \geq \epsilon) \leq \frac{Var(X)}{\epsilon^2}\]

切比雪夫不等式是\textbf{方差定义的目的}：方差之所以被定义为二阶中心矩，正是因为它能给出偏离期望的概率上界。方差越大，这个上界越宽松------你对自己``离中心有多远''的掌控越弱。

\begin{quote}
切比雪夫不等式是最基础的概率集中不等式。它只需要方差存在，不需要知道分布的具体形状。Ch04 将展示一系列更精细的不等式（Hoeffding、Bernstein），它们用更强的假设换取更紧的界。
\end{quote}

\subsection{其他特征数：分布形态的完整素描}\label{ux5176ux4ed6ux7279ux5f81ux6570ux5206ux5e03ux5f62ux6001ux7684ux5b8cux6574ux7d20ux63cf}

\begin{itemize}
\tightlist
\item
  \textbf{k 阶原点矩} \(\mu_k = E(X^k)\)：一阶原点矩是期望。
\item
  \textbf{k 阶中心矩} \(\nu_k = E[(X - E(X))^k]\)：二阶中心矩是方差。
\item
  \textbf{偏度系数} \(\beta_s = \nu_3 / \nu_2^{3/2}\)：度量分布的左右不对称性。\(\beta_s < 0\) 表示左偏（左尾更长），\(\beta_s > 0\) 表示右偏。
\item
  \textbf{峰度系数} \(\beta_k = \nu_4 / \nu_2^2 - 3\)：度量分布的尖峭程度（相对于正态分布）。\(\beta_k > 0\) 表示比正态分布更尖、尾部更厚。
\item
  \textbf{分位数} \(x_p\)：满足 \(F(x_p) = p\)。中位数是 0.5 分位数------它将分布的概率质量一分为二。
\item
  \textbf{变异系数} \(CV = \sigma / \mu\)：消除量纲后的离散度，用于比较不同尺度的变量的波动性。
\end{itemize}

\section{分布的演化树：从伯努利到万物}\label{ux5206ux5e03ux7684ux6f14ux5316ux6811ux4eceux4f2fux52aaux5229ux5230ux4e07ux7269}

分布不是各自孤立的孤岛，而是一个彼此关联的\textbf{演化网络}。理解了这个网络，就理解了一个随机过程如何衍生出不同的观测维度。

\subsection{第一棵树：伯努利过程 ------ 计数与等待}\label{ux7b2cux4e00ux68f5ux6811ux4f2fux52aaux5229ux8fc7ux7a0b-ux8ba1ux6570ux4e0eux7b49ux5f85}

源头是一枚硬币，正面概率为 \(p\)。重复抛掷，每次独立------这就是\textbf{伯努利过程}，所有离散分布之母。

\begin{verbatim}
                  伯努利试验 Be(p)
                       │
          ┌────────────┼────────────┐
          │            │            │
    单次观测：      n次独立和：   等待第r次成功：
    两点分布      二项分布       负二项分布
    X~0-1        X~b(n,p)       X~Nb(r,p)
                    │               │
               n大p小 →          r=1 →
              泊松分布           几何分布
              X~P(λ)            X~Ge(p)
\end{verbatim}

\textbf{两点分布}（0-1 分布）\(P(X=x) = p^x(1-p)^{1-x}, x \in \{0,1\}\) 是最简单的随机变量------一次试验，成功或失败。

\textbf{二项分布} \(P(X=k) = C_n^k p^k (1-p)^{n-k}\)：\(n\) 次独立伯努利试验中成功次数的分布。它的随机变量可以分解为 \(n\) 个独立同分布的两点分布之和。

\[
E(X) = np, \quad Var(X) = np(1-p)
\]

当 \(n\) 很大而 \(p\) 很小时，二项分布的计算量巨大。\textbf{泊松分布}提供了优雅的近似------令 \(\lambda = np\)：

\[\lim_{n \to \infty} C_n^k p_n^k (1-p_n)^{n-k} = \frac{\lambda^k}{k!} e^{-\lambda}\]

泊松分布只有一个参数 \(\lambda\)，且 \(E(X) = Var(X) = \lambda\)------均值和方差相等，这是它的签名。

\begin{quote}
泊松分布常用于单位时间（或单位面积）内随机事件发生次数的建模：一天内进店的顾客数、一页书中的错字数、一平方公里内的陨石坑数。
\end{quote}

\textbf{几何分布} \(P(X=k) = (1-p)^{k-1} p\)：在伯努利过程中，\textbf{首次成功所需的试验次数}。

\[E(X) = \frac{1}{p}, \quad Var(X) = \frac{1-p}{p^2}\]

几何分布拥有\textbf{无记忆性}：\(P(X > m+n \mid X > m) = P(X > n)\)。前面 \(m\) 次已经失败了，后面 \(n\) 次仍然失败的概率，与``从一开始就连续失败 \(n\) 次''完全一样。伯努利过程不记得过去------这是独立性最戏剧性的后果。

\textbf{负二项分布}：等待第 \(r\) 次成功所需的试验次数。它可以分解为 \(r\) 个独立几何分布之和。

\subsection{第二棵树：泊松流 ------ 计数与间隔的双面}\label{ux7b2cux4e8cux68f5ux6811ux6ccaux677eux6d41-ux8ba1ux6570ux4e0eux95f4ux9694ux7684ux53ccux9762}

泊松过程描述的是事件随时间随机发生的模型：事件以恒定速率 \(\lambda\) 独立地发生。

\begin{verbatim}
                    泊松过程 (速率λ)
                         │
            ┌────────────┴────────────┐
            │                         │
        [0,t]内的事件次数         第k-1次到第k次
        N(t) ~ P(λt)              事件的时间间隔
        (计数维度)                T_k ~ Exp(λ)
                                      │
                                 r个间隔之和：
                                Tr ~ Ga(r, λ)
                                (等待时间维度)
\end{verbatim}

\textbf{泊松分布} \(P(N(t) = k) = \frac{(\lambda t)^k}{k!} e^{-\lambda t}\)：在长为 \(t\) 的时间区间内，事件发生 \(k\) 次的概率。

\textbf{指数分布} \(p(t) = \lambda e^{-\lambda t}, t > 0\)：\textbf{相邻两次事件之间的等待时间}。

\[E(T) = \frac{1}{\lambda}, \quad Var(T) = \frac{1}{\lambda^2}\]

指数分布同样拥有无记忆性：\(P(T > s+t \mid T > s) = P(T > t)\)。一个已经使用了 \(s\) 小时没有故障的设备，其剩余寿命的分布和全新设备一样------这是泊松过程的另一个面孔。

\begin{quote}
泊松分布和指数分布从两个角度描述同一个泊松过程。计数维度（离散）与间隔维度（连续）之间的对偶，是随机过程理论中最优美的结构之一。
\end{quote}

\textbf{伽马分布} \(X \sim Ga(\alpha, \lambda)\)：\(\alpha\) 个独立的指数分布之和。

\[p(x) = \frac{\lambda^\alpha}{\Gamma(\alpha)} x^{\alpha-1} e^{-\lambda x}, \quad x \geq 0\]

\[E(X) = \frac{\alpha}{\lambda}, \quad Var(X) = \frac{\alpha}{\lambda^2}\]

当 \(\alpha = 1\) 时，伽马分布退化为指数分布。当 \(\alpha = n/2\)、\(\lambda = 1/2\) 时，伽马分布即为卡方分布 \(\chi^2(n)\)------这是 Ch05 抽样分布的核心角色。

\subsection{第三棵树：极限与共轭 ------ 正态与贝塔的独特地位}\label{ux7b2cux4e09ux68f5ux6811ux6781ux9650ux4e0eux5171ux8f6d-ux6b63ux6001ux4e0eux8d1dux5854ux7684ux72ecux7279ux5730ux4f4d}

\textbf{正态分布}是一条特殊的河流------几乎所有其他分布，在``大量独立叠加''的极限下，都汇入正态之海。这就是中心极限定理（Ch03）的断言。

\[p(x) = \frac{1}{\sqrt{2\pi}\sigma} e^{-\frac{(x-\mu)^2}{2\sigma^2}}, \quad X \sim N(\mu, \sigma^2)\]

\[E(X) = \mu, \quad Var(X) = \sigma^2\]

正态分布有两个参数：\(\mu\) 控制中心位置，\(\sigma\) 控制离散程度。通过标准化变换 \(U = (X - \mu)/\sigma \sim N(0,1)\)，任何正态变量都可以转化为标准正态。

\begin{quote}
高斯在研究天文观测误差时首次使用正态分布刻画误差。他证明了：在``误差的期望为零且独立同分布''的假设下，\textbf{最小二乘法等价于正态误差模型下的最大似然估计}。这个优雅的等价性是回归分析（Ch10）的逻辑起点。
\end{quote}

正态分布还是给定均值和方差下\textbf{熵最大}的分布------在只知道这两个约束时，正态分布是``最不武断''的选择。这个最大熵性质在 Ch08 的贝叶斯先验选择和 Ch20 的变分自编码器中都会再次出现。

\textbf{贝塔分布} \(X \sim Be(a, b)\)，取值在 \((0, 1)\)：

\[p(x) = \frac{\Gamma(a+b)}{\Gamma(a)\Gamma(b)} x^{a-1} (1-x)^{b-1}, \quad x \in (0,1)\]

它的取值天然是\textbf{概率}------贝塔分布在贝叶斯统计中作为二项分布的\textbf{共轭先验}出现：先验是 Be(a,b)，观测到 \(k\) 次成功和 \(n-k\) 次失败后，后验是 Be(a+k, b+n-k)。这是一个干净的贝叶斯更新（Ch08）。

并非所有分布都生长在前面的演化树上。有些分布不依附于特定的随机过程，但因其基础性而在整个概率体系中扮演枢纽角色。

\subsection{均匀分布与超几何分布}\label{ux5747ux5300ux5206ux5e03ux4e0eux8d85ux51e0ux4f55ux5206ux5e03}

\textbf{均匀分布} \(U(a,b)\)：区间内每一点的可能性均等------这是连续版本的``不充分理由原理''。均匀分布还有一个关键地位：若 \(F_X\) 是严格单调的连续分布函数，则 \(Y = F_X(X) \sim U(0,1)\)。这表明\textbf{任何连续分布都可以通过均匀分布生成}------这是所有随机数生成算法的理论基础。

\textbf{超几何分布}：有限总体中的不放回抽样。当样本量 \(n\) 远小于总体量 \(N\) 时，超几何分布可近似为二项分布------不放回近似为放回，这是统计实践中常用的简化。

\section{分布之间的距离：从 KL 散度到交叉熵}\label{ux5206ux5e03ux4e4bux95f4ux7684ux8dddux79bbux4ece-kl-ux6563ux5ea6ux5230ux4ea4ux53c9ux71b5}

前面的熵度量了\textbf{单个分布}的不确定性。接下来一个自然的问题是：\textbf{两个分布之间有多''远''？}

这个问题的答案，跨越了从参数估计到深度学习的所有统计方法。

\subsection{KL 散度：信息论意义上的''距离''}\label{kl-ux6563ux5ea6ux4fe1ux606fux8bbaux610fux4e49ux4e0aux7684ux8dddux79bb}

对于离散分布 \(P\) 和 \(Q\)（定义在同一个样本空间上），\textbf{KL 散度}（Kullback--Leibler Divergence）定义为：

\[D_{KL}(P \parallel Q) = \sum_{k} p_k \log \frac{p_k}{q_k} = \sum_k p_k \log p_k - \sum_k p_k \log q_k\]

连续版本只需将求和换为积分。

拆解开来看：

\begin{itemize}
\tightlist
\item
  第一项 \(\sum p_k \log p_k = -H(P)\)，是 \(P\) 的负熵。
\item
  第二项 \(-\sum p_k \log q_k\)，是\textbf{交叉熵} \(H(P, Q)\)------用 \(Q\) 的编码方案去描述 \(P\) 时，所需的平均比特数。
\end{itemize}

因此：

\[D_{KL}(P \parallel Q) = H(P, Q) - H(P)\]

\begin{quote}
KL 散度 = 交叉熵 − 熵。它的物理含义是：\textbf{用} \(Q\) 近似 \(P\) 时，额外浪费的信息量。
\end{quote}

KL 散度的两个关键性质：

\begin{enumerate}
\def\labelenumi{\arabic{enumi}.}
\item
  \textbf{非负性}：\(D_{KL}(P \parallel Q) \geq 0\)，等号成立当且仅当 \(P = Q\)（几乎处处）。这是吉布斯不等式（Gibbs' inequality）的结论。无论你怎么选 \(Q\)，用 \(Q\) 去编码 \(P\) 永远需要至少 \(H(P)\) 个比特------不可能比用 \(P\) 自己的最优编码更省。
\item
  \textbf{不对称性}：\(D_{KL}(P \parallel Q) \neq D_{KL}(Q \parallel P)\)（一般情况下）。因此 KL 散度不是真正的距离度量，而是\textbf{方向性的信息损失}。\(D_{KL}(P \parallel Q)\) 和 \(D_{KL}(Q \parallel P)\) 回答的是不同的问题：

  \begin{itemize}
  \tightlist
  \item
    \(D_{KL}(P \parallel Q)\)：用 \(Q\) 近似真实分布 \(P\)，损失多少信息？（\textbf{正向 KL}，MLE 的方向）
  \item
    \(D_{KL}(Q \parallel P)\)：用 \(P\) 近似 \(Q\)，损失多少信息？（\textbf{反向 KL}，变分推断的方向，Ch09）
  \end{itemize}
\end{enumerate}

\subsection{从 KL 散度到最大似然估计}\label{ux4ece-kl-ux6563ux5ea6ux5230ux6700ux5927ux4f3cux7136ux4f30ux8ba1}

这是全书最重要的等价关系之一。设有真实数据分布 \(P_{\text{data}}\)（由经验分布 \(\hat{P}_n\) 逼近）和参数模型 \(Q_\theta\)：

\[\hat{\theta}_{\text{MLE}} = \arg\max_\theta \sum_{i=1}^n \log q_\theta(x_i)\]

而最小化 \(D_{KL}(\hat{P}_n \parallel Q_\theta)\)：

\[\arg\min_\theta D_{KL}(\hat{P}_n \parallel Q_\theta) = \arg\min_\theta \left[-\sum_i \frac{1}{n} \log q_\theta(x_i) - H(\hat{P}_n)\right] = \arg\max_\theta \sum_i \log q_\theta(x_i)\]

\(H(\hat{P}_n)\) 与 \(\theta\) 无关，消去。剩下的是对数似然。

\begin{quote}
\textbf{最大似然估计 = 最小化经验分布与模型分布之间的 KL 散度。} 换言之，MLE 在寻找那个与数据分布''信息损失最小''的参数。这是频率学派和贝叶斯学派（通过 KL 先验）共享的逻辑内核，也是 Ch06 与 Ch08 对话的数学基础。
\end{quote}

\subsection{交叉熵损失：深度学习的默认语言}\label{ux4ea4ux53c9ux71b5ux635fux5931ux6df1ux5ea6ux5b66ux4e60ux7684ux9ed8ux8ba4ux8bedux8a00}

在分类问题（Ch18、Ch20）中，设真实标签为 one-hot 向量 \(\mathbf{y}\)，模型预测概率为 \(\hat{\mathbf{y}}\)。交叉熵损失为：

\[\mathcal{L}_{\text{CE}} = -\sum_{c=1}^C y_c \log \hat{y}_c\]

这正是 \(H(\mathbf{y}, \hat{\mathbf{y}})\)------真实分布与预测分布之间的交叉熵。因为 \(\mathbf{y}\) 的熵为零（完全确定），所以 \(D_{KL}(\mathbf{y} \parallel \hat{\mathbf{y}}) = H(\mathbf{y}, \hat{\mathbf{y}})\)。交叉熵损失\textbf{就是 KL 散度在分类问题中的特例}。

\begin{quote}
从伯努利（1713）到香农（1948）到 Kullback--Leibler（1951）到现代深度学习的交叉熵损失------这条 300 年的思想演化线，始于概率，终于信息，贯穿全书。
\end{quote}

\subsection{其他分布距离}\label{ux5176ux4ed6ux5206ux5e03ux8dddux79bb}

除了 KL 散度，还有几类重要的分布距离/散度将在后续章节中出现：

\begin{itemize}
\tightlist
\item
  \textbf{全变差距离}（Total Variation）：\(TV(P, Q) = \frac{1}{2}\sum |p_k - q_k|\)。在假设检验（Ch07）的检验功效分析和概率不等式（Ch04）中自然出现。
\item
  \textbf{Wasserstein 距离}（Earth Mover's Distance）：在深度学习（Ch20，WGAN）和最优传输理论中作为生成模型的训练目标。
\item
  \textbf{Jensen-Shannon 散度}：KL 的对称化版本，是原始 GAN 的理论基础。
\end{itemize}

\section{随机变量函数的分布}\label{ux968fux673aux53d8ux91cfux51fdux6570ux7684ux5206ux5e03}

已知 \(X\) 的分布，\(Y = g(X)\) 的分布是什么？

\textbf{离散情况}：\(Y\) 仍是离散的，枚举 \(X\) 的取值、计算对应的 \(Y\) 值、合并相同 \(Y\) 的概率。

\textbf{连续情况}：若 \(g\) 严格单调，其反函数为 \(h\)，则

\[p_Y(y) = p_X(h(y)) \cdot |h'(y)|\]

直观理解：\(X\) 落在 \([x, x+dx]\) 内的概率质量，被映射到 \(Y\) 的 \([g(x), g(x+dx)]\) 区间。\(|h'(y)|\) 是 \(dx/dy\)------拉伸因子，补偿 \(g\) 对概率质量的压缩或膨胀。

几个重要的特例：

\begin{itemize}
\tightlist
\item
  正态变量的线性变换仍是正态：\(X \sim N(\mu,\sigma^2) \Rightarrow aX+b \sim N(a\mu+b, a^2\sigma^2)\)
\item
  正态变量的指数是对数正态：\(X \sim N(\mu,\sigma^2) \Rightarrow e^X \sim LN(\mu,\sigma^2)\)
\item
  伽马变量的尺度变换：\(X \sim Ga(\alpha,\lambda) \Rightarrow kX \sim Ga(\alpha, \lambda/k)\)
\item
  任意分布的 CDF 变换：\(F_X(X) \sim U(0,1)\)
\end{itemize}

\begin{center}\rule{0.5\linewidth}{0.5pt}\end{center}

\textbf{速查卡片}：

\begin{itemize}
\tightlist
\item
  \textbf{核心问题}：如何用数学语言描述''不知道''？用样本空间枚举可能性，用概率分配信念，用信息熵度量不确定性的程度，用分布函数完整刻画统计规律，用 KL 散度衡量分布之间的差异。
\item
  \textbf{关键概念}：概率公理 → 信息熵 → 条件概率与贝叶斯公式 → 随机变量与分布函数 → 期望与方差 → 分布演化树（伯努利树/泊松流树/正态树）→ KL 散度与交叉熵 → 随机变量函数的分布
\item
  \textbf{R 命令}：\texttt{dbinom()}, \texttt{pnorm()}, \texttt{qexp()}, \texttt{rpois()} 四类函数（\texttt{d}-密度, \texttt{p}-分布, \texttt{q}-分位, \texttt{r}-随机数）；\texttt{entropy::entropy()} 计算经验熵
\item
  \textbf{参见}：Ch02（多维扩展）, Ch03（极限定理）, Ch04（概率不等式）, Ch06（MLE = 最小化 KL 散度）, Ch08（贝叶斯统计）, Ch09（变分推断与反向 KL）, Ch18/Ch20（交叉熵损失）
\end{itemize}

\begin{center}\rule{0.5\linewidth}{0.5pt}\end{center}

\textbf{深度思考}：

\begin{enumerate}
\def\labelenumi{\arabic{enumi}.}
\item
  如果概率是观测者主观信念的度量，那么``客观概率''这个说法还有意义吗？频率方法给出的概率值，是客观的，还是``一群观测者共享了相同的先验经验''？
\item
  独立性假设在统计建模中无处不在------朴素贝叶斯分类器（Ch18）甚至假设所有特征在给定类别下条件独立。为什么一个明显错误的假设，在实践中却常常有效？
\item
  贝叶斯公式将学习描述为 \(P(\theta|\text{data}) \propto P(\text{data}|\theta) \cdot P(\theta)\)。如果先验 \(P(\theta)\) 是主观的，那么科学知识的客观性从何而来？------随着数据量的增加，后验是否会压倒先验？Ch03 和 Ch04 将提供数学层面的答案。
\item
  \(D_{KL}(P \parallel Q) \neq D_{KL}(Q \parallel P)\)：用真实分布近似模型，和用模型近似真实分布，信息损失不一样。在什么场景下你会优先选择正向 KL（MLE 的方向），什么场景下优先选择反向 KL（变分推断的方向）？提示：正向 KL 是''均值追踪''（覆盖所有模式），反向 KL 是''模式追踪''（锁定最可能的模式）。这个差异在 Ch09 的变分推断和 Ch20 的生成模型中会有具体的表现。
\end{enumerate}

\textbf{后续出场预告}：

\begin{longtable}[]{@{}ll@{}}
\toprule\noalign{}
本章概念 & 将在何处再次出现 \\
\midrule\noalign{}
\endhead
\bottomrule\noalign{}
\endlastfoot
信息熵 & Ch18 决策树分裂准则 / Ch09 变分推断的证据下界 \\
贝叶斯公式 & Ch08 贝叶斯统计（完整框架） \\
KL 散度 & Ch06 MLE / Ch09 变分推断 / Ch20 交叉熵损失 \\
交叉熵 & Ch18 分类损失 / Ch20 深度学习默认损失函数 \\
正态分布 & Ch03 CLT / Ch05 抽样 / Ch10 回归 \\
二项分布 & Ch07 符号检验 / Ch12 类别数据 \\
泊松分布 & Ch11 对数线性模型 / 泊松回归 \\
伽马分布 & Ch05 卡方分布 / Ch09 贝叶斯计算 \\
贝塔分布 & Ch08 共轭先验 \\
方差 & Ch04 切比雪夫不等式 / Ch06 估计精度 \\
\end{longtable}

\chapter{多维随机变量及其分布}\label{ux591aux7ef4ux968fux673aux53d8ux91cfux53caux5176ux5206ux5e03}

\begin{quote}
知道每一个，不等于知道全部。这正是多维概率论存在的理由。
\end{quote}

Ch01 的观测者学会了描述单个不确定量。现在他面对身高和体重两个量。他分别给出了身高的分布，又分别给出了体重的分布。他以为任务完成了。

考虑三个世界：

\begin{itemize}
\tightlist
\item
  世界 A：身高和体重同升同降。
\item
  世界 B：身高越高，体重越轻。
\item
  世界 C：两者毫无关联。
\end{itemize}

这三个世界里，身高的边际分布可以完全相同，体重的边际分布也可以完全相同。\textbf{边际分布无法区分它们。} 观测者缺少的，是描述变量之间\textbf{依赖结构}的语言。

Ch01 建立了描述单个不确定量的语言。但单个量只是世界的切片。真实数据天然是高维的------不是多个变量的简单并列，而是变量之间有一条不可见的线。

\begin{itemize}
\tightlist
\item
  一个人的健康报告包含几十个指标。肝功能指标和血脂不是独立的------它们共享同一个代谢系统的信息。
\item
  一张 \(28 \times 28\) 的灰度图像是 784 维空间中的一个点------深度学习（Ch20）本质上是在这个高维流形上做几何变换。
\item
  一个投资组合包含多只股票，它们的涨跌\textbf{几乎从不独立}------协方差矩阵刻画了系统性风险。
\end{itemize}

高维带来了 Ch01 中不存在的新问题：不仅要知道每个分量各自的分布，还要知道它们之间的\textbf{依赖结构}。独立性是特例，相关性是常态；边际分布是投影，联合分布才是全景。

\section{联合分布：高维空间上的概率云}\label{ux8054ux5408ux5206ux5e03ux9ad8ux7ef4ux7a7aux95f4ux4e0aux7684ux6982ux7387ux4e91}

\subsection{联合分布函数}\label{ux8054ux5408ux5206ux5e03ux51fdux6570}

\(n\) 维随机向量 \(\mathbf{X} = (X_1, X_2, \ldots, X_n)\) 是定义在同一个样本空间 \(\Omega\) 上的 \(n\) 个随机变量的组合。关键是``同一个样本空间''------多维变量的各个分量必须共享同一个随机性的源头。

\textbf{联合分布函数}：

\[F(x_1, \ldots, x_n) = P(X_1 \leq x_1, X_2 \leq x_2, \ldots, X_n \leq x_n)\]

这是 \(\mathbf{X}\) 落入由 \(x_1, \ldots, x_n\) 界定的左下矩形区域的概率。对于二维情况：

\[P(a < X \leq b, c < Y \leq d) = F(b,d) - F(a,d) - F(b,c) + F(a,c)\]

联合分布函数同时包含三类信息：

\begin{itemize}
\tightlist
\item
  每个分量的\textbf{边际分布}------单独看 \(X\) 或 \(Y\) 长什么样。
\item
  分量之间的\textbf{关联程度}------\(X\) 大了，\(Y\) 是否也倾向于大？
\item
  给定一个分量的\textbf{条件分布}------知道了 \(X\) 的取值，\(Y\) 的分布如何变化？
\end{itemize}

联合分布函数的结构引出一个根本问题：\textbf{边际和依赖是否可以分开建模？} Sklar 定理（1959）的回答是肯定的：任何联合分布可分解为 \(F(x_1,\ldots,x_p) = C(F_1(x_1),\ldots,F_p(x_p))\)，其中 \textbf{Copula} \(C\) 是剥离掉所有边际分布后剩下的纯依赖结构。这一解耦意味着你可以独立选择边际分布族和依赖结构族------正态边际配合 t-Copula 完全合法，而高斯 Copula 在尾部假设独立（2008 年 CDO 定价危机的技术诱因之一）。Copula 的思想在 Ch13（非参数统计）、Ch14（尾部风险）和 Ch20（生成模型）中回响。

\subsection{离散与连续的双重表达}\label{ux79bbux6563ux4e0eux8fdeux7eedux7684ux53ccux91cdux8868ux8fbe}

\textbf{离散场合}用\textbf{联合分布列}：

\[p_{ij} = P(X = x_i, Y = y_j)\]

所有 \(p_{ij} \geq 0\) 且 \(\sum_i \sum_j p_{ij} = 1\)。

\textbf{连续场合}用\textbf{联合密度函数} \(p(x,y)\)：

\[F(x,y) = \int_{-\infty}^x \int_{-\infty}^y p(u,v) \, du \, dv\]

区域 \(G\) 上的概率为 \(\iint_G p(x,y) \, dx \, dy\)。计算时需要注意两点：积分区域是 \(p(x,y)\) 的非零区域与 \(G\) 的交集；边界线（直线段）的面积为 0，是否包含不影响积分结果。

\subsection{常用多维分布}\label{ux5e38ux7528ux591aux7ef4ux5206ux5e03}

\textbf{多项分布}：将二项分布从``成功/失败''两类推广到 \(r\) 类。\(n\) 次独立试验中，第 \(k\) 类出现 \(n_k\) 次的概率：

\[P(X_1=n_1,\ldots,X_r=n_r) = \frac{n!}{n_1! \cdots n_r!} p_1^{n_1} \cdots p_r^{n_r}\]

这是类别数据分析（Ch12）的底层分布。

\textbf{二维正态分布}：多元分析中最核心的连续分布。

\[p(x,y) = \frac{1}{2\pi\sigma_1\sigma_2\sqrt{1-\rho^2}} \exp\left\{ -\frac{1}{2(1-\rho^2)} \left[ \frac{(x-\mu_1)^2}{\sigma_1^2} - 2\rho\frac{(x-\mu_1)(y-\mu_2)}{\sigma_1\sigma_2} + \frac{(y-\mu_2)^2}{\sigma_2^2} \right] \right\}\]

五个参数：\(\mu_1, \mu_2\)（中心）、\(\sigma_1^2, \sigma_2^2\)（离散度）、\(\rho\)（线性相关）。当 \(\rho = 0\) 时，联合密度分解为两个边际密度的乘积------二元正态中，\textbf{不相关等价于独立}。这是一个特殊性质，并非所有分布都成立。

\textbf{多元正态分布}是二元正态向 \(p\) 维的推广。记 \(\mathbf{X} \sim N_p(\boldsymbol{\mu}, \Sigma)\)，其密度为：

\[p(\mathbf{x}) = \frac{1}{(2\pi)^{p/2} |\Sigma|^{1/2}} \exp\left\{ -\frac{1}{2}(\mathbf{x} - \boldsymbol{\mu})^T \Sigma^{-1} (\mathbf{x} - \boldsymbol{\mu}) \right\}\]

其中 \(\boldsymbol{\mu} \in \mathbb{R}^p\) 为均值向量，\(\Sigma\) 为 \(p \times p\) 正定协方差矩阵。二次型 \((\mathbf{x} - \boldsymbol{\mu})^T \Sigma^{-1} (\mathbf{x} - \boldsymbol{\mu})\) 是马氏距离的平方------它测量 \(\mathbf{x}\) 到中心 \(\boldsymbol{\mu}\) 的距离，单位是各主轴方向上的标准差。

多元正态的三条核心性质贯穿后续几乎所有方法：

\begin{enumerate}
\def\labelenumi{\arabic{enumi}.}
\tightlist
\item
  \textbf{线性变换保持正态}：\(\mathbf{Y} = A\mathbf{X} + \mathbf{b} \sim N(A\boldsymbol{\mu} + \mathbf{b},\, A\Sigma A')\)。Ch05 的抽样分布和 Ch10 的线性回归都依赖这条性质。
\item
  \textbf{边际与条件皆正态}：任意子向量仍是多元正态；给定部分分量后，条件分布也是多元正态，且条件期望是已知分量的\textbf{线性函数}。这是高斯过程回归和高斯图模型的核心。
\item
  \textbf{不相关等价于独立}：\(\Sigma\) 非对角线全为零则各分量独立。这使得 \(\Sigma\) 的稀疏结构直接编码了条件独立------概率图模型（Ch21）的出发点。
\end{enumerate}

\section{边际分布与独立性：在边缘上观察}\label{ux8fb9ux9645ux5206ux5e03ux4e0eux72ecux7acbux6027ux5728ux8fb9ux7f18ux4e0aux89c2ux5bdf}

\subsection{从联合到边际：积分掉你不关心的维度}\label{ux4eceux8054ux5408ux5230ux8fb9ux9645ux79efux5206ux6389ux4f60ux4e0dux5173ux5fc3ux7684ux7ef4ux5ea6}

\textbf{边际分布}：从联合分布中``抹掉''其他分量，只看某一个分量的分布。

连续场合：

\[p_X(x) = \int_{-\infty}^{\infty} p(x,y) \, dy, \quad p_Y(y) = \int_{-\infty}^{\infty} p(x,y) \, dx\]

\begin{quote}
边际分布的运算是\textbf{积分}------把你不关心的维度上的概率质量全部加总起来。这和全概率公式（Ch01）是同一种思维：穷举所有可能，再求和。
\end{quote}

\subsection{独立性：信息通道的断裂}\label{ux72ecux7acbux6027ux4fe1ux606fux901aux9053ux7684ux65adux88c2-1}

\(X\) 和 \(Y\) 独立，当且仅当：

\[F(x,y) = F_X(x) F_Y(y)\]

或等价地，\(p(x,y) = p_X(x) p_Y(y)\)（连续）或 \(P(X=x_i, Y=y_j) = P(X=x_i) P(Y=y_j)\)（离散）。

独立性的多维版本是说：任意两个分量之间都不存在任何统计关联------知道任何一个，不会带来其他任何的信息。

\begin{quote}
独立是数学上的奢侈品。现实中的变量几乎总有关联。但独立假设之所以重要，不仅因为它简化计算，更因为它是\textbf{基准模型}------先假设独立，再看数据在多大程度上偏离独立，这本身就构成了判断关联强度的逻辑起点。
\end{quote}

独立性的多维版还有一个关键的变体：\textbf{条件独立}。\(X\) 和 \(Y\) 在给定 \(Z\) 的条件下独立，记为 \(X \perp\!\!\!\perp Y \mid Z\)，意味着一旦知道了 \(Z\)，\(X\) 和 \(Y\) 之间不再有信息通道------\(p(x,y|z) = p(x|z)p(y|z)\)。条件独立比无条件独立更贴近真实的科学推理：两个看起来相关的变量，常常只是因为共同依赖于第三个变量。控制了这个共同原因后，相关性消失------这正是因果推断（Ch22）的逻辑基础，也是概率图模型（Ch21）用图结构编码条件独立关系的动机。

\section{条件分布与条件期望：已知一部分，推断另一部分}\label{ux6761ux4ef6ux5206ux5e03ux4e0eux6761ux4ef6ux671fux671bux5df2ux77e5ux4e00ux90e8ux5206ux63a8ux65adux53e6ux4e00ux90e8ux5206}

\subsection{条件分布}\label{ux6761ux4ef6ux5206ux5e03}

已知 \(Y = y\)，\(X\) 的\textbf{条件密度函数}：

\[p(x|y) = \frac{p(x,y)}{p_Y(y)}\]

和 Ch01 的条件概率逻辑一致：观测到 \(Y=y\) 后，概率空间从 \(\mathbb{R}^2\) 塌缩到直线 \(Y=y\) 上，\(p(x|y)\) 是这条直线上 \(X\) 的密度。

\begin{quote}
二维正态的条件分布仍是一维正态：\(X \mid Y=y \sim N\left(\mu_1 + \rho\frac{\sigma_1}{\sigma_2}(y-\mu_2), \; \sigma_1^2(1-\rho^2)\right)\)。条件期望是 \(y\) 的线性函数，条件方差小于边际方差（因为 \(Y\) 提供了信息）。这与 Ch01 的条件概率共享同一套逻辑骨架------只是求和换成了积分。
\end{quote}

\subsection{条件期望：统计学最核心的预测工具}\label{ux6761ux4ef6ux671fux671bux7edfux8ba1ux5b66ux6700ux6838ux5fc3ux7684ux9884ux6d4bux5de5ux5177}

观测者知道了 \(Y\) 的取值。他必须更新对 \(X\) 的判断------用的不再是 \(X\) 的边际分布，而是条件分布。条件期望 \(E(X \mid Y=y)\) 是他能做出的、在均方误差意义下最优的预测。

\[E(X \mid Y=y) = \begin{cases} \sum_i x_i \, P(X=x_i \mid Y=y) & \text{离散} \\ \int_{-\infty}^{\infty} x \, p(x|y) \, dx & \text{连续} \end{cases}\]

条件期望 \(E(X \mid Y=y)\) 是 \(y\) 的函数------不同的 \(y\) 给出不同的期望值。把这个函数记为 \(g(y)\)，则 \(E(X \mid Y)\) 本身是一个\textbf{随机变量}------它是 \(Y\) 的变换，其随机性来源于 \(Y\) 的随机性。

\begin{quote}
\textbf{核心洞察}：\(E(Y \mid X)\) 是``已知 \(X\) 时，在最小均方误差（MSE）意义下对 \(Y\) 的最佳预测''。这个事实贯穿全书：

\begin{itemize}
\tightlist
\item
  Ch10 回归分析：用线性函数近似 \(E(Y \mid X)\)
\item
  Ch17 机器学习：任何监督学习的目标函数都是 \(E(Y \mid X)\) 的某种估计
\item
  Ch20 深度学习：神经网络是用梯度下降逼近高度非线性的 \(E(Y \mid X)\)
\end{itemize}
\end{quote}

\subsection{重期望公式：分而治之}\label{ux91cdux671fux671bux516cux5f0fux5206ux800cux6cbbux4e4b}

\[E(X) = E[E(X \mid Y)]\]

\textbf{推导}：由条件期望的定义，

\[E[E(X \mid Y)] = \int E(X \mid Y=y) \, p_Y(y) \, dy = \int \left[ \int x \, p(x|y) \, dx \right] p_Y(y) \, dy = \iint x \, p(x,y) \, dx \, dy = E(X)\]

先按 \(Y\) 的取值把整个空间切成若干片，在每一片上求 \(X\) 的均值，再对片与片之间加权平均------权重是每片的概率。这不仅是计算技巧，更是一种世界观：如果一个问题太大、太复杂，就找一个相关的 \(Y\) 把它拆开，逐个击破。

重期望公式最常见的应用是将方差分解为两部分：

\[\text{Var}(X) = E[\text{Var}(X \mid Y)] + \text{Var}(E[X \mid Y])\]

第一项是''组内方差的平均''（不可由 \(Y\) 解释的部分），第二项是''组间均值的方差''（可由 \(Y\) 解释的部分）。这是在 Ch10 的 \(R^2\) 分解和 Ch09 的 EM 算法 E 步中反复出现的结构。

\section{多维随机变量函数的分布}\label{ux591aux7ef4ux968fux673aux53d8ux91cfux51fdux6570ux7684ux5206ux5e03}

已知 \((X_1, \ldots, X_n)\) 的联合分布，如何求出 \(Y = g(X_1, \ldots, X_n)\) 的分布？

\subsection{离散卷积与可加性}\label{ux79bbux6563ux5377ux79efux4e0eux53efux52a0ux6027}

当 \(X\) 和 \(Y\) 独立时，\(Z = X + Y\) 的分布列由\textbf{卷积}给出：

\[P(Z = k) = \sum_i P(X = i) \, P(Y = k - i)\]

卷积的直观含义是：要得到 \(Z=k\)，枚举 \(X\) 取每一个 \(i\)，相应的 \(Y\) 必须取 \(k-i\)，然后将所有路径的概率累加。这是一种''滑动平均''------将 \(Y\) 的分布翻转后沿 \(X\) 的取值滑过，逐点相乘再求和。连续卷积（后续的小节）将这个过程写成了积分。

卷积是全书最重要的运算之一。Ch03 的特征函数将卷积转化为乘法------这是证明 CLT 的核心技巧；Ch10 的线性回归中，估计量是数据的线性卷积；Ch20 的卷积神经网络直接在像素网格上做离散卷积。

\textbf{可加性}：若同一分布族的两个独立随机变量的和仍属于该族，则称该族具有可加性。

\begin{itemize}
\tightlist
\item
  泊松分布：\(P(\lambda_1) * P(\lambda_2) = P(\lambda_1 + \lambda_2)\)
\item
  二项分布（同 \(p\)）：\(b(n,p) * b(m,p) = b(n+m,p)\)
\item
  正态分布：\(N(\mu_1,\sigma_1^2) * N(\mu_2,\sigma_2^2) = N(\mu_1+\mu_2, \sigma_1^2+\sigma_2^2)\)
\item
  伽马分布（同 \(\lambda\)）：\(Ga(\alpha_1, \lambda) * Ga(\alpha_2, \lambda) = Ga(\alpha_1+\alpha_2, \lambda)\)
\end{itemize}

具有可加性的分布族在独立和运算下\textbf{封闭}------正态分布是唯一在任意线性组合下都保持封闭的分布。可加性是 CLT 的前奏：当你不断对独立同分布变量做卷积，分布形状会逐渐演化，而正态分布是这个演化过程的唯一不动点（Ch03）。

\subsection{最大值与最小值的分布}\label{ux6700ux5927ux503cux4e0eux6700ux5c0fux503cux7684ux5206ux5e03}

\(n\) 个独立同分布变量 \(X_i \sim F\)，令 \(Y = \max\{X_1, \ldots, X_n\}\)，\(Z = \min\{X_1, \ldots, X_n\}\)：

\[F_Y(y) = [F(y)]^n, \quad F_Z(z) = 1 - [1 - F(z)]^n\]

\begin{quote}
最大值分布是极值理论（Ch13 非参数方法相关）的基础，在异常检测中有直接应用。
\end{quote}

\subsection{连续卷积公式}\label{ux8fdeux7eedux5377ux79efux516cux5f0f}

\[p_Z(z) = \int_{-\infty}^{\infty} p_X(z-y) \, p_Y(y) \, dy = \int_{-\infty}^{\infty} p_X(x) \, p_Y(z-x) \, dx\]

\subsection{变量变换法（雅可比方法）}\label{ux53d8ux91cfux53d8ux6362ux6cd5ux96c5ux53efux6bd4ux65b9ux6cd5}

设 \(\mathbf{Y} = \mathbf{g}(\mathbf{X})\) 是一一变换，其逆变换为 \(\mathbf{X} = \mathbf{h}(\mathbf{Y})\)，则：

\[p_{\mathbf{Y}}(\mathbf{y}) = p_{\mathbf{X}}(\mathbf{h}(\mathbf{y})) \cdot |\det J|\]

其中 \(J = \partial \mathbf{h}(\mathbf{y}) / \partial \mathbf{y}\) 是雅可比矩阵，\(|\det J|\) 是变换对体积的伸缩因子------从 \(dx\) 到 \(dy\) 的坐标变换，概率质量必须守恒。

\begin{quote}
多元正态变量经\textbf{线性变换}后仍为多元正态变量。这是多元正态分布最重要的性质，也是它在多元分析（Ch21）和概率图模型中不可替代的原因。
\end{quote}

两个实用的积/商公式。在金融中，收益率之积对应复利累积；在比率估计中，两个随机量的比值自然出现。令 \(U = XY\) 和 \(U = X/Y\)：

\[p_U(u) = \int_{-\infty}^{\infty} p_X\left(\frac{u}{v}\right) p_Y(v) \, \frac{1}{|v|} \, dv\]

\[p_U(u) = \int_{-\infty}^{\infty} p_X(uv) \, p_Y(v) \, |v| \, dv\]

\section{协方差与相关系数：依赖关系的量化}\label{ux534fux65b9ux5deeux4e0eux76f8ux5173ux7cfbux6570ux4f9dux8d56ux5173ux7cfbux7684ux91cfux5316}

\subsection{协方差矩阵：高维分布的结构骨架}\label{ux534fux65b9ux5deeux77e9ux9635ux9ad8ux7ef4ux5206ux5e03ux7684ux7ed3ux6784ux9aa8ux67b6}

\textbf{协方差}是两个随机变量联合偏离各自期望的程度的期望：

\[\text{Cov}(X,Y) = E[(X - EX)(Y - EY)] = E(XY) - E(X)E(Y)\]

\(\text{Cov}(X,Y) > 0\) 意味着 \(X\) 和 \(Y\) 倾向于\textbf{同向偏离}各自的均值；\(\text{Cov}(X,Y) < 0\) 意味着\textbf{反向偏离}。

\begin{quote}
若 \(X\) 与 \(Y\) 独立，则 \(\text{Cov}(X,Y) = 0\)。但\textbf{反之不然}------协方差只捕获线性关联。两个有完美的非线性关系（如 \(Y = X^2\)）的变量，协方差可以为零。
\end{quote}

观测者面对的不再是单变量的波动，而是一个 \(p\) 维数据云。他在问：这个云最伸展的方向在哪里？各方向之间如何彼此牵制？\textbf{协方差矩阵} \(\Sigma\) 回答的就是这些问题。对于 \(n\) 维随机向量 \(\mathbf{X}\)：

\[\Sigma = \text{Cov}(\mathbf{X}) = E[(\mathbf{X} - E\mathbf{X})(\mathbf{X} - E\mathbf{X})'] =
\begin{pmatrix}
\text{Var}(X_1) & \text{Cov}(X_1,X_2) & \cdots & \text{Cov}(X_1,X_n) \\
\text{Cov}(X_2,X_1) & \text{Var}(X_2) & \cdots & \text{Cov}(X_2,X_n) \\
\vdots & \vdots & \ddots & \vdots \\
\text{Cov}(X_n,X_1) & \text{Cov}(X_n,X_2) & \cdots & \text{Var}(X_n)
\end{pmatrix}\]

\(\Sigma\) 是对称的半正定矩阵，对角线是各分量的方差，非对角线是两两间的协方差。\(\Sigma\) 不仅是一张数字表------它是一个\textbf{线性算子}，刻画了数据云的形状。

\textbf{谱分解} \(\Sigma = U \Lambda U^T\) 给出了最直接的几何解读：\(U\) 的列是数据云的\textbf{主轴方向}（特征向量），\(\Lambda\) 的对角线是沿各主轴方向的方差（特征值）。PCA（Ch21）不是一套算法步骤，而是一次坐标旋转------将数据转动到它最自然伸展的坐标系下。

\textbf{马氏距离} \((x - \mu)^T \Sigma^{-1} (x - \mu)\) 以 \(\Sigma^{-1}\) 为度规，沿主轴方向对各向异性做归一化------方差大的方向上''宽容''，方差小的方向上''严格''。这是多元异常检测和 Hotelling \(T^2\) 检验（Ch07、Ch21）的几何核心。

\(\Sigma\) 的三个标量摘要在高维分析中尤为实用：\(\operatorname{tr}(\Sigma) = \sum \lambda_i\)（总方差------数据云包含的总波动能量）、\(|\Sigma| = \prod \lambda_i\)（广义方差------置信椭球的超体积正比于 \(\sqrt{|\Sigma|}\)）、\(\operatorname{rank}(\Sigma)\)（数据云有效占据的维度数）。当 \(p > n\) 时 \(\Sigma\) 必然不满秩------这是高维统计中协方差估计的结构性约束。

\textbf{Cholesky 分解} \(\Sigma = LL^T\) 提供模拟的基础：\(Z \sim N(0,I) \Rightarrow LZ + \mu \sim N(\mu, \Sigma)\)。这一''用确定性变换将噪声映射为目标分布''的思想是变分自编码器重参数化技巧（Ch20）的原型。

\(\Sigma\) 在一些关键应用中超越''汇总''的角色：

\begin{itemize}
\tightlist
\item
  \textbf{精度矩阵} \(\Sigma^{-1}\) 的 \((i,j)\) 元素为零，等价于在给定其他变量的条件下 \(X_i\) 与 \(X_j\) 条件独立------概率图模型（Ch21）用 \(\Sigma^{-1}\) 的稀疏结构编码条件独立性。
\item
  \textbf{深度学习的 BatchNorm}：对每层激活值做标准化，本质上是隐式地对协方差结构做对齐。
\end{itemize}

\begin{quote}
\textbf{高维注记}：当 \(p\) 增大时，协方差矩阵的估计面临 \(p > n\) 的困境（参数 \(O(p^2)\)，样本仅 \(n\)），且高维空间的几何发生质变------任意两点欧氏距离趋于相等（kNN 失效）、高斯测度集中在半径 \(\approx \sqrt{p}\) 的薄球壳上而非均值点。Ch17 的正则化、Ch21 的降维和稀疏 PCA 正是为了在这样的高维世界中找到可工作的简化原理。
\end{quote}

\subsection{方差运算}\label{ux65b9ux5deeux8fd0ux7b97}

\[\text{Var}\left(\sum_{i=1}^n a_i X_i\right) = \sum_{i=1}^n a_i^2 \text{Var}(X_i) + 2\sum_{i<j} a_i a_j \text{Cov}(X_i, X_j)\]

当 \(X_i\) 互不相关时，交叉项消失------方差可加。但当它们相关时，总方差可能被放大（正相关）或被对冲（负相关）。\textbf{分散投资降低风险}的统计学原理，正在于此。

\(n\) 个独立同分布变量的算术平均的方差：

\[\text{Var}\left(\frac{1}{n}\sum_{i=1}^n X_i\right) = \frac{\sigma^2}{n}\]

这 \(1/n\) 的衰减是大数定律（Ch03）的前奏------平均让方差缩小，让估计越来越准。

\subsection{相关系数：消除量纲的纯关联}\label{ux76f8ux5173ux7cfbux6570ux6d88ux9664ux91cfux7eb2ux7684ux7eafux5173ux8054}

\[\text{Corr}(X,Y) = \frac{\text{Cov}(X,Y)}{\sqrt{\text{Var}(X)\,\text{Var}(Y)}} = \frac{\text{Cov}(X,Y)}{\sigma_X \sigma_Y}\]

由施瓦兹不等式保证：\(-1 \leq \text{Corr}(X,Y) \leq 1\)。

\(\text{Corr} = \pm 1\) 的充要条件是 \(X\) 与 \(Y\) 之间\textbf{几乎处处有精确的线性关系}。\(\text{Corr} = 0\) 只意味着``没有线性关系''，不意味``没有关系''。

\begin{quote}
相关系数 \(\rho\) 是二元正态分布的第五个参数，也是整个回归分析（Ch10）的起点------\(\rho^2\) 就是 \(R^2\)（决定系数）的母体，度量''用 \(X\) 解释 \(Y\) 的变异''的解释比例。
\end{quote}

\begin{center}\rule{0.5\linewidth}{0.5pt}\end{center}

\textbf{速查卡片}：

\begin{itemize}
\tightlist
\item
  \textbf{核心问题}：如何描述多个随机变量的联合行为？联合分布给出全景（Copula 揭示边际与依赖可解耦），边际分布给出剪影，条件分布给出局部，协方差矩阵（作为线性算子）刻画数据云的形状与方向。
\item
  \textbf{关键概念}：联合分布函数（含 Copula） → 边际分布 → 条件分布与条件独立 → 条件期望与重期望公式（含方差分解） → 卷积与可加性 → 变量变换法（雅可比） → 协方差矩阵（谱分解/马氏距离/Cholesky） → 高维几何直觉
\item
  \textbf{R 命令}：\texttt{cov()}, \texttt{cor()}, \texttt{mvtnorm::rmvnorm()}, \texttt{outer()} + 密度函数构造联合分布, \texttt{mahalanobis()} 计算马氏距离
\item
  \textbf{参见}：Ch01（一元基础）, Ch03（CLT 与卷积）, Ch08（贝叶斯后验）, Ch10（回归=\(E(Y|X)\)）, Ch09（EM 与重期望/方差分解）, Ch20（VAE/生成模型）, Ch21（PCA/图模型）, Ch22（因果与条件独立）
\end{itemize}

\begin{center}\rule{0.5\linewidth}{0.5pt}\end{center}

\textbf{深度思考}：

\begin{enumerate}
\def\labelenumi{\arabic{enumi}.}
\item
  条件期望 \(E(Y \mid X)\) 是``已知 \(X\) 时对 \(Y\) 的最佳预测''。在什么意义上``最佳''？如果我们换用绝对误差（MAE）而非平方误差（MSE），最佳预测还是条件期望吗？------中位数在 MAE 下的最优性和条件期望在 MSE 下的最优性，揭示了一个深刻的原理：\textbf{损失函数定义了``好''的含义}（Ch06 和 Ch17 将反复回到这个主题）。
\item
  协方差和相关系数只度量线性关联。如果 \(Y = \sin(X)\)，相关系数可以接近零。这是相关系数的缺陷，还是线性性的力量？------在什么场景下，关注线性关联是足够的？
\item
  协方差矩阵 \(\Sigma\) 的维度是 \(p \times p\)，参数数量是 \(O(p^2)\)。当 \(p = 1000\) 时，你需要估计约 50 万个参数------而样本量可能只有几百。这个 \(p > n\) 的困境是高维统计的核心挑战。在你的直觉中，有什么方法可以应对？
\item
  高斯 Copula 在 2008 年金融危机前被广泛用于 CDO 定价------它假设各抵押贷款的违约在尾部是独立的。为什么这个假设在危机中崩溃了？如果一个数据科学家今天要为极端事件建模，他应该优先考虑什么依赖结构？提示：尾部依赖的有无，决定了''百年一遇''的事件是否倾向于扎堆出现。
\end{enumerate}

\textbf{后续出场预告}：

\begin{longtable}[]{@{}
  >{\raggedright\arraybackslash}p{(\columnwidth - 2\tabcolsep) * \real{0.3000}}
  >{\raggedright\arraybackslash}p{(\columnwidth - 2\tabcolsep) * \real{0.7000}}@{}}
\toprule\noalign{}
\begin{minipage}[b]{\linewidth}\raggedright
本章概念
\end{minipage} & \begin{minipage}[b]{\linewidth}\raggedright
将在何处再次出现
\end{minipage} \\
\midrule\noalign{}
\endhead
\bottomrule\noalign{}
\endlastfoot
联合分布 / Copula & Ch13 非参数统计 / Ch14 尾部风险 / Ch20 生成模型 \\
边际分布与独立性 & Ch07 独立性检验 / Ch13 非参数方法 \\
条件独立 & Ch21 概率图模型 / Ch22 因果推断 \\
条件期望 \(E(Y\|X)\) & Ch10 回归 / Ch17 监督学习 / Ch20 神经网络 \\
重期望 / 方差分解 & Ch10 \(R^2\) / Ch09 EM 算法 \\
卷积与可加性 & Ch03 独立和的极限（CLT） \\
协方差矩阵 \(\Sigma\) & Ch05 多元正态 / Ch21 PCA \\
谱分解 / 马氏距离 & Ch07 Hotelling \(T^2\) / Ch21 多元分析 \\
相关系数 \(\rho\) & Ch10 \(R^2\) / Ch14 自相关函数 \\
高维几何直觉 & Ch17 正则化 / Ch20 流形学习 / Ch21 降维 \\
雅可比变换 & Ch16 矩阵计算 \\
\end{longtable}

\chapter{大数定律和中心极限定理}\label{ux5927ux6570ux5b9aux5f8bux548cux4e2dux5fc3ux6781ux9650ux5b9aux7406}

\begin{quote}
频率如何升华为概率？观测如何逼近真实？大数定律和中心极限定理是统计学最深刻的数学担保------它们让``从数据中学习''成为可能，而不仅仅是``凭空猜测''。
\end{quote}

\section{频率如何合法地变成概率？}\label{ux9891ux7387ux5982ux4f55ux5408ux6cd5ux5730ux53d8ux6210ux6982ux7387}

在 Ch01 中，我们绕过了这个问题：用一种叫``概率''的东西去度量不确定性，然后对它进行公理化。但概率从哪里来？你说硬币正面的概率是 \(1/2\)------这句话到底是什么意思？

\textbf{频率学派的回答}：抛很多次之后，正面频率会稳定在 \(1/2\) 附近。\textbf{但这是需要证明的。}

1713 年，雅各布·伯努利在遗著《猜度术》中首次证明了这件事：频率收敛到概率------不是模糊的``趋近''，而是可以用数学严格陈述的极限。这个结果后来被称为\textbf{伯努利大数定律}。

从伯努利开始，概率不再是玄学。它有了一个操作性的定义------你可以通过重复试验去逼近它，验证它，修正它。大数定律是``从经验中学习''的数学许可证。

\begin{quote}
如果世界没有大数定律------频率在每次试验中随机波动、永不收敛------那么一切统计推断、一切机器学习、一切基于数据的科学，都将失去地基。
\end{quote}

中心极限定理则回答了下一个问题：不仅频率会收敛，而且\textbf{收敛过程中的波动形态}也是可知的------它总是趋向正态分布。这意味着我们不仅可以估计概率，还可以量化估计的精度。

两定理合在一起，构成了经典统计推断（Part II）的完整逻辑基础：

\[\text{LLN} \Rightarrow \text{估计的相合性} \quad \text{（数据越多越准）}\] \[\text{CLT} \Rightarrow \text{误差的分布形态} \quad \text{（我们可以量化“多准”）}\]

\section{收敛性：随机序列如何``靠近''极限}\label{ux6536ux655bux6027ux968fux673aux5e8fux5217ux5982ux4f55ux9760ux8fd1ux6781ux9650}

在陈述大数定律之前，必须先定义``随机变量序列收敛''是什么意思。不同于数列的 \(\epsilon\)-\(N\) 语言，随机变量的收敛有多种不同强度的版本。

\subsection{依概率收敛}\label{ux4f9dux6982ux7387ux6536ux655b}

\[P(|X_n - X| \geq \epsilon) \to 0 \quad (n \to \infty), \quad \forall \epsilon > 0\]

记作 \(X_n \xrightarrow{P} X\)。

\textbf{物理直觉}：随着 \(n\) 增大，\(X_n\) 与 \(X\) ``显著偏离''的可能性趋于零。但这不保证对于每一条样本路径，\(X_n\) 最终都会停在 \(X\) 附近------它允许少数``出格的''路径存在，只要这些路径的概率趋于零。

\begin{quote}
依概率收敛是统计中最常用的收敛模式，因为它足够弱------大多数估计量都满足，但又足够强------能保证``大概率情况下估计是准的''。
\end{quote}

\subsection{按分布收敛（弱收敛）}\label{ux6309ux5206ux5e03ux6536ux655bux5f31ux6536ux655b}

\[\lim_{n \to \infty} F_n(x) = F(x) \quad \text{对 } F(x) \text{ 的每一个连续点}\]

记作 \(X_n \xrightarrow{d} X\) 或 \(F_n \xrightarrow{W} F\)。

\textbf{物理直觉}：不看具体的 \(X_n\) 取值，只看它们的\textbf{分布形状}是否越来越像 \(X\) 的分布。这是最弱的收敛模式------它不关心 \(X_n\) 和 \(X\) 是否``接近''，只关心它们的概率分布是否相似。

\begin{quote}
按分布收敛是\textbf{中心极限定理}使用的收敛模式。CLT 说的是标准化和的分布收敛到正态分布------不是标准化和本身``接近''某个正态变量。
\end{quote}

三种收敛的强度关系：

\[X_n \xrightarrow{a.s.} X \quad \Rightarrow \quad X_n \xrightarrow{P} X \quad \Rightarrow \quad X_n \xrightarrow{d} X\]

（几乎必然收敛是最强的------它要求对几乎所有样本路径，\(X_n(\omega)\) 最终稳定在 \(X(\omega)\) 附近。Ch04 将详细展开。）

一个重要特例：\textbf{当极限为常数时，依概率收敛与按分布收敛等价}。这个性质在推导 MLE 的渐近正态性（Ch06）时至关重要。

\section{特征函数：分析分布收敛的数学工具}\label{ux7279ux5f81ux51fdux6570ux5206ux6790ux5206ux5e03ux6536ux655bux7684ux6570ux5b66ux5de5ux5177}

直接分析随机变量序列的分布函数收敛通常很困难。\textbf{特征函数}提供了一条巧妙的迂回路径：将分布函数``翻译''成特征函数（傅里叶变换），在特征函数空间做收敛分析，再用\textbf{连续性定理}翻译回来。

\subsection{定义与性质}\label{ux5b9aux4e49ux4e0eux6027ux8d28}

\[\phi_X(t) = E[e^{itX}] = \begin{cases} \sum_k e^{itx_k} p(x_k) & \text{离散} \\ \int_{-\infty}^{\infty} e^{itx} p(x) dx & \text{连续} \end{cases}\]

特征函数是密度函数的\textbf{傅里叶变换}（差一个符号约定）。它的核心性质：

\begin{enumerate}
\def\labelenumi{\arabic{enumi}.}
\tightlist
\item
  \textbf{唯一性}：分布函数由其特征函数唯一确定。两个随机变量同分布 \(\Leftrightarrow\) 它们的特征函数完全相同。
\item
  \textbf{独立和 → 乘积}：若 \(X\) 与 \(Y\) 独立，则 \(\phi_{X+Y}(t) = \phi_X(t) \cdot \phi_Y(t)\)。卷积变成乘法------这是特征函数最强大的计算优势。
\item
  \textbf{矩的提取}：\(E[X^k] = i^{-k} \phi^{(k)}(0)\)（只要 \(k\) 阶矩存在）。积分变成求导。
\item
  \textbf{有界性}：\(|\phi(t)| \leq \phi(0) = 1\)。
\end{enumerate}

\subsection{连续性定理}\label{ux8fdeux7eedux6027ux5b9aux7406}

\[X_n \xrightarrow{d} X \quad \Leftrightarrow \quad \phi_{X_n}(t) \to \phi_X(t), \quad \forall t\]

这一定理将``随机变量的分布收敛''转化为普通的``函数逐点收敛''------是证明中心极限定理的标准路径。

\begin{quote}
正态分布的特征函数 \(\phi(t) = e^{i\mu t - \sigma^2 t^2 / 2}\) 形式简洁，正是这种简洁性使得 CLT 的证明成为可能------独立和的乘积 \(\prod \phi_{X_i}(t/\sigma\sqrt{n})\) 在极限下自然产生 \(e^{-t^2/2}\)。
\end{quote}

\section{大数定律：不确定性的驯服}\label{ux5927ux6570ux5b9aux5f8bux4e0dux786eux5b9aux6027ux7684ux9a6fux670d}

大数定律（Law of Large Numbers）述说一个最朴素也最深刻的真理：\textbf{大量随机性的平均会呈现出稳定性}。

\subsection{为什么有多个版本？}\label{ux4e3aux4ec0ux4e48ux6709ux591aux4e2aux7248ux672c}

不同的版本对随机变量序列施加了不同的条件。条件越强，结论越容易证明；条件越弱，适用范围越广。理解 LLN 的关键不是背下所有版本，而是理解\textbf{条件递减的逻辑链条}。

\subsection{层次一：最强的条件 ------ 伯努利大数定律}\label{ux5c42ux6b21ux4e00ux6700ux5f3aux7684ux6761ux4ef6-ux4f2fux52aaux5229ux5927ux6570ux5b9aux5f8b}

设有 \(n\) 重伯努利试验，每次成功概率为 \(p\)，\(S_n\) 为成功次数。则对任意 \(\epsilon > 0\)：

\[\lim_{n \to \infty} P\left(\left|\frac{S_n}{n} - p\right| < \epsilon\right) = 1\]

\textbf{这是``频率稳定于概率''的数学表述}。它是雅各布·伯努利在 1713 年给出的，也是历史上第一个被证明的大数定律。

\begin{quote}
蒙特卡洛方法正是伯努利大数定律的直接应用：用随机点的频率 \(\frac{1}{n}\sum I(f(X_i) \leq a)\) 去近似定积分或概率。Ch09 的 MCMC 本质上是在复杂分布上做蒙特卡洛。
\end{quote}

\subsection{层次二：放松同分布 ------ 切比雪夫大数定律}\label{ux5c42ux6b21ux4e8cux653eux677eux540cux5206ux5e03-ux5207ux6bd4ux96eaux592bux5927ux6570ux5b9aux5f8b}

设 \(\{X_n\}\) 两两不相关，且每个 \(X_i\) 的方差存在且有共同上界 \(Var(X_i) \leq c\)，则：

\[\lim_{n \to \infty} P\left(\left|\frac{1}{n}\sum_{i=1}^n X_i - \frac{1}{n}\sum_{i=1}^n E(X_i)\right| < \epsilon\right) = 1\]

\textbf{物理类比}：不需要同分布------不同来源的随机因素可以有不同的分布，只要它们的波动都不超过某个上限。就像交响乐团中各种乐器音色不同，但只要各自音量可控，合奏就能稳定。伯努利大数定律是它的特例。

\subsection{层次三：放松独立 ------ 马尔可夫大数定律}\label{ux5c42ux6b21ux4e09ux653eux677eux72ecux7acb-ux9a6cux5c14ux53efux592bux5927ux6570ux5b9aux5f8b}

只需满足马尔可夫条件：

\[\frac{1}{n^2} Var\left(\sum_{i=1}^n X_i\right) \to 0\]

\textbf{物理类比}：这是最弱的条件------不要求同分布、不要求独立、甚至不要求两两不相关。唯一的要求是``集体行为的波动增长速率慢于 \(n^2\)''。切比雪夫大数定律可由马尔可夫条件推出。

\subsection{层次总结}\label{ux5c42ux6b21ux603bux7ed3}

\begin{longtable}[]{@{}
  >{\raggedright\arraybackslash}p{(\columnwidth - 6\tabcolsep) * \real{0.2500}}
  >{\raggedright\arraybackslash}p{(\columnwidth - 6\tabcolsep) * \real{0.2500}}
  >{\centering\arraybackslash}p{(\columnwidth - 6\tabcolsep) * \real{0.2500}}
  >{\raggedright\arraybackslash}p{(\columnwidth - 6\tabcolsep) * \real{0.2500}}@{}}
\toprule\noalign{}
\begin{minipage}[b]{\linewidth}\raggedright
定理
\end{minipage} & \begin{minipage}[b]{\linewidth}\raggedright
条件
\end{minipage} & \begin{minipage}[b]{\linewidth}\centering
条件强度
\end{minipage} & \begin{minipage}[b]{\linewidth}\raggedright
哲学含义
\end{minipage} \\
\midrule\noalign{}
\endhead
\bottomrule\noalign{}
\endlastfoot
伯努利 & i.i.d. 伯努利 & 最强 & 频率→概率，为蒙特卡洛奠基 \\
辛钦 & i.i.d. + 期望存在 & 强 & 科学的可复制性------只要可重复，就能逼近 \\
切比雪夫 & 两两不相关 + 方差有界 & 中 & 不需要同分布，不同来源可以混合 \\
马尔可夫 & 集体方差增长慢于 \(n^2\) & 最弱 & 几乎不要求个体行为，只约束集体 \\
\end{longtable}

\begin{quote}
辛钦大数定律只需要\textbf{期望存在}，连方差都可以不存在。它的证明使用特征函数方法------先截断（有限期望保证截断后的行为可控），再取极限（截断点趋向无穷）。这一``先截断、后取极限''的技巧是概率论证明的标准范式。
\end{quote}

\subsection{LLN → 统计推断的地基}\label{lln-ux7edfux8ba1ux63a8ux65adux7684ux5730ux57fa}

大数定律的直接推论：

\[\bar{X}_n \xrightarrow{P} \mu\]

样本均值依概率收敛到总体均值。这意味着：

\begin{itemize}
\tightlist
\item
  \textbf{参数估计的相合性}（Ch06）：当 \(n \to \infty\) 时，估计量收敛到真实参数------数据多了就会准。
\item
  \textbf{经验分布函数收敛}（Ch05、Ch13）：\(F_n(x) \xrightarrow{P} F(x)\) ------样本的经验分布逼近真实分布。格利文科定理是这一收敛的强化版（一致收敛）。
\item
  \textbf{交叉验证的合理性}（Ch17、Ch19）：训练误差收敛到泛化误差------验证集可以替代真实分布。
\end{itemize}

\section{中心极限定理：正态分布为什么无处不在}\label{ux4e2dux5fc3ux6781ux9650ux5b9aux7406ux6b63ux6001ux5206ux5e03ux4e3aux4ec0ux4e48ux65e0ux5904ux4e0dux5728}

\subsection{从误差到万物}\label{ux4eceux8befux5deeux5230ux4e07ux7269}

高斯在 1809 年研究天文观测误差时发现：观测误差近似服从正态分布。这个经验规律直到 20 世纪初才被中心极限定理赋予了严格的数学基础。

\textbf{中心极限定理（CLT）的实质}：大量微小、独立的随机因素的叠加效应，在标准化后，其分布趋向标准正态------\textbf{无论原始因素服从什么分布}。

\subsection{独立同分布下的 CLT（Lindeberg--Lévy）}\label{ux72ecux7acbux540cux5206ux5e03ux4e0bux7684-cltlindebergluxe9vy}

设 \(\{X_n\}\) i.i.d.，\(E(X_i) = \mu\)，\(Var(X_i) = \sigma^2 > 0\)。令

\[Y_n^* = \frac{\sum_{i=1}^n X_i - n\mu}{\sigma\sqrt{n}} = \frac{\bar{X}_n - \mu}{\sigma / \sqrt{n}}\]

则对任意实数 \(y\)：

\[\lim_{n \to \infty} P(Y_n^* \leq y) = \Phi(y) = \frac{1}{\sqrt{2\pi}} \int_{-\infty}^y e^{-t^2/2} \, dt\]

CLT 的惊人之处在于：\textbf{原始分布可以是二项、泊松、均匀、指数------无论什么形状，标准化和总是趋向正态。} 正态分布是分布演化的``吸引子''。

\begin{quote}
这一``普适性''源于正态分布是\textbf{稳定分布}：两个独立正态变量之和仍是正态。当大量随机变量被累加时，只有稳定分布能作为极限出现。
\end{quote}

\subsection{二项分布的正态近似（De Moivre--Laplace）}\label{ux4e8cux9879ux5206ux5e03ux7684ux6b63ux6001ux8fd1ux4f3cde-moivrelaplace}

当 \(n\) 很大时，\(b(n,p)\) 近似为 \(N(np, np(1-p))\)：

\[\frac{S_n - np}{\sqrt{np(1-p)}} \xrightarrow{d} N(0,1)\]

这是 CLT 最早被发现的版本（De Moivre, 1733）。实际应用中的经验法则：当 \(np \geq 5\) 且 \(n(1-p) \geq 5\) 时，正态近似通常是可接受的。

\begin{quote}
CLT 的应用三连：\textbf{求概率}（用正态分布近似复杂和的分布）、\textbf{求分位数}（在大样本下构造区间）、\textbf{求样本量}（反解除保证给定精度所需的最小 \(n\)）。
\end{quote}

\subsection{独立不同分布下的 CLT}\label{ux72ecux7acbux4e0dux540cux5206ux5e03ux4e0bux7684-clt}

当 \(X_i\) 的分布不完全相同时，林德伯格条件（Lindeberg Condition）保证了``没有单个变量主导总和''------每个 \(X_i\) 对方差总和的贡献都趋于零。只要这个条件满足，CLT 依然成立。

\subsection{CLT → 统计推断的精度保证}\label{clt-ux7edfux8ba1ux63a8ux65adux7684ux7cbeux5ea6ux4fddux8bc1}

CLT 的直接推论：

\[\sqrt{n}(\bar{X}_n - \mu) \xrightarrow{d} N(0, \sigma^2)\]

这意味着：

\begin{itemize}
\tightlist
\item
  \textbf{置信区间}（Ch06）：\(\bar{X}_n \pm z_{\alpha/2} \cdot \sigma/\sqrt{n}\)------我们可以量化``多准''。
\item
  \textbf{假设检验}（Ch07）：\((\bar{X}_n - \mu_0)/(\sigma/\sqrt{n})\) 在大样本下近似标准正态------我们可以控制第 I 类错误的概率。
\item
  \textbf{MLE 的渐近正态性}（Ch06）：\(\sqrt{n}(\hat{\theta}_{MLE} - \theta_0) \xrightarrow{d} N(0, 1/I(\theta_0))\)------最大似然估计不仅相合（LLN），而且其分布渐近正态（CLT）。
\end{itemize}

\subsection{LLN 与 CLT 的分工}\label{lln-ux4e0e-clt-ux7684ux5206ux5de5}

\begin{longtable}[]{@{}
  >{\raggedright\arraybackslash}p{(\columnwidth - 4\tabcolsep) * \real{0.3333}}
  >{\raggedright\arraybackslash}p{(\columnwidth - 4\tabcolsep) * \real{0.3333}}
  >{\raggedright\arraybackslash}p{(\columnwidth - 4\tabcolsep) * \real{0.3333}}@{}}
\toprule\noalign{}
\begin{minipage}[b]{\linewidth}\raggedright
\end{minipage} & \begin{minipage}[b]{\linewidth}\raggedright
大数定律
\end{minipage} & \begin{minipage}[b]{\linewidth}\raggedright
中心极限定理
\end{minipage} \\
\midrule\noalign{}
\endhead
\bottomrule\noalign{}
\endlastfoot
回答的问题 & 估计会准吗？ & 误差有多大？ \\
收敛类型 & 依概率收敛（到常数） & 按分布收敛（到正态） \\
数学表述 & \(\bar{X}_n \xrightarrow{P} \mu\) & \(\sqrt{n}(\bar{X}_n - \mu) \xrightarrow{d} N(0,\sigma^2)\) \\
统计应用 & 相合性、蒙特卡洛 & 置信区间、假设检验 \\
收敛速率 & 不提供 & \(\sim 1/\sqrt{n}\) \\
\end{longtable}

两者联合起来，构成了频率学派统计推断的完整地基：LLN 保证``方向正确''，CLT 保证``可以度量精度''。

\section{蒙特卡洛方法：大数定律的工程实现}\label{ux8499ux7279ux5361ux6d1bux65b9ux6cd5ux5927ux6570ux5b9aux5f8bux7684ux5de5ux7a0bux5b9eux73b0}

大数定律最直接的工程应用是蒙特卡洛积分：

\[\int_a^b g(x) dx \approx \frac{b-a}{n} \sum_{i=1}^n g(X_i), \quad X_i \sim U(a,b)\]

更一般地：

\[E_f[g(X)] = \int g(x) f(x) dx \approx \frac{1}{n} \sum_{i=1}^n g(X_i), \quad X_i \sim f\]

LLN 保证：当 \(n \to \infty\) 时，这个近似收敛到真实值。CLT 进一步告诉我们：误差以 \(O(1/\sqrt{n})\) 的速率衰减，且与维度无关------这是蒙特卡洛优于数值积分（对高维问题指数爆炸）的核心优势。

\begin{quote}
当 \(f\) 不是标准分布时，如何从中抽样？这是 Ch09（贝叶斯计算）的核心问题：MCMC 方法正是为了解决``从非标准后验分布中抽样''而设计的蒙特卡洛。
\end{quote}

\begin{center}\rule{0.5\linewidth}{0.5pt}\end{center}

\textbf{速查卡片}：

\begin{itemize}
\tightlist
\item
  \textbf{核心问题}：为什么大量的随机性会呈现出稳定性和正态性？这让我们能在不确定的世界中做出可靠的推断。
\item
  \textbf{关键概念}：依概率收敛 → 按分布收敛 → 特征函数（傅里叶变换） → 大数定律四层次（伯努利/辛钦/切比雪夫/马尔可夫） → 中心极限定理（独立同分布/二项正态近似/林德伯格条件） → LLN+CLT → 统计推断的合法性
\item
  \textbf{R 命令}：\texttt{mean()} 大样本验证, \texttt{replicate()} 模拟 CLT, \texttt{qqnorm()} + \texttt{qqline()} 正态性检验, \texttt{runif()} 蒙特卡洛积分
\item
  \textbf{参见}：Ch04（收敛性深化与不等式）, Ch05（抽样分布）, Ch06（MLE 渐近正态性与置信区间）, Ch09（MCMC）
\end{itemize}

\begin{center}\rule{0.5\linewidth}{0.5pt}\end{center}

\textbf{深度思考}：

\begin{enumerate}
\def\labelenumi{\arabic{enumi}.}
\item
  大数定律说 \(\bar{X}_n \xrightarrow{P} \mu\)，但它没有告诉我们\textbf{需要多大的} \(n\) 才能让误差小于某个给定的 \(\epsilon\)。Ch04 的概率不等式将回答这个问题------从``渐近保证''到``有限样本保证''的跃迁，是现代统计学习理论（Ch17）的核心关切。
\item
  中心极限定理保证了标准化和趋向正态分布，但它要求方差存在。如果真实世界的数据生成过程是\textbf{重尾分布}（方差无穷大，如柯西分布），CLT 不适用------样本均值不会收敛到总体均值。面对重尾数据，我们该用什么工具？非参数统计（Ch13）和中位数回归提供了备选方案。
\item
  辛钦大数定律只需要期望存在，而 CLT 需要方差存在。这种``条件递减带来结论递减''的模式------少一点假设，结论就少一点精度------是统计学中无处不在的权衡。你能在后续的学习中识别出更多这样的权衡吗？（提示：Ch06 的无偏性 vs.~MSE，Ch10 的偏差-方差分解，Ch17 的过拟合-欠拟合。）
\end{enumerate}

\textbf{后续出场预告}：

\begin{longtable}[]{@{}ll@{}}
\toprule\noalign{}
本章概念 & 将在何处再次出现 \\
\midrule\noalign{}
\endhead
\bottomrule\noalign{}
\endlastfoot
依概率收敛 \(\xrightarrow{P}\) & Ch04 概率不等式 / Ch06 MLE 相合性 \\
按分布收敛 \(\xrightarrow{d}\) & Ch04 收敛性对比 / Ch06 MLE 渐近正态性 \\
特征函数 & Ch04 连续性定理 / Ch05 卡方分布推导 \\
大数定律 & Ch06 相合性 / Ch09 MCMC 收敛 / Ch17 泛化界 \\
中心极限定理 & Ch06 置信区间 / Ch07 大样本检验 / Ch13 渐近方法 \\
蒙特卡洛 & Ch09 贝叶斯计算 \\
\(O(1/\sqrt{n})\) & Ch04 Hoeffding 的 \(O(e^{-n})\) / Ch06 收敛速率 \\
\end{longtable}

\chapter{概率不等式与收敛性}\label{ux6982ux7387ux4e0dux7b49ux5f0fux4e0eux6536ux655bux6027}

\begin{quote}
Ch03 告诉我们 \(n \to \infty\) 时会发生什么。但现实中的 \(n\) 永远是有限的。概率不等式的任务，是把渐近的承诺转化为\textbf{有限的、可量化的概率陈述}------不仅说``最终会准''，还要说``在给定的 \(n\) 下，有多准''。
\end{quote}

\section{为什么渐近不够？}\label{ux4e3aux4ec0ux4e48ux6e10ux8fd1ux4e0dux591f}

大数定律和中心极限定理是``终极保证''：当样本量趋于无穷时，一切都会完美。但有三个问题它们没有回答：

\begin{enumerate}
\def\labelenumi{\arabic{enumi}.}
\tightlist
\item
  \textbf{\(n\) 要多大才算``够大''？} CLT 说标准化和趋向正态，但 \(n=30\) 够吗？\(n=100\) 够吗？答案取决于原始分布的偏度和尾部------CLT 的渐近结论对有限的 \(n\) 不提供速率。
\item
  \textbf{有限 \(n\) 下的误差有多大？} 你的样本均值是基于 200 个观测计算的。它与真实均值的偏差超过 0.1 的概率是多少？LLN 和 CLT 不直接回答。
\item
  \textbf{如果 \(p\) 相对于 \(n\) 极小呢？} CLT 对二项分布 \(b(n,p)\) 的正态近似要求 \(np \geq 5\) 且 \(n(1-p) \geq 5\)。当 \(p = 0.001\) 时，你需要 \(n \geq 5000\) 才能安全使用正态近似。在 \(n\) 不够大时，怎么办？
\end{enumerate}

概率不等式正是为这些``有限样本''问题而设计的数学工具。

\section{三种收敛性的直观比较}\label{ux4e09ux79cdux6536ux655bux6027ux7684ux76f4ux89c2ux6bd4ux8f83}

Ch03 介绍了两种收敛性，这里我们加入最强的``几乎必然收敛''，形成完整的收敛性谱系。

设 \(\{X_n\}\) 为随机变量序列，\(X\) 为极限（可以是随机变量或常数）。

\subsection{按分布收敛（最弱）}\label{ux6309ux5206ux5e03ux6536ux655bux6700ux5f31}

\[X_n \xrightarrow{d} X \quad \Leftrightarrow \quad F_n(x) \to F(x) \text{ 在 } F \text{ 的连续点上}\]

\textbf{一句话物理直觉}：分布的\textbf{形状}越来越像 \(X\) 的分布。它只看概率分布，不看 \(X_n\) 和 \(X\) 之间的具体数值差异。

\begin{quote}
中心极限定理的结论就是按分布收敛：\(\sqrt{n}(\bar{X}_n - \mu)/\sigma \xrightarrow{d} N(0,1)\)。
\end{quote}

\subsection{依概率收敛（中等）}\label{ux4f9dux6982ux7387ux6536ux655bux4e2dux7b49}

\[X_n \xrightarrow{P} X \quad \Leftrightarrow \quad P(|X_n - X| \geq \epsilon) \to 0, \quad \forall \epsilon > 0\]

\textbf{一句话物理直觉}：\(X_n\) 与 \(X\) ``显著偏离''的概率趋于零------\textbf{大概率情况下，值越来越近}。

\begin{quote}
大数定律的结论是依概率收敛：\(\bar{X}_n \xrightarrow{P} \mu\)。依概率收敛 \(\Rightarrow\) 按分布收敛，但反之不成立。不过当极限为常数时，两者等价------这就是为什么我们能用 CLT 来构造 MLE 的置信区间（Ch06）。
\end{quote}

\subsection{几乎必然收敛（最强）}\label{ux51e0ux4e4eux5fc5ux7136ux6536ux655bux6700ux5f3a}

\[X_n \xrightarrow{a.s.} X \quad \Leftrightarrow \quad P\left(\lim_{n \to \infty} X_n = X\right) = 1\]

\textbf{一句话物理直觉}：除了极少数``倒霉的''样本路径外，\textbf{每一条路径最终都到达目标并且不再离开}。

\begin{quote}
几乎必然收敛要求的是逐路径的收敛------对几乎所有的 \(\omega \in \Omega\)，数列 \(X_n(\omega)\) 在通常的数列意义下收敛到 \(X(\omega)\)。而依概率收敛只要求``偏离的概率趋于零''------它甚至允许每个 \(\omega\) 对应的数列 \(X_n(\omega)\) 永远振荡，只要振荡的幅度逐渐被概率空间吞没。
\end{quote}

\subsection{一句比喻贯通}\label{ux4e00ux53e5ux6bd4ux55bbux8d2fux901a}

想象你在朝一个靶子射箭，\(n\) 是射箭次数：

\begin{itemize}
\tightlist
\item
  \textbf{按分布收敛}：落点的统计分布逐渐集中到靶心周围。不保证任一箭是否偏得离谱------但偏箭越来越``稀有''。
\item
  \textbf{依概率收敛}：脱靶（偏离超过某个固定距离）的概率越来越小。但你无法排除偶尔有一箭脱靶。
\item
  \textbf{几乎必然收敛}：从某一箭开始，\textbf{之后的每一箭都中靶}。只有极少数射手（概率零）会永远间歇性地脱靶。
\end{itemize}

\[
\text{a.s.收敛} \quad \Rightarrow \quad \text{依概率收敛} \quad \Rightarrow \quad \text{按分布收敛}
\]

\begin{quote}
为什么我们没有要求所有估计量都具有几乎必然收敛性？因为依概率收敛已经足够保证统计实践的可靠性------你可以选择足够大的 \(n\) 让偏离概率任意小。而几乎必然收敛的证明往往需要更强的条件。
\end{quote}

\section{特征函数与连续性定理（复习与深化）}\label{ux7279ux5f81ux51fdux6570ux4e0eux8fdeux7eedux6027ux5b9aux7406ux590dux4e60ux4e0eux6df1ux5316}

Ch03 中介绍了特征函数作为 CLT 的证明工具。这里补充它在收敛性分析中的核心地位。

\textbf{连续性定理}是特征函数最重要的应用：\(X_n \xrightarrow{d} X\) 当且仅当 \(\phi_{X_n}(t) \to \phi_X(t)\) 对每个 \(t\) 成立。

这个定理的价值在于\textbf{降低维度}：

\[\text{随机变量序列的分布收敛（概率空间上的问题）}\]
\[\Downarrow\]
\[\text{特征函数的逐点收敛（函数空间上的问题）}\]

一步翻译，把概率论问题变成分析学问题。所有中心极限定理的证明本质上都遵循这条路径：写出标准化和的特征函数 → 泰勒展开 → 证明逐点极限是 \(e^{-t^2/2}\) → 连续性定理给出 \(N(0,1)\) 的收敛。

\section{概率不等式链：用更强的假设换更紧的界}\label{ux6982ux7387ux4e0dux7b49ux5f0fux94feux7528ux66f4ux5f3aux7684ux5047ux8bbeux6362ux66f4ux7d27ux7684ux754c}

概率不等式有一个清晰的光谱------最左端条件最弱但界最松，最右端条件最强但界最紧。

\subsection{第一层：马尔可夫不等式（仅需一阶矩）}\label{ux7b2cux4e00ux5c42ux9a6cux5c14ux53efux592bux4e0dux7b49ux5f0fux4ec5ux9700ux4e00ux9636ux77e9}

设 \(X \geq 0\) 且 \(E[X]\) 存在，则对任意 \(a > 0\)：

\[P(X \geq a) \leq \frac{E[X]}{a}\]

\textbf{证明只用一句话}：\(a \cdot I(X \geq a) \leq X\)，两边取期望。这是所有概率不等式的\textbf{母体}------下面的每一个不等式，都是对特定对象、在附加假设下，反复施用马尔可夫不等式的产物。

\begin{quote}
马尔可夫不等式给出的界是 \(O(1/a)\)。它的美在于惊人地普适------只要是非负随机变量且期望存在即可。它的局限也在于普适：因为几乎不假设什么，所以几乎不保证什么。
\end{quote}

\subsection{第二层：切比雪夫不等式（需要二阶矩）}\label{ux7b2cux4e8cux5c42ux5207ux6bd4ux96eaux592bux4e0dux7b49ux5f0fux9700ux8981ux4e8cux9636ux77e9}

\[P(|X - \mu| \geq \epsilon) \leq \frac{Var(X)}{\epsilon^2}\]

等价形式：\(P(|X - \mu| \geq k\sigma) \leq 1/k^2\)。

\textbf{推导}：令 \(Y = (X - \mu)^2 \geq 0\)，取 \(a = \epsilon^2\)，对 \(Y\) 施用马尔可夫不等式。

\begin{quote}
切比雪夫不等式给出了 \(O(1/\epsilon^2)\) 的衰减------仍然不随样本量 \(n\) 指数衰减。对于 \(\bar{X}_n\)，\(Var(\bar{X}_n) = \sigma^2/n\)，所以 \(P(|\bar{X}_n - \mu| \geq \epsilon) \leq \sigma^2/(n\epsilon^2)\)------衰减速率是 \(1/n\)，而非指数。
\end{quote}

\subsection{第三层：Hoeffding 不等式（有界 + 独立 → 指数衰减）}\label{ux7b2cux4e09ux5c42hoeffding-ux4e0dux7b49ux5f0fux6709ux754c-ux72ecux7acb-ux6307ux6570ux8870ux51cf}

设 \(X_1, \ldots, X_n\) 独立，且 \(X_i \in [a_i, b_i]\)，则：

\[P\left(\left|\sum_{i=1}^n X_i - E\left[\sum_{i=1}^n X_i\right]\right| \geq \epsilon\right) \leq 2\exp\left(-\frac{2\epsilon^2}{\sum_{i=1}^n (b_i - a_i)^2}\right)\]

若所有 \(X_i \in [0, 1]\) 且独立同分布，简化为：

\[P(|\bar{X}_n - \mu| \geq \epsilon) \leq 2e^{-2n\epsilon^2}\]

\textbf{这是质的飞跃}：Hoeffding 界对 \(n\) 呈\textbf{指数衰减}------样本量翻倍，概率上界不是减半，而是\textbf{平方降低}。相比切比雪夫的 \(1/n\) 衰减，指数衰减意味着你可以用少得多的样本达到同样的精度保证。

\begin{quote}
Hoeffding 不等式的代价是\textbf{有界性假设}------\(X_i\) 必须在 \([a_i, b_i]\) 内。这对于很多实际问题是可以接受的（像素值在 {[}0,255{]}、概率在 {[}0,1{]}、标准化特征在有限区间内）。
\end{quote}

\subsection{第四层：Bernstein 不等式（利用方差信息收紧）}\label{ux7b2cux56dbux5c42bernstein-ux4e0dux7b49ux5f0fux5229ux7528ux65b9ux5deeux4fe1ux606fux6536ux7d27}

Hoeffding 只用了有界性，忽略了方差。当 \(Var(X_i) = \sigma^2\) 很小时，真实波动应远小于有界性暗示的范围。Bernstein 不等式弥补了这个缺陷：

\[P(|\bar{X}_n - \mu| \geq \epsilon) \leq 2\exp\left(-\frac{n\epsilon^2}{2\sigma^2 + 2M\epsilon/3}\right)\]

其中 \(|X_i| \leq M\)。

当 \(\epsilon\) 较小时（``小偏差''区域），分母的 \(2M\epsilon/3\) 可忽略，界的行为是 \(\exp(-n\epsilon^2/2\sigma^2)\)------\textbf{由方差驱动}。当 \(\epsilon\) 较大时，\(2M\epsilon/3\) 主导，退化为 Hoeffding 型。

\begin{quote}
Bernstein 不等式揭示了一个更深层的原理：\textbf{分布的尾部行为主要由低阶矩决定}。方差给出了``典型偏差''的尺度，切比雪夫和 Bernstein 分别从不同角度捕捉了这个尺度。
\end{quote}

\subsection{不等式链条总览}\label{ux4e0dux7b49ux5f0fux94feux6761ux603bux89c8}

\begin{longtable}[]{@{}
  >{\raggedright\arraybackslash}p{(\columnwidth - 6\tabcolsep) * \real{0.2500}}
  >{\raggedright\arraybackslash}p{(\columnwidth - 6\tabcolsep) * \real{0.2500}}
  >{\raggedright\arraybackslash}p{(\columnwidth - 6\tabcolsep) * \real{0.2500}}
  >{\raggedright\arraybackslash}p{(\columnwidth - 6\tabcolsep) * \real{0.2500}}@{}}
\toprule\noalign{}
\begin{minipage}[b]{\linewidth}\raggedright
不等式
\end{minipage} & \begin{minipage}[b]{\linewidth}\raggedright
条件
\end{minipage} & \begin{minipage}[b]{\linewidth}\raggedright
衰减速率
\end{minipage} & \begin{minipage}[b]{\linewidth}\raggedright
典型应用
\end{minipage} \\
\midrule\noalign{}
\endhead
\bottomrule\noalign{}
\endlastfoot
马尔可夫 & 非负，一阶矩存在 & \(1/a\) & 推导其他不等式 \\
切比雪夫 & 二阶矩存在 & \(1/\epsilon^2\)（对 \(\bar{X}_n\) 是 \(1/n\)） & 大数定律证明、置信区间 \\
Hoeffding & 有界独立变量 & \(e^{-2n\epsilon^2}\) & 分类器泛化界（Ch17-19） \\
Bernstein & 有界 + 方差已知 & \(e^{-n\epsilon^2/(2\sigma^2)}\) & 精细的集中性分析 \\
\end{longtable}

\textbf{精准度-假设强度的交易}：

\begin{verbatim}
条件弱 ───────────────────────────→ 条件强
马尔可夫        切比雪夫        Hoeffding      Bernstein
界松 ←─────────────────────────── 界紧
\end{verbatim}

每增加一个假设，就获得一个更紧的界。这在统计学中是普遍规律------\textbf{信息不是免费的}（这是 Ch08 最大熵原理和 Ch17 偏差-方差权衡的共同底层逻辑）。

\section{不等式在实际中的应用}\label{ux4e0dux7b49ux5f0fux5728ux5b9eux9645ux4e2dux7684ux5e94ux7528}

\subsection{经验分布函数的集中性（DKW 不等式）}\label{ux7ecfux9a8cux5206ux5e03ux51fdux6570ux7684ux96c6ux4e2dux6027dkw-ux4e0dux7b49ux5f0f}

设 \(F_n(x)\) 为经验分布函数，\(F(x)\) 为真实分布函数。DKW 不等式给出：

\[P\left(\sup_x |F_n(x) - F(x)| > \epsilon\right) \leq 2e^{-2n\epsilon^2}\]

这是格利文科定理（Ch05）的\textbf{非渐近}版本。它说：不仅 \(\sup_x |F_n(x) - F(x)| \to 0\)（渐近），而且对任意有限的 \(n\)，这个最大偏差超过 \(\epsilon\) 的概率以 \(e^{-2n\epsilon^2}\) 被控制。这是非参数统计（Ch13）中置信带的构造基础。

\subsection{样本均值的样本复杂度}\label{ux6837ux672cux5747ux503cux7684ux6837ux672cux590dux6742ux5ea6}

由 Hoeffding 不等式，让 \(P(|\bar{X}_n - \mu| \geq \epsilon) \leq \delta\) 所需的最小样本量为：

\[n \geq \frac{(b-a)^2}{2\epsilon^2} \log\frac{2}{\delta}\]

这个公式是实验设计和机器学习中确定训练集大小的理论基础。它告诉我们：精度要求每提高 10 倍（\(\epsilon \to \epsilon/10\)），样本量需要增加 100 倍------维度灾难在单变量问题中就已经存在。

\subsection{从概率不等式到 PAC 可学习性（伏笔）}\label{ux4eceux6982ux7387ux4e0dux7b49ux5f0fux5230-pac-ux53efux5b66ux4e60ux6027ux4f0fux7b14}

Hoeffding 不等式直接通向 PAC（Probably Approximately Correct）学习理论（Ch17）：

\begin{itemize}
\tightlist
\item
  一个分类器 \(h\) 的经验误差 \(\hat{R}(h)\) 与真实误差 \(R(h)\) 之间的差距，被 Hoeffding 不等式以概率 \(1-\delta\) 控制在 \(\sqrt{\frac{\log(2/\delta)}{2n}}\) 内。
\item
  当假设空间有限（\(|\mathcal{H}| < \infty\)）时，对所有 \(h \in \mathcal{H}\) 做联合界（union bound），得到：要让``对最优 \(h\) 的泛化误差不超 \(\epsilon\)''以概率 \(1-\delta\) 成立，样本量需要 \(n = O(\frac{\log |\mathcal{H}|}{\epsilon^2})\)。
\end{itemize}

\begin{quote}
这就是机器学习中``需要多少数据''这一根本问题的答案的雏形。它的地基就是 Ch04 的概率不等式。从伯努利大数定律（1713）到 PAC 学习理论（Valiant, 1984），跨越了 270 年的数学发展，但逻辑是一脉相承的。
\end{quote}

\subsection{概率不等式在因果推断中的角色（伏笔）}\label{ux6982ux7387ux4e0dux7b49ux5f0fux5728ux56e0ux679cux63a8ux65adux4e2dux7684ux89d2ux8272ux4f0fux7b14}

在因果推断（Ch22）中，处理组和对照组的平均处理效应差异 \(\hat{\tau} = \bar{Y}_T - \bar{Y}_C\) 是一个样本均值之差。Hoeffding 不等式或 Bernstein 不等式给出了 \(\hat{\tau}\) 偏离真实 \(\tau\) 的概率上界，从而可以在有限样本下构造置信区间------这是随机化实验分析的核心工具。

\begin{center}\rule{0.5\linewidth}{0.5pt}\end{center}

\textbf{速查卡片}：

\begin{itemize}
\tightlist
\item
  \textbf{核心问题}：在有限的 \(n\) 下，随机变量的行为有多``集中''？估计的精度如何量化地保证？
\item
  \textbf{关键概念}：三种收敛性（a.s. → 依概率 → 按分布）、特征函数的连续性定理、不等式链（马尔可夫 → 切比雪夫 → Hoeffding → Bernstein）、指数衰减速率、样本复杂度 \(n \geq O(\frac{1}{\epsilon^2} \log\frac{1}{\delta})\)
\item
  \textbf{最重要的不等式}：切比雪夫（普适但松） → Hoeffding（指数衰减但需要有界） → Bernstein（利用方差精细控制）
\item
  \textbf{参见}：Ch03（大数定律与 CLT------渐近保证）, Ch06（MLE 渐近正态性）, Ch13（DKW 与非参数置信带）, Ch17（PAC 学习与泛化界）, Ch22（因果推断中的有限样本推断）
\end{itemize}

\begin{center}\rule{0.5\linewidth}{0.5pt}\end{center}

\textbf{深度思考}：

\begin{enumerate}
\def\labelenumi{\arabic{enumi}.}
\item
  切比雪夫不等式对正态分布 \(N(\mu,\sigma^2)\) 给出 \(P(|X-\mu| \geq 2\sigma) \leq 0.25\)，但正态的实际值是 \(\approx 0.046\)------差距近 6 倍。这种``保守''是 bug 还是 feature？切比雪夫不需要知道分布是正态的------它适用于\textbf{所有}方差存在的分布。``保守''是为了普适性付出的代价。在你的实际工作中，什么时候你愿意接受这种保守，什么时候你必须寻找更紧的界？
\item
  Hoeffding 不等式的指数衰减 \(e^{-2n\epsilon^2}\) 看上去很美，但 \((b-a)^2\) 在分母中意味着：如果变量的范围很大，界会变得很松。假设你在分析股票日收益率（范围可能在 {[}-10\%, +10\%{]}），Hoeffding 界还有用吗？Bernstein 不等式能否救场？
\item
  本章的所有不等式都是\textbf{概率上界}------它们告诉你坏事情发生的概率不超过某个值，但这不意味着坏事情\textbf{不会}发生。在实际决策中，你如何选择 \(\delta\)（容忍的失败概率）？\(\delta = 0.05\) 是传统选择，但如果你的模型如果失败会导致人员伤亡（如自动驾驶），\(\delta = 0.05\) 够吗？
\end{enumerate}

\textbf{后续出场预告}：

\begin{longtable}[]{@{}ll@{}}
\toprule\noalign{}
本章概念 & 将在何处再次出现 \\
\midrule\noalign{}
\endhead
\bottomrule\noalign{}
\endlastfoot
三种收敛性的关系 & Ch05 抽样分布 / Ch06 MLE 性质 \\
连续性定理 & Ch05 卡方分布推导 \\
切比雪夫不等式 & Ch06 区间估计 \\
Hoeffding 不等式 & Ch17 PAC 学习 / Ch18 分类器分析 \\
Bernstein 不等式 & Ch19 模型选择理论 \\
DKW 不等式 & Ch13 非参数统计 / 经验过程 \\
样本复杂度公式 & Ch17-20 机器学习理论 \\
\end{longtable}

\chapter{统计量与抽样分布}\label{ux7edfux8ba1ux91cfux4e0eux62bdux6837ux5206ux5e03}

\section{数理统计前言}\label{ux6570ux7406ux7edfux8ba1ux524dux8a00}

数理统计研究的是一组受到随机性干扰的数据，再加上人们对这组数据的一些认识（即各种假设），就形成数理统计研究的出发点，然后对所考虑的问题做出统计推断或预测。

\begin{quote}
\textbf{统计量作为信息压缩}：样本 \(X_1, X_2, \ldots, X_n\) 是 \(n\) 维空间中的一个点。统计量 \(T(X)\) 的任务是将这 \(n\) 维信息压缩为低维摘要，同时尽可能保留与参数 \(\theta\) 相关的信息。当压缩过程''零损失''时------即 \(T(X)\) 包含了样本中关于 \(\theta\) 的全部信息------我们就称 \(T(X)\) 为充分统计量。这个看似抽象的概念，正是统计推断可以基于低维统计量而非全样本进行的地基。
\end{quote}

统计量的分布称为抽样分布。从正态样本出发构造的某些统计量，可以成为三大抽样分布（卡方分布、t分布、F分布）之一------这是一批精确的抽样分布。

对于一般的抽样分布，我们可以寻求统计量的近似分布。如果某个总体有均值和方差，由中心极限定理可知，即使总体不为正态分布，只要样本容量n较大时，\(\bar{X}\)的分布可用正态分布近似，称为渐近分布：
\[
\bar{X}\sim AN(\mu,\frac{\sigma^2}{n})
\]

统计量具有压缩数据的功能。我们把不损失信息的统计量视为充分统计量，可用因子分解定理方便地判断一个统计量是否为充分统计量。注意，充分统计量的一一变换仍是充分统计量。

指数型分布族概括了大量常用分布族的共同特征，其定义及性质详见Ch01（概率初步）和Ch11（广义线性模型）。

\section{总体与样本}\label{ux603bux4f53ux4e0eux6837ux672c}

1、\textbf{总体与个体}

在一个统计问题中，我们把研究对象的全体称为总体，构成总体的每个成员称为个体。

很多时候我们关注的是总体中个体某方面特征的数量指标值，因此可以说总体就是一个分布，而其数量指标就是服从这个分布的随机变量。``从总体中抽样''和``从某分布中抽样''是同一个意思。

2、\textbf{样本}

为了了解总体的分布，我们从总体中随机抽取n个个体，记其指标值为\(x_1,x_2,...,x_n\)，称其为总体的一个样本，n为样本量，样本中的个体称为样品。

一方面，样本在抽取前无法预知其数值，因此，样本是随机变量，用大写字母表示；另一方面，样本在抽取后有一组确定的观测值，因此样本又是一组数值，用小写字母表示。

样本观测值没有具体数值，只有一个范围的样本称为分组样本，是一种不完全样本。

最常用的简单随机抽样方法有两个要求：样本矩有代表性，样本要有独立性。这样得到的样本总可以看成独立同分布的iid样本。

样本的联合分布函数为:
\[
F(x_1,x_2,...,x_n)=\prod_{i=1}^nF(x_i)
\]

3、\textbf{经验分布函数}

样本观测值由大到小进行排列后得到有序样本，记为\(x_{(1)},x_{(2)},...,x_{(n)}\)。经验分布函数\(F_n(x)\)可表示为：
\[
F_n(x)=\begin{cases}
0,&x<x_{(1)}\\
\frac{k}{n},&x_{(k)}\leq x<x_{(k+1)},k=1,2,...,n-1\\
1,&x\geq x_{(n)}
\end{cases}
\]

格利文科定理：设\(x_1,x_2,...,x_n\)是取自总体分布函数\(F(x)\)的样本，\(F_n(x)\)是其经验分布函数，则有：
\[
P(sup_{-\infty<x<\infty}|F_n(x)-F(x)|\rightarrow0)=1,n\rightarrow\infty
\]

\begin{quote}
此定理表明，当n相当大时，经验分布函数是总体分布函数的一个良好的近似。经典统计学中一切统计推断都已样本为依据，其理由就在于此。
\end{quote}

4、\textbf{频数频率表}

整理样本数据常用的方法是给出频数表或频率表，其具体步骤为：

\begin{itemize}
\item
  对样本进行分组
\item
  确定每组组距
\item
  确定每组组限
\item
  统计样本数据落入每个区间的个数------频数，并列出其频数频率表
\end{itemize}

5、\textbf{样本数据的图形展示}

\begin{itemize}
\item
  直方图
\item
  茎叶图
\end{itemize}

\section{统计量及其分布}\label{ux7edfux8ba1ux91cfux53caux5176ux5206ux5e03}

1、\textbf{统计量与抽样分布}

设\(x_1,x_2,\cdots,x_n\)为取自某总体的样本，若样本函数\(T=T(x_1,x_2,\cdots,x_n)\)中不含有任何未知参数，则称\(T\)为统计量。

统计量的分布称为抽样分布。统计量的分布依赖于未知参数。

2、\textbf{样本均值及其抽样分布}

样本均值用\(\bar{x}\)表示，即
\[
\bar{x}=\frac{x_1+x_2+\cdots+x_n}{n}=\frac{1}{n}\sum_{i=1}^nx_i
\]

样本均值有如下两个性质：

\begin{itemize}
\item
  若把样本中的数据与样本均值之差称为偏差，则样本所有偏差之和为0，即\(\sum_{i=1}^n(x_i-\bar{x})=0\);
\item
  数据观测值与样本均值的偏差平方和最小，即在形如\(\sum(x_i-c)^2\)的函数中，\(\sum(x_i-\bar{x})\)最小。
\end{itemize}

样本均值的抽样分布：

\begin{itemize}
\item
  若总体的分布为\(N(\mu,\sigma^2)\)，则\(\bar{x}\)的精确分布为\(N(\mu,\frac{\sigma^2}{n})\)；
\item
  若总体的分布不是正态分布，且\(E(x)=\mu,Var(x)=\sigma^2\)存在，则n较大时\(\bar{x}\)的渐近分布为\(N(\mu,\frac{\sigma^2}{n})\)
\end{itemize}

3、\textbf{样本方差与样本标准差}

样本方差可以表示为
\[
s_n^2=\frac{1}{n}\sum_{i=1}^n(x_i-\bar{x})^2
\]

在n不大时，常用
\[
s^2=\frac{1}{n-1}\sum_{i=1}^n(x_i-\bar{x})^2
\]
作为样本方差（也叫做无偏方差）。样本标准差为\(s=\sqrt{s^2}\)。

计算样本偏差平方和的三个常用表达式：
\[
\sum(x_i-\bar{x})^2=\sum x_i^2-\frac{(\sum x_i)^2}{n}=\sum x_i^2-n\bar{x}^2
\]

样本均值的期望和方差：
\[
E(\bar{x})=\mu,Var(\bar{x})=\frac{\sigma^2}{n}
\]

样本方差的期望：
\[
E(s^2)=\frac{1}{n-1}E(\sum_{i=1}^n(x_i-\bar{x})^2)=\frac{1}{n-1}(E(x_i^2)-nE(\bar{x}^2)) =\sigma^2
\]

4、\textbf{样本矩及其函数}

样本k阶原点矩：
\[
a_k=\frac{1}{n}\sum_{i=1}^nx_i^k
\]

样本k阶中心矩：
\[
b_k=\frac{1}{n}\sum_{i=1}^n(x_i-\bar{x})^k
\]

特别的，样本一阶原点矩就是样本均值，样本二阶中心矩就是样本方差。

样本偏度反映了样本数据与对称性的偏离程度核偏离方向：
\[
\hat{\beta_s}=\frac{b_3}{b_2^{\frac{3}{2}}}
\]

样本峰度是反映总体分布密度曲线在其峰值附近的陡峭程度和尾部粗细的统计量：
\[
\hat{\beta_k}=\frac{b_4}{b_2^2}-3
\]

5、\textbf{次序统计量及其分布}

称\(x_{(i)}\)是取自总体的样本的第i个次序统计量，其取值是将样本观测值从小到大排列后得到的第i个观测值。

其中，\(x_{(1)}=min\{x_1,x_2,...,x_n\}\)，\(x_{(n)}=max\{x_1,x_2,...,x_n\}\)分别成为最小和最大次序统计量。称\((x_{(1)},x_{(2)},...,x_{(n)})\)为该样本的次序统计量。

\begin{itemize}
\item
  单个次序统计量的分布
  设总体X的密度函数为\(p(x)\)，分布函数为\(F(x)\)，则第k个次序统计量\(x_{(k)}\)的密度函数为
  \[
    p_k(x)=\frac{n!}{{(k-1)!(n-k)!}}(F(x))^{k-1}(1-F(x))^{n-k}p(x)
    \]
  \[
    p_1(x)=n\cdot (1-F(x))^{n-1}p(x),\ \ \ 
    p_n(x)=n\cdot (F(x))^{n-1}p(x)
    \]
\item
  多个次序统计量及其函数的分布
  \[
  p_{ij}(y,z)=\frac{n!}{(i-1)!(j-i-1)!(n-j)!}[F(y)]^{i-1}[F(z)-F(y)]^{j-i-1}[1-F(z)]^{n-j}p(y)p(z)
  \]
\item
  样本中位数
  \[
  m_{0.5}=\begin{cases}
  x_{\frac{n+1}{2}},&n\ is\ odd\\
  \frac{1}{2}(x_{\frac{n}{2}}+x_{\frac{n}{2}+1}),&n\ is\ even
  \end{cases}
  \]
  当数据中含有极端值时，使用中位数比使用均值更好，因为中位数不受极端值的影响。
\end{itemize}

利用次序统计量，可以得到最小观测值、最大观测值、中位数、第一4分位数、第三4分位数。利用这五个数可以绘制箱线图来大致描绘一批数据的轮廓。

\section{三大抽样分布}\label{ux4e09ux5927ux62bdux6837ux5206ux5e03}

1、\(\chi^2\)分布

设\(X_1,X_2,...,X_n\)独立同分布于标准正态分布\(N(0,1)\)，则称：
\[
\chi^2=X_1^2+X_2^2+\cdots+X_n^2\sim\chi^2(n)
\]

\(\chi^2(n)\)分布的密度函数为：
\[
f(x;n)=\frac{1}{2^{n/2}\Gamma(n/2)}x^{n/2-1}e^{-x/2},\quad x>0
\]
这实际上是 \(\text{Gamma}(n/2, 1/2)\) 分布。理解这层关系有助于记忆：\(n\) 个独立标准正态的平方和，其分布恰是形状参数为 \(n/2\)、尺度参数为 \(1/2\) 的 Gamma 分布。
\[
E(\chi^2)=n,\quad Var(\chi^2)=2n
\]

\begin{quote}
该分布的密度函数的图像是一个只取非负值的偏态分布。
\end{quote}

来自正态总体的样本，其样本方差与卡方分布的关系如下：
\[
s^2=\frac{1}{n-1}\sum_{i=1}^n(x_i-\bar{x})^2
\]
\[
\frac{(n-1)s^2}{\sigma^2}\sim \chi^2(n-1)
\]

2、\textbf{F分布}

设随机变量\(X_1\sim\chi^2(m),X_2\sim\chi^2(n)\)，则称
\[
F=\frac{X_1/m}{X_2/n}\sim F(m,n)
\]

来自正态总体的样本，其样本方差与F分布的关系如下：
\[
s_x^2=\frac{1}{m-1}\sum_{i=1}^m(x_i-\bar{x})^2,s_y^2=\frac{1}{n-1}\sum_{i=1}^n(y_i-\bar{y})^2
\]

\[
F=\frac{s_x^2/\sigma_1^2}{s_y^2/\sigma_2^2}\sim F(m-1,n-1)
\]

3、\textbf{t分布}

设随机变量\(X_1\sim N(0,1),X_2\sim\chi^2(n)\)，且 \(X_1\) 与 \(X_2\) 独立，则称
\[
t=\frac{X_1}{\sqrt{X_2/n}}\sim t(n)
\]

t分布的密度函数可通过变量变换法从 \(N(0,1)\) 和 \(\chi^2(n)\) 的联合分布推导得出：
\[
f(t;n)=\frac{\Gamma((n+1)/2)}{\sqrt{n\pi}\,\Gamma(n/2)}\left(1+\frac{t^2}{n}\right)^{-(n+1)/2},\quad -\infty<t<\infty
\]

当 \(n\to\infty\) 时，\(f(t;n)\to \frac{1}{\sqrt{2\pi}}e^{-t^2/2}\)，即 t 分布收敛于标准正态分布。
\[
E(t)=0,\quad Var(t)=\frac{n}{n-2}\ (n>2)
\]

t分布的密度函数的图像关于纵轴对称。与标准正态分布相比，t 分布的峰更低、尾部更厚------这反映了用 \(s\) 估计 \(\sigma\) 带来的额外不确定性。自由度越小，尾部越厚；自由度越大，越接近正态。

\begin{quote}
t分布的发现打破了正态分布一统天下的局面，开创了小样本统计推断的新纪元。
\end{quote}

\textbf{三大分布的关系}：三大抽样分布并非孤立存在。\(t^2(n) \sim F(1,n)\)------当一个 \(F\) 分布的分子自由度为 1 时，其开方即为 \(t\) 分布。而 \(F\) 分布本身是两个独立 \(\chi^2\) 分布之比，统一了所有正态抽样下的精确推断。这三个分布共同构成了 Ch06（区间估计）和 Ch07（假设检验）的数学基础。

\begin{Shaded}
\begin{Highlighting}[]
\CommentTok{\# 模拟验证抽样分布}
\FunctionTok{set.seed}\NormalTok{(}\DecValTok{42}\NormalTok{)}
\NormalTok{n }\OtherTok{\textless{}{-}} \DecValTok{10000}\NormalTok{; df }\OtherTok{\textless{}{-}} \DecValTok{5}

\CommentTok{\# 模拟 χ²(df) = sum of df squared standard normals}
\NormalTok{chi2\_sim }\OtherTok{\textless{}{-}} \FunctionTok{replicate}\NormalTok{(n, }\FunctionTok{sum}\NormalTok{(}\FunctionTok{rnorm}\NormalTok{(df)}\SpecialCharTok{\^{}}\DecValTok{2}\NormalTok{))}
\CommentTok{\# 模拟 t(df)}
\NormalTok{t\_sim }\OtherTok{\textless{}{-}} \FunctionTok{rt}\NormalTok{(n, df)}
\CommentTok{\# 模拟 F(3, 20)}
\NormalTok{f\_sim }\OtherTok{\textless{}{-}} \FunctionTok{rf}\NormalTok{(n, }\DecValTok{3}\NormalTok{, }\DecValTok{20}\NormalTok{)}

\FunctionTok{par}\NormalTok{(}\AttributeTok{mfrow =} \FunctionTok{c}\NormalTok{(}\DecValTok{1}\NormalTok{, }\DecValTok{3}\NormalTok{))}
\FunctionTok{hist}\NormalTok{(chi2\_sim, }\AttributeTok{breaks =} \DecValTok{50}\NormalTok{, }\AttributeTok{freq =} \ConstantTok{FALSE}\NormalTok{, }\AttributeTok{main =} \FunctionTok{expression}\NormalTok{(chi}\SpecialCharTok{\^{}}\DecValTok{2}\SpecialCharTok{*}\NormalTok{(}\DecValTok{5}\NormalTok{)),}
     \AttributeTok{xlab =} \StringTok{""}\NormalTok{, }\AttributeTok{col =} \StringTok{"lightblue"}\NormalTok{)}
\FunctionTok{curve}\NormalTok{(}\FunctionTok{dchisq}\NormalTok{(x, df), }\AttributeTok{add =} \ConstantTok{TRUE}\NormalTok{, }\AttributeTok{col =} \StringTok{"red"}\NormalTok{, }\AttributeTok{lwd =} \DecValTok{2}\NormalTok{)}

\FunctionTok{hist}\NormalTok{(t\_sim, }\AttributeTok{breaks =} \DecValTok{50}\NormalTok{, }\AttributeTok{freq =} \ConstantTok{FALSE}\NormalTok{, }\AttributeTok{main =} \FunctionTok{expression}\NormalTok{(t}\SpecialCharTok{*}\NormalTok{(}\DecValTok{5}\NormalTok{)),}
     \AttributeTok{xlab =} \StringTok{""}\NormalTok{, }\AttributeTok{col =} \StringTok{"lightblue"}\NormalTok{, }\AttributeTok{xlim =} \FunctionTok{c}\NormalTok{(}\SpecialCharTok{{-}}\DecValTok{5}\NormalTok{, }\DecValTok{5}\NormalTok{))}
\FunctionTok{curve}\NormalTok{(}\FunctionTok{dt}\NormalTok{(x, df), }\AttributeTok{add =} \ConstantTok{TRUE}\NormalTok{, }\AttributeTok{col =} \StringTok{"red"}\NormalTok{, }\AttributeTok{lwd =} \DecValTok{2}\NormalTok{)}
\FunctionTok{curve}\NormalTok{(}\FunctionTok{dnorm}\NormalTok{(x), }\AttributeTok{add =} \ConstantTok{TRUE}\NormalTok{, }\AttributeTok{col =} \StringTok{"gray"}\NormalTok{, }\AttributeTok{lty =} \DecValTok{2}\NormalTok{, }\AttributeTok{lwd =} \FloatTok{1.5}\NormalTok{)}
\FunctionTok{legend}\NormalTok{(}\StringTok{"topright"}\NormalTok{, }\FunctionTok{c}\NormalTok{(}\StringTok{"t(5)"}\NormalTok{, }\StringTok{"N(0,1)"}\NormalTok{), }\AttributeTok{col =} \FunctionTok{c}\NormalTok{(}\StringTok{"red"}\NormalTok{, }\StringTok{"gray"}\NormalTok{), }\AttributeTok{lty =} \FunctionTok{c}\NormalTok{(}\DecValTok{1}\NormalTok{, }\DecValTok{2}\NormalTok{))}

\FunctionTok{hist}\NormalTok{(f\_sim, }\AttributeTok{breaks =} \DecValTok{50}\NormalTok{, }\AttributeTok{freq =} \ConstantTok{FALSE}\NormalTok{, }\AttributeTok{main =} \FunctionTok{expression}\NormalTok{(F}\SpecialCharTok{*}\NormalTok{(}\DecValTok{3}\SpecialCharTok{*}\StringTok{","}\SpecialCharTok{*}\DecValTok{20}\NormalTok{)),}
     \AttributeTok{xlab =} \StringTok{""}\NormalTok{, }\AttributeTok{col =} \StringTok{"lightblue"}\NormalTok{, }\AttributeTok{xlim =} \FunctionTok{c}\NormalTok{(}\DecValTok{0}\NormalTok{, }\DecValTok{5}\NormalTok{))}
\FunctionTok{curve}\NormalTok{(}\FunctionTok{df}\NormalTok{(x, }\DecValTok{3}\NormalTok{, }\DecValTok{20}\NormalTok{), }\AttributeTok{add =} \ConstantTok{TRUE}\NormalTok{, }\AttributeTok{col =} \StringTok{"red"}\NormalTok{, }\AttributeTok{lwd =} \DecValTok{2}\NormalTok{)}
\end{Highlighting}
\end{Shaded}

\begin{center}\includegraphics{statistic-book_files/figure-latex/sampling-dist-sim-1} \end{center}

\begin{Shaded}
\begin{Highlighting}[]
\FunctionTok{par}\NormalTok{(}\AttributeTok{mfrow =} \FunctionTok{c}\NormalTok{(}\DecValTok{1}\NormalTok{, }\DecValTok{1}\NormalTok{))}
\end{Highlighting}
\end{Shaded}

来自正态总体的样本，有如下关系：
\[
t=\frac{\sqrt{n}(\bar{x}-\mu)}{s}\sim t(n-1)
\]

\[
\frac{(\bar{x}-\bar{y})-(\mu_1-\mu_2)}{s_w\sqrt{\frac{1}{m}+\frac{1}{n}}}\sim t(m+n-2)
\]

\section{充分统计量}\label{ux5145ux5206ux7edfux8ba1ux91cf}

1、\textbf{充分性的概念}

设\(x_1,x_2,\cdots,x_n\)是来自某个总体的样本，总体分布函数为\(F(x;\theta)\),统计量\(T=T(x_1,x_2,\cdots,x_n)\)称为θ的充分统计量，如果在给定T的取值后，\(x_1,x_2,\cdots,x_n\)的条件分布与θ无关。

\begin{quote}
直观上，我们将''样本加工不损失信息``称为``充分性''。
\end{quote}

2、\textbf{因子分解定理}

统计学中的\textbf{充分性原则}是说，在充分统计量存在的场合，任何统计推断都可以基于充分统计量进行，这可以简化统计推断的程序。

我们使用因子分解定理可以快速判断一个统计量是否充分。

设总体概率密度函数为\(f(x;\theta)\)，则\(T=T(x_1,x_2,\cdots,x_n)\)为充分统计量的充要条件是：存在两个函数\(g(t,\theta)\)，\(h(x_1,x_2,\cdots.x_n)\)使得对任意的θ和任一组观测值\(x_1,x_2,\cdots,x_n\)有：
\[
f(x_1,x_2,\cdots,x_n;\theta)=g(T(x_1,x_2,\cdots,x_n),\theta)h(x_1,x_2,\cdots,x_n)
\]
，其中\(g(t,\theta)\)是通过统计量T的取值而依赖于样本的。

若统计量T是充分统计量，存在某个函数\(h(\cdot)\)，使得T可以表示为\(t=h(s)\)，则统计量S也是充分统计量。

\subsection{指数型分布族与充分统计量}\label{ux6307ux6570ux578bux5206ux5e03ux65cfux4e0eux5145ux5206ux7edfux8ba1ux91cf}

Ch01 中介绍了指数型分布族的一般形式：
\[
f(x;\theta) = h(x)\exp\left(\eta(\theta)^T T(x) - A(\theta)\right)
\]

指数族与充分统计量之间存在天然的对应关系：将密度函数写成指数族形式后，\(T(x)\) 就是 \(\theta\) 的充分统计量。换句话说，\textbf{指数族自动为你提供了充分统计量}。

例如，正态分布 \(N(\mu, \sigma^2)\)（\(\sigma^2\) 已知）属于指数族：
\[
f(x;\mu) = \frac{1}{\sqrt{2\pi\sigma^2}}\exp\left(-\frac{x^2}{2\sigma^2}\right) \cdot \exp\left(\frac{\mu}{\sigma^2}x - \frac{\mu^2}{2\sigma^2}\right)
\]
这里 \(T(x) = x\)，因此样本均值 \(\bar{x}\)（或等价地 \(\sum x_i\)）是 \(\mu\) 的充分统计量------这与因子分解定理的结论一致。

这一联系在 Ch11（广义线性模型）中将发挥核心作用：GLM 假定响应变量来自指数族，从而自然地获得充分统计量和简洁的推断程序。

\begin{center}\rule{0.5\linewidth}{0.5pt}\end{center}

\textbf{速查卡片}：

\begin{itemize}
\tightlist
\item
  \textbf{核心问题}：如何从样本中提炼信息？样本统计量服从什么分布？
\item
  \textbf{关键概念}：总体与样本 → 经验分布函数（格利文科定理） → 样本矩与次序统计量 → 三大抽样分布（χ² / t / F） → 充分统计量与因子分解定理
\item
  \textbf{R命令}：\texttt{mean()}, \texttt{var()}, \texttt{quantile()}, \texttt{ecdf()}, \texttt{dchisq()}, \texttt{dt()}, \texttt{df()}, \texttt{qqplot()}
\item
  \textbf{参见}：Ch03（中心极限定理）, Ch06（参数估计）, Ch07（假设检验）
\end{itemize}

\chapter{参数估计}\label{ux53c2ux6570ux4f30ux8ba1}

\section{估计与优良性}\label{ux4f30ux8ba1ux4e0eux4f18ux826fux6027}

\subsection{点估计}\label{ux70b9ux4f30ux8ba1}

设\(x_1,x_2,...,x_n\)是来自总体的样本，用于估计未知参数\(\theta\)的统计量\(\hat{\theta}=\hat{\theta}(x_1,x_2,...,x_n)\)称为\(\theta\)的估计量、点估计，简称估计。

\subsection{无偏性}\label{ux65e0ux504fux6027}

设\(\hat{\theta}=\hat{\theta}(x_1,x_2,...,x_n)\)为\(\theta\)的一个估计，\(\theta\)的参数空间为\(\Theta\)，若对任意的\(\theta\in \Theta\)，有
\[
E_{\theta}(\hat{\theta})=\theta\ \ \  or \ \ E_{\theta}(\hat{\theta}-\theta)=0
\]
，则称\(\hat{\theta}\)是\(\theta\)的无偏估计。

当总体k阶矩存在时，样本k阶原点矩\(a_k\)是总日k阶原点矩\(\mu_k\)的无偏估计，如样本均值总是总体均值的无偏估计。

对k阶中心矩不一样，比如样本方差\(s_n^2\)就不是总体方差\(\sigma^2\)的无偏估计：
\[
E(s^2)=\sigma^2,\ while\ \ E(s_n^2)=E(\frac{n-1}{n}s^2)=\frac{n-1}{n}\sigma^2
\]
当样本容量趋于无穷时，有\(E(s_n^2)\rightarrow \sigma^2\)，称之为渐进无偏估计。

\begin{quote}
并非所有参数都存在无偏估计，当参数存在无偏估计时，称其为可估的，否则称其为不可估的。
\end{quote}

\subsection{有效性}\label{ux6709ux6548ux6027}

当参数可估时，其无偏估计可以有很多，常用无偏估计的方差的大小作为度量无偏估计优劣的标准，这就是有效性。

设\(\hat{\theta_1},\hat{\theta_2}\)是\(\theta\)的两个无偏估计，若对任意\(\theta\in \Theta\)有：
\[
Var(\hat{\theta_1})\leq Var(\hat{\theta_2})
\]
且至少有一个\(\theta\in \Theta\)使得上述不等号严格成立，则称\(\hat{\theta_1},\hat{\theta_2}\)中前者更有效。

\section{无偏估计}\label{ux65e0ux504fux4f30ux8ba1}

\subsection{均方误差}\label{ux5747ux65b9ux8befux5dee}

在大样本场合下，相合性和渐近正态性是评价估计好坏的两个重要标准。在样本量不是很大时，人们倾向于使用一些基于小样本的评价标准，此时对无偏估计常使用方差，对有偏估计常使用均方误差。

均方误差可以定义如下：
\[
MSE(\hat{\theta})=E(\hat{\theta}-\theta)^2=Var(\hat{\theta})+(E(\hat{\theta})-\theta)^2
\]

可以看出，均方误差由点估计的方差与偏差\(|E(\hat{\theta}-\theta)|^2\)两部分组成。

当\(\hat{\theta}\)不是\(\theta\)的无偏估计时，就要看其均方误差而不只是方差。在均方误差的含义下，有些有偏估计优于无偏估计。

我们先给出一致最小均方误差估计的定义，并断言：一致最小误差估计一般都不存在。
\[
MSE_{\theta}(\hat{\theta})\leq MSE_{\theta}(\tilde{\theta})
\]

\subsection{一致最小方差无偏估计}\label{ux4e00ux81f4ux6700ux5c0fux65b9ux5deeux65e0ux504fux4f30ux8ba1}

当\(\hat{\theta}\)是\(\theta\)的无偏估计时，均方误差就简化为估计的方差。此时一致最小均方误差估计就是一致最小方差无偏估计（UMVUE）：
\[
Var_{\theta}(\hat{\theta})\leq Var_{\theta}(\tilde{\theta})
\]

\subsection{充分性原则}\label{ux5145ux5206ux6027ux539fux5219}

如果充分统计量和UMVUE存在，则UMVUE一定可以表示为充分统计量的函数。

设\(T=T(x_1,x_2,...,x_n)\)是总体的参数\(\theta\)的充分统计量，则对\(\theta\)的任一无偏估计\(\hat{\theta}=\hat{\theta}(x_1,x_2,...,x_n)\)，令\(\tilde{\theta}=E(\hat{\theta}|T)\)，则\(\tilde{\theta}\)也是无偏估计，而且
\[
Var(\tilde{\theta})\leq Var(\hat{\theta})
\]

这说明，如果无偏估计不是充分统计量的函数，将之对充分统计量求条件期望可以得到一个新的无偏估计，该估计的方差比原来的估计的方差要小，从而降低了无偏估计的方差。换言之，考虑\(\theta\)的估计问题只需要在基于充分统计量的函数中进行即可，这对所有的统计推断问题都是成立的。这就是充分性原则。

\section{C-R不等式}\label{c-rux4e0dux7b49ux5f0f}

\subsection{Fisher信息量}\label{fisherux4fe1ux606fux91cf}

费希尔信息量是统计学中的一个基本概念，很多统计结果如MLE的渐近方差、无偏估计的方差下界都与费希尔信息量有关。

计算公式为：
\[
I(\theta)=Var(\frac{\partial}{\partial\theta}ln\ p(x;\theta)) = E[\frac{\partial}{\partial\theta}ln\ p(x;\theta)]^2 = -E\frac{\partial^2ln\ p(x;\theta)}{\partial\theta^2}
\]

\begin{quote}
费希尔信息量越大，可以认为总体分布中包含未知参数的信息越多。
\end{quote}

\subsection{C-R不等式}\label{c-rux4e0dux7b49ux5f0f-1}

设\(T=T(x_1,x_2,...,x_n)\)是\(g(\theta)\)的任一个无偏估计，\(g'(\theta)=\frac{\partial g(\theta)}{\partial\theta}\)存在，则有C-R不等式：
\[
Var(T)\geq \frac{[g'(\theta)]^2}{nI(\theta)}
\]

特别的，对\(\theta\)的无偏估计\(\hat{\theta}\)，有
\[
Var(\hat{\theta})\geq \frac{1}{nI(\theta)}
\]

如果等式成立，则称无偏估计\(T=T(x_1,x_2,...,x_n)\)是\(g(\theta)\)的有效估计，有效估计一定是UMVUE。

能达到C-R下界的无偏估计并不多。大多数场合无偏估计都达不到其C-R下界。

\subsection{有效无偏估计}\label{ux6709ux6548ux65e0ux504fux4f30ux8ba1}

对正则条件下的分布族，有些参数无偏估计的C-R下界可以达到，而有些则不能。通常将达到C-R下界的无偏估计称为有效无偏估计，而把无偏估计的方差与其C-R下界之比的倒数称为该估计的效。

有效无偏估计只在指数型分布族场合才可谈及，而并非指数族中任一可估参数都是有效无偏估计。尽管如此，信息不等式在大样本场合给出了一个很好的方差下界，这个下界是可以逼近的。

\section{估计量的渐近性质}\label{ux4f30ux8ba1ux91cfux7684ux6e10ux8fd1ux6027ux8d28}

\subsection{相合性}\label{ux76f8ux5408ux6027}

如果我们有足够的观测值，根据格利文科定理，随着样本量的不断增大，经验分布函数逼近真实分布函数，因此完全可以要求估计量随着样本量的不断增大而逼近参数真值，这就是相合性。

设\(\theta\in\Theta\)为未知参数，\(\hat{\theta_n}=\hat{\theta_n}(x_1,x_2,...,x_n)\)是\(\theta\)的一个估计量，n是样本容量。若对任意一个\(\epsilon>0\)，有
\[
lim_{n\rightarrow \infty}P(|\hat{\theta_n}-\theta|\geq\epsilon)=0
\]
则称\(\hat{\theta}\)是\(\theta\)的相合估计。

若把依赖于样本量n的估计量\(\hat{\theta_n}\)看作一个随机变量序列，相合性就是\(\hat{\theta_n}\)依概率收敛于\(\theta\)。所以证明估计的相合性可应用依概率收敛的性质以及各种大数定律。

判断相合性时有如下两个判据可以快速判断：
\[
lim_{n\rightarrow \infty}E(\hat{\theta_n})=\theta,\ \ \ \ lim_{n\rightarrow\infty}Var(\hat{\theta_n})=0
\]

另外，相合估计的函数也是相合估计。

最后需要说明，相合性只反映了当\(n\rightarrow \infty\)时估计量的性质，而对任意有限的n，相合性是没有意义的。相合性本身不能说明为了使\(\hat{\theta_n}\)达到一定精度n必须至少为多少，事实上，相合估计可以不止一个，它们之间是有差异的。这种差异往往可由估计量的渐近分布的渐近方差反映出来。最常用的渐近分布是正态分布。

\subsection{渐进正态性}\label{ux6e10ux8fdbux6b63ux6001ux6027}

设\(\hat{\theta_n}=\hat{\theta}(X_1,X_2,...,X_n)\)是\(\theta\)的估计，如果存在\(\sigma_n^2(\theta)\)，满足
\[
\frac{\hat{\theta_n}-\theta}{\sigma_n(\theta)}\xrightarrow{L}N(0,1)
\]
则称\(\hat{\theta_n}\)是\(\theta\)的渐近正态估计，\(\sigma_n^2(\theta)\)称为\(\hat{\theta}\)的渐近方差，记为\(\hat{\theta_n}\sim AN(\theta,\sigma_n^2(\theta))\)

最大似然估计通常有渐近正态性的结论：设\(x_1,x_2,...,x_n\)来自某特定总体（可求导、方差存在）的样本，则存在未知参数\(\theta\)的最大似然估计\(\hat{\theta_n}=\hat{\theta_n}(x_1,x_2,...,x_n)\)，且\(\hat{\theta_n}\)具有相合性和渐近正态性，
\[\hat{\theta_n}\sim AN(\theta,\frac{1}{nI(\theta)})\]
其中，\(I(\theta)\)称为费希尔信息量。

\section{矩估计与替换方法}\label{ux77e9ux4f30ux8ba1ux4e0eux66ffux6362ux65b9ux6cd5}

\subsection{矩估计}\label{ux77e9ux4f30ux8ba1}

矩估计的基本思路是用样本矩及其函数估计相应的总体矩及其函数：

\begin{enumerate}
\def\labelenumi{\arabic{enumi}.}
\item
  设\(X_1,X_2,\cdots,X_n\)是来自某总体的一个样本。用\(\mu_r\)记总体的r阶原点矩，用\(m_r\) 记由样本得到的r阶样本原点矩。即
  \[
  \mu_r=E(X_i^r);m_r=\frac{1}{n}\sum_{i=1}^nX_i^r
  \]
\item
  如果要估计的参数θ可以表示为总体前k阶矩的函数：
  \[
  \theta=g(\mu_1,\mu_2,\cdots,\mu_k)
  \]
\item
  则我们可以用
  \[
  \hat{\theta}(X)=g(m_1,m_2,\cdots,m_k)
  \]
  估计θ，\(\hat{\theta}(X)\) 即称为θ的矩估计。
\item
  在矩估计不唯一时，有两个原则：涉及的矩的阶数尽可能小，所用估计最好是充分统计量的函数。
\end{enumerate}

\subsection{矩估计的特点}\label{ux77e9ux4f30ux8ba1ux7684ux7279ux70b9}

\begin{enumerate}
\def\labelenumi{\arabic{enumi}.}
\item
  矩估计基于经验分布函数，而经验分布函数逼近真实分布的前提条件是样本容量较大。因此矩方法是以大样本为应用对象的。
\item
  矩方法没有用到总体分布的任何信息，本质上讲它是一种非参数方法，对已知的总体分布，矩估计不一定是一个好的估计。
\item
  在诸如二项分布、泊松分布、正态分布等常见分布的均值，其矩估计具有许多优良的性质，它是UMVUE。
\item
  在大样本下，矩估计是相合估计。
\end{enumerate}

\section{最大似然估计}\label{ux6700ux5927ux4f3cux7136ux4f30ux8ba1}

\subsection{MLE}\label{mle}

在概率统计中，密度函数扮演了重要的角色。当\(\theta\)已知时，\(p(x;\theta)\) 显示密度怎样随x变化；而当样本给定后，可以考虑对不同的\(\theta\)，密度如何变化，它反映了对x的解释能力，这便是\textbf{似然}。

样本的似然函数：
\[
L(\theta)=L(\theta;x_1,x_2,...,x_n)=p(x_1;\theta)p(x_2;\theta)\cdots p(x_n;\theta)
\]
参数的极大似然估计要满足：
\[
L(\hat{\theta})=max_{\theta\in\Theta}L(\theta)
\]

由于密度函数大多具有指数函数形式，采用似然函数的对数通常更为简便。

当极大似然估计（MLE）存在时，寻找MLE的常用方法是求似然方程的解：

\[
\frac{\partial l(\theta;x)}{\partial \theta_i}=0,i=1,2,...,k
\]

\begin{quote}
从信息论视角看，MLE 等价于最小化经验分布与模型分布之间的 KL 散度------你是在用数据''压缩''对真实分布的无知。这一点在 Ch01 §分布之间的距离中做了完整的推导。
\end{quote}

\subsection{MLE的性质}\label{mleux7684ux6027ux8d28}

MLE常是充分统计量的函数。

\[
l(\theta;x)=lnp(x;\theta)=ln(g(\theta,T(x))h(x))=ln g(\theta,T(x))+ln h(x)
\]

MLE通常具有渐近正态性。

\begin{center}\rule{0.5\linewidth}{0.5pt}\end{center}

\textbf{速查卡片}：

\begin{itemize}
\tightlist
\item
  \textbf{核心问题}：如何用样本数据估计未知参数？怎样的估计是''好''的？
\item
  \textbf{关键概念}：无偏性 → 有效性 → 相合性 → 渐近正态性 → 矩估计 → 最大似然估计（MLE） → Fisher信息量与C-R下界 → 充分性原则 → UMVUE
\item
  \textbf{R命令}：\texttt{fitdistrplus::fitdist()}, \texttt{MASS::fitdistr()}, \texttt{optim()} 自定义似然函数
\item
  \textbf{参见}：Ch05（充分统计量）, Ch07（假设检验）, Ch08（贝叶斯估计）
\end{itemize}

\chapter{参数假设检验和方差分析}\label{ux53c2ux6570ux5047ux8bbeux68c0ux9a8cux548cux65b9ux5deeux5206ux6790}

\section{假设检验的基本思想和概念}\label{ux5047ux8bbeux68c0ux9a8cux7684ux57faux672cux601dux60f3ux548cux6982ux5ff5}

\subsection{假设}\label{ux5047ux8bbe}

设样本空间中样本点x是我们的观测值，除了参数估计，统计推断中另一类重要问题是，检验观测值x与给定的假设是否存在矛盾。这类问题称为假设检验问题。

在一个假设检验问题中涉及两个假设，所要检验的假设称为原假设，记为\(H_0\)，与\(H_0\)不相容的假设称为备择假设，记为\(H_1\)。参数假设检验问题的原假设和备择假设可记为：

\[
H_0:\theta\in\Theta_0\ \ vs\ \ H_1:\theta\in\Theta_1
\]

\subsection{检验、拒绝域与检验统计量}\label{ux68c0ux9a8cux62d2ux7eddux57dfux4e0eux68c0ux9a8cux7edfux8ba1ux91cf}

在检验问题\((H_0,H_1)\)中，所谓检验法则就是设法把样本空间划分为互不相交的可测集：
\[
\mathcal{X}=W+\bar{W}
\] 并规定：
当观测值\(x\in W\)时，就拒绝原假设\(H_0\)，认为备择假设\(H_1\)成立；
当观测值\(x\notin W\)时，就不拒绝原假设\(H_0\)。这里\(W\)称为检验的拒绝域。

这样，选定了检验法，就是确定了拒绝域；选定了拒绝域，也就确定了检验法。

为了确定拒绝域，往往首先由问题的直观背景出发，寻找一个统计量，使得在原假设\(H_0\)成立时和在备择假设\(H_1\)成立时，该统计量的值有差异。
从而使我们能够根据这个统计量的值的大小选定拒绝域。我们\textbf{称这个能从样本空间中划分出拒绝域的统计量称为检验统计量}。

\subsection{两类错误}\label{ux4e24ux7c7bux9519ux8bef}

在进行检验时，由于样本的随机性，我们可能做出正确的判断，也可能做出错误的判断。为了对检验法的好换给出一个合理的评选系统，需要考察一个检验法可能犯错误的概率。

当原假设\(H_0\)成立时，样本观测值却落在拒绝域\(W\)中，从而拒绝了原假设。这种错误称为\textbf{第一类错误}。犯第一类错误的概率是：
\[
\alpha(\theta)=P_{\theta}(X\in W),\ \ \  \theta\in\Theta_0
\]

当原假设\(H_0\)不成立时，样本观测值却没有落在拒绝域\(W\)中，从而没有拒绝原假设，这种错误称为\textbf{第二类错误}，犯第二类错误的概率为：
\[
\beta(\theta)=P_{\theta}(X\notin W)=1-P_{\theta}(X\in W),\ \ \ \ \theta\in\Theta_1
\]

\begin{quote}
\textbf{两种范式的视角差异}：频率学派控制 \(P(\text{数据} \mid H_0)\)------``如果原假设为真，这种数据出现的概率是多少？''贝叶斯学派（Ch08）则直接计算 \(P(H_0 \mid \text{数据})\)------``给定这些数据，原假设为真的概率是多少？''前者的答案是关于数据的，后者的答案是关于假设本身的。两种回答各有其适用场景，但问的是不同的问题。
\end{quote}

\subsection{势函数}\label{ux52bfux51fdux6570}

称样本观测值落在拒绝域的概率为检验的势函数，记为 \[
g(\theta)=P_{\theta}(X\in W),\theta\in\Theta\\=\begin{cases}
\alpha(\theta),&\theta\in\Theta_0\\
1-\beta(\theta),& \theta\in\Theta_1
\end{cases}
\]

\subsection{检验的水平}\label{ux68c0ux9a8cux7684ux6c34ux5e73}

一般情况下，当样本容量n固定时，不可能使得犯这两类错误的概率都减小。基于此，我们使犯第一类错误的概率\(\alpha(\theta)\)的概率限制在某一个范围内，然后寻找使犯第二类错误的概率尽可能小的检验。

这样，寻找一个好的检验，就是对选定的一个较小的数\(\alpha(0<\alpha<1)\)，在满足
\[
g(\theta)=P_{\theta}(X\in W)\leq\alpha,\theta\in\Theta_0
\]
的检验中，寻找这样的检验，使得在\(\theta\in\Theta_1\)时，\(g(\theta)\)尽可能大（样本观测值尽可能落在拒绝域中）。

在\(\theta\in\Theta_0\)时，\(g(\theta)\leq\alpha\)的检验称为显著性水平为\(\alpha\)的检验。\(\alpha\)是事先选定的，由于\(\alpha\)的大小反映了检验犯第一类错误的概率的大小，所以常取一个较小的数。比如说，若根据以往经验，非常相信原假设是真的，那么可以把\(\alpha\)取小一点，这样第一类错误的犯错概率会减小。

由于检验犯第二类错误的概率没有限制，所以在备择假设为真时，样本观测值没有落入拒绝域可能不是一个小概率事件，因此接受原假设的理由可能是不充足的，所以在这种情况下，\textbf{较为科学的说法是不能拒绝原假设}。

\subsection{假设检验的一般步骤}\label{ux5047ux8bbeux68c0ux9a8cux7684ux4e00ux822cux6b65ux9aa4}

\begin{itemize}
\item
  根据实际问题，建立统计假设\(H_0\ vs\ H_1\)
\item
  选择一个合适的检验统计量\(T(X)\)，使当\(H_0\)成立时（或\(H_0\)中某个具体参数下），\(T\)的分布完全已知，并根据\(H_0\)和\(H_1\)的特点，确定拒绝域\(W\)的形式。
\item
  选择合适的显著性水平\(\alpha\)，确定具体的拒绝域\(W\).
\item
  由样本观测值\(x_1,x_2,...,x_n\)，计算检验统计量的\(T(x_1,x_2,...,x_n)\)，由\(T(x_1,x_2,...,x_n)\)是否属于\(W\)，做出最终判断。
\end{itemize}

\subsection{检验的p值}\label{ux68c0ux9a8cux7684pux503c}

在一个假设检验问题中，利用样本观测值能够做出拒绝原假设的最小显著性水平称为检验的p值。

实际中，p很小时即可拒绝原假设，p很大时即可接收原假设。只有当p与\(\alpha\)接近时才需要比较，以减少争论。

\section{正态总体参数假设检验}\label{ux6b63ux6001ux603bux4f53ux53c2ux6570ux5047ux8bbeux68c0ux9a8c}

\subsection{关于正态总体均值的检验}\label{ux5173ux4e8eux6b63ux6001ux603bux4f53ux5747ux503cux7684ux68c0ux9a8c}

\textbf{1、}\(\sigma=\sigma_0\)\textbf{已知时的u检验}

检验统计量： \[
u=\frac{\bar{x}-\mu_0}{\sigma_0/\sqrt{n}}\sim N(0,1)
\]

双侧检验的拒绝域可以写为：

\[
W=\{|u|\geq u_{1-\alpha/2}\}
\]

\textbf{2、}\(\sigma\)\textbf{未知时的t检验}

检验统计量： \[
t=\frac{\bar{x}-\mu}{s/\sqrt{n}}\sim t(n-1),\text{s为样本标准差}
\]

双侧检验的拒绝域可以写为：

\[
W=\{|t|\geq t_{1-\alpha/2}(n-1)\}
\]

\textbf{3、}\(\sigma_1,\sigma_2\)\textbf{已知时的两样本}\(u\)\textbf{检验}

假设检验问题：

\[
H_0:\mu_1-\mu_2\leq0\ or\ \mu_1-\mu_2\geq0 \ or \ \mu_1=\mu_2
\]

检验统计量（简化形式，\(\mu_1=\mu_2\)）：

\[
u=\frac{\bar{x}-\bar{y}}{\sqrt{\frac{\sigma_1^2}{m}+\frac{\sigma_2^2}{n}}}\sim N(0,1)
\]

双侧检验的拒绝域：

\[
W=\{|u|\geq u_{1-\alpha/2}\}
\]

\textbf{4、}\(\sigma_1=\sigma_2=\sigma\)\textbf{但未知时的两样本}\(t\)\textbf{检验}

记： \[
s_w^2=\frac{1}{m+n-2}[\sum_{i=1}^m(x_i-\bar{x})+\sum_{i=1}^n(y_i-\bar{y})]
\] 检验统计量： \[
t=\frac{(\bar{x}-\bar{y})-(\mu_1-\mu_2)}{s_w\sqrt{\frac{1}{m}+\frac{1}{n}}}\sim t(m+n-2)
\]

双边检验的拒绝域： \[
W=\{|t|\geq t_{1-\alpha/2}(m+n-2)\}
\]

\textbf{5、大样本下的u检验和小样本下的近似t检验}

当总体不是正态分布时，对应大样本和小样本下的检验统计量为：

\[
u=\frac{\bar{x}-\bar{y}}{\sqrt{\frac{s_x^2}{m}+\frac{s_y^2}{n}}}\sim N(0,1)
\]

\[
t=\frac{\bar{x}-\bar{y}}{\sqrt{\frac{s_x^2}{m}+\frac{s_y^2}{n}}}\sim t(l)
\]

其中t分布的自由度由公式给出： \[
l=\frac{(\frac{s_x^2}{m}+\frac{s_y^2}{n})^2}{\frac{s_x^4}{m^2(m-1)}+\frac{s_y^4}{n^2(n-1)}}
\]

\textbf{6、成对数据检验}

如果要比较两个总体均值时，数据是成对出现的，此时采用两样本的t检验并不准确，应该将成对数据做差，转换为单样本的t检验问题。

\subsection{关于正态总体方差的检验}\label{ux5173ux4e8eux6b63ux6001ux603bux4f53ux65b9ux5deeux7684ux68c0ux9a8c}

\textbf{1、单个正态总体方差的}\(\chi^2\)\textbf{检验}

检验统计量（当\(\sigma^2_0=\sigma^2\)）： \[
\chi^2=\frac{(n-1)s^2}{\sigma_0^2}\sim \chi^2(n-1)
\]

双侧检验的拒绝域： \[
W=\{\chi^2\leq\chi^2_{\alpha/2}(n-1)\text{或}\chi^2\geq\chi^2_{1-\alpha/2}(n-1)\}
\]

\textbf{2、两个正态总体方差比的F检验}

检验统计量（当\(\sigma^2_1=\sigma_2^2\)）： \[
F=\frac{s_x^2}{s_y^2}\sim F(m-1,n-1)
\]

双边检验的拒绝域：

\[
W=\{F\leq F_{\alpha/2}\ or \ F\geq F_{1-\alpha/2} \}
\]

\subsection{指数分布参数的假设检验}\label{ux6307ux6570ux5206ux5e03ux53c2ux6570ux7684ux5047ux8bbeux68c0ux9a8c}

记： \[
X\sim Exp(\frac{1}{\theta}),\theta\text{为均值}
\]

检验问题： \[
H_0:\theta\leq \theta_0 \ \ \ \ H_1:\theta>\theta_0
\]

检验统计量： \[
\chi^2 = \frac{2n\bar{x}}{\theta_0}\sim \chi^2(2n)
\]

单侧检验拒绝域： \[
W=\{\chi^2\geq \chi^2_{1-\alpha}(2n) \}
\]

\subsection{大样本检验}\label{ux5927ux6837ux672cux68c0ux9a8c}

设\(x_1,x_2,...,x_n\)是来自总体分布\(F(x;\theta)\)的样本，均值为\(\theta\)，方差为\(\theta\)的函数，记为\(\sigma^2(\theta)\)。

利用中心极限定理，构建检验统计量： \[
u=\frac{\sqrt{n}(\bar{x}-\theta_0)}{\sqrt{\sigma^2(\hat{\theta})}}\sim AN(0,1)
\]

其中\(\hat{\theta}\)为\(\theta\)的极大似然估计。

拒绝域为：

\[
W=\{u \geq u_{1-\alpha} \}
\]

\section{方差分析}\label{ux65b9ux5deeux5206ux6790}

\subsection{方差分析的本质及误差分解}\label{ux65b9ux5deeux5206ux6790ux7684ux672cux8d28ux53caux8befux5deeux5206ux89e3}

方差分析由费希尔在进行农业试验设计时为解释试验数据首次提出。本质上，方差分析研究分类变量对数值变量的影响，其含义是通过检验总体的均值是否相等来判断分类变量对数值变量是否有显著影响。

总误差（SST）：反映全部数据误差大小的平方和：

\[
SST = \sum_{i=1}^4\sum_{j=1}^{n_i}(x_{ij}-\bar{x})^2
\]

组内误差（SSE）：反映每个样本内部误差大小的平方和

\[
SSE=\sum_{i=1}^4\sum_{j=1}^{n_i}(x_{ij}-\bar{x_i})^2
\]

组间误差（SSA）：反映样本间误差大小的平方和

\[
SSA=\sum_{i=1}^4\sum_{j=1}^{n_i}(\bar{x_i}-\bar{x})^2
\]

误差分解： \[
SST = SSE + SSA
\]

\subsection{单因子方差分析}\label{ux5355ux56e0ux5b50ux65b9ux5deeux5206ux6790}

在单因子试验中，记因子为A，设其有r个水平，记为\(A_1,A_2,...,A_r\)，在每一水平下考察的指标可以看成个总体，共有r个总体。

假定：

\begin{itemize}
\item
  每一总体均为正态总体，记为\(N(\mu,\sigma_i^2),i=1,2,...,r\)
\item
  各总体的方差相同，记为\(\sigma_1^2=\sigma_2^2=\cdots=\sigma_r^2\)
\item
  从每一总体中抽取的样本是相互独立的，即所有的试验结果\(y_{ij}\)都相互独立。
\end{itemize}

检验内容是比较各水平下的均值是否相同，即： \[
H_0:\mu_1=\mu_2=\cdots=\mu_r
\]

如果\(H_0\)成立，称因子A的r个水平间没有显著差异，简称因子A不显著。反之称因子A显著。

统计模型：表示处理i的第j个观测值\(x_{ij}\)等于不考虑处理影响时的总均值\(\mu\)加上处理为i时的附加效应\(\alpha_i\)，及其处理i的第j个观测值的随机误差\(\epsilon_{ij}\)
\[
x_{ij}=\mu+\alpha_i+\epsilon_{ij},\sum_{i=1}^kn_i\alpha_i=0
\] 可以据此做出第二种假设检验： \[
H_0:\alpha_1=\alpha_2=\cdots=\alpha_k=0
\]

构造检验统计量（\(H_0\)为真时）：

\[
F=\frac{SSA/(k-1)}{SSE/(n-k)} = \frac{MSA}{MSE}\sim F(k-1,n-k)
\]

拒绝域： \[
W = \{F>F_{1-\alpha}(k-1,n-k) \}
\]

\begin{longtable}[]{@{}llllll@{}}
\caption{单因素方差分析表}\tabularnewline
\toprule\noalign{}
误差来源 & 平方和 & 自由度 & 均方（MS） & F值 & p值 \\
\midrule\noalign{}
\endfirsthead
\toprule\noalign{}
误差来源 & 平方和 & 自由度 & 均方（MS） & F值 & p值 \\
\midrule\noalign{}
\endhead
\bottomrule\noalign{}
\endlastfoot
处理效应 & SSA & k-1 & MSA=SSA/(k-1) & F=MSA/MSE & \\
误差 & SSE & n-k & MSE=SSE/(n-k) & & \\
总效应 & SST=SSA+SSE & n-1 & & & \\
\end{longtable}

\begin{center}\rule{0.5\linewidth}{0.5pt}\end{center}

\textbf{速查卡片}：

\begin{itemize}
\tightlist
\item
  \textbf{核心问题}：观测到的效应是真实的还是随机波动？多组均值是否有显著差异？
\item
  \textbf{关键概念}：两类错误（第I类/第II类） → 势函数 → p值 → 显著性水平 → u检验/t检验/F检验/χ²检验 → 方差分析（ANOVA）：SST = SSA + SSE
\item
  \textbf{R命令}：\texttt{t.test()}, \texttt{var.test()}, \texttt{aov()}, \texttt{anova()}, \texttt{oneway.test()}, \texttt{pairwise.t.test()}
\item
  \textbf{参见}：Ch06（参数估计）, Ch10（回归的F检验）, Ch13（非参数检验）
\end{itemize}

\chapter{贝叶斯统计}\label{ux8d1dux53f6ux65afux7edfux8ba1}

\begin{quote}
频率学派问：``如果参数固定，数据出现的概率是多少？''贝叶斯学派问：``给定数据，参数取某个值的信念应该多少？''两个问题看似对称，却通向截然不同的统计哲学。
\end{quote}

\section{贝叶斯范式}\label{ux8d1dux53f6ux65afux8303ux5f0f}

\subsection{频率学派 vs 贝叶斯学派}\label{ux9891ux7387ux5b66ux6d3e-vs-ux8d1dux53f6ux65afux5b66ux6d3e}

前面几章（Ch05-07）介绍的参数估计和假设检验，本质上属于\textbf{频率学派}（Frequentist）的框架。在频率学派看来：

\begin{itemize}
\tightlist
\item
  参数 \(\theta\) 是固定的、未知的常数。
\item
  数据是随机的，来自 \(P(X \mid \theta)\)。
\item
  推断基于\textbf{重复抽样}的性质：无偏性、相合性、显著性水平等。
\end{itemize}

\textbf{贝叶斯学派}（Bayesian）则采取了完全不同的立场：

\begin{itemize}
\tightlist
\item
  参数 \(\theta\) 本身也是\textbf{随机的}------它服从一个先验分布 \(P(\theta)\)。
\item
  数据一旦被观测到，就是固定的。
\item
  推断基于\textbf{给定数据下参数的条件分布} \(P(\theta \mid \text{data})\)，即后验分布。
\end{itemize}

这两种范式的差异，根源在于对''概率是什么''这一元问题的不同回答。频率学派认为概率是长期频率的极限；贝叶斯学派则认为概率是信念（或知识状态）的量化。

\subsection{贝叶斯定理回顾}\label{ux8d1dux53f6ux65afux5b9aux7406ux56deux987e}

贝叶斯定理在 Ch01 中已经出现。将其应用于统计推断：

\[p(\theta \mid x) = \frac{p(x \mid \theta) \, p(\theta)}{p(x)} = \frac{p(x \mid \theta) \, p(\theta)}{\int p(x \mid \theta) \, p(\theta) \, d\theta}\]

其中：

\begin{itemize}
\tightlist
\item
  \(p(\theta)\)：\textbf{先验分布}------在观测数据之前对 \(\theta\) 的信念。
\item
  \(p(x \mid \theta)\)：\textbf{似然函数}------在参数 \(\theta\) 下观测到数据 \(x\) 的概率。
\item
  \(p(\theta \mid x)\)：\textbf{后验分布}------在观测数据之后对 \(\theta\) 的更新信念。
\item
  \(p(x)\)：\textbf{边际似然}（或证据）------归一化常数，确保后验是合法的概率分布。
\end{itemize}

由于 \(p(x)\) 不依赖于 \(\theta\)，后验分布常常写为比例形式：

\[p(\theta \mid x) \propto p(x \mid \theta) \, p(\theta)\]

\begin{quote}
后验 ∝ 似然 × 先验。这就是贝叶斯学习的核心公式。
\end{quote}

\section{先验分布的选择}\label{ux5148ux9a8cux5206ux5e03ux7684ux9009ux62e9}

\subsection{共轭先验}\label{ux5171ux8f6dux5148ux9a8c}

如果先验分布与后验分布属于同一分布族，则称该先验为\textbf{共轭先验}（Conjugate Prior）。共轭先验大大简化了后验的计算。

\begin{longtable}[]{@{}
  >{\raggedright\arraybackslash}p{(\columnwidth - 6\tabcolsep) * \real{0.2500}}
  >{\raggedright\arraybackslash}p{(\columnwidth - 6\tabcolsep) * \real{0.2500}}
  >{\raggedright\arraybackslash}p{(\columnwidth - 6\tabcolsep) * \real{0.2500}}
  >{\raggedright\arraybackslash}p{(\columnwidth - 6\tabcolsep) * \real{0.2500}}@{}}
\toprule\noalign{}
\begin{minipage}[b]{\linewidth}\raggedright
似然（数据分布）
\end{minipage} & \begin{minipage}[b]{\linewidth}\raggedright
参数
\end{minipage} & \begin{minipage}[b]{\linewidth}\raggedright
共轭先验
\end{minipage} & \begin{minipage}[b]{\linewidth}\raggedright
后验
\end{minipage} \\
\midrule\noalign{}
\endhead
\bottomrule\noalign{}
\endlastfoot
二项 \(B(n, \theta)\) & \(\theta\) & Beta\((\alpha, \beta)\) & Beta\((\alpha + x, \beta + n - x)\) \\
泊松 \(P(\lambda)\) & \(\lambda\) & Gamma\((\alpha, \beta)\) & Gamma\((\alpha + \sum x_i, \beta + n)\) \\
正态（均值，方差已知） & \(\mu\) & Normal\((\mu_0, \tau_0^2)\) & Normal（加权平均） \\
正态（方差，均值已知） & \(\sigma^2\) & Inv-Gamma\((\alpha, \beta)\) & Inv-Gamma（更新参数） \\
多项分布 & \(\mathbf{p}\) & Dirichlet\((\boldsymbol{\alpha})\) & Dirichlet\((\boldsymbol{\alpha} + \mathbf{n})\) \\
\end{longtable}

\textbf{例子}：硬币的公平性

设我们抛一枚硬币 \(n = 20\) 次，观察到 \(x = 13\) 次正面。设先验为 \(\theta \sim \text{Beta}(2, 2)\)（相当于事先''看过''2 正 2 反）。

后验：\(\theta \mid x \sim \text{Beta}(2 + 13, 2 + 7) = \text{Beta}(15, 9)\)

后验均值：\(E[\theta \mid x] = \frac{15}{15+9} = 0.625\)。而 MLE 为 \(\hat{\theta} = 13/20 = 0.65\)。先验的''收缩''效应将后验均值拉向先验均值 \(0.5\)。

\begin{Shaded}
\begin{Highlighting}[]
\NormalTok{n }\OtherTok{\textless{}{-}} \DecValTok{20}\NormalTok{; x }\OtherTok{\textless{}{-}} \DecValTok{13}
\NormalTok{alpha\_prior }\OtherTok{\textless{}{-}} \DecValTok{2}\NormalTok{; beta\_prior }\OtherTok{\textless{}{-}} \DecValTok{2}
\NormalTok{alpha\_post }\OtherTok{\textless{}{-}}\NormalTok{ alpha\_prior }\SpecialCharTok{+}\NormalTok{ x; beta\_post }\OtherTok{\textless{}{-}}\NormalTok{ beta\_prior }\SpecialCharTok{+}\NormalTok{ n }\SpecialCharTok{{-}}\NormalTok{ x}

\NormalTok{theta }\OtherTok{\textless{}{-}} \FunctionTok{seq}\NormalTok{(}\DecValTok{0}\NormalTok{, }\DecValTok{1}\NormalTok{, }\AttributeTok{length.out =} \DecValTok{200}\NormalTok{)}
\FunctionTok{plot}\NormalTok{(theta, }\FunctionTok{dbeta}\NormalTok{(theta, alpha\_post, beta\_post), }\AttributeTok{type =} \StringTok{"l"}\NormalTok{, }\AttributeTok{lwd =} \DecValTok{2}\NormalTok{,}
     \AttributeTok{xlab =} \FunctionTok{expression}\NormalTok{(theta), }\AttributeTok{ylab =} \StringTok{"密度"}\NormalTok{, }\AttributeTok{main =} \StringTok{"先验、似然与后验"}\NormalTok{)}
\FunctionTok{lines}\NormalTok{(theta, }\FunctionTok{dbeta}\NormalTok{(theta, alpha\_prior, beta\_prior), }\AttributeTok{lty =} \DecValTok{2}\NormalTok{, }\AttributeTok{col =} \StringTok{"blue"}\NormalTok{)}
\FunctionTok{lines}\NormalTok{(theta, }\FunctionTok{dbeta}\NormalTok{(theta, x }\SpecialCharTok{+} \DecValTok{1}\NormalTok{, n }\SpecialCharTok{{-}}\NormalTok{ x }\SpecialCharTok{+} \DecValTok{1}\NormalTok{), }\AttributeTok{lty =} \DecValTok{3}\NormalTok{, }\AttributeTok{col =} \StringTok{"red"}\NormalTok{)}
\FunctionTok{legend}\NormalTok{(}\StringTok{"topright"}\NormalTok{, }\FunctionTok{c}\NormalTok{(}\StringTok{"后验"}\NormalTok{, }\StringTok{"先验 Beta(2,2)"}\NormalTok{, }\StringTok{"似然"}\NormalTok{), }\AttributeTok{lty =} \DecValTok{1}\SpecialCharTok{:}\DecValTok{3}\NormalTok{, }\AttributeTok{col =} \FunctionTok{c}\NormalTok{(}\DecValTok{1}\NormalTok{, }\StringTok{"blue"}\NormalTok{, }\StringTok{"red"}\NormalTok{))}
\FunctionTok{abline}\NormalTok{(}\AttributeTok{v =}\NormalTok{ alpha\_post }\SpecialCharTok{/}\NormalTok{ (alpha\_post }\SpecialCharTok{+}\NormalTok{ beta\_post), }\AttributeTok{lty =} \DecValTok{4}\NormalTok{)}
\end{Highlighting}
\end{Shaded}

\begin{center}\includegraphics{statistic-book_files/figure-latex/unnamed-chunk-1-1} \end{center}

\subsection{无信息先验}\label{ux65e0ux4fe1ux606fux5148ux9a8c}

当我们缺乏先验知识时，可以使用\textbf{无信息先验}（Non-informative Prior），让数据''自己说话''。常见的无信息先验包括：

\begin{itemize}
\tightlist
\item
  \textbf{Laplace 先验}（均匀分布）：\(p(\theta) \propto 1\)。适用于参数空间有界的情况。
\item
  \textbf{Jeffreys 先验}：\(p(\theta) \propto \sqrt{I(\theta)}\)，其中 \(I(\theta)\) 为 Fisher 信息量。Jeffreys 先验在一对一变换下保持不变，具有良好的理论性质。
\item
  \textbf{位置参数的均匀先验}：\(p(\mu) \propto 1\)。
\item
  \textbf{尺度参数的倒数先验}：\(p(\sigma) \propto 1/\sigma\)。
\end{itemize}

\begin{quote}
注意：无信息先验可能是\textbf{非正常}（improper）的------即积分不为 1。但只要后验是正常的，就不会出问题。
\end{quote}

\subsection{层次先验}\label{ux5c42ux6b21ux5148ux9a8c}

当先验本身含有未知参数时，可以对这些超参数再设定先验，形成\textbf{层次模型}（Hierarchical Model）：

\[p(\theta, \phi \mid x) \propto p(x \mid \theta) \, p(\theta \mid \phi) \, p(\phi)\]

层次模型自然地处理了''不同组的参数之间存在相似性''的问题，在随机效应模型、多臂老虎机问题、贝叶斯优化等场景中应用广泛。

\section{贝叶斯推断}\label{ux8d1dux53f6ux65afux63a8ux65ad}

\subsection{点估计}\label{ux70b9ux4f30ux8ba1-1}

后验分布是一个完整的分布，而不是一个单点。但如果需要点估计，常用的选择是：

\begin{itemize}
\tightlist
\item
  \textbf{后验均值}：\(\hat{\theta}_{\text{mean}} = E[\theta \mid x]\)------最小化平方损失下的最优估计。
\item
  \textbf{后验中位数}：\(\hat{\theta}_{\text{med}}\)------最小化绝对损失下的最优估计。
\item
  \textbf{最大后验估计（MAP）}：\(\hat{\theta}_{\text{MAP}} = \arg\max_\theta p(\theta \mid x)\)------最大化后验密度。
\end{itemize}

MAP 估计与 MLE 的关系：由于 \(p(\theta \mid x) \propto p(x \mid \theta) p(\theta)\)，有

\[\log p(\theta \mid x) = \log p(x \mid \theta) + \log p(\theta) + \text{const}\]

因此，MAP 等价于在似然函数上加了一个来自先验的\textbf{正则化项}。这是连接频率学派与贝叶斯学派的重要桥梁。

\subsection{区间估计：可信区间}\label{ux533aux95f4ux4f30ux8ba1ux53efux4fe1ux533aux95f4}

贝叶斯的区间估计是\textbf{可信区间}（Credible Interval），与频率学派的置信区间有本质区别：

\begin{itemize}
\tightlist
\item
  \textbf{置信区间}：\(P(\text{区间包含} \, \theta) = 0.95\)------这个概率是对''区间''的随机性而言的（重复抽样下 95\% 的置信区间包含 \(\theta\)）。
\item
  \textbf{可信区间}：\(P(\theta \in \text{区间} \mid \text{data}) = 0.95\)------这个概率直接是对''参数''的信念而言的（给定数据，\(\theta\) 有 95\% 的概率落在该区间内）。
\end{itemize}

贝叶斯可以直接说：``我有 95\% 的把握认为参数在 \([a, b]\) 之间。''频率学派则不能说这句话------这是两种范式最直观的差异。

常用的可信区间：
- \textbf{等尾可信区间}：取后验的 \(\alpha/2\) 和 \(1-\alpha/2\) 分位数。
- \textbf{最高后验密度（HPD）区间}：在给定概率下，取后验密度最高的区域。HPD 区间是最短的可信区间。

\subsection{假设检验：贝叶斯因子}\label{ux5047ux8bbeux68c0ux9a8cux8d1dux53f6ux65afux56e0ux5b50}

贝叶斯假设检验不是通过 p 值，而是通过比较两个模型的\textbf{后验概率比}来进行：

\[\frac{P(H_1 \mid x)}{P(H_0 \mid x)} = \frac{P(x \mid H_1)}{P(x \mid H_0)} \cdot \frac{P(H_1)}{P(H_0)}\]

其中，\(BF_{10} = \frac{P(x \mid H_1)}{P(x \mid H_0)} = \frac{\int p(x \mid \theta_1) p(\theta_1) d\theta_1}{\int p(x \mid \theta_0) p(\theta_0) d\theta_0}\) 称为\textbf{贝叶斯因子}（Bayes Factor）。

\begin{longtable}[]{@{}lll@{}}
\toprule\noalign{}
\(\log_{10} BF_{10}\) & \(BF_{10}\) & 证据强度 \\
\midrule\noalign{}
\endhead
\bottomrule\noalign{}
\endlastfoot
\(< 0\) & \(< 1\) & 支持 \(H_0\) \\
\(0\) --- \(0.5\) & \(1\) --- \(3.2\) & 弱 \\
\(0.5\) --- \(1\) & \(3.2\) --- \(10\) & 中等 \\
\(1\) --- \(2\) & \(10\) --- \(100\) & 强 \\
\(> 2\) & \(> 100\) & 决定性 \\
\end{longtable}

\begin{quote}
贝叶斯因子自动''惩罚''模型复杂度------复杂模型的先验预测分布更分散，因此需要更强的数据支持才能获得大的贝叶斯因子。这体现了\textbf{奥卡姆剃刀}原则。
\end{quote}

\section{贝叶斯预测}\label{ux8d1dux53f6ux65afux9884ux6d4b}

贝叶斯方法自然地提供了预测分布。给定观测数据 \(x\)，新观测 \(x^*\) 的预测分布为：

\[p(x^* \mid x) = \int p(x^* \mid \theta) \, p(\theta \mid x) \, d\theta\]

这被称为\textbf{后验预测分布}（Posterior Predictive Distribution）。它\textbf{不是}用一个点估计 \(\hat{\theta}\) 代入 \(p(x^* \mid \hat{\theta})\)，而是\textbf{对参数的不确定性做积分平均}。这一区别意味着贝叶斯预测自然地包含了参数不确定性，通常比基于点估计的预测更稳健。

\section{贝叶斯统计的优缺点}\label{ux8d1dux53f6ux65afux7edfux8ba1ux7684ux4f18ux7f3aux70b9}

\textbf{优点}：
1. 推断解释直觉自然（``参数有 95\% 概率落在区间内''）。
2. 自然地融入先验知识。
3. 自动的正则化------先验防止过拟合。
4. 预测分布自然地包含了参数不确定性。
5. 在小样本下表现通常优于频率学派方法。

\textbf{缺点}：
1. 先验的选择可能主观，不同先验可能导致不同结论（敏感性分析是必要的）。
2. 计算复杂------后验往往没有解析形式，需要 MCMC（Ch09）等数值方法。
3. 高维问题中的计算挑战。

\section{与后续章节的联系}\label{ux4e0eux540eux7eedux7ae0ux8282ux7684ux8054ux7cfb}

\begin{itemize}
\tightlist
\item
  Ch09（贝叶斯计算）：当后验无法解析求解时，如何使用 MCMC 和变分推断进行数值计算------紧接本章形成完整的贝叶斯路径。
\item
  Ch11（广义线性模型）：贝叶斯 GLM。
\item
  Ch17（机器学习基础）：贝叶斯视角下的正则化与模型选择。
\item
  Ch22（因果推断）：贝叶斯因果推断。
\end{itemize}

\begin{center}\rule{0.5\linewidth}{0.5pt}\end{center}

\textbf{速查卡片}：

\begin{itemize}
\tightlist
\item
  \textbf{核心问题}：如何用概率量化信念，并用数据更新信念？
\item
  \textbf{核心公式}：后验 \(\propto\) 似然 × 先验
\item
  \textbf{关键概念}：共轭先验、MAP估计、可信区间、贝叶斯因子、后验预测分布
\item
  \textbf{与频率学派的核心差异}：参数是随机的（贝叶斯）vs 固定的（频率学派）；概率是信念 vs 长期频率
\item
  \textbf{R命令}：\texttt{rstanarm::stan\_glm()}, \texttt{brms::brm()}, \texttt{LearnBayes::} 系列函数
\item
  \textbf{参见}：Ch01（贝叶斯公式）, Ch06（参数估计）, Ch09（贝叶斯计算）
\end{itemize}

\chapter{贝叶斯计算}\label{ux8d1dux53f6ux65afux8ba1ux7b97}

\begin{quote}
Ch08 介绍了贝叶斯统计的理论框架。本章回答一个更实际的问题：当后验分布没有解析形式时，如何用计算机从后验分布中抽样？
\end{quote}

\section{贝叶斯思维}\label{ux8d1dux53f6ux65afux601dux7ef4}

\subsubsection{Beyes rule}\label{beyes-rule}

The process of using Bayes' rule to update a probability based on an event affecting it, is called Bayes' updating.

The what one tries to update can be considered `prior' information, sometimes simply called the \textbf{prior}.

Then, updating this prior using Bayes' rule gives the information conditional on the data, also known as the \textbf{posterior}, as in the information after having seen the data.

Going from the prior to the posterior is \textbf{Bayes learning} or Bayes updating.

不同于频率学派对概率的定义，贝叶斯学派认为概率反映的是我们先验的信念，因此\(P(E)\)可以是任何可能的分布。

我们的一生又何尝不是一场旷日持久的贝叶斯更新。

A Beyesian credible interval is a range for which the Bayesian thinks that the probability of including the true value is , say, 0.95.

Thus a Bayesian can say that there is a 95 percent chance that the credible interval contains the true parameter value.

Beyesian inference is to modify our beliefs to account for new data.

贝叶斯推断就是更新我们的信念去解释新的数据。

\[
p(\theta|data) = \frac{p(data|\theta)p(\theta)}{p(data)}\propto p(data|\theta)\cdot p(\theta)
\]

如何选择先验？

There is no ``true'' prior. In some cases expert opinion can be used to specify an informative prior. It would be a waste to discard this information.

If prior information is unavailable, then the prior should be uninformative. The prior is best viewed as an initial value to a statistical procedure.

只要先验选择是合理的，随着数据不断收集和贝叶斯学习不断的进行，后验会集中围绕在真值附近。但是在有限的样本中，先验是会产生影响的。

\subsubsection{Example1}\label{example1}

下面是一个关于睡眠时长分布的例子。

\begin{Shaded}
\begin{Highlighting}[]
\CommentTok{\# 描述先验分布}
\NormalTok{p }\OtherTok{=} \FunctionTok{seq}\NormalTok{(}\FloatTok{0.05}\NormalTok{, }\FloatTok{0.95}\NormalTok{, }\AttributeTok{by =} \FloatTok{0.1}\NormalTok{) }\CommentTok{\#p的可能取值}
\NormalTok{prior }\OtherTok{=} \FunctionTok{c}\NormalTok{(}\DecValTok{1}\NormalTok{, }\FloatTok{5.2}\NormalTok{, }\DecValTok{8}\NormalTok{, }\FloatTok{7.2}\NormalTok{, }\FloatTok{4.6}\NormalTok{, }\FloatTok{2.1}\NormalTok{, }\FloatTok{0.7}\NormalTok{, }\FloatTok{0.1}\NormalTok{, }\DecValTok{0}\NormalTok{, }\DecValTok{0}\NormalTok{)}
\NormalTok{prior }\OtherTok{=}\NormalTok{ prior}\SpecialCharTok{/}\FunctionTok{sum}\NormalTok{(prior) }\CommentTok{\#权重}
\FunctionTok{plot}\NormalTok{(p, prior, }\AttributeTok{ylim =} \FunctionTok{c}\NormalTok{(}\DecValTok{0}\NormalTok{,}\DecValTok{1}\NormalTok{), }\AttributeTok{type =} \StringTok{"h"}\NormalTok{, }\AttributeTok{ylab=}\StringTok{"Prior"}\NormalTok{)}
\end{Highlighting}
\end{Shaded}

\begin{center}\includegraphics{statistic-book_files/figure-latex/unnamed-chunk-2-1} \end{center}

\begin{Shaded}
\begin{Highlighting}[]
\CommentTok{\# 计算后验分布}
\NormalTok{s }\OtherTok{\textless{}{-}} \DecValTok{11}\NormalTok{; f }\OtherTok{\textless{}{-}} \DecValTok{16}
\NormalTok{like }\OtherTok{\textless{}{-}}\NormalTok{ p}\SpecialCharTok{\^{}}\NormalTok{s}\SpecialCharTok{*}\NormalTok{(}\DecValTok{1}\SpecialCharTok{{-}}\NormalTok{p)}\SpecialCharTok{\^{}}\NormalTok{f }\CommentTok{\# 计算似然}
\NormalTok{post }\OtherTok{\textless{}{-}}\NormalTok{ prior}\SpecialCharTok{*}\NormalTok{like }\CommentTok{\# 计算后验}
\NormalTok{post }\OtherTok{\textless{}{-}}\NormalTok{ post}\SpecialCharTok{/}\FunctionTok{sum}\NormalTok{(post) }\CommentTok{\#归一化}
\FunctionTok{plot}\NormalTok{(p,post,}\AttributeTok{type =} \StringTok{"h"}\NormalTok{,}\AttributeTok{lty =} \DecValTok{1}\NormalTok{,}\AttributeTok{ylim =} \FunctionTok{c}\NormalTok{(}\DecValTok{0}\NormalTok{,}\DecValTok{1}\NormalTok{),}\AttributeTok{ylab =} \StringTok{"Posterior"}\NormalTok{)}
\end{Highlighting}
\end{Shaded}

\begin{center}\includegraphics{statistic-book_files/figure-latex/unnamed-chunk-2-2} \end{center}

下面我们把参数的取值指定为一个beta分布。现在的问题是，在已知一些先验知识的前提下，如何确定一个beta分布的形状参数来最好的匹配先验的知识？

\begin{Shaded}
\begin{Highlighting}[]
\NormalTok{quantile2}\OtherTok{=}\FunctionTok{list}\NormalTok{(}\AttributeTok{p=}\FloatTok{0.9}\NormalTok{,}\AttributeTok{x=}\FloatTok{0.5}\NormalTok{)}\CommentTok{\#  p的0.9分位数是0.5}
\NormalTok{quantile1}\OtherTok{=}\FunctionTok{list}\NormalTok{(}\AttributeTok{p=}\FloatTok{0.5}\NormalTok{,}\AttributeTok{x=}\FloatTok{0.3}\NormalTok{)}\CommentTok{\#  p的中位数是0.3}
\NormalTok{LearnBayes}\SpecialCharTok{::}\FunctionTok{beta.select}\NormalTok{(quantile1,quantile2)}\CommentTok{\# 找到最符合的形状参数取值}
\end{Highlighting}
\end{Shaded}

\begin{verbatim}
## [1] 3.26 7.19
\end{verbatim}

将先验和似然函数相乘，可以得到后验分布： \[
g(p|data)\propto p^{a+s-1}(1-p)^{b+f-1}
\]

可以看出，它们同属于beta分布族。

\begin{Shaded}
\begin{Highlighting}[]
\CommentTok{\# Plot beta prior, likelihood and posterior}
\NormalTok{a }\OtherTok{=} \FloatTok{3.26}\NormalTok{; b }\OtherTok{=} \FloatTok{7.19}\NormalTok{; s }\OtherTok{=} \DecValTok{11}\NormalTok{; f }\OtherTok{=} \DecValTok{16}
\FunctionTok{curve}\NormalTok{(}\FunctionTok{dbeta}\NormalTok{(x,a}\SpecialCharTok{+}\NormalTok{s,b}\SpecialCharTok{+}\NormalTok{f), }\AttributeTok{from=}\DecValTok{0}\NormalTok{, }\AttributeTok{to=}\DecValTok{1}\NormalTok{,}
      \AttributeTok{xlab=}\StringTok{"p"}\NormalTok{,}\AttributeTok{ylab=}\StringTok{"Density"}\NormalTok{,}\AttributeTok{lty=}\DecValTok{1}\NormalTok{,}\AttributeTok{lwd=}\DecValTok{4}\NormalTok{)  }\CommentTok{\# 后验分布}
\FunctionTok{curve}\NormalTok{(}\FunctionTok{dbeta}\NormalTok{(x,s}\SpecialCharTok{+}\DecValTok{1}\NormalTok{,f}\SpecialCharTok{+}\DecValTok{1}\NormalTok{),}\AttributeTok{add=}\ConstantTok{TRUE}\NormalTok{,}\AttributeTok{lty=}\DecValTok{2}\NormalTok{,}\AttributeTok{lwd=}\DecValTok{4}\NormalTok{) }\CommentTok{\# 似然分布}
\FunctionTok{curve}\NormalTok{(}\FunctionTok{dbeta}\NormalTok{(x,a,b),}\AttributeTok{add=}\ConstantTok{TRUE}\NormalTok{,}\AttributeTok{lty=}\DecValTok{3}\NormalTok{,}\AttributeTok{lwd=}\DecValTok{4}\NormalTok{)}\CommentTok{\# 先验分布}
\FunctionTok{legend}\NormalTok{(.}\DecValTok{7}\NormalTok{,}\DecValTok{4}\NormalTok{,}\FunctionTok{c}\NormalTok{(}\StringTok{"Prior"}\NormalTok{,}\StringTok{"Likelihood"}\NormalTok{,}\StringTok{"Posterior"}\NormalTok{),}
       \AttributeTok{lty=}\FunctionTok{c}\NormalTok{(}\DecValTok{3}\NormalTok{,}\DecValTok{2}\NormalTok{,}\DecValTok{1}\NormalTok{),}\AttributeTok{lwd=}\FunctionTok{c}\NormalTok{(}\DecValTok{3}\NormalTok{,}\DecValTok{3}\NormalTok{,}\DecValTok{3}\NormalTok{))}
\end{Highlighting}
\end{Shaded}

\begin{center}\includegraphics{statistic-book_files/figure-latex/unnamed-chunk-4-1} \end{center}

\subsubsection{Example2}\label{example2}

下面我们要考虑的是均值已知，方差未知的正态分布。我们用美国足球成绩来估计未知方差。观测值为\(d_1,d_2,\cdots,d_n\)。

似然函数：\(L(\sigma^2)\propto (\sigma^2)^{-n/2}exp\{-\sum_{i=1}^nd_i^2/(2\sigma^2) \}\)

先验分布（使用非信息先验）：\(p(\sigma^2)\propto 1/\sigma^2\)

后验分布：\(g(\sigma^2|data)\propto (\sigma^2)^{-n/2-1}exp\{-v/(2\sigma^2)\},v=\sum_{i=1}^nd_i^2\)

定义\(Y=1/\sigma^2\)，给出Y的分布，以及\(\sigma\)的点估计和0.95置信区间。

\subsubsection{Example3}\label{example3}

下面我们考虑双参数未知的正态分布数据。我们使用的数据是纽约马拉松完赛时间，共有20个观测值\(y_1,\cdots,y_{20}\)，观测值都来自正态分布\(N(\mu,\sigma^2)\)

非信息先验分布：\(g(\mu,\sigma^2)\propto 1/\sigma^2\).

似然函数：\(L(y|\mu,\sigma^2) \propto \sigma^{-n}exp\{-\frac{1}{2\sigma^2}\sum_{i=1}^n(y_i-\mu)^2\}\)

联合后验分布： \[
g(\mu,\sigma^2|y)\propto \sigma^{-n-2}exp\{-\frac{1}{2\sigma^2}\sum_{i=1}^n(y_i-\mu)^2 \}
\]

联合后验密度可以分解为条件后验密度和边际后验密度的乘积。如何做到？ \[
g(\mu,\sigma^2|y)=g(\mu|\sigma^2,y)\times p(\sigma^2|y)
\]

条件后验分布：\(\mu|\sigma^2,y \sim N(\bar{y},\sigma^2/n)\).

边际后验分布：\(p(\sigma^2|y) \sim Scale-Inv-\chi^2(n-1,s^2)\).

将双参数的联合抽样转化为``先抽\(\sigma^2\)，再抽\(\mu\)''的两步抽样。

\begin{Shaded}
\begin{Highlighting}[]
\FunctionTok{library}\NormalTok{(LearnBayes)}
\FunctionTok{data}\NormalTok{(marathontimes)}
\FunctionTok{attach}\NormalTok{(marathontimes)}
\CommentTok{\# 绘制联合后验等高线图，等高线图的每一条线代表联合后验密度相等的点。}
\NormalTok{d }\OtherTok{=} \FunctionTok{mycontour}\NormalTok{(normchi2post, }\FunctionTok{c}\NormalTok{(}\DecValTok{220}\NormalTok{, }\DecValTok{330}\NormalTok{, }\DecValTok{500}\NormalTok{, }\DecValTok{9000}\NormalTok{), time, }\AttributeTok{xlab=}\StringTok{"mean"}\NormalTok{,}\AttributeTok{ylab=}\StringTok{"variance"}\NormalTok{)}

\CommentTok{\# simulation}
\NormalTok{S }\OtherTok{=} \FunctionTok{sum}\NormalTok{((time }\SpecialCharTok{{-}} \FunctionTok{mean}\NormalTok{(time))}\SpecialCharTok{\^{}}\DecValTok{2}\NormalTok{)}
\NormalTok{n }\OtherTok{=} \FunctionTok{length}\NormalTok{(time)}
\NormalTok{sigma2 }\OtherTok{=}\NormalTok{ S}\SpecialCharTok{/}\FunctionTok{rchisq}\NormalTok{(}\DecValTok{1000}\NormalTok{, n }\SpecialCharTok{{-}} \DecValTok{1}\NormalTok{)}\CommentTok{\#从边际后验抽方差}
\NormalTok{mu }\OtherTok{=} \FunctionTok{rnorm}\NormalTok{(}\DecValTok{1000}\NormalTok{, }\AttributeTok{mean =} \FunctionTok{mean}\NormalTok{(time), }\AttributeTok{sd =} \FunctionTok{sqrt}\NormalTok{(sigma2)}\SpecialCharTok{/}\FunctionTok{sqrt}\NormalTok{(n))}\CommentTok{\# 从条件后验抽均值}
\CommentTok{\# 验证抽样有效性，将抽样样本堆叠到等高线图上}
\NormalTok{d }\OtherTok{=} \FunctionTok{mycontour}\NormalTok{(normchi2post, }\FunctionTok{c}\NormalTok{(}\DecValTok{220}\NormalTok{, }\DecValTok{330}\NormalTok{, }\DecValTok{500}\NormalTok{, }\DecValTok{9000}\NormalTok{), time, }\AttributeTok{xlab=}\StringTok{"mean"}\NormalTok{,}\AttributeTok{ylab=}\StringTok{"variance"}\NormalTok{)}
\FunctionTok{points}\NormalTok{(mu, sigma2)}
\end{Highlighting}
\end{Shaded}

\begin{center}\includegraphics{statistic-book_files/figure-latex/unnamed-chunk-5-1} \end{center}

\begin{Shaded}
\begin{Highlighting}[]
\CommentTok{\#inference}
\FunctionTok{quantile}\NormalTok{(mu, }\FunctionTok{c}\NormalTok{(}\FloatTok{0.025}\NormalTok{, }\FloatTok{0.975}\NormalTok{))}
\end{Highlighting}
\end{Shaded}

\begin{verbatim}
##     2.5%    97.5% 
## 255.3482 300.3790
\end{verbatim}

\begin{Shaded}
\begin{Highlighting}[]
\FunctionTok{quantile}\NormalTok{(}\FunctionTok{sqrt}\NormalTok{(sigma2), }\FunctionTok{c}\NormalTok{(}\FloatTok{0.025}\NormalTok{, }\FloatTok{0.975}\NormalTok{))}
\end{Highlighting}
\end{Shaded}

\begin{verbatim}
##     2.5%    97.5% 
## 37.60029 73.70613
\end{verbatim}

使用网格估计的方法重新做一遍：

\begin{Shaded}
\begin{Highlighting}[]
\CommentTok{\# Data}
\FunctionTok{library}\NormalTok{(LearnBayes)}
\NormalTok{Y }\OtherTok{\textless{}{-}}\NormalTok{ marathontimes}\SpecialCharTok{$}\NormalTok{time}
\NormalTok{n }\OtherTok{\textless{}{-}} \FunctionTok{length}\NormalTok{(Y)}
\CommentTok{\# Create a grid}
\NormalTok{npts    }\OtherTok{\textless{}{-}} \DecValTok{100}
\NormalTok{mu      }\OtherTok{\textless{}{-}} \FunctionTok{seq}\NormalTok{( }\DecValTok{250}\NormalTok{, }\DecValTok{350}\NormalTok{,}\AttributeTok{length=}\NormalTok{npts)}
\NormalTok{s2      }\OtherTok{\textless{}{-}} \FunctionTok{seq}\NormalTok{( }\DecValTok{30}\SpecialCharTok{\^{}}\DecValTok{2}\NormalTok{,}\DecValTok{80}\SpecialCharTok{\^{}}\DecValTok{2}\NormalTok{,}\AttributeTok{length=}\NormalTok{npts)}
\NormalTok{grid\_mu }\OtherTok{\textless{}{-}} \FunctionTok{matrix}\NormalTok{(mu,npts,npts,}\AttributeTok{byrow=}\ConstantTok{FALSE}\NormalTok{)}
\NormalTok{grid\_s2 }\OtherTok{\textless{}{-}} \FunctionTok{matrix}\NormalTok{(s2,npts,npts,}\AttributeTok{byrow=}\ConstantTok{TRUE}\NormalTok{)}

\CommentTok{\#先验}
\NormalTok{prior }\OtherTok{\textless{}{-}} \FunctionTok{matrix}\NormalTok{(}\DecValTok{1}\NormalTok{,npts,npts)}
\ControlFlowTok{for}\NormalTok{(i }\ControlFlowTok{in} \DecValTok{1}\SpecialCharTok{:}\NormalTok{n)\{}
\NormalTok{  prior }\OtherTok{\textless{}{-}} \DecValTok{1}\SpecialCharTok{/}\NormalTok{grid\_s2}
\NormalTok{\}}
\NormalTok{prior }\OtherTok{\textless{}{-}}\NormalTok{ prior}\SpecialCharTok{/}\FunctionTok{sum}\NormalTok{(prior)}
\FunctionTok{library}\NormalTok{(fields)}
\FunctionTok{image.plot}\NormalTok{(mu,s2,prior,}\AttributeTok{main=}\StringTok{"Prior"}\NormalTok{)}
\end{Highlighting}
\end{Shaded}

\begin{center}\includegraphics{statistic-book_files/figure-latex/unnamed-chunk-6-1} \end{center}

\begin{Shaded}
\begin{Highlighting}[]
\CommentTok{\#似然}
\NormalTok{like   }\OtherTok{\textless{}{-}} \FunctionTok{matrix}\NormalTok{(}\DecValTok{1}\NormalTok{,npts,npts)}
\ControlFlowTok{for}\NormalTok{(i }\ControlFlowTok{in} \DecValTok{1}\SpecialCharTok{:}\NormalTok{n)\{}
\NormalTok{  like }\OtherTok{\textless{}{-}}\NormalTok{ like }\SpecialCharTok{*} \FunctionTok{dnorm}\NormalTok{(Y[i],grid\_mu,}\FunctionTok{sqrt}\NormalTok{(grid\_s2))}
\NormalTok{\}}
\NormalTok{like }\OtherTok{\textless{}{-}}\NormalTok{ like}\SpecialCharTok{/}\FunctionTok{sum}\NormalTok{(like)}
\FunctionTok{image.plot}\NormalTok{(mu,s2,like,}\AttributeTok{main=}\StringTok{"Likelihood"}\NormalTok{)}
\end{Highlighting}
\end{Shaded}

\begin{center}\includegraphics{statistic-book_files/figure-latex/unnamed-chunk-6-2} \end{center}

\begin{Shaded}
\begin{Highlighting}[]
\CommentTok{\# 后验}
\NormalTok{post }\OtherTok{\textless{}{-}}\NormalTok{ prior }\SpecialCharTok{*}\NormalTok{ like}
\NormalTok{post }\OtherTok{\textless{}{-}}\NormalTok{ post}\SpecialCharTok{/}\FunctionTok{sum}\NormalTok{(post)}
\FunctionTok{image.plot}\NormalTok{(mu,s2,post,}\AttributeTok{main=}\StringTok{"Posterior"}\NormalTok{)}
\end{Highlighting}
\end{Shaded}

\begin{center}\includegraphics{statistic-book_files/figure-latex/unnamed-chunk-6-3} \end{center}

\begin{Shaded}
\begin{Highlighting}[]
\CommentTok{\#mu的边际分布}
\NormalTok{post\_mu }\OtherTok{\textless{}{-}}\FunctionTok{colSums}\NormalTok{(post)}
\NormalTok{post\_mu }\OtherTok{\textless{}{-}}\NormalTok{ post\_mu}\SpecialCharTok{/}\FunctionTok{sum}\NormalTok{(post\_mu)}
\FunctionTok{plot}\NormalTok{(mu,post\_mu,}\AttributeTok{ylab=}\StringTok{"f(mu|Y)"}\NormalTok{,}\AttributeTok{type=}\StringTok{"l"}\NormalTok{)}
\FunctionTok{abline}\NormalTok{(}\AttributeTok{v=}\FunctionTok{mean}\NormalTok{(Y))}
\end{Highlighting}
\end{Shaded}

\begin{center}\includegraphics{statistic-book_files/figure-latex/unnamed-chunk-6-4} \end{center}

\begin{Shaded}
\begin{Highlighting}[]
\CommentTok{\#方差的边际分布}
\NormalTok{post\_s2 }\OtherTok{\textless{}{-}}\FunctionTok{rowSums}\NormalTok{(post)}
\NormalTok{post\_s2 }\OtherTok{\textless{}{-}}\NormalTok{ post\_s2}\SpecialCharTok{/}\FunctionTok{sum}\NormalTok{(post\_s2)}
\FunctionTok{plot}\NormalTok{(s2,post\_s2,}\AttributeTok{ylab=}\StringTok{"f(sigma2|Y)"}\NormalTok{,}\AttributeTok{type=}\StringTok{"l"}\NormalTok{)}
\FunctionTok{abline}\NormalTok{(}\AttributeTok{v=}\FunctionTok{var}\NormalTok{(Y))}
\end{Highlighting}
\end{Shaded}

\begin{center}\includegraphics{statistic-book_files/figure-latex/unnamed-chunk-6-5} \end{center}

\begin{Shaded}
\begin{Highlighting}[]
\CommentTok{\# 整合结果}
\NormalTok{post\_mean\_mu }\OtherTok{\textless{}{-}} \FunctionTok{sum}\NormalTok{(mu}\SpecialCharTok{*}\NormalTok{post\_mu)}
\NormalTok{post\_sd\_mu   }\OtherTok{\textless{}{-}} \FunctionTok{sum}\NormalTok{(mu}\SpecialCharTok{*}\NormalTok{mu}\SpecialCharTok{*}\NormalTok{post\_mu) }\SpecialCharTok{{-}}\NormalTok{ post\_mean\_mu}\SpecialCharTok{\^{}}\DecValTok{2}
\NormalTok{post\_q025\_mu }\OtherTok{\textless{}{-}} \FunctionTok{max}\NormalTok{(mu[}\FunctionTok{cumsum}\NormalTok{(post\_mu)}\SpecialCharTok{\textless{}}\FloatTok{0.025}\NormalTok{])}
\NormalTok{post\_q975\_mu }\OtherTok{\textless{}{-}} \FunctionTok{max}\NormalTok{(mu[}\FunctionTok{cumsum}\NormalTok{(post\_mu)}\SpecialCharTok{\textless{}}\FloatTok{0.975}\NormalTok{])}
\NormalTok{post\_mean\_s2 }\OtherTok{\textless{}{-}} \FunctionTok{sum}\NormalTok{(s2}\SpecialCharTok{*}\NormalTok{post\_s2)}
\NormalTok{post\_sd\_s2   }\OtherTok{\textless{}{-}} \FunctionTok{sum}\NormalTok{(s2}\SpecialCharTok{*}\NormalTok{s2}\SpecialCharTok{*}\NormalTok{post\_s2) }\SpecialCharTok{{-}}\NormalTok{ post\_mean\_s2}\SpecialCharTok{\^{}}\DecValTok{2}
\NormalTok{post\_q025\_s2 }\OtherTok{\textless{}{-}} \FunctionTok{max}\NormalTok{(s2[}\FunctionTok{cumsum}\NormalTok{(post\_s2)}\SpecialCharTok{\textless{}}\FloatTok{0.025}\NormalTok{])}
\NormalTok{post\_q975\_s2 }\OtherTok{\textless{}{-}} \FunctionTok{max}\NormalTok{(s2[}\FunctionTok{cumsum}\NormalTok{(post\_s2)}\SpecialCharTok{\textless{}}\FloatTok{0.975}\NormalTok{])}
\NormalTok{out1 }\OtherTok{\textless{}{-}} \FunctionTok{c}\NormalTok{(post\_mean\_mu,post\_sd\_mu,post\_q025\_mu,post\_q975\_mu)}
\NormalTok{out2 }\OtherTok{\textless{}{-}} \FunctionTok{c}\NormalTok{(post\_mean\_s2,post\_sd\_s2,post\_q025\_s2,post\_q975\_s2)}
\NormalTok{out  }\OtherTok{\textless{}{-}} \FunctionTok{rbind}\NormalTok{(out1,out2)}
\FunctionTok{colnames}\NormalTok{(out) }\OtherTok{\textless{}{-}} \FunctionTok{c}\NormalTok{(}\StringTok{"Mean"}\NormalTok{,}\StringTok{"SD"}\NormalTok{,}\StringTok{"Q025"}\NormalTok{,}\StringTok{"Q975"}\NormalTok{)}
\FunctionTok{rownames}\NormalTok{(out) }\OtherTok{\textless{}{-}} \FunctionTok{c}\NormalTok{(}\StringTok{"mu"}\NormalTok{,}\StringTok{"sigma2"}\NormalTok{)}
\FunctionTok{print}\NormalTok{(out)}
\end{Highlighting}
\end{Shaded}

\begin{verbatim}
##             Mean          SD      Q025      Q975
## mu      282.7072    270.5767  258.0808  322.7273
## sigma2 2435.5151 380807.4824 1177.7778 3622.2222
\end{verbatim}

\subsection{贝叶斯计算}\label{ux8d1dux53f6ux65afux8ba1ux7b97-1}

Bear in mind that given the prior and data, the posterior is fixed and a Bayesian analysis really boils down to summarizing the posterior.

Summarizing a p-dimensional posterior distribution is challenging for large p.

\subsubsection{MAP 估计}\label{map-ux4f30ux8ba1}

MAP与极大似然估计很像，但是添加了先验。

\[
\hat{\theta}=argmax_{\theta}\ p(\theta|Y) = argmax_{\theta} log[f(Y|\theta)]+log[p(\theta)]
\]

一个具体的例子，我们想研究一支队伍的胜率，它在上一个赛季的38场比赛中赢了30场。

我们先用MLE估计，写出似然函数，极大化就好。 \[
P(X=k|\theta)=C_n^k\theta^k(1-\theta)^{n-k}
\]

\[
L(\theta)=p(data|\theta) = C_{38}^{30}\theta^{30}(1-\theta)^8
\] 求导解得极大似然估计的参数值：

\[
\hat{\theta}_{MLE}=\frac{k}{n}=\frac{30}{38}
\]

当样本足够多的时候，MLE是很好的方法。但是样本数据不足的情况下呢？MAP就使用了先验的知识，辅助我们估计参数，使之更贴合实际。

先验分布：\(\theta\sim Beta(10,10),p(\theta)\propto \theta^9(1-\theta)^9\)

后验分布：\(p(\theta|data)\propto \theta^{39}(1-\theta)^{17}\)

求导解得MAP估计的参数值：

\[
\hat{\theta}_{MAP}=\frac{k+\alpha-1}{n+\alpha+\beta-2}=\frac{30+10-1}{38+10+10-2}=\frac{39}{56}
\]

\subsubsection{后验分布高斯近似}\label{ux540eux9a8cux5206ux5e03ux9ad8ux65afux8fd1ux4f3c}

贝叶斯中心极限定理：当样本量非常大的时候， \[
\theta|Y\sim Normal[\hat{\theta}_{MAP},I(\hat{\theta}_{MAP})^{-1}]
\]

\[
I_{jk}=-\frac{\partial^2}{\partial\theta_j\partial\theta_k}log[p(\theta|Y)]
\]

练习：Say \(Y|\theta\sim B(n,\theta),\theta\sim Beta(0.5,0.5)\),find the Gaussian approximation for \(p(\theta|Y)\)

后验分布：\(\theta|Y\sim Beta(Y+1/2,n-Y+1/2)\)

MAP估计：\(\hat{\theta}=A/(A+B),A=Y-1/2,B=n-Y-1/2\)

后验方差估计：\((\frac{A}{\hat{\theta}^2}+\frac{B}{(1-\hat{\theta})^2})^{-1}\)

\begin{Shaded}
\begin{Highlighting}[]
\NormalTok{theta }\OtherTok{\textless{}{-}} \FunctionTok{seq}\NormalTok{(}\FloatTok{0.001}\NormalTok{,}\FloatTok{0.999}\NormalTok{,.}\DecValTok{001}\NormalTok{) }\CommentTok{\# Grid of thetas for plotting}
\CommentTok{\#Y \textless{}{-} 2; n \textless{}{-} 5\# 小样本下估计}
\NormalTok{Y }\OtherTok{\textless{}{-}} \DecValTok{30}\NormalTok{; n  }\OtherTok{\textless{}{-}} \DecValTok{100}\CommentTok{\#大样本下估计}
\CommentTok{\# Compute the posterior mean and Fisher information matrix}
\NormalTok{A }\OtherTok{\textless{}{-}}\NormalTok{ Y}\FloatTok{{-}0.5}\NormalTok{; B }\OtherTok{\textless{}{-}}\NormalTok{ n}\SpecialCharTok{{-}}\NormalTok{Y}\FloatTok{{-}0.5}
\NormalTok{theta\_MAP }\OtherTok{\textless{}{-}}\NormalTok{ A}\SpecialCharTok{/}\NormalTok{(A}\SpecialCharTok{+}\NormalTok{B)}
\NormalTok{Info  }\OtherTok{\textless{}{-}}\NormalTok{ A}\SpecialCharTok{/}\NormalTok{theta\_MAP}\SpecialCharTok{\^{}}\DecValTok{2}\SpecialCharTok{+}\NormalTok{B}\SpecialCharTok{/}\NormalTok{(}\DecValTok{1}\SpecialCharTok{{-}}\NormalTok{theta\_MAP)}\SpecialCharTok{\^{}}\DecValTok{2}
\CommentTok{\# Plot the true and approximate posteriors}
\NormalTok{post1 }\OtherTok{\textless{}{-}} \FunctionTok{dbinom}\NormalTok{(Y,n,theta)}\SpecialCharTok{*}\FunctionTok{dbeta}\NormalTok{(theta,}\FloatTok{0.5}\NormalTok{,}\FloatTok{0.5}\NormalTok{)}
\NormalTok{post1 }\OtherTok{\textless{}{-}}\NormalTok{ post1}\SpecialCharTok{/}\FunctionTok{sum}\NormalTok{(post1)}
\NormalTok{post2 }\OtherTok{\textless{}{-}} \FunctionTok{dnorm}\NormalTok{(theta,theta\_MAP,}\FunctionTok{sqrt}\NormalTok{(}\DecValTok{1}\SpecialCharTok{/}\NormalTok{Info))}
\NormalTok{post2 }\OtherTok{\textless{}{-}}\NormalTok{ post2}\SpecialCharTok{/}\FunctionTok{sum}\NormalTok{(post2)}
\FunctionTok{plot}\NormalTok{(theta,post1,}\AttributeTok{type=}\StringTok{"l"}\NormalTok{,}\AttributeTok{lwd=}\DecValTok{2}\NormalTok{,}
     \AttributeTok{xlab=}\FunctionTok{expression}\NormalTok{(theta),}\AttributeTok{ylab=}\StringTok{"Posterior"}\NormalTok{)}
\FunctionTok{abline}\NormalTok{(}\AttributeTok{v=}\NormalTok{theta\_MAP,}\AttributeTok{col=}\DecValTok{3}\NormalTok{,}\AttributeTok{lwd=}\DecValTok{2}\NormalTok{)}
\FunctionTok{lines}\NormalTok{(theta,post2,}\AttributeTok{col=}\DecValTok{2}\NormalTok{,}\AttributeTok{lwd=}\DecValTok{2}\NormalTok{)}
\FunctionTok{legend}\NormalTok{(}\StringTok{"topright"}\NormalTok{,}\FunctionTok{c}\NormalTok{(}\StringTok{"Exact Posterior "}\NormalTok{,}\StringTok{"CLT"}\NormalTok{,}\StringTok{"MAP"}\NormalTok{),}\AttributeTok{bty=}\StringTok{"n"}\NormalTok{,}\AttributeTok{col=}\DecValTok{1}\SpecialCharTok{:}\DecValTok{3}\NormalTok{,}\AttributeTok{cex=}\FloatTok{1.5}\NormalTok{,}\AttributeTok{lwd=}\DecValTok{2}\NormalTok{)}
\end{Highlighting}
\end{Shaded}

\begin{center}\includegraphics{statistic-book_files/figure-latex/unnamed-chunk-7-1} \end{center}

\subsection{独立蒙特卡洛}\label{ux72ecux7acbux8499ux7279ux5361ux6d1b}

\subsubsection{直接抽样法}\label{ux76f4ux63a5ux62bdux6837ux6cd5}

直接抽样要求目标分布有CDF且能求逆。

Suppose that we are given \(U \sim U(0,1)\) and want to simulate a continuous r.v.X with cdf \(F_x\). Put \(Y=F_X^{-1}(U)\), we have \[
F_Y(y)=P(Y\leq y)=P(F_X^{-1}(U)\leq y)=P(U\leq F_X(y))=F_X(y)
\] For any continuous r.v.X,\(F_X(X)\sim U(0,1)\),Why?

来自指数分布的例子。

\begin{Shaded}
\begin{Highlighting}[]
\CommentTok{\#小样本可视化映射关系}
\NormalTok{lambda }\OtherTok{\textless{}{-}} \DecValTok{1}
\NormalTok{u }\OtherTok{\textless{}{-}} \FunctionTok{runif}\NormalTok{(}\DecValTok{20}\NormalTok{)}\CommentTok{\# 生成20个U(0,1)均匀分布的样本}
\NormalTok{x }\OtherTok{\textless{}{-}} \SpecialCharTok{{-}}\DecValTok{1}\SpecialCharTok{/}\NormalTok{lambda }\SpecialCharTok{*} \FunctionTok{log}\NormalTok{(}\DecValTok{1} \SpecialCharTok{{-}}\NormalTok{ u)}\CommentTok{\# 逆CDF变换：将均匀样本u转化为指数分布样本x}
\CommentTok{\# 绘制指数分布的CDF曲线}
\FunctionTok{curve}\NormalTok{(}\DecValTok{1} \SpecialCharTok{{-}} \FunctionTok{exp}\NormalTok{(}\SpecialCharTok{{-}}\NormalTok{x), }\DecValTok{0}\NormalTok{, }\DecValTok{5}\NormalTok{, }\AttributeTok{ylab =} \StringTok{"u"}\NormalTok{, }\AttributeTok{main =} \StringTok{"Inverse CDF transformation for exp(1)"}\NormalTok{)}
\CommentTok{\# 为每个(u[i], x[i])画虚线，展示“u→x”的映射关系}
\ControlFlowTok{for}\NormalTok{ (i }\ControlFlowTok{in} \DecValTok{1}\SpecialCharTok{:}\FunctionTok{length}\NormalTok{(x)) \{}
    \FunctionTok{lines}\NormalTok{(}\FunctionTok{c}\NormalTok{(x[i], x[i]), }\FunctionTok{c}\NormalTok{(u[i], }\DecValTok{0}\NormalTok{), }\AttributeTok{lty =} \StringTok{"dashed"}\NormalTok{)}
    \FunctionTok{lines}\NormalTok{(}\FunctionTok{c}\NormalTok{(x[i], }\DecValTok{0}\NormalTok{), }\FunctionTok{c}\NormalTok{(u[i], u[i]), }\AttributeTok{lty =} \StringTok{"dashed"}\NormalTok{)}
\NormalTok{\}}
\end{Highlighting}
\end{Shaded}

\begin{center}\includegraphics{statistic-book_files/figure-latex/unnamed-chunk-8-1} \end{center}

\begin{Shaded}
\begin{Highlighting}[]
\CommentTok{\# 大样本验证用逆CDF变换的抽样结果}
\NormalTok{lambda }\OtherTok{\textless{}{-}} \DecValTok{1}
\NormalTok{u }\OtherTok{\textless{}{-}} \FunctionTok{runif}\NormalTok{(}\DecValTok{1000}\NormalTok{)}\CommentTok{\# 生成1000个U(0,1)样本}
\NormalTok{x }\OtherTok{\textless{}{-}} \SpecialCharTok{{-}}\DecValTok{1}\SpecialCharTok{/}\NormalTok{lambda }\SpecialCharTok{*} \FunctionTok{log}\NormalTok{(}\DecValTok{1} \SpecialCharTok{{-}}\NormalTok{ u)}\CommentTok{\# 同样用逆CDF变换生成指数分布样本}
\FunctionTok{hist}\NormalTok{(x)}\CommentTok{\# 绘制样本x的直方图}
\end{Highlighting}
\end{Shaded}

\begin{center}\includegraphics{statistic-book_files/figure-latex/unnamed-chunk-8-2} \end{center}

\subsubsection{拒绝抽样}\label{ux62d2ux7eddux62bdux6837}

\textbf{What if we can not write the inverse?}

Although we cannot easily sample from \(f(\theta|y)\), there exists another density \(g(\theta)\)(called proposal density), like a Normal distribution or perhaps a t-distribution, from which it is easy for us to sample.

Then we can sample from \(g(\theta)\) directly and then ``reject'' the samples in a strategic way to make the resulting ``non-rejected'' samples look like they came from \(f(\theta|y)\) .

\begin{longtable}[]{@{}
  >{\raggedright\arraybackslash}p{(\columnwidth - 2\tabcolsep) * \real{0.5000}}
  >{\raggedright\arraybackslash}p{(\columnwidth - 2\tabcolsep) * \real{0.5000}}@{}}
\toprule\noalign{}
\begin{minipage}[b]{\linewidth}\raggedright
符号
\end{minipage} & \begin{minipage}[b]{\linewidth}\raggedright
含义
\end{minipage} \\
\midrule\noalign{}
\endhead
\bottomrule\noalign{}
\endlastfoot
\(f(\theta|y)\) & 目标分布的概率密度函数 \\
\(g(\theta)\) & 提议分布的概率密度函数 \\
\(M\) & 包络常数，满足\(M = max_{\theta}\frac{f(\theta|y)}{g(\theta)}\), 作用是让\(M\cdot g(\theta)\geq f(\theta|y)\) ，即放大后的提议分布能包裹目标分布 \\
\(U\) & 从均匀分布\(Unif(0,1)\)中抽取的随机数，保证阈值是随机的（而不能是固定的） \\
\(\theta\) & 提议分布抽样得到的候选样本 \\
\(\alpha = \frac{f(\theta|y)}{Mg(\theta)}\in [0,1]\) & 样本被接受概率 \\
\end{longtable}

\textbf{拒绝抽样步骤：}

\begin{enumerate}
\def\labelenumi{\arabic{enumi}.}
\item
  Sample \(U \sim Unif(0,1)\).
\item
  Sample \(\theta\) at randon from the probability density proportional to \(g(\theta)\).
\item
  if \(U\leq \frac{f(\theta|y)}{Mg(\theta)}\),accept \(\theta\) as a draw from f; if the drown \(\theta\) is rejected, return to step 1.
\end{enumerate}

重复此步骤，直到接受的样本足够多。

\textbf{直观理解：}

拒绝抽样的几何直观（以一维为例）：

\begin{itemize}
\tightlist
\item
  上方曲线为放大后的提议分布\(M \cdot g(\theta)\)，下方曲线为目标分布\(f(\theta|y)\)。两条曲线之间围成的区域代表了被拒绝的样本空间。
\item
  从\(g(\theta)\)中抽取一个候选值\(\theta^*\)（对应横轴上的点），在两条曲线之间随机取一个高度\(U \cdot M g(\theta^*)\)。若该高度落在下方曲线之下（即\(U \leq f(\theta^*|y) / [M g(\theta^*)]\)），则接受\(\theta^*\)。
\item
  接受的样本比例恰好等于\(1/M\)，即下方曲线面积与上方曲线面积之比。因此\(M\)越接近1，抽样效率越高。
\end{itemize}

The top curve is an approximation function, \(Mg(\theta)\) （注意这是放大后的提议分布）, and the button curve is the target density, \(f(\theta|y)\).

As required, \(Mg(\theta)\geq f(\theta|y)\) for all \(\theta\). The vertical line indicates a single random draw \(\theta\) from the density proportional to g. The probability that a sampled draw \(\theta\) is accepted is the ratio of the height of the lower curve to the height of the highter curve at the value \(\theta\).

正是通过设置随机阈值，当迭代次数足够多时，目标分布密度高的区域，接受概率也高，这样接受的样本数量多。反之目标分布密度低的区域接受额样本数量就会少，最终接受样本的分布会严格服从目标分布！

\textbf{如何选择}\(g(\theta)\)\textbf{？}

A simple choice is choose \(g(\theta)\) that heavier tails.（厚尾分布优先）

For example, if we want to get samples from \(N(0,1)\), we could use the \(t_2\) distribution as our proposal density as it has heavier tails than the Normal.

\begin{Shaded}
\begin{Highlighting}[]
\FunctionTok{curve}\NormalTok{(}\FunctionTok{dnorm}\NormalTok{(x), }\SpecialCharTok{{-}}\DecValTok{6}\NormalTok{, }\DecValTok{6}\NormalTok{, }\AttributeTok{xlab =} \FunctionTok{expression}\NormalTok{(theta), }\AttributeTok{ylab =} \StringTok{"Density"}\NormalTok{, }\AttributeTok{n =} \DecValTok{200}\NormalTok{)}
\FunctionTok{curve}\NormalTok{(}\FunctionTok{dt}\NormalTok{(x, }\DecValTok{2}\NormalTok{), }\SpecialCharTok{{-}}\DecValTok{6}\NormalTok{, }\DecValTok{6}\NormalTok{, }\AttributeTok{add =} \ConstantTok{TRUE}\NormalTok{, }\AttributeTok{col =} \DecValTok{4}\NormalTok{, }\AttributeTok{n =} \DecValTok{200}\NormalTok{)}
\FunctionTok{legend}\NormalTok{(}\StringTok{"topright"}\NormalTok{, }\FunctionTok{c}\NormalTok{(}\StringTok{"Normal density"}\NormalTok{, }\StringTok{"t density"}\NormalTok{), }\AttributeTok{col =} \FunctionTok{c}\NormalTok{(}\DecValTok{1}\NormalTok{, }\DecValTok{4}\NormalTok{), }\AttributeTok{bty =} \StringTok{"n"}\NormalTok{, }\AttributeTok{lty =} \DecValTok{1}\NormalTok{)}
\NormalTok{x }\OtherTok{\textless{}{-}} \FunctionTok{rt}\NormalTok{(}\DecValTok{200}\NormalTok{, }\DecValTok{2}\NormalTok{)}
\FunctionTok{rug}\NormalTok{(x, }\AttributeTok{col =} \DecValTok{4}\NormalTok{)}
\NormalTok{y }\OtherTok{\textless{}{-}} \FunctionTok{rnorm}\NormalTok{(}\DecValTok{200}\NormalTok{)}
\FunctionTok{rug}\NormalTok{(y, }\AttributeTok{col =} \DecValTok{1}\NormalTok{)}
\end{Highlighting}
\end{Shaded}

\begin{center}\includegraphics{statistic-book_files/figure-latex/unnamed-chunk-9-1} \end{center}

样本\(\theta\) 被接受的概率：

\[
P(\theta_{accept})=P(U\leq \frac{f(\theta|y)}{Mg(\theta)})=\frac{1}{M}
\]

由此，我们应该选取尽可能接近目标分布的提议分布g。

\textbf{拒绝抽样的问题}

For example, if g has heavier tails than f, then there will be too many extreme values generated from g. Rejection sampling simply thins out those extreme values to obtain the right proportion. But what if we could take those rejected values and, instead of discarding them, simply downweight or upweight them in a specific way?

\subsubsection{重要性抽样}\label{ux91cdux8981ux6027ux62bdux6837}

重要性抽样是拒绝抽样的进阶版，核心是用``样本加权''替代``样本拒绝''，直接计算目标分布f下函数h的期望。

target:

\[
E_f[h(\theta)|y]=E_g[\frac{f(\theta|y)}{g(\theta)}h(\theta|y)]
\]

if \(\theta_1.\cdots,\theta_n\sim g\), we get:

\[
\mu_n'=\frac{1}{n}\sum_{i=1}^n\frac{f(\theta_i|y)}{g(\theta_i)}h(\theta_i|y)=\frac{1}{n}\sum_{i=1}^nw_ih(\theta_i|y)\approx E_f[h(\theta|)]
\]

\(w_i\) are referred to as the importance weights.

\subsection{马尔科夫链蒙特卡洛方法MCMC}\label{ux9a6cux5c14ux79d1ux592bux94feux8499ux7279ux5361ux6d1bux65b9ux6cd5mcmc}

\textbf{What is Markov chain?}

A Markov chain is a sequence of random variables \(\theta^1,\theta^2,\cdots,\) for which, for any t, the distribution of \(\theta^t\) given all previous \(\theta'\)s depends only on the most recent value, \(\theta^{t-1}\).

The key to MCMC's success is that the approximate distributions are improved at each step in simulation (in the sense of converging to the target distribution).

\subsubsection{The Metropolis algorithm}\label{the-metropolis-algorithm}

The aim of Metropolis algorithm is to simulate \(x\sim p(x)\) where \(p(x)\) is called the target density, and we will need a proposal density \(q(x^{[new]}|x^{[old]})\).

An important property of q in the Metropolis algorithm is symmetry of the proposal:

\[
q(x^{[new]}|x^{[old]})=q(x^{[old]}|x^{[new]})
\]

To make sure this Markov chain converges to our target, we need to include the following steps.

\begin{enumerate}
\def\labelenumi{\arabic{enumi}.}
\item
  At step i+1, set \(x^{[old]}=x^{[i]}\).
\item
  Generate \(x^{[new]}\sim N(x^{[old]},1)\), (or \(x^{[new]}=x^{[old]}+\epsilon\) random walk proposal) and compute

  \[
  \alpha = min(1,\frac{p(x^{[new]})}{p(x^{[old]})})
  \]
\item
  Set \(x^{[i+1]}=x^{[new]}\) with probability \(\alpha\) (accept) ; set \(x^{[i+1]}=x^{[old]}\) with probability \(1-\alpha\) (reject).
\end{enumerate}

Example.

\begin{Shaded}
\begin{Highlighting}[]
\CommentTok{\#1. 抽样的目标分布：混合正态分布}
\NormalTok{p }\OtherTok{\textless{}{-}} \FloatTok{0.4}\CommentTok{\# 第一个正态分布的权重}
\NormalTok{mu }\OtherTok{\textless{}{-}} \FunctionTok{c}\NormalTok{(}\SpecialCharTok{{-}}\DecValTok{1}\NormalTok{, }\DecValTok{2}\NormalTok{)}
\NormalTok{sd }\OtherTok{\textless{}{-}} \FunctionTok{c}\NormalTok{(.}\DecValTok{5}\NormalTok{, }\DecValTok{2}\NormalTok{)}
\NormalTok{f }\OtherTok{\textless{}{-}} \ControlFlowTok{function}\NormalTok{(x)\{}
\NormalTok{    p     }\SpecialCharTok{*} \FunctionTok{dnorm}\NormalTok{(x, mu[}\DecValTok{1}\NormalTok{], sd[}\DecValTok{1}\NormalTok{]) }\SpecialCharTok{+}
\NormalTok{    (}\DecValTok{1}\SpecialCharTok{{-}}\NormalTok{p) }\SpecialCharTok{*} \FunctionTok{dnorm}\NormalTok{(x, mu[}\DecValTok{2}\NormalTok{], sd[}\DecValTok{2}\NormalTok{])}
\NormalTok{\} }
\FunctionTok{curve}\NormalTok{(}\FunctionTok{f}\NormalTok{(x), }\SpecialCharTok{{-}}\DecValTok{4}\NormalTok{, }\DecValTok{8}\NormalTok{, }\AttributeTok{n=}\DecValTok{301}\NormalTok{, }\AttributeTok{las=}\DecValTok{1}\NormalTok{)}
\end{Highlighting}
\end{Shaded}

\begin{center}\includegraphics{statistic-book_files/figure-latex/unnamed-chunk-10-1} \end{center}

\begin{Shaded}
\begin{Highlighting}[]
\CommentTok{\#2. 定义对称提议分布}
\NormalTok{q }\OtherTok{\textless{}{-}} \ControlFlowTok{function}\NormalTok{(x) }\FunctionTok{rnorm}\NormalTok{(}\DecValTok{1}\NormalTok{, x, }\DecValTok{4}\NormalTok{)}

\CommentTok{\#3. 定义单步转移函数}
\NormalTok{step }\OtherTok{\textless{}{-}} \ControlFlowTok{function}\NormalTok{(x, f, q) \{}
  \DocumentationTok{\#\# 从提议分布生成候选状态xp}
\NormalTok{  xp }\OtherTok{\textless{}{-}} \FunctionTok{q}\NormalTok{(x)}
  \DocumentationTok{\#\# 计算接受概率α}
\NormalTok{  alpha }\OtherTok{\textless{}{-}} \FunctionTok{min}\NormalTok{(}\DecValTok{1}\NormalTok{, }\FunctionTok{f}\NormalTok{(xp) }\SpecialCharTok{/} \FunctionTok{f}\NormalTok{(x))}
  \DocumentationTok{\#\# 随机判断是否接受候选状态}
  \ControlFlowTok{if}\NormalTok{ (}\FunctionTok{runif}\NormalTok{(}\DecValTok{1}\NormalTok{) }\SpecialCharTok{\textless{}}\NormalTok{ alpha)}
\NormalTok{    x }\OtherTok{\textless{}{-}}\NormalTok{ xp}
  \DocumentationTok{\#\# 返回当前状态}
\NormalTok{  x}
\NormalTok{\}}
\CommentTok{\#4. 定义迭代抽样函数}
\NormalTok{run }\OtherTok{\textless{}{-}} \ControlFlowTok{function}\NormalTok{(x, f, q, nsteps) \{}
  \CommentTok{\# 初始化矩阵：存储nsteps次迭代的状态}
\NormalTok{  res }\OtherTok{\textless{}{-}} \FunctionTok{matrix}\NormalTok{(}\ConstantTok{NA}\NormalTok{, nsteps, }\FunctionTok{length}\NormalTok{(x))}
  \CommentTok{\# 循环nsteps次，每次调用step()更新状态，并存储结果}
  \ControlFlowTok{for}\NormalTok{ (i }\ControlFlowTok{in} \FunctionTok{seq\_len}\NormalTok{(nsteps))}
\NormalTok{    res[i,] }\OtherTok{\textless{}{-}}\NormalTok{ x }\OtherTok{\textless{}{-}} \FunctionTok{step}\NormalTok{(x, f, q)}
  \CommentTok{\# 返回所有迭代的状态序列（马尔可夫链）}
  \FunctionTok{return}\NormalTok{(res)}
\NormalTok{\}}
\CommentTok{\# 执行抽样：初始状态x={-}10，迭代1100次}
\NormalTok{res }\OtherTok{\textless{}{-}} \FunctionTok{run}\NormalTok{(}\SpecialCharTok{{-}}\DecValTok{10}\NormalTok{, f, q, }\DecValTok{1100}\NormalTok{)}
\CommentTok{\# 绘制烧脑期后样本的直方图}
\FunctionTok{hist}\NormalTok{(res[}\DecValTok{101}\SpecialCharTok{:}\DecValTok{1100}\NormalTok{], }\DecValTok{50}\NormalTok{, }\AttributeTok{freq=}\ConstantTok{FALSE}\NormalTok{, }\AttributeTok{main=}\StringTok{""}\NormalTok{, }\AttributeTok{ylim=}\FunctionTok{c}\NormalTok{(}\DecValTok{0}\NormalTok{, .}\DecValTok{4}\NormalTok{), }\AttributeTok{las=}\DecValTok{1}\NormalTok{,}
     \AttributeTok{xlab=}\StringTok{"x"}\NormalTok{, }\AttributeTok{ylab=}\StringTok{"Probability density"}\NormalTok{)}
\NormalTok{z }\OtherTok{\textless{}{-}} \FunctionTok{integrate}\NormalTok{(f, }\SpecialCharTok{{-}}\ConstantTok{Inf}\NormalTok{, }\ConstantTok{Inf}\NormalTok{)}\SpecialCharTok{$}\NormalTok{value}
\FunctionTok{curve}\NormalTok{(}\FunctionTok{f}\NormalTok{(x) }\SpecialCharTok{/}\NormalTok{ z, }\AttributeTok{add=}\ConstantTok{TRUE}\NormalTok{, }\AttributeTok{col=}\StringTok{"red"}\NormalTok{, }\AttributeTok{n=}\DecValTok{200}\NormalTok{)}
\end{Highlighting}
\end{Shaded}

\begin{center}\includegraphics{statistic-book_files/figure-latex/unnamed-chunk-10-2} \end{center}

\begin{Shaded}
\begin{Highlighting}[]
\CommentTok{\# 绘制迹图验证收敛性}
\FunctionTok{plot}\NormalTok{(res[,}\DecValTok{1}\NormalTok{], }\AttributeTok{type=}\StringTok{"l"}\NormalTok{)}
\end{Highlighting}
\end{Shaded}

\begin{center}\includegraphics{statistic-book_files/figure-latex/unnamed-chunk-10-3} \end{center}

\subsubsection{The Convergence efficiency of MCMC}\label{the-convergence-efficiency-of-mcmc}

我们可以计算样本均值方差来表示 Monte Carlo error.

对于MCMC样本，由于样本之间存在自相关，样本均值的方差不能简单用\(Var(X)/M\)计算。考虑自相关后的正确方差为：

\[\text{Var}(\bar{x}) = \frac{\sigma^2}{M} \left(1 + 2\sum_{k=1}^{\infty} \rho_k \right)\]

其中\(\sigma^2 = Var(X)\)，\(\rho_k\)为滞后\(k\)的自相关系数。定义积分自相关时间（IACT）\(\tau = 1 + 2\sum_{k=1}^{\infty} \rho_k\)，则\(\text{Var}(\bar{x}) = \sigma^2 \tau / M\)。

ACF（自相关强度）

ACF（自相关函数）图展示了马尔可夫链在不同滞后距离下的自相关强度\(\rho_k = Corr(x_t, x_{t+k})\)。如果ACF随滞后快速衰减至零，说明链的混合性能好；如果ACF缓慢衰减，说明链的自相关性强，有效样本量（ESS）低。ESS定义为\(M_{ESS} = M / \tau\)，其中\(\tau\)为积分自相关时间。ESS越高，蒙特卡洛估计越精确。

Effective sample size (ESS):

\[
Var(\frac{1}{M}\sum_{i=1}^Mx_i)=\frac{Var(X)}{M}+\frac{2}{M}\sum\sum cov(x_i,x_j) = \frac{Var(X)}{M_{ESS}}
\]

\subsubsection{The Metropolis Hastings algorithm}\label{the-metropolis-hastings-algorithm}

MH是马尔可夫蒙特卡洛的核心算法，用于从无法直接采样的目标分布中生成样本。

Metropolis-Hastings算法的接受概率为：

\[\alpha(x^{[old]} \rightarrow x^{[new]}) = \min\left(1, \frac{p(x^{[new]}) \cdot q(x^{[old]} | x^{[new]})}{p(x^{[old]}) \cdot q(x^{[new]} | x^{[old]})}\right)\]

关键点：
- 当提议分布\(q\)为对称分布时（即\(q(x'|x) = q(x|x')\)），接受率退化为\(\min(1, p(x^{[new]})/p(x^{[old]}))\)，此时退化为Metropolis算法。
- MH算法通过引入提议分布比值\(q(x^{[old]}|x^{[new]})/q(x^{[new]}|x^{[old]})\)，纠正了非对称提议带来的采样偏差，使马尔可夫链的平稳分布恰好为目标分布\(p(x)\)。
- 该修正因子被称为Hastings比率，是MH算法区别于Metropolis算法的核心。

2d example

\begin{Shaded}
\begin{Highlighting}[]
\CommentTok{\# 构造多元正态分布的概率密度函数}
\NormalTok{make.mvn }\OtherTok{\textless{}{-}} \ControlFlowTok{function}\NormalTok{(mean, vcv) \{}
\NormalTok{  logdet }\OtherTok{\textless{}{-}} \FunctionTok{as.numeric}\NormalTok{(}\FunctionTok{determinant}\NormalTok{(vcv, }\ConstantTok{TRUE}\NormalTok{)}\SpecialCharTok{$}\NormalTok{modulus) }\CommentTok{\# 协方差矩阵的对数行列式}
\NormalTok{  tmp }\OtherTok{\textless{}{-}} \FunctionTok{length}\NormalTok{(mean) }\SpecialCharTok{*} \FunctionTok{log}\NormalTok{(}\DecValTok{2} \SpecialCharTok{*}\NormalTok{ pi) }\SpecialCharTok{+}\NormalTok{ logdet }\CommentTok{\# 多元正态PDF的常数项部分}
\NormalTok{  vcv.i }\OtherTok{\textless{}{-}} \FunctionTok{solve}\NormalTok{(vcv) }\CommentTok{\# 协方差矩阵的逆矩阵}
  \CommentTok{\# 返回闭包（输入样本x，输出PDF值）}
  \ControlFlowTok{function}\NormalTok{(x) \{}
\NormalTok{    dx }\OtherTok{\textless{}{-}}\NormalTok{ x }\SpecialCharTok{{-}}\NormalTok{ mean}
    \FunctionTok{exp}\NormalTok{(}\SpecialCharTok{{-}}\NormalTok{(tmp }\SpecialCharTok{+} \FunctionTok{rowSums}\NormalTok{((dx }\SpecialCharTok{\%*\%}\NormalTok{ vcv.i) }\SpecialCharTok{*}\NormalTok{ dx))}\SpecialCharTok{/}\DecValTok{2}\NormalTok{)}
\NormalTok{  \}}
\NormalTok{\}}

\CommentTok{\# 定义目标分布：双峰混合双变量正态分布}
\NormalTok{mu1 }\OtherTok{\textless{}{-}} \FunctionTok{c}\NormalTok{(}\SpecialCharTok{{-}}\DecValTok{1}\NormalTok{, }\DecValTok{1}\NormalTok{) }\CommentTok{\#峰1均值}
\NormalTok{mu2 }\OtherTok{\textless{}{-}} \FunctionTok{c}\NormalTok{(}\DecValTok{2}\NormalTok{, }\SpecialCharTok{{-}}\DecValTok{2}\NormalTok{) }\CommentTok{\#峰2均值}
\NormalTok{vcv1 }\OtherTok{\textless{}{-}} \FunctionTok{matrix}\NormalTok{(}\FunctionTok{c}\NormalTok{(}\DecValTok{1}\NormalTok{, .}\DecValTok{25}\NormalTok{, .}\DecValTok{25}\NormalTok{, }\FloatTok{1.5}\NormalTok{), }\DecValTok{2}\NormalTok{, }\DecValTok{2}\NormalTok{) }\CommentTok{\#峰1协方差矩阵}
\NormalTok{vcv2 }\OtherTok{\textless{}{-}} \FunctionTok{matrix}\NormalTok{(}\FunctionTok{c}\NormalTok{(}\DecValTok{2}\NormalTok{, }\SpecialCharTok{{-}}\NormalTok{.}\DecValTok{5}\NormalTok{, }\SpecialCharTok{{-}}\NormalTok{.}\DecValTok{5}\NormalTok{, }\DecValTok{2}\NormalTok{), }\DecValTok{2}\NormalTok{, }\DecValTok{2}\NormalTok{) }\CommentTok{\#峰2协方差矩阵}
\NormalTok{f1 }\OtherTok{\textless{}{-}} \FunctionTok{make.mvn}\NormalTok{(mu1, vcv1) }\CommentTok{\# 峰1概率密度函数}
\NormalTok{f2 }\OtherTok{\textless{}{-}} \FunctionTok{make.mvn}\NormalTok{(mu2, vcv2) }\CommentTok{\# 峰2概率密度函数}
\NormalTok{f }\OtherTok{\textless{}{-}} \ControlFlowTok{function}\NormalTok{(x)}\FunctionTok{f1}\NormalTok{(x) }\SpecialCharTok{+} \FunctionTok{f2}\NormalTok{(x) }\CommentTok{\# 目标分布}

\CommentTok{\# 可视化目标分布}
\NormalTok{x }\OtherTok{\textless{}{-}} \FunctionTok{seq}\NormalTok{(}\SpecialCharTok{{-}}\DecValTok{5}\NormalTok{, }\DecValTok{6}\NormalTok{, }\AttributeTok{length=}\DecValTok{71}\NormalTok{) }\CommentTok{\#x轴网格点}
\NormalTok{y }\OtherTok{\textless{}{-}} \FunctionTok{seq}\NormalTok{(}\SpecialCharTok{{-}}\DecValTok{7}\NormalTok{, }\DecValTok{6}\NormalTok{, }\AttributeTok{length=}\DecValTok{61}\NormalTok{) }\CommentTok{\#y轴网格点}
\NormalTok{xy }\OtherTok{\textless{}{-}} \FunctionTok{expand.grid}\NormalTok{(}\AttributeTok{x=}\NormalTok{x, }\AttributeTok{y=}\NormalTok{y) }\CommentTok{\#生成所有网格点组合}
\NormalTok{z }\OtherTok{\textless{}{-}} \FunctionTok{matrix}\NormalTok{(}\FunctionTok{apply}\NormalTok{(}\FunctionTok{as.matrix}\NormalTok{(xy), }\DecValTok{1}\NormalTok{, f), }\FunctionTok{length}\NormalTok{(x), }\FunctionTok{length}\NormalTok{(y))}
\CommentTok{\# 计算每个网格点的目标分布密度}
\FunctionTok{image}\NormalTok{(x, y, z, }\AttributeTok{las=}\DecValTok{1}\NormalTok{) }\CommentTok{\# 画热力图}
\FunctionTok{contour}\NormalTok{(x, y, z, }\AttributeTok{add=}\ConstantTok{TRUE}\NormalTok{) }\CommentTok{\# 叠加等高线}
\end{Highlighting}
\end{Shaded}

\begin{center}\includegraphics{statistic-book_files/figure-latex/unnamed-chunk-11-1} \end{center}

\begin{Shaded}
\begin{Highlighting}[]
\CommentTok{\# 定义提议分布：对称随机游走}
\NormalTok{q }\OtherTok{\textless{}{-}} \ControlFlowTok{function}\NormalTok{(x, }\AttributeTok{d=}\DecValTok{8}\NormalTok{) x }\SpecialCharTok{+} \FunctionTok{runif}\NormalTok{(}\FunctionTok{length}\NormalTok{(x), }\SpecialCharTok{{-}}\NormalTok{d}\SpecialCharTok{/}\DecValTok{2}\NormalTok{, d}\SpecialCharTok{/}\DecValTok{2}\NormalTok{)}

\CommentTok{\# MH采样与可视化}
\NormalTok{x0 }\OtherTok{\textless{}{-}} \FunctionTok{c}\NormalTok{(}\SpecialCharTok{{-}}\DecValTok{4}\NormalTok{, }\SpecialCharTok{{-}}\DecValTok{4}\NormalTok{) }\CommentTok{\#采样初始值}
\FunctionTok{set.seed}\NormalTok{(}\DecValTok{1}\NormalTok{)}
\CommentTok{\# 迭代1000次采样（短序列，看轨迹）}
\NormalTok{samples }\OtherTok{\textless{}{-}} \FunctionTok{run}\NormalTok{(x0, f, q, }\DecValTok{1000}\NormalTok{)}
\FunctionTok{image}\NormalTok{(x, y, z, }\AttributeTok{xlim=}\FunctionTok{range}\NormalTok{(x, samples[,}\DecValTok{1}\NormalTok{]), }\AttributeTok{ylim=}\FunctionTok{range}\NormalTok{(x, samples[,}\DecValTok{2}\NormalTok{]))}
\FunctionTok{contour}\NormalTok{(x, y, z, }\AttributeTok{add=}\ConstantTok{TRUE}\NormalTok{)}
\FunctionTok{lines}\NormalTok{(samples[,}\DecValTok{1}\NormalTok{], samples[,}\DecValTok{2}\NormalTok{], }\AttributeTok{col=}\StringTok{"\#00000088"}\NormalTok{)}
\end{Highlighting}
\end{Shaded}

\begin{center}\includegraphics{statistic-book_files/figure-latex/unnamed-chunk-11-2} \end{center}

\begin{Shaded}
\begin{Highlighting}[]
\CommentTok{\# 迭代10000次采样（长序列，验证收敛）}
\NormalTok{samples }\OtherTok{\textless{}{-}} \FunctionTok{run}\NormalTok{(x0, f, q, }\DecValTok{100000}\NormalTok{)}
\FunctionTok{smoothScatter}\NormalTok{(samples)}
\FunctionTok{contour}\NormalTok{(x, y, z, }\AttributeTok{add=}\ConstantTok{TRUE}\NormalTok{)}
\end{Highlighting}
\end{Shaded}

\begin{center}\includegraphics{statistic-book_files/figure-latex/unnamed-chunk-11-3} \end{center}

\textbf{间接抽样法和MH方法的异同点比较}

相同点：都依赖提议分布；提议分布的适配性决定采样效率；均含有接受-拒绝机制。

不同点：间接抽样法产生独立样本，MH方法产生相关样本；间接抽样法要求存在上界M；MH更适合高维向量情况

\subsubsection{Gibbs sampler}\label{gibbs-sampler}

吉布斯采样相比于MH算法的优点在于，将高维采样问题拆解为一系列低维条件分布的采样，避免设计复杂的高维提议分布。

吉布斯采样的工作流程（以二维为例）：

设目标分布为\(p(x, y)\)，已知条件分布\(p(x|y)\)和\(p(y|x)\)均可直接采样。

\begin{enumerate}
\def\labelenumi{\arabic{enumi}.}
\item
  初始化：给定初始值\(y^{(0)}\)。
\item
  第\(t\)次迭代：

  \begin{itemize}
  \tightlist
  \item
    从\(p(x | y^{(t-1)})\)中采样得到\(x^{(t)}\)
  \item
    从\(p(y | x^{(t)})\)中采样得到\(y^{(t)}\)
  \end{itemize}
\item
  重复步骤2，经过烧脑期后，样本对\((x^{(t)}, y^{(t)})\)近似来自联合分布\(p(x, y)\)。
\end{enumerate}

Gibbs采样在参数空间中沿坐标轴方向移动，每一步只改变一个分量。在高维问题中，这种''分而治之''的策略避免了直接处理高维提议分布的困难。

Gibbs采样是MH算法的特例，其提议分布是目标分布的条件分布，因此接受率恒等于1。

\[
\alpha=min(1,\frac{f(x^{[i+1]},y)f(x^{[i]}|y)}{f(x^{[i]},y)f(x^{[i+1]}|y)})=1
\]

\begin{itemize}
\item
  优点：无需调优提议分布；提议总是被接受
\item
  缺点：依赖条件分布的可采样性；参数高度相关时采样极慢
\end{itemize}

\begin{Shaded}
\begin{Highlighting}[]
\FunctionTok{set.seed}\NormalTok{(}\DecValTok{123}\NormalTok{)}
\NormalTok{rho }\OtherTok{\textless{}{-}} \FloatTok{0.8} \CommentTok{\# 相关系数}
\NormalTok{n }\OtherTok{\textless{}{-}} \DecValTok{5000} \CommentTok{\# Gibbs迭代次数}
\NormalTok{X }\OtherTok{\textless{}{-}}\NormalTok{ Y }\OtherTok{\textless{}{-}} \FunctionTok{numeric}\NormalTok{(n) }\CommentTok{\# 初始化存储采样结果的向量}
\CommentTok{\# 初始值}
\NormalTok{X[}\DecValTok{1}\NormalTok{] }\OtherTok{\textless{}{-}} \DecValTok{0}
\NormalTok{Y[}\DecValTok{1}\NormalTok{] }\OtherTok{\textless{}{-}} \DecValTok{0}
\NormalTok{sd }\OtherTok{\textless{}{-}} \FunctionTok{sqrt}\NormalTok{(}\DecValTok{1} \SpecialCharTok{{-}}\NormalTok{ rho}\SpecialCharTok{\^{}}\DecValTok{2}\NormalTok{) }\CommentTok{\# 条件分布的标准差}

\CommentTok{\# Gibbs 核心迭代循环}
\ControlFlowTok{for}\NormalTok{ (t }\ControlFlowTok{in} \DecValTok{2}\SpecialCharTok{:}\NormalTok{n) \{}
\NormalTok{  X[t] }\OtherTok{\textless{}{-}} \FunctionTok{rnorm}\NormalTok{(}\DecValTok{1}\NormalTok{, }\AttributeTok{mean =}\NormalTok{ rho }\SpecialCharTok{*}\NormalTok{ Y[t}\DecValTok{{-}1}\NormalTok{], }\AttributeTok{sd =}\NormalTok{ sd) }\CommentTok{\# Sample X | Y}
\NormalTok{  Y[t] }\OtherTok{\textless{}{-}} \FunctionTok{rnorm}\NormalTok{(}\DecValTok{1}\NormalTok{, }\AttributeTok{mean =}\NormalTok{ rho }\SpecialCharTok{*}\NormalTok{ X[t], }\AttributeTok{sd =}\NormalTok{ sd)  }\CommentTok{\# Sample Y | X}
\NormalTok{\}}

\CommentTok{\# 可视化采样结果}
\FunctionTok{plot}\NormalTok{(X, Y, }\AttributeTok{pch =} \DecValTok{20}\NormalTok{, }\AttributeTok{col =} \FunctionTok{rgb}\NormalTok{(}\DecValTok{0}\NormalTok{,}\DecValTok{0}\NormalTok{,}\DecValTok{1}\NormalTok{,}\FloatTok{0.4}\NormalTok{),}
     \AttributeTok{main =} \StringTok{"Joint Samples (X, Y) with Contour"}\NormalTok{,}
     \AttributeTok{xlab =} \StringTok{"X"}\NormalTok{, }\AttributeTok{ylab =} \StringTok{"Y"}\NormalTok{)}

\CommentTok{\# 与真实二维正态分布对比}
\FunctionTok{library}\NormalTok{(mvtnorm)}
\CommentTok{\# Grid for contour plot}
\NormalTok{xg }\OtherTok{\textless{}{-}} \FunctionTok{seq}\NormalTok{(}\FunctionTok{min}\NormalTok{(X)}\SpecialCharTok{{-}}\DecValTok{1}\NormalTok{, }\FunctionTok{max}\NormalTok{(X)}\SpecialCharTok{+}\DecValTok{1}\NormalTok{, }\AttributeTok{length=}\DecValTok{80}\NormalTok{)}
\NormalTok{yg }\OtherTok{\textless{}{-}} \FunctionTok{seq}\NormalTok{(}\FunctionTok{min}\NormalTok{(Y)}\SpecialCharTok{{-}}\DecValTok{1}\NormalTok{, }\FunctionTok{max}\NormalTok{(Y)}\SpecialCharTok{+}\DecValTok{1}\NormalTok{, }\AttributeTok{length=}\DecValTok{80}\NormalTok{)}
\NormalTok{grid }\OtherTok{\textless{}{-}} \FunctionTok{expand.grid}\NormalTok{(}\AttributeTok{x=}\NormalTok{xg, }\AttributeTok{y=}\NormalTok{yg)}
\CommentTok{\# True bivariate normal density}
\NormalTok{Sigma }\OtherTok{\textless{}{-}} \FunctionTok{matrix}\NormalTok{(}\FunctionTok{c}\NormalTok{(}\DecValTok{1}\NormalTok{, rho, rho, }\DecValTok{1}\NormalTok{), }\DecValTok{2}\NormalTok{, }\DecValTok{2}\NormalTok{)}
\NormalTok{dens }\OtherTok{\textless{}{-}} \FunctionTok{dmvnorm}\NormalTok{(grid, }\AttributeTok{mean=}\FunctionTok{c}\NormalTok{(}\DecValTok{0}\NormalTok{,}\DecValTok{0}\NormalTok{), }\AttributeTok{sigma=}\NormalTok{Sigma)}
\CommentTok{\# Add contour}
\FunctionTok{contour}\NormalTok{(xg, yg, }\FunctionTok{matrix}\NormalTok{(dens, }\DecValTok{80}\NormalTok{, }\DecValTok{80}\NormalTok{),}
        \AttributeTok{add=}\ConstantTok{TRUE}\NormalTok{, }\AttributeTok{drawlabels=}\ConstantTok{FALSE}\NormalTok{,}
        \AttributeTok{col=}\StringTok{"red"}\NormalTok{, }\AttributeTok{lwd=}\DecValTok{2}\NormalTok{)}
\CommentTok{\# Add legend}
\FunctionTok{legend}\NormalTok{(}\StringTok{"bottomright"}\NormalTok{,}
       \AttributeTok{legend=}\FunctionTok{c}\NormalTok{(}\StringTok{"Gibbs Samples"}\NormalTok{, }\StringTok{"True Contour"}\NormalTok{),}
       \AttributeTok{col=}\FunctionTok{c}\NormalTok{(}\FunctionTok{rgb}\NormalTok{(}\DecValTok{0}\NormalTok{,}\DecValTok{0}\NormalTok{,}\DecValTok{1}\NormalTok{,}\FloatTok{0.4}\NormalTok{), }\StringTok{"red"}\NormalTok{),}
       \AttributeTok{pch=}\FunctionTok{c}\NormalTok{(}\DecValTok{20}\NormalTok{, }\ConstantTok{NA}\NormalTok{),}
       \AttributeTok{lty=}\FunctionTok{c}\NormalTok{(}\ConstantTok{NA}\NormalTok{, }\DecValTok{1}\NormalTok{),}
       \AttributeTok{lwd=}\FunctionTok{c}\NormalTok{(}\ConstantTok{NA}\NormalTok{, }\DecValTok{2}\NormalTok{))}
\end{Highlighting}
\end{Shaded}

\begin{center}\includegraphics{statistic-book_files/figure-latex/unnamed-chunk-12-1} \end{center}

\textbf{Metropolis within Gibbs}

即使单个条件分布难以模拟采样，也可在每个 Gibbs 步骤中使用 Metropolis-Hastings 算法。这样既保留了Gibbs高维拆低维的核心优势，又用MH弥补Gibbs条件分布需可采样的缺陷，适配现实中复杂的目标分布。

\subsubsection{贝叶斯线性回归模型}\label{ux8d1dux53f6ux65afux7ebfux6027ux56deux5f52ux6a21ux578b}

线性回归模型：

\[
y_i=\beta_0+\beta_1x_i+\epsilon_i
\]

贝叶斯法则：

\[
p(\beta_0,\beta_1,\sigma^2|y,x)\propto p(y|\beta_0,\beta_1,\sigma^2,x)p(\beta_0,\beta_1,\sigma^2)
\]

Gibbs采样：

贝叶斯线性回归的Gibbs采样步骤如下：

给定先验\(p(\beta, \sigma^2) \propto 1/\sigma^2\)，后验条件分布为：

\begin{enumerate}
\def\labelenumi{\arabic{enumi}.}
\item
  \(\beta | \sigma^2, y, X \sim N(\hat{\beta}_{OLS}, \sigma^2 (X'X)^{-1})\)
\item
  \(\sigma^2 | \beta, y, X \sim \text{Scale-Inv-}\chi^2(n, s^2)\)，其中\(s^2 = \frac{1}{n}\sum_{i=1}^n (y_i - x_i'\beta)^2\)
\end{enumerate}

Gibbs迭代：
- 步骤1：从条件后验\(\beta | \sigma^{2(t-1)}, y, X\)中抽取\(\beta^{(t)}\)
- 步骤2：从条件后验\(\sigma^2 | \beta^{(t)}, y, X\)中抽取\(\sigma^{2(t)}\)
- 交替重复，收敛后获得联合后验样本\((\beta^{(t)}, \sigma^{2(t)})\)。

\chapter{回归分析}\label{ux56deux5f52ux5206ux6790}

回归分析是研究变量间相关关系的一种统计方法。``回归''一词由英国统计学家高尔顿（Galton）在研究父子身高关系时首次提出：他发现高个子父亲的儿子平均身高倾向于低于父亲，而矮个子父亲的儿子平均身高倾向于高于父亲，即子代身高有''回归''到总体均值的趋势。如今，回归分析已成为统计学中应用最广泛的方法之一。

\section{一元线性回归模型}\label{ux4e00ux5143ux7ebfux6027ux56deux5f52ux6a21ux578b}

\subsection{模型形式与基本假设}\label{ux6a21ux578bux5f62ux5f0fux4e0eux57faux672cux5047ux8bbe}

设我们关心响应变量\(Y\)与一个自变量\(X\)之间的关系。一元线性回归模型假定：

\[
Y = \beta_0 + \beta_1 X + \varepsilon
\]

其中\(\beta_0, \beta_1\)是未知的回归系数，\(\varepsilon\)是随机误差项。

给定\(n\)组独立观测值\((x_1, y_1), (x_2, y_2), \ldots, (x_n, y_n)\)，模型可写为：

\[
y_i = \beta_0 + \beta_1 x_i + \varepsilon_i, \quad i = 1, 2, \ldots, n
\]

经典的线性回归模型通常包含以下\textbf{高斯-马尔可夫（Gauss-Markov）假设}：

\begin{enumerate}
\def\labelenumi{\arabic{enumi}.}
\tightlist
\item
  \textbf{线性性}：\(E(Y|X) = \beta_0 + \beta_1 X\)，即响应变量的条件期望是自变量的线性函数。
\item
  \textbf{独立性}：\(\varepsilon_1, \varepsilon_2, \ldots, \varepsilon_n\)相互独立，\(Cov(\varepsilon_i, \varepsilon_j) = 0 \ (i \neq j)\)。
\item
  \textbf{同方差性}：\(Var(\varepsilon_i) = \sigma^2\)对所有的\(i\)都相同。
\item
  \textbf{零均值}：\(E(\varepsilon_i) = 0\)。
\end{enumerate}

若进一步假设\(\varepsilon_i \sim N(0, \sigma^2)\)且相互独立，则得到正态线性回归模型。该假设主要用于区间估计和假设检验。

\subsection{参数的最小二乘估计}\label{ux53c2ux6570ux7684ux6700ux5c0fux4e8cux4e58ux4f30ux8ba1}

我们希望找到\(\beta_0, \beta_1\)的估计\(\hat{\beta}_0, \hat{\beta}_1\)，使得拟合直线\(\hat{y} = \hat{\beta}_0 + \hat{\beta}_1 x\)``最好地''拟合数据。常用的准则是\textbf{最小二乘法}（OLS），即极小化残差平方和：

\[
Q(\beta_0, \beta_1) = \sum_{i=1}^n (y_i - \beta_0 - \beta_1 x_i)^2
\]

对\(Q\)分别求关于\(\beta_0\)和\(\beta_1\)的偏导数并令其为零，得到正规方程：

\[
\begin{cases}
\sum_{i=1}^n (y_i - \hat{\beta}_0 - \hat{\beta}_1 x_i) = 0 \\
\sum_{i=1}^n (y_i - \hat{\beta}_0 - \hat{\beta}_1 x_i) x_i = 0
\end{cases}
\]

求解得：

\[
\hat{\beta}_1 = \frac{\sum_{i=1}^n (x_i - \bar{x})(y_i - \bar{y})}{\sum_{i=1}^n (x_i - \bar{x})^2} = \frac{S_{xy}}{S_{xx}}
\]

\[
\hat{\beta}_0 = \bar{y} - \hat{\beta}_1 \bar{x}
\]

其中\(\bar{x} = \frac{1}{n}\sum x_i\)，\(\bar{y} = \frac{1}{n}\sum y_i\)。

记\(\hat{y}_i = \hat{\beta}_0 + \hat{\beta}_1 x_i\)为第\(i\)个拟合值，\(e_i = y_i - \hat{y}_i\)为第\(i\)个残差。残差是误差\(\varepsilon_i\)的观测体现，是回归诊断的重要工具。

\subsection{最小二乘估计的性质}\label{ux6700ux5c0fux4e8cux4e58ux4f30ux8ba1ux7684ux6027ux8d28}

在高斯-马尔可夫假设下，OLS估计具有以下性质：

\textbf{1、线性性}

\(\hat{\beta}_1\)和\(\hat{\beta}_0\)都是观测值\(y_i\)的线性组合：

\[
\hat{\beta}_1 = \sum_{i=1}^n \frac{x_i - \bar{x}}{S_{xx}} y_i, \quad \hat{\beta}_0 = \sum_{i=1}^n \left[\frac{1}{n} - \frac{\bar{x}(x_i - \bar{x})}{S_{xx}}\right] y_i
\]

\textbf{2、无偏性}

\[
E(\hat{\beta}_0) = \beta_0, \quad E(\hat{\beta}_1) = \beta_1
\]

\textbf{3、方差}

\[
Var(\hat{\beta}_0) = \sigma^2 \left(\frac{1}{n} + \frac{\bar{x}^2}{S_{xx}}\right), \quad Var(\hat{\beta}_1) = \frac{\sigma^2}{S_{xx}}
\]

\[
Cov(\hat{\beta}_0, \hat{\beta}_1) = -\frac{\bar{x} \sigma^2}{S_{xx}}
\]

\textbf{4、高斯-马尔可夫定理}

在所有的线性无偏估计中，OLS估计具有最小方差。也就是说，OLS估计是\textbf{最佳线性无偏估计}（BLUE）。

\begin{quote}
BLUE性质不依赖于正态性假设，仅需Gauss-Markov条件。这是OLS估计被广泛使用的重要理论基础。
\end{quote}

\textbf{5、误差方差的估计}

\(\sigma^2\)的无偏估计为：

\[
\hat{\sigma}^2 = \frac{1}{n-2} \sum_{i=1}^n (y_i - \hat{y}_i)^2 = \frac{SSE}{n-2}
\]

这里除以\(n-2\)而非\(n\)，是因为我们用了两个估计量\(\hat{\beta}_0, \hat{\beta}_1\)代替真实的\(\beta_0, \beta_1\)，损失了两个自由度。

\subsection{回归系数的区间估计与假设检验}\label{ux56deux5f52ux7cfbux6570ux7684ux533aux95f4ux4f30ux8ba1ux4e0eux5047ux8bbeux68c0ux9a8c}

在正态性假设\(\varepsilon_i \sim N(0, \sigma^2)\)下：

\[
\hat{\beta}_1 \sim N\left(\beta_1, \frac{\sigma^2}{S_{xx}}\right), \quad \frac{\hat{\beta}_1 - \beta_1}{\hat{\sigma} / \sqrt{S_{xx}}} \sim t(n-2)
\]

因此，\(\beta_1\)的置信水平为\(1-\alpha\)的置信区间为：

\[
\hat{\beta}_1 \pm t_{1-\alpha/2}(n-2) \cdot \frac{\hat{\sigma}}{\sqrt{S_{xx}}}
\]

对假设检验问题\(H_0: \beta_1 = 0 \ \ vs \ \ H_1: \beta_1 \neq 0\)，检验统计量为：

\[
t = \frac{\hat{\beta}_1}{\hat{\sigma} / \sqrt{S_{xx}}} \sim t(n-2) \quad (\text{当}H_0\text{成立时})
\]

拒绝域为\(W = \{|t| \geq t_{1-\alpha/2}(n-2)\}\)。若拒绝\(H_0\)，则说明自变量\(X\)对响应变量\(Y\)有显著的线性影响。

类似地，对截距项\(\beta_0\)也可以构造t检验：

\[
t = \frac{\hat{\beta}_0}{\hat{\sigma} \sqrt{\frac{1}{n} + \frac{\bar{x}^2}{S_{xx}}}} \sim t(n-2)
\]

\subsection{回归方程的显著性检验：F检验}\label{ux56deux5f52ux65b9ux7a0bux7684ux663eux8457ux6027ux68c0ux9a8cfux68c0ux9a8c}

除了检验单个系数，我们通常还需要检验整个回归方程是否显著。这等价于检验：

\[
H_0: \beta_1 = 0 \quad vs \quad H_1: \beta_1 \neq 0
\]

在第六章方差分析中我们学习了误差分解的思想，这里同样适用。

\textbf{平方和分解}

记\(\bar{y} = \frac{1}{n}\sum y_i\)，总离差平方和（SST）可做如下分解：

\[
\begin{aligned}
SST &= \sum_{i=1}^n (y_i - \bar{y})^2 \\
    &= \sum_{i=1}^n (\hat{y}_i - \bar{y})^2 + \sum_{i=1}^n (y_i - \hat{y}_i)^2 \\
    &= SSR + SSE
\end{aligned}
\]

其中：
- \(SST\)（Total Sum of Squares）：总离差平方和，反映\(y\)的总变差。
- \(SSR\)（Regression Sum of Squares）：回归平方和，反映回归方程能够解释的变差部分。
- \(SSE\)（Error Sum of Squares）：残差平方和，反映回归方程不能解释的变差部分。

在\(H_0\)成立且正态假设下，有：

\[
\frac{SSR}{\sigma^2} \sim \chi^2(1), \quad \frac{SSE}{\sigma^2} \sim \chi^2(n-2)
\]

且\(SSR\)与\(SSE\)相互独立。构造F检验统计量：

\[
F = \frac{SSR / 1}{SSE / (n-2)} = \frac{MSR}{MSE} \sim F(1, n-2)
\]

拒绝域为\(W = \{F \geq F_{1-\alpha}(1, n-2)\}\)。

在一元线性回归中，F检验与\(\beta_1\)的t检验是等价的，因为\(F = t^2\)。

\begin{longtable}[]{@{}lllll@{}}
\caption{一元线性回归的方差分析表}\tabularnewline
\toprule\noalign{}
误差来源 & 平方和 & 自由度 & 均方 & F值 \\
\midrule\noalign{}
\endfirsthead
\toprule\noalign{}
误差来源 & 平方和 & 自由度 & 均方 & F值 \\
\midrule\noalign{}
\endhead
\bottomrule\noalign{}
\endlastfoot
回归 & SSR & 1 & MSR = SSR/1 & F = MSR/MSE \\
残差 & SSE & n-2 & MSE = SSE/(n-2) & \\
总和 & SST & n-1 & & \\
\end{longtable}

\subsection{\texorpdfstring{判定系数 \(R^2\)}{判定系数 R\^{}2}}\label{ux5224ux5b9aux7cfbux6570-r2}

回归方程对数据的拟合优度可以用\textbf{判定系数}来衡量：

\[
R^2 = \frac{SSR}{SST} = 1 - \frac{SSE}{SST}
\]

\(R^2\)取值在\([0, 1]\)之间，越接近1表示回归方程对\(y\)的变差解释得越好。可以证明，在一元线性回归中，\(R^2\)恰好等于\(x\)与\(y\)的样本相关系数的平方：

\[
R^2 = r_{xy}^2 = \left( \frac{S_{xy}}{\sqrt{S_{xx} S_{yy}}} \right)^2
\]

\begin{quote}
注意：\(R^2\)大并不一定意味着模型设定正确。一个很高的\(R^2\)可能来自伪回归，需要结合残差图和其他诊断工具综合判断。
\end{quote}

\subsection{预测}\label{ux9884ux6d4b}

回归方程的一个重要应用是对新观测值进行预测。

\textbf{点预测}：给定新的自变量值\(x_0\)，响应变量的预测值为：

\[
\hat{y}_0 = \hat{\beta}_0 + \hat{\beta}_1 x_0
\]

\textbf{区间预测}分两种：

\begin{enumerate}
\def\labelenumi{\arabic{enumi}.}
\tightlist
\item
  \textbf{均值\(E(Y|x_0)\)的置信区间}
\end{enumerate}

\[
\hat{y}_0 \pm t_{1-\alpha/2}(n-2) \cdot \hat{\sigma} \sqrt{\frac{1}{n} + \frac{(x_0 - \bar{x})^2}{S_{xx}}}
\]

这给出的是回归直线在\(x_0\)处的平均响应值的区间估计。

\begin{enumerate}
\def\labelenumi{\arabic{enumi}.}
\setcounter{enumi}{1}
\tightlist
\item
  \textbf{新观测值\(y_0\)的预测区间}
\end{enumerate}

\[
\hat{y}_0 \pm t_{1-\alpha/2}(n-2) \cdot \hat{\sigma} \sqrt{1 + \frac{1}{n} + \frac{(x_0 - \bar{x})^2}{S_{xx}}}
\]

预测区间比均值置信区间更宽，因为它不仅包含了估计回归直线的误差，还包含了单个观测值本身的随机波动。当\(x_0\)远离\(\bar{x}\)时，两种区间的宽度都会增大，这是因为在数据中心附近回归直线的估计最为精确。

\subsection{残差分析}\label{ux6b8bux5deeux5206ux6790}

残差\(e_i = y_i - \hat{y}_i\)是模型诊断的核心工具。常见的残差图包括：

\begin{itemize}
\tightlist
\item
  \textbf{残差-拟合值图}：以\(\hat{y}_i\)为横轴，\(e_i\)为纵轴。如果模型假设正确，残差应随机散布在\(y=0\)附近，无明显趋势。若呈漏斗形，提示方差不齐；若呈曲线趋势，提示模型形式设定不当。
\item
  \textbf{残差正态QQ图}：检验误差正态性假设。若点大致落在直线上，则正态性合理。
\item
  \textbf{残差-自变量图}：检查是否遗漏了自变量的非线性效应。
\end{itemize}

\section{多元线性回归模型}\label{ux591aux5143ux7ebfux6027ux56deux5f52ux6a21ux578b}

\subsection{模型与矩阵表示}\label{ux6a21ux578bux4e0eux77e9ux9635ux8868ux793a}

实际问题中，响应变量往往受多个自变量的影响。设\(Y\)与\(p-1\)个自变量\(X_1, X_2, \ldots, X_{p-1}\)有关，多元线性回归模型为：

\[
Y = \beta_0 + \beta_1 X_1 + \beta_2 X_2 + \cdots + \beta_{p-1} X_{p-1} + \varepsilon
\]

对\(n\)组独立观测值，采用矩阵记号可以大大简化表达。令：

\[
\mathbf{y} = \begin{bmatrix} y_1 \\ y_2 \\ \vdots \\ y_n \end{bmatrix}, \quad
\mathbf{X} = \begin{bmatrix}
1 & x_{11} & \cdots & x_{1,p-1} \\
1 & x_{21} & \cdots & x_{2,p-1} \\
\vdots & \vdots & \ddots & \vdots \\
1 & x_{n1} & \cdots & x_{n,p-1}
\end{bmatrix}, \quad
\boldsymbol{\beta} = \begin{bmatrix} \beta_0 \\ \beta_1 \\ \vdots \\ \beta_{p-1} \end{bmatrix}, \quad
\boldsymbol{\varepsilon} = \begin{bmatrix} \varepsilon_1 \\ \varepsilon_2 \\ \vdots \\ \varepsilon_n \end{bmatrix}
\]

则模型可简洁地写为：

\[
\mathbf{y} = \mathbf{X} \boldsymbol{\beta} + \boldsymbol{\varepsilon}
\]

其中\(\mathbf{X}\)是\(n \times p\)的设计矩阵（第一列全为1对应于截距项）。

基本假设为：
1. \(E(\boldsymbol{\varepsilon}) = \mathbf{0}\)
2. \(Cov(\boldsymbol{\varepsilon}) = \sigma^2 \mathbf{I}_n\)（独立同方差）
3. 若进一步假设\(\boldsymbol{\varepsilon} \sim N(\mathbf{0}, \sigma^2 \mathbf{I}_n)\)，则得到正态线性模型。

\subsection{参数的最小二乘估计}\label{ux53c2ux6570ux7684ux6700ux5c0fux4e8cux4e58ux4f30ux8ba1-1}

目标仍然是极小化残差平方和：

\[
Q(\boldsymbol{\beta}) = (\mathbf{y} - \mathbf{X}\boldsymbol{\beta})'(\mathbf{y} - \mathbf{X}\boldsymbol{\beta}) = \|\mathbf{y} - \mathbf{X}\boldsymbol{\beta}\|^2
\]

对\(\boldsymbol{\beta}\)求导并令其为零，得到正规方程：

\[
\mathbf{X}'\mathbf{X} \hat{\boldsymbol{\beta}} = \mathbf{X}'\mathbf{y}
\]

当\(\mathbf{X}'\mathbf{X}\)可逆时（即\(\mathbf{X}\)满列秩，自变量之间不存在完全的线性相关），OLS估计为：

\[
\hat{\boldsymbol{\beta}} = (\mathbf{X}'\mathbf{X})^{-1} \mathbf{X}'\mathbf{y}
\]

拟合值向量和残差向量分别为：

\[
\hat{\mathbf{y}} = \mathbf{X}\hat{\boldsymbol{\beta}} = \mathbf{X}(\mathbf{X}'\mathbf{X})^{-1}\mathbf{X}'\mathbf{y} \triangleq \mathbf{H}\mathbf{y}
\]

\[
\mathbf{e} = \mathbf{y} - \hat{\mathbf{y}} = (\mathbf{I} - \mathbf{H})\mathbf{y}
\]

其中\(\mathbf{H} = \mathbf{X}(\mathbf{X}'\mathbf{X})^{-1}\mathbf{X}'\)称为帽子矩阵（hat matrix），它将观测值\(\mathbf{y}\)投影到\(\mathbf{X}\)的列空间上。\(\mathbf{H}\)是对称幂等矩阵，满足\(\mathbf{H}^2 = \mathbf{H}\)。

\subsection{OLS估计的统计性质}\label{olsux4f30ux8ba1ux7684ux7edfux8ba1ux6027ux8d28}

在高斯-马尔可夫假设下：

\textbf{1、无偏性}

\[
E(\hat{\boldsymbol{\beta}}) = \boldsymbol{\beta}
\]

\textbf{2、协方差矩阵}

\[
Cov(\hat{\boldsymbol{\beta}}) = \sigma^2 (\mathbf{X}'\mathbf{X})^{-1}
\]

\textbf{3、高斯-马尔可夫定理}

\(\hat{\boldsymbol{\beta}}\)是\(\boldsymbol{\beta}\)的最佳线性无偏估计（BLUE）：对任意\(p\)维向量\(\mathbf{c}\)，在所有线性无偏估计中，\(\mathbf{c}'\hat{\boldsymbol{\beta}}\)具有最小方差。

\textbf{4、误差方差的无偏估计}

\[
\hat{\sigma}^2 = \frac{\mathbf{e}'\mathbf{e}}{n-p} = \frac{SSE}{n-p}
\]

\textbf{5、正态假设下的分布}

若\(\boldsymbol{\varepsilon} \sim N(\mathbf{0}, \sigma^2 \mathbf{I}_n)\)，则：

\[
\hat{\boldsymbol{\beta}} \sim N(\boldsymbol{\beta}, \sigma^2 (\mathbf{X}'\mathbf{X})^{-1})
\]

\[
\frac{(n-p)\hat{\sigma}^2}{\sigma^2} = \frac{SSE}{\sigma^2} \sim \chi^2(n-p)
\]

且\(\hat{\boldsymbol{\beta}}\)与\(\hat{\sigma}^2\)相互独立。

\subsection{回归方程的显著性检验}\label{ux56deux5f52ux65b9ux7a0bux7684ux663eux8457ux6027ux68c0ux9a8c}

对于多元线性回归，首先检验回归方程整体是否显著：

\[
H_0: \beta_1 = \beta_2 = \cdots = \beta_{p-1} = 0 \quad vs \quad H_1: \text{至少有一个} \beta_j \neq 0
\]

平方和分解同样成立：

\[
SST = SSR + SSE
\]

\[
F = \frac{SSR / (p-1)}{SSE / (n-p)} = \frac{MSR}{MSE} \sim F(p-1, n-p)
\]

\begin{longtable}[]{@{}lllll@{}}
\caption{多元线性回归的方差分析表}\tabularnewline
\toprule\noalign{}
误差来源 & 平方和 & 自由度 & 均方 & F值 \\
\midrule\noalign{}
\endfirsthead
\toprule\noalign{}
误差来源 & 平方和 & 自由度 & 均方 & F值 \\
\midrule\noalign{}
\endhead
\bottomrule\noalign{}
\endlastfoot
回归 & SSR & p-1 & MSR = SSR/(p-1) & F = MSR/MSE \\
残差 & SSE & n-p & MSE = SSE/(n-p) & \\
总和 & SST & n-1 & & \\
\end{longtable}

若\(F \geq F_{1-\alpha}(p-1, n-p)\)，则拒绝\(H_0\)，认为回归方程在整体上是显著的。

\subsection{回归系数的显著性检验}\label{ux56deux5f52ux7cfbux6570ux7684ux663eux8457ux6027ux68c0ux9a8c}

对单个回归系数\(\beta_j\)进行检验：

\[
H_0: \beta_j = 0 \quad vs \quad H_1: \beta_j \neq 0
\]

检验统计量为：

\[
t_j = \frac{\hat{\beta}_j}{\hat{\sigma} \sqrt{(\mathbf{X}'\mathbf{X})^{-1}_{jj}}} \sim t(n-p)
\]

其中\((\mathbf{X}'\mathbf{X})^{-1}_{jj}\)是\((\mathbf{X}'\mathbf{X})^{-1}\)的第\(j\)个对角元。

拒绝域为\(W = \{|t_j| \geq t_{1-\alpha/2}(n-p)\}\)。

\(\beta_j\)的\(1-\alpha\)置信区间为：

\[
\hat{\beta}_j \pm t_{1-\alpha/2}(n-p) \cdot \hat{\sigma} \sqrt{(\mathbf{X}'\mathbf{X})^{-1}_{jj}}
\]

\subsection{判定系数与调整判定系数}\label{ux5224ux5b9aux7cfbux6570ux4e0eux8c03ux6574ux5224ux5b9aux7cfbux6570}

多元回归中仍用\(R^2 = SSR/SST\)衡量拟合优度。但\(R^2\)的一个缺陷是：随着自变量的增加，\(R^2\)永远不会减小，即使加入的是无关变量。为此引入\textbf{调整\(R^2\)}（Adjusted \(R^2\)）：

\[
R^2_{adj} = 1 - \frac{SSE / (n-p)}{SST / (n-1)} = 1 - \frac{n-1}{n-p}(1 - R^2)
\]

调整\(R^2\)对模型中参数个数进行了''惩罚''。当增加一个对模型贡献不大的变量时，\(SSE\)减少不多但\(p\)增大，\(R^2_{adj}\)可能会下降。因此，\(R^2_{adj}\)常用于比较不同自变量个数的模型。

\subsection{多重共线性}\label{ux591aux91cdux5171ux7ebfux6027}

当自变量之间存在近似的线性关系时，称模型存在\textbf{多重共线性}。多重共线性会导致\(|\mathbf{X}'\mathbf{X}| \approx 0\)，使得\((\mathbf{X}'\mathbf{X})^{-1}\)的对角元很大，从而OLS估计的方差膨胀，估计变得不稳定。

常用的诊断指标是\textbf{方差膨胀因子}（VIF）：

\[
VIF_j = \frac{1}{1 - R_j^2}
\]

其中\(R_j^2\)是以\(X_j\)为响应变量，其余\(p-2\)个自变量为解释变量做回归所得的判定系数。

经验上，当\(VIF_j > 10\)（即\(R_j^2 > 0.9\)）时，认为存在严重的多重共线性问题。

处理多重共线性的常用方法包括：删除相关性高的变量、增大样本量、使用岭回归或有偏估计方法、主成分回归等。

\subsection{变量选择}\label{ux53d8ux91cfux9009ux62e9}

当备选的自变量较多时，需要从中选出一个''最优''子集。常用的准则包括：

\begin{itemize}
\tightlist
\item
  \textbf{调整\(R^2\)}：选择使\(R^2_{adj}\)最大的模型。
\item
  \textbf{AIC准则}（Akaike Information Criterion）：\(AIC = n \ln(SSE/n) + 2p\)，选择AIC最小的模型。
\item
  \textbf{BIC准则}（Bayesian Information Criterion）：\(BIC = n \ln(SSE/n) + p \ln n\)，比AIC对模型复杂度惩罚更重。
\end{itemize}

常用的变量选择策略有：

\begin{itemize}
\tightlist
\item
  \textbf{向前选择}（Forward Selection）：从空模型开始，每次加入一个最显著的自变量，直到没有变量显著为止。
\item
  \textbf{向后剔除}（Backward Elimination）：从全模型开始，每次剔除一个最不显著的变量，直到所有变量都显著为止。
\item
  \textbf{逐步回归}（Stepwise Regression）：结合向前选择和向后剔除，每加入一个变量后检查是否需剔除已入选的变量。
\end{itemize}

\begin{quote}
注意：完全依赖自动变量选择可能产生有偏的结果（系数估计绝对值偏大、p值偏小），应结合实际问题背景和专业知识来选定最终模型。
\end{quote}

\subsection{其他常见问题}\label{ux5176ux4ed6ux5e38ux89c1ux95eeux9898}

\textbf{1、异方差性}

若\(Var(\varepsilon_i) \neq Var(\varepsilon_j)\)，即存在异方差，OLS估计仍是无偏的，但不再具有最小方差。此时可采用加权最小二乘法（WLS），对观测值按其方差的倒数加权，或使用异方差稳健的标准误（如White's sandwich estimator）。

\textbf{2、自相关性}

若误差项之间存在相关性（尤其常见于时间序列数据），OLS估计仍无偏但效率降低，且标准误的估计有偏。Durbin-Watson检验是常用的自相关诊断方法。处理自相关的常用方法包括广义最小二乘法（GLS）和引入滞后项。

\textbf{3、异常值与影响点}

某些观测点可能对回归结果产生不成比例的影响。常用诊断统计量包括：

\begin{itemize}
\tightlist
\item
  \textbf{标准化残差}：\(r_i = \frac{e_i}{\hat{\sigma}\sqrt{1 - h_{ii}}}\)，其中\(h_{ii}\)是帽子矩阵的第\(i\)个对角元（杠杆值）。\(|r_i| > 2\)或\(3\)的点值得关注。
\item
  \textbf{Cook距离}：衡量删除某个观测点后回归系数估计的变化程度。\(D_i > 4/n\)时需警惕。
\end{itemize}

\section{非线性回归模型}\label{ux975eux7ebfux6027ux56deux5f52ux6a21ux578b}

在实际问题中，\(Y\)与\(X\)之间的关系往往不是线性的。根据能否通过变量变换转化为线性形式，分为两类。

\subsection{可线性化的模型}\label{ux53efux7ebfux6027ux5316ux7684ux6a21ux578b}

许多非线性模型可以通过适当的变换转化为线性模型后，仍用OLS估计。

\textbf{1、多项式回归}

\[
Y = \beta_0 + \beta_1 X + \beta_2 X^2 + \cdots + \beta_k X^k + \varepsilon
\]

令\(Z_1 = X, Z_2 = X^2, \ldots, Z_k = X^k\)，即可化为多元线性回归。

\begin{quote}
多项式回归仍属于线性模型的范畴，因为它是关于\textbf{参数}\(\beta_j\)线性的，而非关于自变量\(X\)线性。当\(k\)选择过大时容易出现过拟合，通常\(k\)取2或3即可。
\end{quote}

\textbf{2、对数模型（幂函数模型）}

\[
Y = \alpha X^{\beta_1} e^{\varepsilon}
\]

两边取对数：\(\ln Y = \ln \alpha + \beta_1 \ln X + \varepsilon\)，化为线性模型。

\textbf{3、指数模型}

\[
Y = \alpha e^{\beta_1 X} \cdot \varepsilon
\]

两边取对数：\(\ln Y = \ln \alpha + \beta_1 X + \ln \varepsilon\)。

\textbf{4、倒数变换}

\[
Y = \beta_0 + \frac{\beta_1}{X} + \varepsilon
\]

令\(Z = 1/X\)即可。

\textbf{5、Box-Cox变换}

Box-Cox变换是一族幂变换，用于对响应变量进行变换以改善正态性和方差齐性：

\[
Y^{(\lambda)} = \begin{cases}
\dfrac{Y^\lambda - 1}{\lambda}, & \lambda \neq 0 \\
\ln Y, & \lambda = 0
\end{cases}
\]

其中\(\lambda\)是需要估计的变换参数，可通过最大似然方法估计。

\subsection{一般非线性回归}\label{ux4e00ux822cux975eux7ebfux6027ux56deux5f52}

对于本质上无法线性化的模型（如\(Y = \frac{\alpha}{1 + e^{\beta_0 + \beta_1 X}} + \varepsilon\)），需要采用非线性最小二乘（NLS）方法。

非线性模型的一般形式为：

\[
y_i = f(x_i; \boldsymbol{\theta}) + \varepsilon_i
\]

残差平方和为\(Q(\boldsymbol{\theta}) = \sum [y_i - f(x_i; \boldsymbol{\theta})]^2\)。由于\(f\)关于\(\boldsymbol{\theta}\)是非线性的，通常没有解析解，需要使用迭代优化方法。

\textbf{Gauss-Newton迭代法}是最常用的方法之一。其基本思想是将\(f\)在当前的参数估计值附近做一阶泰勒展开，将非线性问题近似为线性问题求解，然后反复迭代。

记第\(m\)步的参数值为\(\boldsymbol{\theta}^{(m)}\)，将\(f(x_i; \boldsymbol{\theta})\)在\(\boldsymbol{\theta}^{(m)}\)处展开：

\[
f(x_i; \boldsymbol{\theta}) \approx f(x_i; \boldsymbol{\theta}^{(m)}) + \left[ \frac{\partial f(x_i; \boldsymbol{\theta}^{(m)})}{\partial \boldsymbol{\theta}} \right]' (\boldsymbol{\theta} - \boldsymbol{\theta}^{(m)})
\]

代入残差平方和并求极小，得到迭代公式：

\[
\boldsymbol{\theta}^{(m+1)} = \boldsymbol{\theta}^{(m)} + [\mathbf{J}'^{(m)} \mathbf{J}^{(m)}]^{-1} \mathbf{J}'^{(m)} (\mathbf{y} - \mathbf{f}^{(m)})
\]

其中\(\mathbf{J}^{(m)}\)是\(n \times p\)的雅可比矩阵，第\((i, j)\)元为\(\partial f(x_i; \boldsymbol{\theta}) / \partial \theta_j\)在\(\boldsymbol{\theta}^{(m)}\)处的值；\(\mathbf{f}^{(m)} = (f(x_1; \boldsymbol{\theta}^{(m)}), \ldots, f(x_n; \boldsymbol{\theta}^{(m)}))'\)。

迭代进行直至收敛。非线性回归的一个重要问题是参数初值的选择：好的初值能确保算法收敛到全局最优而非局部最优。通常可以借助先验知识、可线性化的近似模型，或网格搜索来选取初值。

\begin{quote}
值得注意的是，从第十章可以看到，广义线性模型（GLM）将许多看似非线性的回归问题统一在线性模型的框架下，通过连接函数和指数族分布来处理非正态响应变量，这是现代回归分析的重要发展。
\end{quote}

\begin{center}\rule{0.5\linewidth}{0.5pt}\end{center}

\textbf{速查卡片}：

\begin{itemize}
\tightlist
\item
  \textbf{核心问题}：如何用自变量解释和预测响应变量？如何诊断和改进模型？
\item
  \textbf{关键概念}：OLS估计 → Gauss-Markov定理（BLUE） → R²与调整R² → F检验与t检验 → 残差诊断 → 多重共线性（VIF） → 变量选择（AIC/BIC/逐步回归） → 非线性回归（多项式/Box-Cox/Gauss-Newton）
\item
  \textbf{R命令}：\texttt{lm()}, \texttt{summary()}, \texttt{plot(x)} 诊断四图, \texttt{AIC()}, \texttt{step()}, \texttt{nls()}, \texttt{car::vif()}
\item
  \textbf{参见}：Ch06（估计理论）, Ch07（方差分析）, Ch11（GLM扩展）, Ch15（优化方法）
\end{itemize}

\chapter{广义线性回归}\label{ux5e7fux4e49ux7ebfux6027ux56deux5f52}

\section{引言}\label{ux5f15ux8a00}

\section{统一方法}\label{ux7edfux4e00ux65b9ux6cd5}

\subsection{指数族 EF}\label{ux6307ux6570ux65cf-ef}

\textbf{1、指数族的定义}

设\(\mathcal{F} =\{f(x,\theta):\theta\in\Theta\}\)是分布族，若其概率函数\(f(x,\theta)\)可以表示成如下形式 \[
f(x,\theta) = A(\theta)exp\{\sum_{i=1}^k\ Q_i(\theta)\ T_i(x)\}h(x)
\] 则称此分布族为指数型分布族（简称指数族），其中\(A(\theta),Q_i(\theta)\)是定义在参数空间上的函数，\(h(x),T_i(x)\)是定义在样本空间上的函数。

由概率函数的积分为1，可知：

\[
0<\int_{\mathcal{X}}exp\{\sum_{i=1}^k\ Q_i(\theta)\ T_i(x) \}h(x)dx = \frac{1}{A(\theta)}<\infty
\]

称\(G(x)\)为概率函数\(f(x,\theta)\)的支撑集，若\(G(x) = \{x:f(x,\theta)>0\}\)。指数族中所有分布具有相同的支撑集。

\textbf{2、指数族的自然形式}

\[
f(x,\theta) = C(\theta)exp\{\sum_{i=1}^k\theta_iT_i(x) \}h(x)
\]

自然参数空间表示如下：

\[
\Theta^* = \{(\theta_1,\theta_2,\cdots,\theta_k):\int_{\mathcal{X}}exp\{\sum_{i=1}^k\theta_iT_i(x) \}h(x)dx<\infty \}
\]

\textbf{3、正态分布族和二项分布族}

正态分布族的例子：

\[
f(x,\theta) = \frac{1}{\sqrt{2\pi}\sigma} e^{-\frac{(x-\mu)^2}{2\sigma^2}} = \frac{1}{\sqrt{2\pi}\sigma} e^{-\frac{\mu^2}{2\sigma^2}} e^{\frac{\mu}{\sigma^2}x-\frac{1}{2\sigma^2}x^2}
\]

令\(\phi_1 = \frac{\mu}{\sigma^2},\phi_2 =-\frac{1}{2\sigma^2}\)，则 \[
\frac{1}{\sqrt{2\pi}\sigma}e^{-\frac{\mu^2}{2\sigma^2}} = \sqrt{-\frac{\phi_2}{\pi}}e^{\frac{\phi_1^2}{4\phi_2}} = C(\phi)
\] 因此，正态分布的自然形式为：

\[
f(x,\phi) = \sqrt{-\frac{\phi_2}{\pi}}e^{\frac{\phi_1^2}{4\phi_2}} e^{\phi_1x + \phi_2x^2} = C(\phi)e^{\phi_1t_1(x)+\phi_2t_2(x)}h(x)
\]

二项分布族的例子：

\[
f(x,\theta) = C_n^x\theta^x(1-\theta)^{1-x} = (1-\theta)^n\cdot e^{log(\frac{\theta}{1-\theta}\cdot x)}\cdot C_n^x
\]

令\(log(\frac{\theta}{1-\theta}) = \phi\)，则

\[
f(x,\phi) = \frac{1}{(1+e^\phi)^n}\cdot e^{\phi x}\cdot C_n^x = C(\phi)\cdot e^{\phi t_1(x)}h(x)
\]

值得注意，均匀分布和t分布不是指数族分布。

\textbf{4、指数族的性质}

\begin{itemize}
\item
  性质1：在指数族的自然形式下，自然参数空间为凸集。
\item
  性质2：设指数族的自然形式中，自然参数空间有内点，\(g(x)\)是任一有界可积函数，则对 \(G(\theta) = \int g(x)exp\{\sum\theta_jT_j(x) \}h(x)dx\)，积分和求导顺序可以互换。 \[
    \frac{\partial^mG(\theta)}{\partial\theta_1^{m_1}\cdots \partial\theta_k^{m_k}} = \int\frac{\partial^m}{\partial\theta_1^{m_1}\cdots \partial\theta_k^{m_k}}[g(x)exp\{\sum\theta_jT_j(x) \}h(x)]dx
  \]
\item
  推论1：以\(\Theta_0\)记自然参数空间的内点集，则指数族自然形式表达式中的\(C(\theta)\)在\(\Theta_0\)上解析，其对各\(\theta_j\)的偏导数可通过在积分号下求导算出。
\item
  推论2：若\(\theta\in\Theta_0\)，则在分布\(P_{\theta}\)下，\(T_j(x)\)的各阶矩都存在有限，且可通过积分号下求导算出。
\end{itemize}

\subsection{指数分散族 EDF}\label{ux6307ux6570ux5206ux6563ux65cf-edf}

\textbf{1、指数分散族定义}

对\textbf{单一兴趣参数}的指数族分布，其随机变量Y的概率函数可表示为：

\[
p(y,\theta)=exp[\eta(\theta)t(y)+log(s(\theta))+log(h(y))]
\] 经过转换，可以得到：

\[
p(y,\theta) = exp[\theta y-b(\theta)+c(y)]
\]

如果还有冗余参数不能简单忽略，则可以进一步改写为更一般的形式：

\[
p(y;\theta,\phi)=exp[\frac{y\theta-b(\theta)}{\phi}+c(y,\phi)]
\]

其中：

\begin{itemize}
\item
  \(\theta\)为典型参数，常为兴趣参数或均值参数的函数。与Y的均值密切关联。
\item
  \(b(\theta)\)称为累积量函数，因其不含Y也常称为归一化常数，在计算分布的矩方面有重要作用。
\item
  \(\phi\)为分散参数，其常与Y的方差密切相关。
\end{itemize}

\textbf{2、指数分散族的性质}

设指数族随机变量Y的概率函数表示为\(p(y;\theta,\phi)=exp[\frac{y\theta-b(\theta)}{\phi}+c(y,\phi)]\)，对数似然函数\(l(\theta,\phi;y)= \frac{y\theta-b(\theta)}{\phi}+c(y,\phi)\)，则有：

\[
a.\ E(\frac{\partial l}{\partial\theta})=0;\ \ b.E(\frac{\partial^2l}{\partial\theta^2})+E((\frac{\partial l }{\partial\theta})^2)=0;\ \ c.Var(\frac{\partial l }{\partial\theta})=-E[\frac{\partial^2 l }{\partial\theta^2}]
\]

指数分散族得分函数的性质推导如下：

对对数似然函数求导：\(\frac{\partial l}{\partial \theta} = \frac{y - b'(\theta)}{\phi}\)。

由\(E(\frac{\partial l}{\partial \theta}) = 0\)可得\(E(Y) = b'(\theta)\)。

由\(\frac{\partial^2 l}{\partial \theta^2} = -\frac{b''(\theta)}{\phi}\)可得\(Var(\frac{\partial l}{\partial \theta}) = -E(\frac{\partial^2 l}{\partial \theta^2}) = \frac{b''(\theta)}{\phi}\)，因此\(Var(Y) = \phi b''(\theta)\)。

进一步，\(b(\theta)\)与归一化常数的关系为：\(b(\theta) = \phi \log \int \exp\left(\frac{y\theta}{\phi} + c(y,\phi)\right) dy\)。对\(b(\theta)\)求导：\(b'(\theta) = E(Y)\)，\(b''(\theta) = Var(Y)/\phi\)。\(b''(\theta)\)被称为方差函数，记为\(V(\mu)\)，它刻画了均值-方差关系。

\[
(a)\frac{b(\theta)}{\phi}=log[\int exp(\frac{y\theta}{\phi}+c(y,\phi))dy];\ \ (b)E(Y)=b'(\theta);\ \ (c)Var(Y)=\phi b''(\theta)
\]

\textbf{指数分散族常见分布汇总表：}

\begin{longtable}[]{@{}
  >{\raggedright\arraybackslash}p{(\columnwidth - 10\tabcolsep) * \real{0.1667}}
  >{\raggedright\arraybackslash}p{(\columnwidth - 10\tabcolsep) * \real{0.1667}}
  >{\raggedright\arraybackslash}p{(\columnwidth - 10\tabcolsep) * \real{0.1667}}
  >{\raggedright\arraybackslash}p{(\columnwidth - 10\tabcolsep) * \real{0.1667}}
  >{\raggedright\arraybackslash}p{(\columnwidth - 10\tabcolsep) * \real{0.1667}}
  >{\raggedright\arraybackslash}p{(\columnwidth - 10\tabcolsep) * \real{0.1667}}@{}}
\toprule\noalign{}
\begin{minipage}[b]{\linewidth}\raggedright
分布
\end{minipage} & \begin{minipage}[b]{\linewidth}\raggedright
\(\theta\)
\end{minipage} & \begin{minipage}[b]{\linewidth}\raggedright
\(\phi\)
\end{minipage} & \begin{minipage}[b]{\linewidth}\raggedright
\(b(\theta)\)
\end{minipage} & \begin{minipage}[b]{\linewidth}\raggedright
\(b'(\theta) = \mu\)
\end{minipage} & \begin{minipage}[b]{\linewidth}\raggedright
\(b''(\theta) = V(\mu)\)
\end{minipage} \\
\midrule\noalign{}
\endhead
\bottomrule\noalign{}
\endlastfoot
正态 \(N(\mu, \sigma^2)\) & \(\mu\) & \(\sigma^2\) & \(\theta^2/2\) & \(\theta\) & \(1\) \\
二项 \(B(m, p)/m\) & \(\log\frac{p}{1-p}\) & \(1/m\) & \(\log(1+e^\theta)\) & \(\frac{e^\theta}{1+e^\theta}\) & \(\mu(1-\mu)\) \\
泊松 \(P(\lambda)\) & \(\log\lambda\) & \(1\) & \(e^\theta\) & \(e^\theta\) & \(\mu\) \\
伽马 \(G(\mu, \nu)\) & \(-1/\mu\) & \(1/\nu\) & \(-\log(-\theta)\) & \(-1/\theta\) & \(\mu^2\) \\
逆高斯 & \(-1/(2\mu^2)\) & \(\phi\) & \(-\sqrt{-2\theta}\) & \(1/\sqrt{-2\theta}\) & \(\mu^3\) \\
\end{longtable}

\section{模型简介}\label{ux6a21ux578bux7b80ux4ecb}

\subsection{定义}\label{ux5b9aux4e49}

广义线性模型具有以下三个部分：

\begin{itemize}
\item
  随机成分：被解释变量\(Y_1,Y_2,\cdots,Y_n\) 来自指数族的同一分布族，具有相同形式的概率函数
\item
  系统成分：线性结构\(\eta_i=x_i'\beta\) 是线性预测的系统成分。
\item
  连接函数：单调连接函数\(g(\cdot)\) 连接随机部分的均值参数或兴趣参数和系统成分
\end{itemize}

\subsection{典型连接}\label{ux5178ux578bux8fdeux63a5}

典型连接\(g=b'^{-1}\) 使得线性系统成分\(\eta\)直接等于典型参数\(\theta\)。

\begin{enumerate}
\def\labelenumi{\arabic{enumi}.}
\tightlist
\item
  二项分布的例子

  \begin{enumerate}
  \def\labelenumii{\arabic{enumii}.}
  \tightlist
  \item
    \[
    f(y;p,m)=C_m^yp^y(1-p)^{m-y}=exp[ylog(\frac{p}{1-p})+mlog(1-p)+log(C_m^y)]
    \]
  \item
    典型参数\(\theta = log(\frac{p}{1-p})\) ；分散参数\(\phi=1\) ；累计函数\(b(\theta)=-mlog(1-p)=mlog(1+e^{\theta})\)
  \item
    均值\(\mu = b'(\theta) = m\frac{e^{\theta}}{1+e^{\theta}}=mp\) ；方差函数\(v(\mu)=b''(\theta)=mp(1-p)\) ；方差\(Var(Y)=mp(1-p)\)
  \item
    典型连接：\(\mu=\theta=b'^{-1}(\mu)=log(\frac{p}{1-p})=logit(p)\)；
  \item
    \[
    Y_i\sim B(m_i,p_i),l(\beta)=\sum_{i=1}^n[y_i\eta_i-m_ilog(1+e_i^{\eta})+log(C_{m_i}^{y_i})]
    \]
  \end{enumerate}
\end{enumerate}

\section{估计}\label{ux4f30ux8ba1}

\subsection{Newton-Raphson算法}\label{newton-raphsonux7b97ux6cd5}

迭代方程：

\[
\beta^{(m)}=\beta^{(m-1)}+[J(\beta^{m-1})]^{-1}U(\beta^{m-1})
\]

其中，\(U(\beta)=l'(\beta)\) 是对数似然函数关于参数向量的一阶导数（得分函数，记分函数）;\(J(\beta)=-E[U'(\beta)]\) 为信息矩阵

\subsection{极大似然估计}\label{ux6781ux5927ux4f3cux7136ux4f30ux8ba1}

对于所有观测值的对数似然函数为：

\[
l=\sum_{i=1}^nl_i=\sum_{i=1}^n[\frac{y_i\theta_i-b(\theta_i)}{\phi}+c(y_i,\phi)]
\]

对数似然函数关于参数\(\beta_j\) 的得分函数为：

\[
U_j=\frac{\partial l}{\partial \beta_j}=\sum_{i=1}^n\frac{\partial l_i}{\partial\beta_j}=\sum_{i=1}^n\frac{\partial l_i}{\partial\theta_i}\frac{\partial\theta_i}{\partial\mu_i}\frac{\partial\mu_i}{\partial \eta_i}\frac{\partial\eta_i}{\partial\beta_j}=\sum_{i=1}^n\frac{y_i-\mu_i}{\phi b''(\theta_i)}\frac{d\mu_i}{d\eta_i}x_{ij}
\]

得分向量的期望为零，则信息矩阵可表示为：

\[
J_{jk}=E(U_jU_k)=\sum_{i=1}^n\frac{x_{ij}x_{ik}}{\phi b''(\theta_i)}(\frac{d\mu_i}{d\eta_i})^2
\]

进一步，可写成矩阵形式：

\[
U=\frac{1}{\phi}X'WD\ \ ;\ \ J=\frac{1}{\phi}X'WX
\]

其中，\(W\)为对角权重矩阵，其对角元素为\(w_i = \frac{1}{\phi b''(\theta_i)} \left(\frac{d\mu_i}{d\eta_i}\right)^2\)；\(D\)为偏差向量，其分量为\(D_i = (y_i - \mu_i) \frac{d\eta_i}{d\mu_i}\)。注意\(W\)依赖于当前的参数估计值，因此需要迭代求解。

得到最终的迭代方程：

\[
\beta^{(m)}=(X'W^{m-1}X)^{-1}X'W^{m-1}Z^{m-1},Z^{m-1}=X\beta^{m-1}+D^{m-1}
\]

该迭代过程即为\textbf{迭代再加权最小二乘法（IRLS）}。其核心思想是：每一步迭代，使用当前的参数估计值计算加权矩阵\(W\)和调整响应变量\(Z\)，然后求解加权最小二乘问题更新\(\beta\)。该过程不断重复，直到参数估计收敛。这一算法的优雅之处在于将非线性极大似然估计问题转化为一系列加权线性回归问题。

\section{推断}\label{ux63a8ux65ad}

\subsection{Wald检验}\label{waldux68c0ux9a8c}

问题：对于指数分散族形式的概率函数，分散参数𝜙已知，如何基于近似得分函数得到极大似然估计的抽样分布？

\(\hat{\beta}\)的渐进分布为：\((\hat{\beta}-\beta)'J(\beta)(\hat{\beta}-\beta)\sim \chi^2(p)\)

\subsection{偏差统计量}\label{ux504fux5deeux7edfux8ba1ux91cf}

首先引入兴趣模型和饱和模型的概念。定义对数似然比：

\[
log\lambda = l(\hat{\beta}_m;y)-l(\hat{\beta};y)
\]

显然，\textbf{饱和模型对数据的描述最全面}，对数似然比越大说明兴趣模型对数据的描述相对更差。

含分散参数\(\phi\)的定标偏差定义为（偏差D不含\(\phi\)）：

\[
D^*=\frac{D}{\phi}=2log(\lambda) = 2(l(\hat{\beta}_m;y)-l(\hat{\beta};y))
\]

定标偏差服从非中心卡方分布：

\[
D^*\sim \chi^2(n-p,v),D\sim \phi \chi^2(n-p,v)
\]

\subsection{偏差之差}\label{ux504fux5deeux4e4bux5dee}

对于两个嵌套关系模型，其中一个模型\(M_0\)是另一个更一般化模型\(M_1\)的特殊形式（它们有共同得饱和模型），可以通过偏差之差判定哪个模型对数据拟合得更好。

定标偏差之差： \[
\Delta D^*=D_0^*-D_1^*=2(l(\hat{\beta}_m;y)-l(\hat{\beta_0};y))-2(l(\hat{\beta}_m;y)-l(\hat{\beta_1};y))=2(l(\hat{\beta}_0;y)-l(\hat{\beta_1};y))
\]

统计假设：

\[
H_0:\beta_{q+1}=\beta_{q+2}=\cdots=\beta_p=0
\]

如果模型\(M_0\) 和\(M_1\) 对数据描述都较好，一般选择对应原假设\(H_0\) 的简洁模型\(M_0\)。

如果模型\(M_0\) 对数据描述不好，同时模型\(M_1\) 对数据描述较好，此时认为模型\(M_1\) 明显更适合描述数据。

\subsection{残差}\label{ux6b8bux5dee}

\subsubsection{Pearson残差}\label{pearsonux6b8bux5dee}

\[
r_i^P=\frac{y_i-\hat{\mu_i}}{\sqrt{v(\hat{\mu}_i)}}=\frac{y_i-b'(\hat{\theta}_i)}{\sqrt{b''(\hat{\theta}_i)}}
\]

正态分布：

二项分布：

泊松分布：

负二项分布：

伽马分布：

\subsubsection{偏差残差}\label{ux504fux5deeux6b8bux5dee}

\[
r_i^D=sign(y_i-\hat{\mu}_i)\sqrt{D_i}
\]

其中，\(sign(y_i-\hat{\mu}_i)=\frac{y_i-\hat{\mu}_i}{|y_i-\hat{\mu}_i|}\)，

各常见分布的偏差残差分量\(D_i\)（满足\(D = \sum D_i\)）的计算公式：

\begin{itemize}
\tightlist
\item
  \textbf{正态分布}：\(D_i = (y_i - \hat{\mu}_i)^2\)
\item
  \textbf{二项分布}：\(D_i = 2\left[y_i \log\frac{y_i}{\hat{\mu}_i} + (m_i - y_i) \log\frac{m_i - y_i}{m_i - \hat{\mu}_i}\right]\)
\item
  \textbf{泊松分布}：\(D_i = 2\left[y_i \log\frac{y_i}{\hat{\mu}_i} - (y_i - \hat{\mu}_i)\right]\)
\item
  \textbf{负二项分布}：\(D_i = 2\left[y_i \log\frac{y_i}{\hat{\mu}_i} - (y_i + \kappa) \log\frac{y_i + \kappa}{\hat{\mu}_i + \kappa}\right]\)
\item
  \textbf{伽马分布}：\(D_i = 2\left[-\log\frac{y_i}{\hat{\mu}_i} + \frac{y_i - \hat{\mu}_i}{\hat{\mu}_i}\right]\)
\end{itemize}

\section{模型应用}\label{ux6a21ux578bux5e94ux7528}

常用广义线性模型包括二分类响应，计数响应，比例响应以及多分类响应模型。

\subsection{二分类响应}\label{ux4e8cux5206ux7c7bux54cdux5e94}

\subsubsection{二点响应}\label{ux4e8cux70b9ux54cdux5e94}

二点响应模型（Bernoulli响应）：设\(Y_i \in \{0, 1\}\)，\(P(Y_i = 1) = p_i\)。其概率函数可写为指数分散族形式：

\[f(y_i; p_i) = p_i^{y_i}(1-p_i)^{1-y_i} = \exp\left[y_i \log\frac{p_i}{1-p_i} + \log(1-p_i)\right]\]

其中，典型参数\(\theta_i = \log\frac{p_i}{1-p_i}\)，\(b(\theta_i) = \log(1+e^{\theta_i})\)，\(\phi = 1\)。均值\(\mu_i = p_i\)，方差函数\(V(\mu_i) = p_i(1-p_i)\)。

\subsubsection{二项响应}\label{ux4e8cux9879ux54cdux5e94}

二项响应模型：设\(Y_i \sim B(m_i, p_i)\)，其中\(m_i\)为第\(i\)次观测的试验次数。其概率函数为：

\[f(y_i; p_i, m_i) = \binom{m_i}{y_i} p_i^{y_i} (1-p_i)^{m_i - y_i} = \exp\left[y_i \log\frac{p_i}{1-p_i} + m_i \log(1-p_i) + \log\binom{m_i}{y_i}\right]\]

此时\(\theta_i = \log\frac{p_i}{1-p_i}\)，\(b(\theta_i) = m_i \log(1+e^{\theta_i})\)，\(\phi = 1\)。均值\(\mu_i = m_i p_i\)，方差函数\(V(\mu_i) = m_i p_i (1-p_i) = \mu_i(1 - \mu_i/m_i)\)。在实际建模中，通常对比例\(\hat{p}_i = Y_i/m_i\)建模，使用权重\(m_i\)处理不等精度。

挑战者算例回顾。

\begin{Shaded}
\begin{Highlighting}[]
\CommentTok{\# 数据导入}
\FunctionTok{data}\NormalTok{(orings, }\AttributeTok{package =} \StringTok{"faraway"}\NormalTok{)}
\NormalTok{orings}\SpecialCharTok{$}\NormalTok{total }\OtherTok{\textless{}{-}} \DecValTok{6}  \CommentTok{\#定义每个观测的 O 型环总数为 6}
\CommentTok{\# 拟合二项式GLM}
\NormalTok{lmod }\OtherTok{\textless{}{-}} \FunctionTok{glm}\NormalTok{(}\FunctionTok{cbind}\NormalTok{(damage, total}\SpecialCharTok{{-}}\NormalTok{damage) }\SpecialCharTok{\textasciitilde{}} \DecValTok{1} \SpecialCharTok{+}\NormalTok{ temp, }\AttributeTok{family=}\NormalTok{binomial, orings)}
\FunctionTok{summary}\NormalTok{(lmod)}\CommentTok{\#输出系数、标准误、z 统计量、p 值，以及残差偏差、AIC 等}
\end{Highlighting}
\end{Shaded}

\begin{verbatim}
## 
## Call:
## glm(formula = cbind(damage, total - damage) ~ 1 + temp, family = binomial, 
##     data = orings)
## 
## Coefficients:
##             Estimate Std. Error z value Pr(>|z|)    
## (Intercept) 11.66299    3.29626   3.538 0.000403 ***
## temp        -0.21623    0.05318  -4.066 4.78e-05 ***
## ---
## Signif. codes:  0 '***' 0.001 '**' 0.01 '*' 0.05 '.' 0.1 ' ' 1
## 
## (Dispersion parameter for binomial family taken to be 1)
## 
##     Null deviance: 38.898  on 22  degrees of freedom
## Residual deviance: 16.912  on 21  degrees of freedom
## AIC: 33.675
## 
## Number of Fisher Scoring iterations: 6
\end{verbatim}

\begin{Shaded}
\begin{Highlighting}[]
\FunctionTok{cat}\NormalTok{(}\StringTok{"\#\#\#\#\#\#\#\#\#\#"}\NormalTok{)}
\end{Highlighting}
\end{Shaded}

\begin{verbatim}
## ##########
\end{verbatim}

\begin{Shaded}
\begin{Highlighting}[]
\CommentTok{\# 手动计算拟合值}
\NormalTok{beta }\OtherTok{\textless{}{-}} \FunctionTok{coef}\NormalTok{(lmod)}\CommentTok{\#提取系数}
\NormalTok{x }\OtherTok{\textless{}{-}} \FunctionTok{matrix}\NormalTok{(}\DecValTok{1}\NormalTok{, }\AttributeTok{nrow =} \DecValTok{23}\NormalTok{, }\AttributeTok{ncol =} \DecValTok{2}\NormalTok{)}\CommentTok{\#构造设计矩阵}
\NormalTok{x[ ,}\DecValTok{2}\NormalTok{] }\OtherTok{\textless{}{-}}\NormalTok{ orings}\SpecialCharTok{$}\NormalTok{temp}
\NormalTok{eta }\OtherTok{\textless{}{-}}\NormalTok{ x}\SpecialCharTok{\%*\%}\NormalTok{beta}
\NormalTok{phat }\OtherTok{\textless{}{-}} \FunctionTok{exp}\NormalTok{(eta)}\SpecialCharTok{/}\NormalTok{(}\DecValTok{1} \SpecialCharTok{+} \FunctionTok{exp}\NormalTok{(eta)) }\DocumentationTok{\#\# fitted(lmod)}
\CommentTok{\# 计算偏差残差}
\NormalTok{d1 }\OtherTok{\textless{}{-}}\NormalTok{ orings}\SpecialCharTok{$}\NormalTok{damage}\SpecialCharTok{*}\FunctionTok{log}\NormalTok{(}\FunctionTok{ifelse}\NormalTok{(orings}\SpecialCharTok{$}\NormalTok{damage}\SpecialCharTok{!=}\DecValTok{0}\NormalTok{,orings}\SpecialCharTok{$}\NormalTok{damage}\SpecialCharTok{/}\NormalTok{(orings}\SpecialCharTok{$}\NormalTok{total}\SpecialCharTok{*}\NormalTok{phat), }\DecValTok{1}\NormalTok{))}
\NormalTok{d2 }\OtherTok{\textless{}{-}}\NormalTok{ (orings}\SpecialCharTok{$}\NormalTok{total}\SpecialCharTok{{-}}\NormalTok{orings}\SpecialCharTok{$}\NormalTok{damage)}\SpecialCharTok{*}\FunctionTok{log}\NormalTok{(}\FunctionTok{ifelse}\NormalTok{(}
\NormalTok{orings}\SpecialCharTok{$}\NormalTok{damage}\SpecialCharTok{!=}\NormalTok{orings}\SpecialCharTok{$}\NormalTok{total,}
\NormalTok{(orings}\SpecialCharTok{$}\NormalTok{total}\SpecialCharTok{{-}}\NormalTok{orings}\SpecialCharTok{$}\NormalTok{damage)}\SpecialCharTok{/}\NormalTok{(orings}\SpecialCharTok{$}\NormalTok{total}\SpecialCharTok{{-}}\NormalTok{orings}\SpecialCharTok{$}\NormalTok{total}\SpecialCharTok{*}\NormalTok{phat), }\DecValTok{1}\NormalTok{))}
\NormalTok{devres }\OtherTok{\textless{}{-}} \FunctionTok{sign}\NormalTok{(orings}\SpecialCharTok{$}\NormalTok{damage }\SpecialCharTok{{-}}\NormalTok{ orings}\SpecialCharTok{$}\NormalTok{total}\SpecialCharTok{*}\NormalTok{phat) }\SpecialCharTok{*} \FunctionTok{sqrt}\NormalTok{(}\DecValTok{2} \SpecialCharTok{*}\NormalTok{ (d1}\SpecialCharTok{+}\NormalTok{d2))}
\CommentTok{\# 计算偏差残差的分位数}
\FunctionTok{quantile}\NormalTok{(devres, }\FunctionTok{c}\NormalTok{(}\DecValTok{0}\NormalTok{, }\FloatTok{0.25}\NormalTok{, }\FloatTok{0.5}\NormalTok{, }\FloatTok{0.75}\NormalTok{, }\DecValTok{1}\NormalTok{))}
\end{Highlighting}
\end{Shaded}

\begin{verbatim}
##         0%        25%        50%        75%       100% 
## -0.9529154 -0.7345040 -0.4392853 -0.2079417  1.9565043
\end{verbatim}

\begin{Shaded}
\begin{Highlighting}[]
\CommentTok{\# 计算总偏差（所有偏差残差平方和）}
\NormalTok{dev }\OtherTok{\textless{}{-}} \FunctionTok{sum}\NormalTok{(devres}\SpecialCharTok{\^{}}\DecValTok{2}\NormalTok{)}
\NormalTok{dev}
\end{Highlighting}
\end{Shaded}

\begin{verbatim}
## [1] 16.91228
\end{verbatim}

\begin{Shaded}
\begin{Highlighting}[]
\CommentTok{\# 手动计算系数p值（z检验）}
\NormalTok{pvalue }\OtherTok{\textless{}{-}} \DecValTok{2} \SpecialCharTok{*}\NormalTok{ (}\DecValTok{1} \SpecialCharTok{{-}} \FunctionTok{pnorm}\NormalTok{(}\FunctionTok{abs}\NormalTok{(}\FunctionTok{summary}\NormalTok{(lmod)}\SpecialCharTok{$}\NormalTok{coef[ ,}\DecValTok{1}\NormalTok{]}\SpecialCharTok{/}\FunctionTok{summary}\NormalTok{(lmod)}\SpecialCharTok{$}\NormalTok{coef[ ,}\DecValTok{2}\NormalTok{])))}
\NormalTok{pvalue}
\end{Highlighting}
\end{Shaded}

\begin{verbatim}
##  (Intercept)         temp 
## 4.027947e-04 4.776581e-05
\end{verbatim}

\begin{Shaded}
\begin{Highlighting}[]
\CommentTok{\# 手动计算AIC（模型选择指标）}
\NormalTok{aic }\OtherTok{\textless{}{-}} \SpecialCharTok{{-}}\DecValTok{2}\SpecialCharTok{*}\FunctionTok{sum}\NormalTok{(orings}\SpecialCharTok{$}\NormalTok{damage}\SpecialCharTok{*}\NormalTok{eta }\SpecialCharTok{{-}}\NormalTok{ orings}\SpecialCharTok{$}\NormalTok{total}\SpecialCharTok{*}\FunctionTok{log}\NormalTok{(}\DecValTok{1}\SpecialCharTok{+}\FunctionTok{exp}\NormalTok{(eta))}
\SpecialCharTok{+} \FunctionTok{log}\NormalTok{(}\FunctionTok{choose}\NormalTok{(orings}\SpecialCharTok{$}\NormalTok{total,orings}\SpecialCharTok{$}\NormalTok{damage))) }\SpecialCharTok{+} \DecValTok{2} \SpecialCharTok{*} \DecValTok{2}
\NormalTok{aic}
\end{Highlighting}
\end{Shaded}

\begin{verbatim}
## [1] 33.67479
\end{verbatim}

\subsubsection{公差分布}\label{ux516cux5deeux5206ux5e03}

公差分布模型的基本设定为：存在一个潜在的连续公差变量\(Y^*\)，当\(Y^*\)超过某个阈值时，响应为1。设\(Y^*\)的CDF为\(F\)，则有

\[P(Y=1 | x) = P(Y^* > \gamma | x) = 1 - F(\gamma - x^T\beta)\]

若\(F\)为对称分布，可写为\(P(Y=1 | x) = F(x^T\beta - \gamma)\)。将\(\gamma\)吸收进截距项，得：

\[g(p) = \eta = x^T\beta\]

其中\(g = F^{-1}\)为连接函数。不同\(F\)的选择对应不同的连接函数。

下面是三种常见的，\textbf{基于不同公差分布}的二分类模型连接函数。

\begin{enumerate}
\def\labelenumi{\arabic{enumi}.}
\item
  Logit连接

  \[
  f(s)=\frac{e^s}{(1+e^s)^2},\eta=F^{-1}(p)=log(\frac{p}{1-p})
  \]
\item
  Probit连接

  \[
  f(s)=\frac{1}{\sqrt{2\pi}}e^{-s^2/2} ,\eta=\Phi^{-1}(p)
  \]
\item
  Cloglog连接

  \[
  f(s)=exp(s-exp(s)),\eta=F^{-1}(p)=log(-log(1-p))
  \]
\end{enumerate}

\subsubsection{潜变量}\label{ux6f5cux53d8ux91cf}

除了公差分布，也可以用潜变量构造二元响应概率和线性系统成分之间的连接函数。

潜变量方法假设存在一个不可观测的连续变量\(Y^* = x^T\beta + \epsilon\)，其中\(\epsilon\)的分布决定了连接函数：

\begin{enumerate}
\def\labelenumi{\arabic{enumi}.}
\item
  \textbf{Logit连接}：\(\epsilon\)服从标准Logistic分布，其CDF为\(F(s) = \frac{e^s}{1+e^s}\)，PDF为\(f(s) = \frac{e^s}{(1+e^s)^2}\)。
\item
  \textbf{Probit连接}：\(\epsilon\)服从标准正态分布，其CDF为\(F(s) = \Phi(s)\)，PDF为\(f(s) = \frac{1}{\sqrt{2\pi}} e^{-s^2/2}\)。
\item
  \textbf{互补双对数(Cloglog)连接}：\(\epsilon\)服从极值分布，其CDF为\(F(s) = 1 - \exp(-e^s)\)，PDF为\(f(s) = \exp(s - e^s)\)。
\end{enumerate}

观察二分类响应\(Y = I(Y^* > \gamma)\)，则\(P(Y=1|x) = P(Y^* > \gamma|x) = F(x^T\beta - \gamma)\)。

\begin{enumerate}
\def\labelenumi{\arabic{enumi}.}
\tightlist
\item
  Logit连接
\item
  Probit连接
\item
  互补双对数连接
\end{enumerate}

下面是一个昆虫试剂浓度的算例。

\begin{Shaded}
\begin{Highlighting}[]
\CommentTok{\# 加载数据：五种不同浓度下昆虫死亡数和存活数}
\FunctionTok{data}\NormalTok{(bliss, }\AttributeTok{package =} \StringTok{"faraway"}\NormalTok{)}
\CommentTok{\# 3中连接函数下的拟合}
\NormalTok{mlogit }\OtherTok{\textless{}{-}} \FunctionTok{glm}\NormalTok{(}\FunctionTok{cbind}\NormalTok{(dead, alive) }\SpecialCharTok{\textasciitilde{}}\NormalTok{ conc, }\AttributeTok{family =}\NormalTok{ binomial, bliss)}
\NormalTok{mprobit }\OtherTok{\textless{}{-}} \FunctionTok{glm}\NormalTok{(}\FunctionTok{cbind}\NormalTok{(dead, alive) }\SpecialCharTok{\textasciitilde{}}\NormalTok{ conc, }\AttributeTok{family =} \FunctionTok{binomial}\NormalTok{(}\AttributeTok{link=}\NormalTok{probit), bliss)}
\NormalTok{mcloglog }\OtherTok{\textless{}{-}} \FunctionTok{glm}\NormalTok{(}\FunctionTok{cbind}\NormalTok{(dead, alive) }\SpecialCharTok{\textasciitilde{}}\NormalTok{ conc, }\AttributeTok{family =} \FunctionTok{binomial}\NormalTok{(}\AttributeTok{link=}\NormalTok{cloglog), bliss)}
\CommentTok{\# 提取并整理各模型的拟合值}
\NormalTok{predval }\OtherTok{\textless{}{-}} \FunctionTok{sapply}\NormalTok{(}\FunctionTok{list}\NormalTok{(mlogit,mprobit,mcloglog), fitted)}
\FunctionTok{dimnames}\NormalTok{(predval) }\OtherTok{\textless{}{-}} \FunctionTok{list}\NormalTok{(}\DecValTok{1}\SpecialCharTok{:}\DecValTok{5}\NormalTok{, }\FunctionTok{c}\NormalTok{(}\StringTok{"logit"}\NormalTok{,}\StringTok{"probit"}\NormalTok{,}\StringTok{"cloglog"}\NormalTok{))}
\FunctionTok{round}\NormalTok{(predval, }\DecValTok{3}\NormalTok{)}\CommentTok{\# 不同浓度下，三种连接函数拟合的昆虫死亡概率。}
\end{Highlighting}
\end{Shaded}

\begin{verbatim}
##   logit probit cloglog
## 1 0.089  0.084   0.127
## 2 0.238  0.245   0.250
## 3 0.500  0.498   0.455
## 4 0.762  0.752   0.722
## 5 0.911  0.914   0.933
\end{verbatim}

\subsubsection{过度分散}\label{ux8fc7ux5ea6ux5206ux6563}

现实数据经常发生过度分散，表现为数据方差膨胀，其对结果可解释性造成负面影响，常见原因是存在模型未能捕捉到的变异性，从而观测值方差大于模型的估计方差。

下面是一个鳟鱼卵箱的算例。

\begin{Shaded}
\begin{Highlighting}[]
\FunctionTok{data}\NormalTok{(troutegg, }\AttributeTok{package =} \StringTok{"faraway"}\NormalTok{)}
\CommentTok{\#ftable(xtabs(cbind(survive,total) \textasciitilde{} location + period, troutegg))}
\NormalTok{bmod }\OtherTok{\textless{}{-}} \FunctionTok{glm}\NormalTok{(}\FunctionTok{cbind}\NormalTok{(survive, total}\SpecialCharTok{{-}}\NormalTok{survive) }\SpecialCharTok{\textasciitilde{}}\NormalTok{ location }\SpecialCharTok{+}\NormalTok{ period,}
\AttributeTok{family =}\NormalTok{ binomial, troutegg)}
\FunctionTok{summary}\NormalTok{(bmod)}
\end{Highlighting}
\end{Shaded}

\begin{verbatim}
## 
## Call:
## glm(formula = cbind(survive, total - survive) ~ location + period, 
##     family = binomial, data = troutegg)
## 
## Coefficients:
##             Estimate Std. Error z value Pr(>|z|)    
## (Intercept)   4.6358     0.2813  16.479  < 2e-16 ***
## location2    -0.4168     0.2461  -1.694   0.0903 .  
## location3    -1.2421     0.2194  -5.660 1.51e-08 ***
## location4    -0.9509     0.2288  -4.157 3.23e-05 ***
## location5    -4.6138     0.2502 -18.439  < 2e-16 ***
## period7      -2.1702     0.2384  -9.103  < 2e-16 ***
## period8      -2.3256     0.2429  -9.573  < 2e-16 ***
## period11     -2.4500     0.2341 -10.466  < 2e-16 ***
## ---
## Signif. codes:  0 '***' 0.001 '**' 0.01 '*' 0.05 '.' 0.1 ' ' 1
## 
## (Dispersion parameter for binomial family taken to be 1)
## 
##     Null deviance: 1021.469  on 19  degrees of freedom
## Residual deviance:   64.495  on 12  degrees of freedom
## AIC: 157.03
## 
## Number of Fisher Scoring iterations: 5
\end{verbatim}

\begin{Shaded}
\begin{Highlighting}[]
\NormalTok{beta }\OtherTok{\textless{}{-}} \FunctionTok{coef}\NormalTok{(bmod)}
\NormalTok{x }\OtherTok{\textless{}{-}} \FunctionTok{matrix}\NormalTok{(}\DecValTok{1}\NormalTok{, }\AttributeTok{nrow =} \DecValTok{20}\NormalTok{, }\AttributeTok{ncol =} \DecValTok{8}\NormalTok{)}
\NormalTok{x[ ,}\DecValTok{2}\NormalTok{] }\OtherTok{\textless{}{-}}\NormalTok{ troutegg}\SpecialCharTok{$}\NormalTok{location }\SpecialCharTok{==} \DecValTok{2}
\NormalTok{x[ ,}\DecValTok{3}\NormalTok{] }\OtherTok{\textless{}{-}}\NormalTok{ troutegg}\SpecialCharTok{$}\NormalTok{location }\SpecialCharTok{==} \DecValTok{3}
\NormalTok{x[ ,}\DecValTok{4}\NormalTok{] }\OtherTok{\textless{}{-}}\NormalTok{ troutegg}\SpecialCharTok{$}\NormalTok{location }\SpecialCharTok{==} \DecValTok{4}
\NormalTok{x[ ,}\DecValTok{5}\NormalTok{] }\OtherTok{\textless{}{-}}\NormalTok{ troutegg}\SpecialCharTok{$}\NormalTok{location }\SpecialCharTok{==} \DecValTok{5}
\NormalTok{x[ ,}\DecValTok{6}\NormalTok{] }\OtherTok{\textless{}{-}}\NormalTok{ troutegg}\SpecialCharTok{$}\NormalTok{period }\SpecialCharTok{==} \DecValTok{7}
\NormalTok{x[ ,}\DecValTok{7}\NormalTok{] }\OtherTok{\textless{}{-}}\NormalTok{ troutegg}\SpecialCharTok{$}\NormalTok{period }\SpecialCharTok{==} \DecValTok{8}
\NormalTok{x[ ,}\DecValTok{8}\NormalTok{] }\OtherTok{\textless{}{-}}\NormalTok{ troutegg}\SpecialCharTok{$}\NormalTok{period }\SpecialCharTok{==} \DecValTok{11}
\NormalTok{phat }\OtherTok{\textless{}{-}} \FunctionTok{exp}\NormalTok{(x}\SpecialCharTok{\%*\%}\NormalTok{beta)}\SpecialCharTok{/}\NormalTok{(}\DecValTok{1} \SpecialCharTok{+} \FunctionTok{exp}\NormalTok{(x}\SpecialCharTok{\%*\%}\NormalTok{beta))}
\NormalTok{peares }\OtherTok{\textless{}{-}}\NormalTok{ (troutegg}\SpecialCharTok{$}\NormalTok{survive}\SpecialCharTok{{-}}\NormalTok{troutegg}\SpecialCharTok{$}\NormalTok{total}\SpecialCharTok{*}\NormalTok{phat)}\SpecialCharTok{/}\FunctionTok{sqrt}\NormalTok{(troutegg}\SpecialCharTok{$}\NormalTok{total}\SpecialCharTok{*}\NormalTok{phat}\SpecialCharTok{*}\NormalTok{(}\DecValTok{1}\SpecialCharTok{{-}}\NormalTok{phat))}
\CommentTok{\#residuals(bmod, type = "pearson")}
\NormalTok{phi }\OtherTok{\textless{}{-}} \FunctionTok{sum}\NormalTok{(peares}\SpecialCharTok{\^{}}\DecValTok{2}\NormalTok{)}\SpecialCharTok{/}\NormalTok{(}\DecValTok{20{-}8}\NormalTok{)}
\NormalTok{phi}
\end{Highlighting}
\end{Shaded}

\begin{verbatim}
## [1] 5.330322
\end{verbatim}

\begin{Shaded}
\begin{Highlighting}[]
\FunctionTok{summary}\NormalTok{(bmod, }\AttributeTok{dispersion =}\NormalTok{ phi)}
\end{Highlighting}
\end{Shaded}

\begin{verbatim}
## 
## Call:
## glm(formula = cbind(survive, total - survive) ~ location + period, 
##     family = binomial, data = troutegg)
## 
## Coefficients:
##             Estimate Std. Error z value Pr(>|z|)    
## (Intercept)   4.6358     0.6495   7.138 9.49e-13 ***
## location2    -0.4168     0.5682  -0.734   0.4632    
## location3    -1.2421     0.5066  -2.452   0.0142 *  
## location4    -0.9509     0.5281  -1.800   0.0718 .  
## location5    -4.6138     0.5777  -7.987 1.39e-15 ***
## period7      -2.1702     0.5504  -3.943 8.05e-05 ***
## period8      -2.3256     0.5609  -4.146 3.38e-05 ***
## period11     -2.4500     0.5405  -4.533 5.82e-06 ***
## ---
## Signif. codes:  0 '***' 0.001 '**' 0.01 '*' 0.05 '.' 0.1 ' ' 1
## 
## (Dispersion parameter for binomial family taken to be 5.330322)
## 
##     Null deviance: 1021.469  on 19  degrees of freedom
## Residual deviance:   64.495  on 12  degrees of freedom
## AIC: 157.03
## 
## Number of Fisher Scoring iterations: 5
\end{verbatim}

\begin{center}\rule{0.5\linewidth}{0.5pt}\end{center}

\textbf{速查卡片}：

\begin{itemize}
\tightlist
\item
  \textbf{核心问题}：当响应变量不是连续正态变量时，如何建立回归模型？
\item
  \textbf{关键概念}：指数分散族（EDF） → 连接函数与典型连接 → IRLS算法 → 偏差统计量与偏差之差 → Pearson残差与偏差残差 → Logit/Probit/Cloglog连接 → 过度分散与准似然
\item
  \textbf{R命令}：\texttt{glm()}, \texttt{family\ =\ binomial()\ /\ poisson()\ /\ Gamma()}, \texttt{MASS::glm.nb()}, \texttt{faraway::troutegg} (示例)
\item
  \textbf{参见}：Ch10（线性回归基础）, Ch06（MLE）, Ch18（二分类方法）
\end{itemize}

\chapter{类别数据分析}\label{ux7c7bux522bux6570ux636eux5206ux6790}

\begin{quote}
人类天生就会分门别类。性别、职业、是否患病、是否点击------类别数据遍布社会科学、医学和商业分析的每一个角落。Ch11（GLM）中的 logit/probit 模型已经为二分类响应提供了建模工具，本章将从列联表出发，系统构建类别数据的推断方法，并走完从''交叉表''到''对数线性模型''的认知升级。
\end{quote}

\section{类别数据的关联分析}\label{ux7c7bux522bux6570ux636eux7684ux5173ux8054ux5206ux6790}

类别数据主要分析如何衡量分类（定性）数据之间关联关系。

\subsection{列联表和卡方独立性检验}\label{ux5217ux8054ux8868ux548cux5361ux65b9ux72ecux7acbux6027ux68c0ux9a8c}

研究两个类别变量时，每个变量有多个类别，通常将两个变量多个类别的频数用交叉表的形式表示出来个变量放在行(row)的位置，称为行变量，其类别数(行数)用r衣示另一个变量放在列(column)的位置，称为列变量，其类别数(列数)用c表示。这种由两个或两个以上类别变量交叉分类的频数分布表称为列联表。

列联表分析的基本问题:考察属性之间有无关联,即判别两属性是否独立.

在应用卡方检验时，\textbf{要样本量足够大，特别是每个单元格的期望频数不能太小，否则应用检验时可能会出现错误的结论}。

\(\chi^2\) 独立性检验：行和列代表两个类别变量，度量两者之间是否存在关联关系。

\[
H_0:p_{ij} = p_i\cdot p_j
\]

\subsection{优势比（Odds Ratio）}\label{ux4f18ux52bfux6bd4odds-ratio}

在 \(2 \times 2\) 列联表中，除了卡方检验，\textbf{优势比}（Odds Ratio, OR）是度量关联强度的核心指标。设行变量为暴露因素（有/无），列变量为结果（阳性/阴性）：

\begin{longtable}[]{@{}lll@{}}
\toprule\noalign{}
& 阳性 & 阴性 \\
\midrule\noalign{}
\endhead
\bottomrule\noalign{}
\endlastfoot
暴露组 & \(a\) & \(b\) \\
非暴露组 & \(c\) & \(d\) \\
\end{longtable}

暴露组的优势（odds）为 \(a/b\)，非暴露组的优势为 \(c/d\)，优势比定义为：
\[
OR = \frac{a/b}{c/d} = \frac{ad}{bc}
\]

\begin{itemize}
\tightlist
\item
  \(OR = 1\)：暴露与结果无关。
\item
  \(OR > 1\)：暴露组的结果优势高于非暴露组（正关联）。
\item
  \(OR < 1\)：暴露组的结果优势低于非暴露组（负关联，保护因素）。
\end{itemize}

优势比的置信区间通常通过对数变换构造：
\[
\text{Var}(\log OR) \approx \frac{1}{a} + \frac{1}{b} + \frac{1}{c} + \frac{1}{d}
\]
\[
95\%\ \text{CI}: \exp\left(\log OR \pm 1.96 \sqrt{\text{Var}(\log OR)}\right)
\]

\begin{quote}
优势比与相对风险（Relative Risk, \(RR = \frac{a/(a+b)}{c/(c+d)}\)）的区别在于：RR 比较的是概率之比，OR 比较的是优势之比。当结果罕见时（低发病率），OR 近似等于 RR。
\end{quote}

\subsection{卡方齐性检验}\label{ux5361ux65b9ux9f50ux6027ux68c0ux9a8c}

齐性检验：行表示所要研究的变量或者问题，列表示不同区组，检验行变量比例分布在各个区组之间是否一致。

提出原假设和备择假设：

\[
\forall i = 1,2,\cdots,r , H_0:p_{i1} = \cdots = p_{ic}=p_{i.}
\]

\(p_{ij}\) 表示第i个变量在第j个区组出现的频率。在原假设下这些概率和区组无关，都等于\(p_{i.}\)

检验统计量：

\[
\chi^2 = \sum_{i=1}^r\sum_{j=1}^c\frac{(o_{ij}-e_{ij})^2}{e_{ij}}\sim \chi^2((r-1)(c-1))
\]

\textbf{自由度为原假设下和备择假设下自由参数之差}。

卡方齐性检验常用于\textbf{抄袭判定}。其基本做法为：若具有抄袭争议作品与该作者其他时期的用词习惯不同，同时具有抄袭争议与被抄袭作品的情节高度相似，则可以判定有抄袭争议作品是抄袭其他作品。

\subsection{Fisher精确性检验（二维列联表）}\label{fisherux7cbeux786eux6027ux68c0ux9a8cux4e8cux7ef4ux5217ux8054ux8868}

Fisher精确性检验适用于二维列联表或小样本下的\(\chi^2\)检验。本质上是在原假设下使用超几何分布的精确分布（小样本下）或其正态近似分布（大样本），而不是用 \(\chi^2\) 拟合统计量。

\begin{longtable}[]{@{}llll@{}}
\toprule\noalign{}
& \(B_1\) & \(B_2\) & 总和 \\
\midrule\noalign{}
\endhead
\bottomrule\noalign{}
\endlastfoot
\(A_1\) & \(n_{11}\) & \(n_{12}\) & \(n_1.\) \\
\(A_2\) & \(n_{21}\) & \(n_{22}\) & \(n_2.\) \\
总和 & \(n_{.1}\) & \(n_{.2}\) & \(n_{..}\) \\
\end{longtable}

\[
n_{11}= \frac{C_{n.1}^{n11}C_{n.2}^{n12}}{C_{n..}^{n1.}} \sim hyper(n_{.1},n_{.2},n_{1.})
\]

Fisher精确检验的原假设为行变量和列变量独立。在原假设下，\(n_{11}\)服从超几何分布：

\[P(n_{11} = k) = \frac{\binom{n_{.1}}{k}\binom{n_{.2}}{n_{1.} - k}}{\binom{n_{..}}{n_{1.}}}, \quad k = 0, 1, \cdots, \min(n_{.1}, n_{1.})\]

双边检验的p值为：\(p = 2 \cdot \min\{P(n_{11} \leq k_0), P(n_{11} \geq k_0)\}\)。在R中，使用\texttt{fisher.test()}函数进行检验。

\subsection{Mantel-Haenszel检验（分层卡方检验）}\label{mantel-haenszelux68c0ux9a8cux5206ux5c42ux5361ux65b9ux68c0ux9a8c}

Mantel-Haenszel检验主要用于分析分层2×2列联表的数据，目的是评估在不同分层条件下两个分类变量之间的关联是否一致。【多张2×2列联表，剥离混杂因素】辛普森悖论！！！

基本做法: 回答处理与反应结果之间是否独立时，\textbf{需要先按层计算差异，再将各层的差异进行综合比较做出判断}。

Mantel-Haenszel提出\(Q_{MH}\) 统计量：\\
\[
Q_{MH}=\frac{[\sum_{h=1}^k(n_{11}^{(h)}-E(n_{11}^{(h)}))]^2}{\sum_{h=1}^kVar(n_{11}^{(h)})} \sim \chi^2(1)
\]

Mantel-Haenszel方法消除了层次因素的干扰而提高了检验功效（检验出变量关联性的可靠性），可以很好地解决辛普森悖论的出现。

\begin{quote}
\textbf{辛普森悖论的数值例子}：某医院有两种肾结石手术方案 A 和 B。汇总数据显示：A 的成功率为 78\%（273/350），B 的成功率为 83\%（289/350）------似乎 B 更好。但按结石大小分层后：

\begin{longtable}[]{@{}llll@{}}
\toprule\noalign{}
小结石 & 成功 & 失败 & 成功率 \\
\midrule\noalign{}
\endhead
\bottomrule\noalign{}
\endlastfoot
方案 A & 81 & 6 & 93\% \\
方案 B & 234 & 36 & 87\% \\
\end{longtable}

\begin{longtable}[]{@{}llll@{}}
\toprule\noalign{}
大结石 & 成功 & 失败 & 成功率 \\
\midrule\noalign{}
\endhead
\bottomrule\noalign{}
\endlastfoot
方案 A & 192 & 71 & 73\% \\
方案 B & 55 & 25 & 69\% \\
\end{longtable}

每一层中 A 都优于 B，但汇总后 B 反而更好------因为病情重的患者（大结石）更多分配给了 A。辛普森悖论的根源在于\textbf{混杂变量}（结石大小）同时影响治疗分配和结果。分层分析和 Mantel-Haenszel 检验正是为检测和纠正这一问题而生。
\end{quote}

\subsection{Ridit检验法}\label{riditux68c0ux9a8cux6cd5}

Ridit检验是用于处理有序分类数据的非参数统计方法，主要用于比较不同组在有序分类变量上的分布情况。其\textbf{核心思想是将有序类别转化为比例得分，然后比较不同组之间的得分差异}。【参照单元分析法】

\begin{longtable}[]{@{}ll@{}}
\toprule\noalign{}
符号 & 解释 \\
\midrule\noalign{}
\endhead
\bottomrule\noalign{}
\endlastfoot
行向量A & 表示不同比较组或不同处理的分类变量 \\
列向量B & 表示顺序尺度变量 \\
Oij & 回答第i处理在第j个顺序类上的响应数 \\
\end{longtable}

原理：取一个样本量较多的组或者将几组数据汇总成参照组，根据参照组的样本结构将原来各组响应数变换为参考得分-Ridit得分，利用变换后的Ridit得分进行各处理之间强弱的比较。

\begin{enumerate}
\def\labelenumi{\arabic{enumi}.}
\tightlist
\item
  选取标准组

  \begin{enumerate}
  \def\labelenumii{\arabic{enumii}.}
  \tightlist
  \item
    利用已知的标准分布作为标准组
  \item
    以例数最多的一组作为标准组
  \item
    把各组观察结果合并起来作为标准组（常用）
  \end{enumerate}
\item
  计算标准组的Ridit得分

  \begin{enumerate}
  \def\labelenumii{\arabic{enumii}.}
  \item
    第i个处理的得分取值在0-1之间 ， \[
    R_i=\sum_{j=1}^sr_j\frac{p_{ij}}{p_{i.}}=\sum_{j=1}^sr_j\frac{O_{ij}}{O_{i.}}
    \]
  \item
    Ridit 得分选择用累积概率得分表示各顺序真实的强弱顺序

    \[
    r_j = \sum_{k=1}^{j-1}p_{.k}+\frac{p_{.j}}{2}=\frac{F_{j-1}^B+F_{j}^B}{2}
    \] 用频数估计概率： \[
    \hat{r_j}=\frac{\frac{O_{.j}}{2}+\sum_{k=1}^{j-1}O_{.k}}{O_{..}}
    \]
  \item
    如上定义的Ridit得分，满足：标准组的得分恒等于0.5. \[
    R=\sum_{j=1}^sr_jp_{.j}=\frac{1}{2}
    \]
  \end{enumerate}
\item
  利用对照组计算各组的平均Ridit得分
\item
  做出判断：卡方检验或者置信区间判别法

  \begin{enumerate}
  \def\labelenumii{\arabic{enumii}.}
  \item
    Ridit检验统计量：（对整体的检验）

    \[
    W = 12\sum_{i=1}^rO_{i.}(R_i-0.5)^2 \sim \chi^2(r-1)
    \]
  \item
    \(R_i\)的0.95置信区间为：

    \[
    R_i\pm1.96\hat{\sigma_{R_i}}
    \]

    可以近似处理为：

    \[
    R_i \pm \frac{1}{\sqrt{3O_i}}
    \]

    检查0.5是否在这个置信区间内，如果不在，则认为有差异性。
  \end{enumerate}
\end{enumerate}

\section{对数线性模型}\label{ux5bf9ux6570ux7ebfux6027ux6a21ux578b}

列联表分析的一个关键局限是：每次只能考察少数几个变量之间的关系。当涉及三个或更多类别变量时，\textbf{对数线性模型}（Log-linear Model）提供了系统化的建模框架。

对于一个二维列联表，设单元格 \((i,j)\) 的期望频数为 \(\mu_{ij}\)，对数线性模型将 \(\log \mu_{ij}\) 分解为主效应和交互效应：
\[
\log \mu_{ij} = \lambda + \lambda_i^A + \lambda_j^B + \lambda_{ij}^{AB}
\]

\begin{itemize}
\tightlist
\item
  \(\lambda\)：总平均效应
\item
  \(\lambda_i^A\)：行变量 \(A\) 的主效应
\item
  \(\lambda_j^B\)：列变量 \(B\) 的主效应
\item
  \(\lambda_{ij}^{AB}\)：\(A\) 与 \(B\) 的交互效应
\end{itemize}

两变量的独立性 \(H_0: p_{ij} = p_{i\cdot} p_{\cdot j}\) 等价于 \(\lambda_{ij}^{AB} = 0\)（对所有 \(i,j\)）。因此，卡方独立性检验可以视为检验''交互效应为零''这一假设。这一框架自然地纳入 Ch11（GLM）的体系------对数线性模型是 Poisson GLM 的特例，其中 \(\log \mu_{ij}\) 是由分类协变量 \(A\) 和 \(B\) 决定的线性预测子。

对于三维列联表，对数线性模型可以包含所有二阶和三阶交互项，从而系统分析变量之间的条件独立性关系。

\section{有序类别数据的处理}\label{ux6709ux5e8fux7c7bux522bux6570ux636eux7684ux5904ux7406}

当类别变量具有天然顺序时（如''不满意/一般/满意''），前述方法忽略了顺序信息。\textbf{累积 logit 模型}（比例优势模型）是最常用的处理方法。

设有 \(J\) 个有序类别，记 \(P(Y \leq j \mid X)\) 为累积概率。模型假设：
\[
\text{logit}(P(Y \leq j \mid X)) = \alpha_j + \beta^T X, \quad j = 1, 2, \ldots, J-1
\]

\begin{itemize}
\tightlist
\item
  \(\alpha_j\)：截距项，随类别而变化（\(\alpha_1 < \alpha_2 < \cdots < \alpha_{J-1}\)）
\item
  \(\beta\)：回归系数，\textbf{对所有 \(j\) 都相同}------这是''比例优势''假设的核心
\end{itemize}

这一假设意味着：协变量 \(X\) 对''从 \(\leq j\) 跳到 \(>j\)``的影响对所有分割点 \(j\) 都是相同的。比例优势模型是 Ch11 中 GLM 框架的自然延伸，在 R 中使用 \texttt{MASS::polr()} 或 \texttt{ordinal::clm()} 拟合。

\begin{quote}
有序类别模型桥接了类别数据分析（本章）与广义线性模型（Ch11）------它告诉我们，类别不是孤岛，可以在统一的 GLM 基础设施上进行处理。
\end{quote}

\begin{center}\rule{0.5\linewidth}{0.5pt}\end{center}

\textbf{速查卡片}：

\begin{itemize}
\tightlist
\item
  \textbf{核心问题}：如何量化、检验和建模类别变量之间的关联关系？
\item
  \textbf{关键概念}：列联表 → 卡方独立性检验/齐性检验 → 优势比（OR）→ Fisher 精确检验 → Mantel-Haenszel 分层检验（辛普森悖论）→ Ridit 检验 → 对数线性模型 → 累积 logit 模型
\item
  \textbf{R命令}：\texttt{chisq.test()}, \texttt{fisher.test()}, \texttt{mantelhaen.test()}, \texttt{glm(family\ =\ poisson)}, \texttt{MASS::polr()}
\item
  \textbf{参见}：Ch07（假设检验基础）, Ch11（GLM 框架）, Ch13（非参数统计------Ridit 的有序视角）
\end{itemize}

\chapter{非参数统计}\label{ux975eux53c2ux6570ux7edfux8ba1}

\begin{quote}
前面几章的方法都要求你假定数据来自某个分布族------正态、二项、泊松。非参数统计则说：让数据自己说话。你不需要假定分布的形式，代价是需要更多的数据。这是一种''少假设、多证据''的哲学。
\end{quote}

本章涵盖非参数统计的三个核心板块：\textbf{非参数假设检验}（不依赖分布假设的推断）、\textbf{非参数密度估计}（让数据直接描绘自己的形状），以及\textbf{分位数回归}（不只看均值，看整个条件分布）。它们共同构成了''不做不必要假设''的方法论工具箱。

\section{单一样本的推断问题}\label{ux5355ux4e00ux6837ux672cux7684ux63a8ux65adux95eeux9898}

\subsection{符号检验}\label{ux7b26ux53f7ux68c0ux9a8c}

符号检验只关心观测值相对某个参考值的符号信息（+和-的个数），用于检验某个总体的中位数是否等于某个特定的值，以及配对数据的中位数是否为零。

符号检验无需假定数据服从正态分布（对比t检验）,对异常值不敏感，适用于小样本场景。但会损失数据信息，降低检验效率。

t检验是对总体均值的检验，符号检验是对总体中位数的检验。对非对称分布总体来说，\textbf{中位数较均值而言是对总体中心位置更稳健的估计。}

\begin{enumerate}
\def\labelenumi{\arabic{enumi}.}
\item
  双边符号检验过程

  \begin{enumerate}
  \def\labelenumii{\arabic{enumii}.}
  \item
    明确检验问题：\[
    H_0:M_e=M_0\ \ vs\ \ H1:M_e\neq M_0
    \]
  \item
    定义检验统计量：\[
    S^+ = \sum_{i=1}^nI(X_i>M_0),S^-=\sum_{i=1}^nI(X_i<M_0),S^++S^-=n'\leq n 
    \]

    在双边检验下，可以取\[
    K=min\{S^+,S^-\}
    \]
  \item
    给出显著性水平下的拒绝域：\[
    W = \{K\leq k\},where\ 2P_{binom}(K\leq k|K\sim b(n',0.5))\leq \alpha
    \]
  \item
    计算p值：\[
    p = 2P_{binom}(K\leq k_0|K\sim b(n',0.5))
    \]
  \end{enumerate}
\item
  在R中，使用binom.test()进行符号检验

\begin{Shaded}
\begin{Highlighting}[]
\NormalTok{build.price}\OtherTok{\textless{}{-}}\FunctionTok{c}\NormalTok{(}\DecValTok{36}\NormalTok{,}\DecValTok{32}\NormalTok{,}\DecValTok{31}\NormalTok{,}\DecValTok{25}\NormalTok{,}\DecValTok{28}\NormalTok{,}\DecValTok{36}\NormalTok{,}\DecValTok{40}\NormalTok{,}\DecValTok{32}\NormalTok{,}\DecValTok{41}\NormalTok{,}\DecValTok{26}\NormalTok{,}\DecValTok{35}\NormalTok{,}\DecValTok{35}\NormalTok{,}\DecValTok{32}\NormalTok{,}\DecValTok{87}\NormalTok{,}\DecValTok{33}\NormalTok{,}\DecValTok{35}\NormalTok{)}
\FunctionTok{binom.test}\NormalTok{(}\FunctionTok{sum}\NormalTok{(build.price}\SpecialCharTok{\textgreater{}}\DecValTok{37}\NormalTok{),}\FunctionTok{length}\NormalTok{(build.price),}\FloatTok{0.5}\NormalTok{)}
\end{Highlighting}
\end{Shaded}

\begin{verbatim}
## 
##  Exact binomial test
## 
## data:  sum(build.price > 37) and length(build.price)
## number of successes = 3, number of trials = 16, p-value = 0.02127
## alternative hypothesis: true probability of success is not equal to 0.5
## 95 percent confidence interval:
##  0.04047373 0.45645655
## sample estimates:
## probability of success 
##                 0.1875
\end{verbatim}
\item
  符号检验的大样本正态近似

  符号检验的精确分布为二项分布，当样本量大时可采用正态分布近似。记\[
  Z_0 = \frac{S_0^+-n'/2+C}{\sqrt{n'/4}}
  \]

  其中\(S^+>n'/2,C=-1/2;S^+_0<n'/2,C=1/2\)

  其中，连续性修正的规则为：当\(S^+ > n'/2\)时取\(C = -1/2\)，当\(S^+ < n'/2\)时取\(C = +1/2\)。在大样本下，\(Z_0\)近似服从标准正态分布\(N(0,1)\)。
\item
  符号检验在配对样本比较中的应用

  设有两组样本A,B，将A中小于B的样本记为-，大于B的样本记为+，将原问题转化为配对样本中+和-的比例是否相等，再应用符号检验。
\item
  广义符号检验

  本质上是将关于中位数的符号检验转换为p0分位数的符号检验。注意，在双边检验原假设成立下，有\[
  S^+\sim b(n',1-p_0),S^-\sim b(n',p_0)
  \]

  在R中，可以使用\texttt{binom.test()}函数进行广义符号检验，只需将参数\texttt{p}设定为\(1-p_0\)（若检验\(S^+\)）或\(p_0\)（若检验\(S^-\)）。
\end{enumerate}

\subsection{Cox-Stuart趋势存在性检验}\label{cox-stuartux8d8bux52bfux5b58ux5728ux6027ux68c0ux9a8c}

趋势存在性检验的基本思路是将数据分为前后两部分，若数据有上升趋势，则排在后面的数据的值要比排在前面的数据的值显著的大。反之，若数据有下降的趋势，那么排在后面的数据的值要比排在前面的数据的值显著的小。利用前后两个时期不同数据的差值正负来判断数据总的变化趋势。

Cox-Stuart提出最优的拆分点是数列的中间位置的数（中位数）。

Cox-Stuart趋势检验的具体步骤为：设观测序列为\(x_1, x_2, \cdots, x_n\)，取\(c = \lfloor n/2 \rfloor\)为中间拆分点。构造配对\((x_i, x_{i+c})\)，\(i = 1, 2, \cdots, c\)。计算差值\(D_i = x_{i+c} - x_i\)，统计\(D_i > 0\)的个数\(S^+\)。在原假设（无趋势）下，\(S^+ \sim B(c, 0.5)\)，可使用符号检验进行判断。

\subsection{Wilcoxon符号秩检验}\label{wilcoxonux7b26ux53f7ux79e9ux68c0ux9a8c}

对于配对两样本数据的检验提出的Wilcoxon符号秩检验不仅考虑差值的正负符号，还考虑了差值的大小。

使用Wilcoxon符号秩检验的\textbf{前提假设是要求数据的差值具有对称性}，即正负差值的分布应该大致相同。对于单样本的位置检验，要求数据的分布具有对称性。

工作原理：如果数据关于零点对称，则对称中心两侧数据的疏密程度应该大致相同。将样本的绝对值排序，那么原来取正值的数据和取负值的数据在绝对值样本中的秩和应该近似相等。

\begin{enumerate}
\def\labelenumi{\arabic{enumi}.}
\item
  Wilcoxon符号秩统计量的定义

  \[
  W^+ = \sum_{j=1}^nR_j^+S(X_j) = \sum_{j=1}^njW_j
  \]

  其中，\(R^+_j\)表示样本在绝对值样本中的秩，\(S(x)\)判定原始数据的正负号，\(W^+\)表示正样本在绝对值样本中的秩和。
\item
  Wilcoxon符号秩统计量的分布

  \begin{enumerate}
  \def\labelenumii{\arabic{enumii}.}
  \tightlist
  \item
    \(W^+\in (0,\frac{n(n+1)}{2})\)
  \item
    均值： \(E(W^+) = E(\sum_{j=1}^n jW_j) = \frac{n(n+1)}{4}\)
  \item
    方差：\(Var(W^+)=Var(\sum_{j=1}^njW_j)=\frac{n(n+1)(2n+1)}{24}\)
  \item
    大样本渐近正态：\[
    Z = \frac{W^+-\frac{n(n+1)}{4}}{\sqrt{\frac{n(n+1)(2n+1)}{24}}}
    \sim N(0,1)
    \]
  \item
    小样本的连续性修正：\[
    Z = \frac{W^+-\frac{n(n+1)}{4}\pm \frac{1}{2}}{\sqrt{\frac{n(n+1)(2n+1)}{24}}}
    \]
  \end{enumerate}
\item
  Wilcoxon符号秩检验过程

  Wilcoxon符号秩检验的步骤如下：

  \begin{enumerate}
  \def\labelenumii{\arabic{enumii}.}
  \item
    计算观测值与参考值\(M_0\)的差值\(D_i = X_i - M_0\)，剔除\(D_i = 0\)的样本，得到有效样本量\(n'\)。
  \item
    对差值的绝对值\(|D_i|\)进行升幂排序，得到秩\(R_i\)。
  \item
    计算正差值的秩和\(W^+ = \sum_{D_i > 0} R_i\)。
  \item
    在\(H_0\)成立时，\(W^+\)的分布关于\(\frac{n'(n'+1)}{4}\)对称。小样本查Wilcoxon符号秩检验临界值表；大样本使用正态近似\(Z = \frac{W^+ - n'(n'+1)/4 \pm 1/2}{\sqrt{n'(n'+1)(2n'+1)/24}} \sim N(0,1)\)。
  \item
    计算p值，做出判断。
  \end{enumerate}

  在R中，使用\texttt{wilcox.test()}函数进行Wilcoxon符号秩检验：

\begin{Shaded}
\begin{Highlighting}[]
\CommentTok{\# 单样本Wilcoxon符号秩检验}
\NormalTok{x }\OtherTok{\textless{}{-}} \FunctionTok{c}\NormalTok{(}\DecValTok{36}\NormalTok{, }\DecValTok{32}\NormalTok{, }\DecValTok{31}\NormalTok{, }\DecValTok{25}\NormalTok{, }\DecValTok{28}\NormalTok{, }\DecValTok{36}\NormalTok{, }\DecValTok{40}\NormalTok{, }\DecValTok{32}\NormalTok{, }\DecValTok{41}\NormalTok{, }\DecValTok{26}\NormalTok{, }\DecValTok{35}\NormalTok{, }\DecValTok{35}\NormalTok{, }\DecValTok{32}\NormalTok{, }\DecValTok{87}\NormalTok{, }\DecValTok{33}\NormalTok{, }\DecValTok{35}\NormalTok{)}
\FunctionTok{wilcox.test}\NormalTok{(x, }\AttributeTok{mu =} \DecValTok{37}\NormalTok{, }\AttributeTok{alternative =} \StringTok{"two.sided"}\NormalTok{)}
\end{Highlighting}
\end{Shaded}

\begin{verbatim}
## 
##  Wilcoxon signed rank test with continuity correction
## 
## data:  x
## V = 29.5, p-value = 0.04904
## alternative hypothesis: true location is not equal to 37
\end{verbatim}
\item
  Hodges-Lemmann估计量

  \begin{enumerate}
  \def\labelenumii{\arabic{enumii}.}
  \item
    Walsh平均值：对一组样本，计算任意两个样本的平均值得到一组长度为\(n(n+1)/2\)的新样本，称为Walsh平均值，又称为配对平均数。
  \item
    Hodges-Lemmann估计量：假设样本独立同分布取自\(F(x-\theta)\)，若\(F(\cdot)\)对称，定义Walsh平均值的中位数为：\[
    \hat{\theta}=median\{\frac{X_i+X_j}{2},i\leq j,i,j = 1,\cdots,n \}
    \]

    称\(\hat{\theta}\) 为\(\theta\) 的Hodges-Lemmann估计量。从均值，中位数、众数、H-L估计量，它们对数据中心位置的估计是越来越稳健的。
  \end{enumerate}
\end{enumerate}

\subsection{正态记分检验}\label{ux6b63ux6001ux8bb0ux5206ux68c0ux9a8c}

基本想法：选取正态分布作为参考，将秩转化为相应的正态分布的分位点，这样依赖于秩的检验，可转化为分位点大小的检验。而且正态记分检验对尾部数据做了差距放大处理，对中间数据做了差距压缩处理，\textbf{强调了尾部数据对位置判断的影响，更加稳健}。

\begin{enumerate}
\def\labelenumi{\arabic{enumi}.}
\item
  正态记分检验

  \begin{enumerate}
  \def\labelenumii{\arabic{enumii}.}
  \tightlist
  \item
    将升幂排列的秩Ri用升幂排列的正态分位点\(\Phi^{-1}(\frac{R_i}{n+1})\) 代替；
  \item
    为保证变换后为正，第\(R_i\) 的正态记分一般记为：\[
    s(i) = \Phi^{-1}(\frac{1}{2}(1+\frac{R_i}{n+1}))
    \]
  \end{enumerate}
\item
  正态记分检验步骤与案例：

  正态记分检验的步骤如下：

  \begin{enumerate}
  \def\labelenumii{\arabic{enumii}.}
  \item
    计算差值\(D_i = X_i - M_0\)，剔除零值，得到有效样本量\(n'\)。
  \item
    对\(|D_i|\)排序得到秩\(R_i\)，计算正态记分\(s_i = \Phi^{-1}\left(\frac{1}{2}\left(1 + \frac{R_i}{n'+1}\right)\right)\)。
  \item
    计算检验统计量\(T = \sum_{D_i > 0} s_i\)，即正差值对应正态记分之和。
  \item
    在原假设下，\(E(T) = 0\)，\(Var(T) = \sum_{i=1}^{n'} s_i^2\)。大样本时\(Z = T / \sqrt{Var(T)} \sim N(0,1)\)。
  \end{enumerate}

  案例：某地区抽取15户家庭，调查其月用电量（单位：度）。检验该地区家庭月用电量的中位数是否等于300度。计算差值\(D_i = X_i - 300\)，对\(|D_i|\)排序得秩，计算正态记分，求得\(T\)统计量并与标准正态分布临界值比较，得出检验结论。具体计算可通过R中的\texttt{qnorm()}函数配合秩次完成。
\end{enumerate}

\subsection{单组数据的位置参数置信区间估计}\label{ux5355ux7ec4ux6570ux636eux7684ux4f4dux7f6eux53c2ux6570ux7f6eux4fe1ux533aux95f4ux4f30ux8ba1}

\begin{enumerate}
\def\labelenumi{\arabic{enumi}.}
\item
  顺序统计量构造分位数置信区间：

  设\(X_{(1)} \leq X_{(2)} \leq \cdots \leq X_{(n)}\)为样本的顺序统计量。对于给定的置信水平\(1-\alpha\)，寻找两个顺序统计量\(X_{(r)}\)和\(X_{(s)}\)（\(r < s\)），使得

  \[P(X_{(r)} \leq M \leq X_{(s)}) = \sum_{k=r}^{s-1} \binom{n}{k} (0.5)^n \geq 1-\alpha\]

  其中\(M\)为总体中位数。通常选择\(r\)和\(s\)使得\(r + s = n + 1\)（对称选取），此时置信区间为\([X_{(r)}, X_{(n-r+1)}]\)。在大样本下，可取\(r \approx \frac{n}{2} - z_{\alpha/2}\frac{\sqrt{n}}{2}\)，\(s \approx \frac{n}{2} + z_{\alpha/2}\frac{\sqrt{n}}{2}\)。
\item
  方差估计的Bootstrap方法

  \begin{enumerate}
  \def\labelenumii{\arabic{enumii}.}
  \item
    从经验分布\(\hat{F_n}\)中重采样\(X_1^*,X_2^*,\cdots,X_m^*\) ，m为重抽样的样本量；
  \item
    计算统计量\(T^*_m=g(X_1^*,X_2^*,\cdots,X_m^*)\) ；
  \item
    重复步骤1和2共B次，得到\(T^*_{m,1},T^*_{m,2},\cdots,T^*_{m,B}\) ；
  \item
    计算B个统计量的样本方差：\[
    v_{boot}=\frac{1}{B}\sum_{b=1}^B(T^*_{m,b}-\frac{1}{B}\sum_{b=1}^BT^*_{m,b})^2
    \]
  \item
    Bootstrap中位数方差估计

    对于中位数\(T_n = \text{median}(X_1, \cdots, X_n)\)，Bootstrap方差估计的具体做法为：从原样本中重采样\(m\)次（通常\(m = n\)）得到Bootstrap样本，计算该样本的中位数\(T^*_{m,b}\)，重复\(B\)次（通常\(B \geq 1000\)），计算\(v_{boot} = \frac{1}{B}\sum_{b=1}^B (T^*_{m,b} - \bar{T}^*_m)^2\)，其中\(\bar{T}^*_m = \frac{1}{B}\sum_{b=1}^B T^*_{m,b}\)。该\(v_{boot}\)即为Bootstrap估计的中位数方差。
  \end{enumerate}
\item
  基于方差估计法的位置参数置信区间估计

  \begin{enumerate}
  \def\labelenumii{\arabic{enumii}.}
  \item
    正态置信区间：当\(T_n\) 的分布接近正态分布时，\(1-\alpha\) Bootstrap正态置信区间为\[
    T_n\pm z_{1-\frac{\alpha}{2}}\hat{SD}_{boot}
    \]
  \item
    枢轴量置信区间：当无法确定\(T_n\) 的分布是否正态或不是正态时，可以用枢轴量方法给出Bootstrap\(T_n\) 的置信区间，其中\(1-\alpha\) Bootstrap枢轴量置信区间\[
    (2\hat{\theta_n}-\hat{\theta^*_{1-\frac{\alpha}{2}}},2\hat{\theta_n}-\hat{\theta^*_{\frac{\alpha}{2}}})
    \]

    核心逻辑：原始估计 \(\hat{\theta_n}\) 与真实\(\theta\) 的距离，近似等于Bootstrap 估计\(\hat{\theta^*}\) 与\(\hat{\theta_n}\) 的距离。
  \item
    分位数置信区间：\(1-\alpha\)Bootstrap分位数置信区间\[
    (T^*_{B,\frac{\alpha}{2}},T^*_{B,1-\frac{\alpha}{2}})
    \]
  \item
    相关符号的含义

    \begin{longtable}[]{@{}
      >{\raggedright\arraybackslash}p{(\columnwidth - 4\tabcolsep) * \real{0.2674}}
      >{\raggedright\arraybackslash}p{(\columnwidth - 4\tabcolsep) * \real{0.3372}}
      >{\raggedright\arraybackslash}p{(\columnwidth - 4\tabcolsep) * \real{0.3837}}@{}}
    \toprule\noalign{}
    \begin{minipage}[b]{\linewidth}\raggedright
    来源
    \end{minipage} & \begin{minipage}[b]{\linewidth}\raggedright
    参数估计/统计量
    \end{minipage} & \begin{minipage}[b]{\linewidth}\raggedright
    分位数表示
    \end{minipage} \\
    \midrule\noalign{}
    \endhead
    \bottomrule\noalign{}
    \endlastfoot
    原始样本（n个） & \(
    \hat{\theta_n} \ ,\
     T_n
    \) & 无分位数，仅单个估计值 \\
    Bootstrap样本（B个） & \(
    \hat{\theta^*_b}\ ,\ T^*_b
    \) & \(
    \hat{\theta^*_q}\ ,\ T^*_{B,q}
    \) B个样本排序后的q分位数 \\
    \end{longtable}
  \end{enumerate}
\end{enumerate}

\section{两样本独立数据的位置和尺度推断}\label{ux4e24ux6837ux672cux72ecux7acbux6570ux636eux7684ux4f4dux7f6eux548cux5c3aux5ea6ux63a8ux65ad}

\subsection{Brown-Mood中位数检验}\label{brown-moodux4e2dux4f4dux6570ux68c0ux9a8c}

检验原理：在原假设成立时，两组数据的中位数相同，则将两组数据混合后，两组数据的混合中位数 \(med_{XY}\) 与\(med_{X}\) 和\(med_{Y}\) 相等，两组数据应该比较均匀的分布在\(med_{XY}\) 的两边。

检验步骤：

\begin{enumerate}
\def\labelenumi{\arabic{enumi}.}
\item
  找到混合数据的样本中位数，将X和Y按照分布在\(med_{XY}\) 的左右侧分为4类，计数得到2×2联表。

  \begin{longtable}[]{@{}llll@{}}
  \toprule\noalign{}
  & X & Y & Total \\
  \midrule\noalign{}
  \endhead
  \bottomrule\noalign{}
  \endlastfoot
  \(>med_{XY}\) & \(A\) & \(B\) & \(t\) \\
  \(<med_{XY}\) & \(C\) & \(D\) & \((m+n)-t\) \\
  \(Total\) & \(m\) & \(n\) & \(m+n=A+B+C+D\) \\
  \end{longtable}
\item
  检验统计量A服从超几何分布：\(A\sim hyper(m,n,t)\)

  \[
  P(A=k)=\frac{C_m^kC_n^{t-k}}{C_{m+n}^t},k = 0,1,\cdots, min\{m,t\}
  \]
\item
  在原假设成立下，当A过大或过小时可以考虑拒绝原假设。对应的p值为

  \[
  p = 2 min \{ P_{hyper}(A\leq a),P_{hyper}(A\geq a)\}
  \]
\item
  大样本情形下，在原假设成立时，可利用超几何分布的渐近正态分布进行检验

  \[
  Z=\frac{A-mt/(m+n)\pm C}{\sqrt{mnt(m+n-t)/(m+n)^3}}\sim N(0,1)
  \]
\end{enumerate}

\subsection{Wilcoxon-Mann-Whitney秩和检验}\label{wilcoxon-mann-whitneyux79e9ux548cux68c0ux9a8c}

基本思想：如果两组样本来自同一总体，那么它们在混合样本中排序后的秩次应该是随机分布的，两组样本在混合样本中的秩和应该大致相等。若两组样本的秩和相差很大，就有理由怀疑它们来自不同的总体。这种方法\textbf{通过比较秩和的大小来判断两组样本是否有显著差异}，避免了对数据分布的严格假设，具有较强的稳定性。

一些符号说明：

\begin{longtable}[]{@{}
  >{\raggedright\arraybackslash}p{(\columnwidth - 4\tabcolsep) * \real{0.2571}}
  >{\raggedright\arraybackslash}p{(\columnwidth - 4\tabcolsep) * \real{0.3071}}
  >{\raggedright\arraybackslash}p{(\columnwidth - 4\tabcolsep) * \real{0.4286}}@{}}
\toprule\noalign{}
\begin{minipage}[b]{\linewidth}\raggedright
符号
\end{minipage} & \begin{minipage}[b]{\linewidth}\raggedright
定义
\end{minipage} & \begin{minipage}[b]{\linewidth}\raggedright
补充说明
\end{minipage} \\
\midrule\noalign{}
\endhead
\bottomrule\noalign{}
\endlastfoot
\(R_i\) & 样本\(Y_i\) 在混合样本中的秩 & 包括了比\(Y_i\)小的X样本的个数和Y样本自身的个数，两部分构成 \\
\(W_Y = \sum_{i=1}^n R_i\) & 样本\(Y_i\)的Wilcoxon秩和统计量 & \(W_Y+W_X=\frac{(m+n)(m+n+1)}{2}\) \\
\(W_{XY} = W_Y - \frac{n(n+1)}{2}\) & 混合样本中\(Y\) 观测值大于\(X\) 观测值的个数 & \(W_{XY}+W_{YX}=mn\) \\
\end{longtable}

称\(W_{XY}\)或者\(W_{YX}\) 为Mann-Whitney检验。由于检验等价，\textbf{Wilcoxon秩和检验（1945提出，等样本）有时也称为Mann-Whitney检验（1947提出，不等样本），或者Wilcoxon-Mann-Whitney检验}。

我们可以推导出WMW秩和检验的统计量的均值和方差，进而提出大样本近似检验。

\[
E(W_{XY})=\frac{mn}{2},\ \ Var(W_{XY})=\frac{mn(m+n+1)}{12}
\]

\[
Z = \frac{W_{XY}-\frac{mn}{2}\pm C}{\sqrt{\frac{mn(m+n+1)}{12}}}\sim N(0,1)
\]

\subsection{Mood检验}\label{moodux68c0ux9a8c}

Mood检验关注两组样本的离散情况是否一样，需要\textbf{假定两位置参数相等}，不失一般假定为零。

\begin{enumerate}
\def\labelenumi{\arabic{enumi}.}
\item
  检验问题：样本\(X_1,\cdots,X_m\sim F(\frac{x}{\sigma_1}),Y_1,\cdots,Y_n\sim F(\frac{y}{\sigma_2})\)，两两独立，且F连续，则两样本刻度检验问题为\[
  H_0:\sigma_1^2 = \sigma_2^2\ \ vs\ \ H_1:\sigma_1^2\neq \sigma_2^2
  \]
\item
  检验统计量：样本X在混合样本中的秩的离差平方和\[
  M=\sum_{i=1}^m(R_i-\frac{m+n+1}{2})^2
  \]
\item
  检验统计量的性质和大样本近似（小样本时用正态连续性修正）：\[
  E(M)=\frac{m(m+n+1)(m+n-1)}{12}\ ,\ Var(M)=\frac{mn(m+n+1)(m+n+2)(m+n-2)}{180}
  \]

  \[
  Z = \frac{M-E(M)\pm C}{\sqrt{Var(M)}} \sim AN(0,1)
  \]
\end{enumerate}

\subsection{Moses方差检验}\label{mosesux65b9ux5deeux68c0ux9a8c}

Moses方差检验用于检验两总体方差是否相等，\textbf{无需事先假定两分布位置参数相等}。

Moses方差检验的步骤如下：

\begin{enumerate}
\def\labelenumi{\arabic{enumi}.}
\item
  将两组样本分别随机分为若干个子组（每组大小通常为2至4）。
\item
  对每个子组计算离差平方和（即组内各观测值减子组均值的平方和）。
\item
  将两组子组的离差平方和混合后，使用Wilcoxon-Mann-Whitney秩和检验比较两组离差平方和的分布。若两组离差平方和有显著差异，则说明两总体的方差不同。
\item
  由于随机分组会影响结果，建议多次随机分组进行检验，综合各次结果做出判断。
\end{enumerate}

一个具体的例子：

举例：设有两组样本，A组为\([2.1, 3.5, 4.0, 2.8, 3.2, 2.6]\)，B组为\([5.2, 4.8, 6.1, 5.5, 4.3, 5.0]\)。将每组各分为3个子组（每组2个观测值），计算各子组的离差平方和，然后对离差平方和应用Wilcoxon秩和检验，判断两组方差是否有显著差异。通过多次随机分组取平均p值，得到更稳健的结论。

Moses方差检验\textbf{在不同的分组情况下，检验的结果会不同}，可以多次随机分组，总和每次检验的结果。另外，Moses方差检验通过分组计算离差平方和，保留了两组样本的独立性，但\textbf{每组样本的数量会变少}，导致后续的Wilcoxon秩和检验不可靠。

\section{多组数据的位置参数检验}\label{ux591aux7ec4ux6570ux636eux7684ux4f4dux7f6eux53c2ux6570ux68c0ux9a8c}

在参数检验中，多组数据的位置推断需要使用方差分析，方差分析需要假定各组数据满足正态分布假定。当假定不成立时，我们需要用非参数方差分析。

\subsection{Kruskal-Wallis单因素方差分析}\label{kruskal-wallisux5355ux56e0ux7d20ux65b9ux5deeux5206ux6790}

Kruskal-Wallis单因素方差分析是将两样本W-M-W检验推广到三个或更多组检验的方法。采用类似的处理方法，将多个样本混合起来求秩，如果遇到打结的情况，采用平均值，然后再按照样本组求秩和。

符号说明：

\begin{longtable}[]{@{}
  >{\raggedright\arraybackslash}p{(\columnwidth - 4\tabcolsep) * \real{0.5338}}
  >{\raggedright\arraybackslash}p{(\columnwidth - 4\tabcolsep) * \real{0.1149}}
  >{\raggedright\arraybackslash}p{(\columnwidth - 4\tabcolsep) * \real{0.3446}}@{}}
\toprule\noalign{}
\begin{minipage}[b]{\linewidth}\raggedright
符号
\end{minipage} & \begin{minipage}[b]{\linewidth}\raggedright
定义
\end{minipage} & \begin{minipage}[b]{\linewidth}\raggedright
补充说明
\end{minipage} \\
\midrule\noalign{}
\endhead
\bottomrule\noalign{}
\endlastfoot
\(n = \sum_{j=1}^kn_j\) & 总样本数 & \(n_j\) 为不同处理下的样本数 \\
\(R_{..}=\frac{n(n+1)}{2}\) & 混合数据的秩和 & \\
\(R_{.j}=\sum_{i=1}^{n_j}R_{ij}\) & 处理内的秩和 & \\
\(\bar{R_{.j}}=\frac{R_{.j}}{n_j}\) & 处理内的平均秩 & 在\(H_0\)下成立\(E(\bar{R_{.j}})=\frac{n+1}{2}\)

\(Var(\bar{R_{.j}}) = \frac{(n-n_j)(n+1)}{12n_j}\) \\
\(MST=\frac{SST}{n-1}=\frac{n(n+1)}{12}\) & 总均方 & \\
\(SSt=\sum_{j=1}^k\frac{R_{.j}^2}{n_j}-\frac{n(n+1)^2}{4}\) & 处理间平方和 & \\
\(H=\frac{SSt}{MST}=\frac{12}{n(n+1)}\sum_{j=1}^k\frac{R_{.j}^2}{n_j}-3(n+1)\) & H统计量 & \(H\sim \chi^2(k-1)\) \\
\end{longtable}

Kruskal-Wallis检验的结果可以告知各组之间是否存在差异，但不会告知哪些组与其他组不同。为了确定哪些组与其他组不同，可以进行事后检验。常见的时候检验是Dunn检验。

\subsection{Jonckheere-Terpstra检验}\label{jonckheere-terpstraux68c0ux9a8c}

检验问题：样本的位置呈现出上升或下降的趋势，这种趋势从统计上来讲是否显著？

基本思想：与Mannundefined-Whitney检验类似，如果一个样本观测值小于另外一个样本观测值得个数较多或较少，则可以考虑两总体的位置之间存在大小关系。

\begin{longtable}[]{@{}
  >{\raggedright\arraybackslash}p{(\columnwidth - 2\tabcolsep) * \real{0.4368}}
  >{\raggedright\arraybackslash}p{(\columnwidth - 2\tabcolsep) * \real{0.5632}}@{}}
\toprule\noalign{}
\begin{minipage}[b]{\linewidth}\raggedright
符号
\end{minipage} & \begin{minipage}[b]{\linewidth}\raggedright
定义
\end{minipage} \\
\midrule\noalign{}
\endhead
\bottomrule\noalign{}
\endlastfoot
\(W_{ij}=\sum \sum I(W_{iu}<W_{jv})\) & 样本i中观测值小于样本j中观测值的个数 \\
\(J = \sum_{i<j}W_{ij}\) & J统计量，若取值过大，则认为k个总体有上升趋势。 \\
\end{longtable}

J检验统计量的期望和方差（公式会给出）：

\[
E(J)=\frac{1}{4}(N^2-\sum_{i=1}^kn_i^2)
\]

\[
Var(J)=\frac{1}{72}[N^2(2N+3)-\sum_{i=1}^kn_i^2(2n_i+3)]
\]

样本量较大时采用正态近似：

\[
Z = \frac{J-E(J)}{\sqrt{Var(J)}}\sim N(0,1)
\]

\begin{Shaded}
\begin{Highlighting}[]
\NormalTok{JTnorm }\OtherTok{\textless{}{-}} \ControlFlowTok{function}\NormalTok{(x, group) \{}
\NormalTok{  N }\OtherTok{\textless{}{-}} \FunctionTok{length}\NormalTok{(x)}
\NormalTok{  index }\OtherTok{\textless{}{-}} \FunctionTok{unique}\NormalTok{(group)}
\NormalTok{  k }\OtherTok{\textless{}{-}} \FunctionTok{length}\NormalTok{(index)}
\NormalTok{  ns }\OtherTok{\textless{}{-}} \ConstantTok{NULL}
  \ControlFlowTok{for}\NormalTok{ (i }\ControlFlowTok{in} \DecValTok{1}\SpecialCharTok{:}\NormalTok{k) \{}
\NormalTok{    ns }\OtherTok{\textless{}{-}} \FunctionTok{c}\NormalTok{(ns, }\FunctionTok{sum}\NormalTok{(group }\SpecialCharTok{==}\NormalTok{ index[i]))}
\NormalTok{  \}}
\NormalTok{  Diffval }\OtherTok{\textless{}{-}} \ConstantTok{NULL}
  \ControlFlowTok{for}\NormalTok{ (i }\ControlFlowTok{in} \DecValTok{1}\SpecialCharTok{:}\NormalTok{(k}\DecValTok{{-}1}\NormalTok{)) \{}
\NormalTok{    xi }\OtherTok{\textless{}{-}}\NormalTok{ x[}\FunctionTok{which}\NormalTok{(group }\SpecialCharTok{==}\NormalTok{ index[i])]}
    \ControlFlowTok{for}\NormalTok{ (j }\ControlFlowTok{in}\NormalTok{ (i}\SpecialCharTok{+}\DecValTok{1}\NormalTok{)}\SpecialCharTok{:}\NormalTok{k) \{}
\NormalTok{      xj }\OtherTok{\textless{}{-}}\NormalTok{ x[}\FunctionTok{which}\NormalTok{(group }\SpecialCharTok{==}\NormalTok{ index[j])]}
\NormalTok{      Diffvali }\OtherTok{\textless{}{-}} \DecValTok{0}
      \ControlFlowTok{for}\NormalTok{ (l }\ControlFlowTok{in} \DecValTok{1}\SpecialCharTok{:}\FunctionTok{length}\NormalTok{(xj)) \{}
\NormalTok{        Diffvali }\OtherTok{\textless{}{-}}\NormalTok{ Diffvali }\SpecialCharTok{+} \FunctionTok{sum}\NormalTok{(xi }\SpecialCharTok{{-}}\NormalTok{ xj[l] }\SpecialCharTok{\textless{}} \DecValTok{0}\NormalTok{)}
\NormalTok{      \}}
\NormalTok{      Diffval }\OtherTok{\textless{}{-}} \FunctionTok{c}\NormalTok{(Diffval, Diffvali)}
\NormalTok{    \}}
\NormalTok{  \}}
\NormalTok{  m.val }\OtherTok{\textless{}{-}}\NormalTok{ (N}\SpecialCharTok{\^{}}\DecValTok{2} \SpecialCharTok{{-}} \FunctionTok{sum}\NormalTok{(ns}\SpecialCharTok{\^{}}\DecValTok{2}\NormalTok{)) }\SpecialCharTok{/} \DecValTok{4}
\NormalTok{  sd.val }\OtherTok{\textless{}{-}} \FunctionTok{sqrt}\NormalTok{((N}\SpecialCharTok{\^{}}\DecValTok{2} \SpecialCharTok{*}\NormalTok{ (}\DecValTok{2}\SpecialCharTok{*}\NormalTok{N }\SpecialCharTok{+} \DecValTok{3}\NormalTok{) }\SpecialCharTok{{-}} \FunctionTok{sum}\NormalTok{(ns}\SpecialCharTok{\^{}}\DecValTok{2} \SpecialCharTok{*}\NormalTok{ (}\DecValTok{2}\SpecialCharTok{*}\NormalTok{ns }\SpecialCharTok{+} \DecValTok{3}\NormalTok{))) }\SpecialCharTok{/} \DecValTok{72}\NormalTok{)}
\NormalTok{  zval }\OtherTok{\textless{}{-}}\NormalTok{ (}\FunctionTok{sum}\NormalTok{(Diffval) }\SpecialCharTok{{-}}\NormalTok{ m.val) }\SpecialCharTok{/}\NormalTok{ sd.val}
\NormalTok{  pval }\OtherTok{\textless{}{-}} \FunctionTok{pnorm}\NormalTok{(zval, }\DecValTok{0}\NormalTok{, }\DecValTok{1}\NormalTok{, }\AttributeTok{lower.tail =}\NormalTok{ F)}
  \FunctionTok{list}\NormalTok{(}
    \AttributeTok{ns =}\NormalTok{ ns,}
    \AttributeTok{k =}\NormalTok{ k,}
    \AttributeTok{Diffval =}\NormalTok{ Diffval,}
    \AttributeTok{J =} \FunctionTok{sum}\NormalTok{(Diffval),}
    \AttributeTok{m.val =}\NormalTok{ m.val,}
    \AttributeTok{sd.val =}\NormalTok{ sd.val,}
    \AttributeTok{zval =}\NormalTok{ zval,}
    \AttributeTok{pval =}\NormalTok{ pval}
\NormalTok{  )}
\NormalTok{\}}

\NormalTok{jump }\OtherTok{\textless{}{-}} \FunctionTok{c}\NormalTok{(}\DecValTok{125}\NormalTok{, }\DecValTok{136}\NormalTok{, }\DecValTok{116}\NormalTok{, }\DecValTok{101}\NormalTok{, }\DecValTok{105}\NormalTok{, }\DecValTok{109}\NormalTok{, }\DecValTok{122}\NormalTok{, }\DecValTok{114}\NormalTok{, }\DecValTok{131}\NormalTok{, }\DecValTok{120}\NormalTok{, }\DecValTok{119}\NormalTok{, }\DecValTok{127}\NormalTok{, }\DecValTok{128}\NormalTok{, }\DecValTok{142}\NormalTok{, }\DecValTok{128}\NormalTok{, }\DecValTok{134}\NormalTok{, }\DecValTok{135}\NormalTok{, }\DecValTok{131}\NormalTok{, }\DecValTok{140}\NormalTok{, }\DecValTok{129}\NormalTok{)}
\NormalTok{gr.jump }\OtherTok{\textless{}{-}} \FunctionTok{factor}\NormalTok{(}\FunctionTok{rep}\NormalTok{(}\DecValTok{1}\SpecialCharTok{:}\DecValTok{3}\NormalTok{, }\FunctionTok{c}\NormalTok{(}\DecValTok{6}\NormalTok{, }\DecValTok{6}\NormalTok{, }\DecValTok{8}\NormalTok{)), }\AttributeTok{labels =} \FunctionTok{c}\NormalTok{(}\StringTok{\textquotesingle{}G1\textquotesingle{}}\NormalTok{, }\StringTok{\textquotesingle{}G2\textquotesingle{}}\NormalTok{, }\StringTok{\textquotesingle{}G3\textquotesingle{}}\NormalTok{))}
\FunctionTok{JTnorm}\NormalTok{(jump, gr.jump)}
\end{Highlighting}
\end{Shaded}

\begin{verbatim}
## $ns
## [1] 6 6 8
## 
## $k
## [1] 3
## 
## $Diffval
## [1] 25 42 44
## 
## $J
## [1] 111
## 
## $m.val
## [1] 66
## 
## $sd.val
## [1] 14.38749
## 
## $zval
## [1] 3.127716
## 
## $pval
## [1] 0.000880851
\end{verbatim}

\section{非参数密度估计}\label{ux975eux53c2ux6570ux5bc6ux5ea6ux4f30ux8ba1}

\subsection{直方图 Histogram}\label{ux76f4ux65b9ux56fe-histogram}

构造与估计方法：

\[
B_j = [x_o + (j-1)h , x_0+jh],j\in z
\]

\[
n_j = \sum_{i=1}^nI(X_i\in B_j),\ \ f_j = \frac{n_j}{nh}
\]

\[
\hat{f_h}(x)=\frac{1}{nh}\sum_{i=1}^n\sum_j I(X_i\in B_j)I(x\in B_j)
\]

\[
\hat{f_h}(m_j) = \frac{1}{nh}\sum_{i=1}^nI(X_i\in B_j)\ ,\ B_j=[m_j-h/2,m_j+h/2]
\]

怎么选合适的窗宽？使得偏差bias和方差variance都尽可能小。

\[
Bias = E[ \hat{f_h}(x)-f(x)] \approx f'(m_j)(m_j-x) \propto f'(m_j)\frac{h}{2}
\]

\[
Variance =Var(\hat{f_h}(x))\approx \frac{1}{nh}f(x)\propto \frac{1}{nh}
\]

\[
MSE = E[\hat{f_h}(x)-f(x)]^2 = Var(\hat{f_h}(x))+Bias^2(\hat{f_h}(x))
\]

均方误趋于零，估计是相合的。

MSE对x积分，得到MISE。

\[
MISE = E[\int_{-\infty}^{\infty}\{\hat{f_h}(x)-f(x) \}^2dx] = \int MSEdx
\]

\[
AMISE(\hat{f_h}) = \frac{1}{nh}+\frac{h^2}{12}||f'||_2^2
\]

\[
||f'||_2^2 = \int f^{'2}(x)dx
\]

\[
h_{opt}=(\frac{6}{n||f'||_2^2})^{1/3}\sim n^{-1/3}
\]

极小化AMISE得到最有窗宽，再标准正态下，可以取\(h_{opt} = 3.5n^{-1/3}\)

直方图还受到原点位置的影响，可以用多次求平均的方法处理。

\subsection{核密度估计}\label{ux6838ux5bc6ux5ea6ux4f30ux8ba1}

要算哪个点，就以该点为中心构造区间，计算密度。

构造均匀核：（积分为1，偶函数。）

\[
K(u)=\frac{1}{2}I(|u|\leq 1)
\]

\[
\hat{f_h}(x) = \frac{1}{2nh}\sum_{i=1}^nI(|\frac{x_i-x}{h}|\leq 1)= \frac{1}{nh}\sum_{i=1}^n k(\frac{x_i-x}{h})
\]

均匀核对待估点附近所有点权重一样，我们希望离x越近权重越大。下面展示了一些常用的核函数

常用的核函数包括：

\begin{longtable}[]{@{}
  >{\raggedright\arraybackslash}p{(\columnwidth - 4\tabcolsep) * \real{0.3333}}
  >{\raggedright\arraybackslash}p{(\columnwidth - 4\tabcolsep) * \real{0.3333}}
  >{\raggedright\arraybackslash}p{(\columnwidth - 4\tabcolsep) * \real{0.3333}}@{}}
\toprule\noalign{}
\begin{minipage}[b]{\linewidth}\raggedright
核函数名称
\end{minipage} & \begin{minipage}[b]{\linewidth}\raggedright
表达式 \(K(u)\)
\end{minipage} & \begin{minipage}[b]{\linewidth}\raggedright
支撑集
\end{minipage} \\
\midrule\noalign{}
\endhead
\bottomrule\noalign{}
\endlastfoot
均匀核 (Uniform) & \(\frac{1}{2} I(|u| \leq 1)\) & \([-1, 1]\) \\
高斯核 (Gaussian) & \(\frac{1}{\sqrt{2\pi}} e^{-u^2/2}\) & \(\mathbb{R}\) \\
Epanechnikov核 & \(\frac{3}{4}(1 - u^2) I(|u| \leq 1)\) & \([-1, 1]\) \\
三角核 (Triangular) & \((1 - |u|) I(|u| \leq 1)\) & \([-1, 1]\) \\
四次核 (Quartic/Biweight) & \(\frac{15}{16}(1 - u^2)^2 I(|u| \leq 1)\) & \([-1, 1]\) \\
\end{longtable}

其中Epanechnikov核在均方误差意义下是最优的，但实践中高斯核因其光滑性和数学便利性最为常用。

核函数要求：积分为1，函数要对称，方差存在。（一阶矩为零，二阶矩小于无穷）

\[
\int k(u)du=1\ \ ;\int uk(u)du=0
\]

核密度估计的表达式：

\[
\hat{f_h}(x) = \frac{1}{n}\sum_{i=1}^nK_h(x-X_i)=\frac{1}{n}\sum_{i=1}^nK(\frac{x-X_i}{h})
\]

得到估计后，需要了解估计的性质。

\[
bias(\hat{f_h}(x))=E(\hat{f_h}(x))-f(x)= \frac{1}{2}h^2f''(x)\mu_2(K)+o(h^2) \propto h^2
\]

偏差的正负取决于\(f''(x)\) 的符号！！

核密度估计偏差的符号取决于\(f(x)\)的二阶导数\(f''(x)\)的符号：

\begin{itemize}
\tightlist
\item
  当\(f''(x) > 0\)（局部凸/峰谷）：\(\hat{f_h}(x)\)倾向于\textbf{高估}真实密度（正偏差）。
\item
  当\(f''(x) < 0\)（局部凹/峰顶）：\(\hat{f_h}(x)\)倾向于\textbf{低估}真实密度（负偏差）。
\item
  当\(f''(x) = 0\)（拐点）：偏差最小。
\end{itemize}

这意味着核密度估计在密度峰值处会''削峰''，在密度谷底处会''填谷''，且窗宽\(h\)越大，这种平滑效应越明显。

\[
Var\{\hat{f_h}(x) \}=\frac{1}{nh}||K||^2_2f(x)+o(\frac{1}{nh})
\]

K的二范数大小决定了方差的大小。

推导过程要掌握！！！

下面求MSE。

\[
MSE(\hat{f_h}(x))=Bias^2(\hat{f_h}(x))+Var(\hat{f_h}(x))
\]

\[
MISE(\hat{f_h})=\int MSE(\hat{f_h}(x))dx
\]

忽略高阶无穷小，得到：

\[
AMISE(\hat{f_h})=\frac{1}{nh}||K||^2_2+\frac{h^4}{4}[\mu_2(K)]^2||f''||^2_2
\]

\[
h_{opt}=(\frac{||K||^2_2}{||f''||^2_2[\mu_2(K)]^2n})^{1/5}\sim n^{-1/5}
\]

相合性：

渐近正态性：

\subsection{实际窗宽选取法则}\label{ux5b9eux9645ux7a97ux5bbdux9009ux53d6ux6cd5ux5219}

\[
h_{opt}=(\frac{||K||^2_2}{||f''||^2_2[\mu_2(K)]^2n})^{1/5}\sim n^{-1/5}
\]

假设1：\(f\) 服从正态分布，\(||f''||^2=\sigma^{-5}\frac{3}{8\sqrt{\pi}}\) ,\$\textbackslash hat\{\textbackslash sigma\}=\textbackslash sqrt\{\textbackslash frac\{1\}\{n-1\}\textbackslash sum(X\_i=\textbackslash bar(X))\^{}2\} \$ 假设2：核函数为正态核，\(K(u)=\frac{1}{\sqrt{2\pi}}e^{-u^2/2}\) 此时，窗宽选取的经验法则为：

\[
\hat{h}_{rot}\approx 1.06\hat{\sigma}n^{-1/5}
\]

更稳健的方法：

\[
\hat{h}_{rot}=1.06 min\{\hat{\sigma^2},\frac{\hat{R}}{1.34} \}n^{-1/5}
\]

\(\hat{R}\)是样本的3/4分位数与1/4分位数的差。

\section{分位数回归}\label{ux5206ux4f4dux6570ux56deux5f52}

\subsection{引入}\label{ux5f15ux5165}

（1）条件期望回归 ( condition mean regression ) 的例子：

\begin{itemize}
\item
  CLM：\(E(Y|X)=X^T\beta\)
\item
  Logistics Regression: \(logit[P(Y=1|X)]=logit[(E(Y|X))]=X^T\beta\)
\end{itemize}

这些回归都假设同方差性。但有时候X对Y的影响可能是\textbf{异质的}，如图：

例如，在研究家庭收入（X）与食品支出（Y）的关系时，散点图通常呈现''扇形''或''喇叭形''分布：在低收入水平，食品支出的变异性较小；在高收入水平，食品支出的变异性显著增大。这说明，条件均值回归（如OLS）仅能刻画中心趋势，而分位数回归可以进一步描述不同消费层次家庭的支出模式，从而揭示X对Y的异质效应。

（2）考虑第二种模型：Location-Scale model 位置刻度模型

考虑了给定X下的均值和方差：

\[
E(y_i|x_i)=\mu(x_i)=x_i^T\beta
\]

\[
\sqrt{Var(y_i|x_i)}=\sigma(x_i)=exp(x_i^T\gamma)
\]

这两个模型刻画的都是条件分布最主要的特征。

（3）分位数回归

\[
F_{Y|X}[Y\leq q_{\tau}(Y|X)|X]=\tau,\tau\in(0,1)
\]

一般选取特殊的分位点进行分位数回归。比如中位数回归，往往比均值回归\textbf{更加稳健}。

\textbf{分位数回归不假设给定X下Y的分布是什么}（比如均值回归是正态分布假设）。

\textbf{定义}

\[
Q_X(\tau)=F_X^{-1}(\tau)=inf\{x|F(x)\geq \tau \}
\]

回顾均值的定义：

\[
\mu=EX=argmin_{\theta}E[(x-\theta)]^2
\]

\[
\bar{x}=\hat{\mu}=argmin_{\theta}\frac{1}{n}\sum(x_i-\theta)^2
\]

\[
\mu(x)=E[Y|X=x]=argmin_{m}E_{Y|X=x}(Y-m(X))^2
\]

均值可以通过求一个\textbf{平方损失函数}的最小值求得。同理，分位数回归中也存在类似的损失函数。

\[
\alpha_{\tau}=F^{-1}(\tau)=argmin_{u}E[\rho_{\tau}(y-u)],\ where\ \rho_{\tau}=u(\tau-I(u<0))
\]

接下来，用经验分布函数代替。

\[
y_{[n\tau]}=inf\{y:F_n(y)\geq\tau\}=argmin_{\epsilon}\frac{1}{n}\sum_{i=1}^n\rho_{\tau}(y_i-\epsilon)\]

条件分位数：

\[
\alpha_{\tau}(x)=argmin_{\alpha}E_{Y|X=x}[\rho_{\tau}(Y-\alpha(X))]
\]

\textbf{分位数回归模型}

线性分位数回归模型（关注的模型）

\[
Q_{Y|X=x}(\tau)=X^T\beta_{\tau}
\]

参数非线性分位数回归模型

\[
Q_{Y|X=x}(\tau)=g(X;\beta_{\tau})
\]

非参数分位数回归模型

半参数分位数回归模型

\subsection{线性分位数回归模型}\label{ux7ebfux6027ux5206ux4f4dux6570ux56deux5f52ux6a21ux578b}

\subsection{代码实现分位数回归}\label{ux4ee3ux7801ux5b9eux73b0ux5206ux4f4dux6570ux56deux5f52}

\begin{Shaded}
\begin{Highlighting}[]
\FunctionTok{library}\NormalTok{(quantreg)}
\FunctionTok{data}\NormalTok{(}\StringTok{\textquotesingle{}engel\textquotesingle{}}\NormalTok{)}
\FunctionTok{attach}\NormalTok{(engel)}
\NormalTok{fit1 }\OtherTok{=} \FunctionTok{rq}\NormalTok{(foodexp }\SpecialCharTok{\textasciitilde{}}\NormalTok{ income, }\AttributeTok{tau =} \FloatTok{0.5}\NormalTok{, }\AttributeTok{data =}\NormalTok{ engel) }
\FunctionTok{summary}\NormalTok{(fit1) }
\end{Highlighting}
\end{Shaded}

\begin{verbatim}
## 
## Call: rq(formula = foodexp ~ income, tau = 0.5, data = engel)
## 
## tau: [1] 0.5
## 
## Coefficients:
##             coefficients lower bd  upper bd 
## (Intercept)  81.48225     53.25915 114.01156
## income        0.56018      0.48702   0.60199
\end{verbatim}

\chapter{时间序列分析}\label{ux65f6ux95f4ux5e8fux5217ux5206ux6790}

\begin{quote}
前面各章讨论的数据，一个隐含的假设是观测值之间相互独立。但时间序列数据最本质的特征恰恰在于：\textbf{观测值之间存在时间上的依赖关系}。昨天的股价影响今天的判断，上个月的销售额预示本月的趋势------忽视这种依赖会导致错误的推断。
\end{quote}

\section{时间序列的基本概念}\label{ux65f6ux95f4ux5e8fux5217ux7684ux57faux672cux6982ux5ff5}

\subsection{什么是时间序列}\label{ux4ec0ux4e48ux662fux65f6ux95f4ux5e8fux5217}

\textbf{时间序列}（Time Series）是按照时间顺序排列的一组观测值，记为 \(\{X_t\}_{t=1}^T\)。与截面数据不同，时间序列数据的两个基本特征是：

\begin{enumerate}
\def\labelenumi{\arabic{enumi}.}
\tightlist
\item
  \textbf{自相关（Autocorrelation）}：\(X_t\) 与 \(X_{t-k}\) 之间不是独立的。
\item
  \textbf{非平稳性（Non-stationarity）}：数据的统计特性（均值、方差、自协方差）可能随时间变化。
\end{enumerate}

\subsection{平稳性}\label{ux5e73ux7a33ux6027}

一个时间序列称为\textbf{严平稳}（Strictly Stationary），如果其任意有限维分布对时间的平移不变：

\[(X_{t_1}, \ldots, X_{t_k}) \stackrel{d}{=} (X_{t_1+h}, \ldots, X_{t_k+h}), \quad \forall h\]

在实际应用中，我们通常只要求\textbf{弱平稳}（Weak Stationarity）：

\begin{enumerate}
\def\labelenumi{\arabic{enumi}.}
\tightlist
\item
  \(E[X_t] = \mu\)（均值恒定）。
\item
  \(Var[X_t] = \sigma^2\)（方差恒定）。
\item
  \(Cov(X_t, X_{t+k})\) 仅依赖于滞后 \(k\)，不依赖于 \(t\)。
\end{enumerate}

\begin{quote}
对于正态过程，弱平稳与严平稳等价。
\end{quote}

\subsection{自相关函数（ACF）与偏自相关函数（PACF）}\label{ux81eaux76f8ux5173ux51fdux6570acfux4e0eux504fux81eaux76f8ux5173ux51fdux6570pacf}

\textbf{自协方差函数}（Autocovariance Function）：

\[\gamma(k) = Cov(X_t, X_{t+k}) = E[(X_t - \mu)(X_{t+k} - \mu)]\]

\textbf{自相关函数}（Autocorrelation Function, ACF）：

\[\rho(k) = \frac{\gamma(k)}{\gamma(0)}\]

其中 \(\rho(k) \in [-1, 1]\)，\(\rho(0) = 1\)。

\textbf{偏自相关函数}（Partial Autocorrelation Function, PACF）：\(\phi_{kk}\) 度量了在消除中间 \(X_{t+1}, \ldots, X_{t+k-1}\) 的线性影响后，\(X_t\) 与 \(X_{t+k}\) 之间的净相关性。

ACF 和 PACF 是识别 ARMA 模型阶数的核心工具。

\section{白噪声过程}\label{ux767dux566aux58f0ux8fc7ux7a0b}

\textbf{白噪声}（White Noise）是最简单的时间序列模型：

\[\varepsilon_t \sim WN(0, \sigma^2)\]

即：\(E[\varepsilon_t] = 0\)，\(Var[\varepsilon_t] = \sigma^2\)，且对任意 \(k \neq 0\)，\(Cov(\varepsilon_t, \varepsilon_{t+k}) = 0\)。

白噪声是构造复杂时间序列模型的基本构件。ACF 图若显示所有非零滞后都落在置信区间内，则可初步判断序列为白噪声。

\section{ARMA 模型族}\label{arma-ux6a21ux578bux65cf}

ARMA（Autoregressive Moving Average）模型是平稳时间序列建模的核心框架。

\subsection{自回归模型 AR(p)}\label{ux81eaux56deux5f52ux6a21ux578b-arp}

AR(p) 模型假设当前值是其过去 \(p\) 个值的线性组合加上白噪声：

\[X_t = \phi_1 X_{t-1} + \phi_2 X_{t-2} + \cdots + \phi_p X_{t-p} + \varepsilon_t\]

其中 \(\varepsilon_t \sim WN(0, \sigma^2)\)。

\textbf{平稳性条件}：特征方程 \(1 - \phi_1 z - \phi_2 z^2 - \cdots - \phi_p z^p = 0\) 的根全部落在单位圆外。对于 AR(1)，即 \(|\phi_1| < 1\)。

AR(1) 过程的性质：

\begin{itemize}
\tightlist
\item
  \(E[X_t] = 0\)（若均值为零）
\item
  \(Var[X_t] = \frac{\sigma^2}{1 - \phi_1^2}\)
\item
  \(\rho(k) = \phi_1^{|k|}\)（指数衰减）
\item
  PACF：\(\phi_{11} = \phi_1\)，\(\phi_{kk} = 0\) 对 \(k > 1\)（\textbf{截尾}）
\end{itemize}

\textbf{识别特征}：ACF 拖尾（指数衰减），PACF 在滞后 \(p\) 处截尾。

\subsection{移动平均模型 MA(q)}\label{ux79fbux52a8ux5e73ux5747ux6a21ux578b-maq}

MA(q) 模型将当前值表示为当前和过去 \(q\) 个白噪声的线性组合：

\[X_t = \varepsilon_t + \theta_1 \varepsilon_{t-1} + \theta_2 \varepsilon_{t-2} + \cdots + \theta_q \varepsilon_{t-q}\]

MA(q) 过程总是平稳的（有限个 \(\varepsilon\) 的和）。

\textbf{可逆性条件}：特征方程 \(1 + \theta_1 z + \cdots + \theta_q z^q = 0\) 的根全部落在单位圆外。

MA(1) 过程的性质：

\begin{itemize}
\tightlist
\item
  \(\rho(1) = \frac{\theta_1}{1 + \theta_1^2}\)，\(\rho(k) = 0\) 对 \(k > 1\)（ACF \textbf{截尾}）
\item
  PACF 拖尾（指数衰减）
\end{itemize}

\textbf{识别特征}：ACF 在滞后 \(q\) 处截尾，PACF 拖尾。

\subsection{ARMA(p, q) 模型}\label{armap-q-ux6a21ux578b}

ARMA(p, q) 结合了 AR 和 MA 两个成分：

\[X_t = \sum_{i=1}^p \phi_i X_{t-i} + \varepsilon_t + \sum_{j=1}^q \theta_j \varepsilon_{t-j}\]

\textbf{识别特征}：ACF 和 PACF 均拖尾。模型阶数选择常借助 AIC/BIC 准则。

\section{非平稳时间序列与 ARIMA}\label{ux975eux5e73ux7a33ux65f6ux95f4ux5e8fux5217ux4e0e-arima}

现实中的时间序列往往是非平稳的。最常见的情况是\textbf{差分平稳}------序列本身不平稳，但一阶或多阶差分后平稳。

\subsection{ARIMA(p, d, q)}\label{arimap-d-q}

ARIMA 模型在 ARMA 的基础上增加了差分阶数 \(d\)：

\[\nabla^d X_t = (1 - B)^d X_t\]

其中 \(B\) 是后移算子 \(B X_t = X_{t-1}\)，\(\nabla^d\) 是 \(d\) 阶差分。

ARIMA(p, d, q) 的建模步骤（Box-Jenkins 方法）：

\begin{enumerate}
\def\labelenumi{\arabic{enumi}.}
\tightlist
\item
  \textbf{识别}：观察时序图和 ACF/PACF 图，判断是否需要差分（\(d\) 的选择），初步判断 \(p\) 和 \(q\) 的范围。
\item
  \textbf{估计}：用最大似然法或条件最小二乘估计参数。
\item
  \textbf{诊断}：检验残差是否为白噪声（Ljung-Box 检验），若不通过则回到步骤 1。
\item
  \textbf{预测}：利用拟合模型进行区间预测。
\end{enumerate}

\begin{Shaded}
\begin{Highlighting}[]
\CommentTok{\# ARIMA 建模示例}
\FunctionTok{library}\NormalTok{(forecast)}
\CommentTok{\# 模拟 ARIMA(1,1,1) 过程}
\FunctionTok{set.seed}\NormalTok{(}\DecValTok{123}\NormalTok{)}
\NormalTok{y }\OtherTok{\textless{}{-}} \FunctionTok{arima.sim}\NormalTok{(}\AttributeTok{n =} \DecValTok{200}\NormalTok{, }\FunctionTok{list}\NormalTok{(}\AttributeTok{order =} \FunctionTok{c}\NormalTok{(}\DecValTok{1}\NormalTok{, }\DecValTok{1}\NormalTok{, }\DecValTok{1}\NormalTok{), }\AttributeTok{ar =} \FloatTok{0.7}\NormalTok{, }\AttributeTok{ma =} \SpecialCharTok{{-}}\FloatTok{0.3}\NormalTok{))}
\CommentTok{\# 自动选择 ARIMA 模型}
\NormalTok{fit }\OtherTok{\textless{}{-}} \FunctionTok{auto.arima}\NormalTok{(y)}
\FunctionTok{summary}\NormalTok{(fit)}
\end{Highlighting}
\end{Shaded}

\begin{verbatim}
## Series: y 
## ARIMA(2,1,2) 
## 
## Coefficients:
##           ar1     ar2     ma1      ma2
##       -0.1311  0.5963  0.5142  -0.4236
## s.e.   0.1539  0.1053  0.1698   0.1490
## 
## sigma^2 = 0.8612:  log likelihood = -267.14
## AIC=544.27   AICc=544.58   BIC=560.76
## 
## Training set error measures:
##                       ME      RMSE       MAE      MPE     MAPE      MASE
## Training set -0.01203173 0.9163854 0.7374012 22.28646 42.07816 0.9165299
##                     ACF1
## Training set -0.01183265
\end{verbatim}

\begin{Shaded}
\begin{Highlighting}[]
\CommentTok{\# 残差诊断}
\FunctionTok{checkresiduals}\NormalTok{(fit)}
\end{Highlighting}
\end{Shaded}

\begin{center}\includegraphics{statistic-book_files/figure-latex/unnamed-chunk-20-1} \end{center}

\begin{verbatim}
## 
##  Ljung-Box test
## 
## data:  Residuals from ARIMA(2,1,2)
## Q* = 2.6934, df = 6, p-value = 0.8462
## 
## Model df: 4.   Total lags used: 10
\end{verbatim}

\begin{Shaded}
\begin{Highlighting}[]
\CommentTok{\# 预测未来 20 期}
\FunctionTok{forecast}\NormalTok{(fit, }\AttributeTok{h =} \DecValTok{20}\NormalTok{) }\SpecialCharTok{|\textgreater{}} \FunctionTok{autoplot}\NormalTok{()}
\end{Highlighting}
\end{Shaded}

\begin{center}\includegraphics{statistic-book_files/figure-latex/unnamed-chunk-20-2} \end{center}

\subsection{单位根检验}\label{ux5355ux4f4dux6839ux68c0ux9a8c}

在建模之前，需要先判断序列是否需要差分。常用的检验包括：

\begin{itemize}
\tightlist
\item
  \textbf{ADF 检验}（Augmented Dickey-Fuller）：\(H_0\)：存在单位根（非平稳） vs \(H_1\)：平稳。
\item
  \textbf{KPSS 检验}：\(H_0\)：平稳 vs \(H_1\)：存在单位根。
\end{itemize}

两种检验的 \(H_0\) 方向相反，结合使用可以提高判断的可靠性。

\section{GARCH 模型与波动率建模}\label{garch-ux6a21ux578bux4e0eux6ce2ux52a8ux7387ux5efaux6a21}

金融时间序列的一个重要特征是\textbf{波动率聚集}（Volatility Clustering）------大的波动往往跟随大的波动（无论方向），因此虽然收益率序列本身可能是白噪声，但其\textbf{平方序列}或\textbf{绝对值序列}却呈现显著的自相关。

\subsection{ARCH(m) 模型}\label{archm-ux6a21ux578b}

ARCH（Autoregressive Conditional Heteroskedasticity）模型假设：

\[\varepsilon_t = \sigma_t z_t, \quad z_t \sim N(0, 1)\]
\[\sigma_t^2 = \alpha_0 + \alpha_1 \varepsilon_{t-1}^2 + \cdots + \alpha_m \varepsilon_{t-m}^2\]

\subsection{GARCH(m, s) 模型}\label{garchm-s-ux6a21ux578b}

GARCH 模型在 ARCH 基础上加入了条件方差自身的滞后项：

\[\sigma_t^2 = \alpha_0 + \sum_{i=1}^m \alpha_i \varepsilon_{t-i}^2 + \sum_{j=1}^s \beta_j \sigma_{t-j}^2\]

最常用的是 GARCH(1,1)：\(\sigma_t^2 = \alpha_0 + \alpha_1 \varepsilon_{t-1}^2 + \beta_1 \sigma_{t-1}^2\)

\begin{quote}
持久性条件：\(\alpha_1 + \beta_1 < 1\) 保证了 \(\sigma_t^2\) 的平稳性。\(\alpha_1 + \beta_1\) 越接近 1，波动率的持续性越强。
\end{quote}

\begin{Shaded}
\begin{Highlighting}[]
\CommentTok{\# GARCH 建模示例}
\FunctionTok{library}\NormalTok{(rugarch)}
\CommentTok{\# 手动模拟 GARCH(1,1) 数据}
\FunctionTok{set.seed}\NormalTok{(}\DecValTok{123}\NormalTok{)}
\NormalTok{n }\OtherTok{\textless{}{-}} \DecValTok{500}
\NormalTok{omega }\OtherTok{\textless{}{-}} \FloatTok{0.1}\NormalTok{; alpha }\OtherTok{\textless{}{-}} \FloatTok{0.1}\NormalTok{; beta }\OtherTok{\textless{}{-}} \FloatTok{0.8}
\NormalTok{sigma2 }\OtherTok{\textless{}{-}} \FunctionTok{numeric}\NormalTok{(n); y }\OtherTok{\textless{}{-}} \FunctionTok{numeric}\NormalTok{(n)}
\NormalTok{epsilon }\OtherTok{\textless{}{-}} \FunctionTok{rnorm}\NormalTok{(n)}
\NormalTok{sigma2[}\DecValTok{1}\NormalTok{] }\OtherTok{\textless{}{-}}\NormalTok{ omega }\SpecialCharTok{/}\NormalTok{ (}\DecValTok{1} \SpecialCharTok{{-}}\NormalTok{ alpha }\SpecialCharTok{{-}}\NormalTok{ beta)}
\NormalTok{y[}\DecValTok{1}\NormalTok{] }\OtherTok{\textless{}{-}} \FunctionTok{sqrt}\NormalTok{(sigma2[}\DecValTok{1}\NormalTok{]) }\SpecialCharTok{*}\NormalTok{ epsilon[}\DecValTok{1}\NormalTok{]}
\ControlFlowTok{for}\NormalTok{ (t }\ControlFlowTok{in} \DecValTok{2}\SpecialCharTok{:}\NormalTok{n) \{}
\NormalTok{  sigma2[t] }\OtherTok{\textless{}{-}}\NormalTok{ omega }\SpecialCharTok{+}\NormalTok{ alpha }\SpecialCharTok{*}\NormalTok{ y[t}\DecValTok{{-}1}\NormalTok{]}\SpecialCharTok{\^{}}\DecValTok{2} \SpecialCharTok{+}\NormalTok{ beta }\SpecialCharTok{*}\NormalTok{ sigma2[t}\DecValTok{{-}1}\NormalTok{]}
\NormalTok{  y[t] }\OtherTok{\textless{}{-}} \FunctionTok{sqrt}\NormalTok{(sigma2[t]) }\SpecialCharTok{*}\NormalTok{ epsilon[t]}
\NormalTok{\}}
\CommentTok{\# 拟合 GARCH}
\NormalTok{spec }\OtherTok{\textless{}{-}} \FunctionTok{ugarchspec}\NormalTok{(}
  \AttributeTok{variance.model =} \FunctionTok{list}\NormalTok{(}\AttributeTok{model =} \StringTok{"sGARCH"}\NormalTok{, }\AttributeTok{garchOrder =} \FunctionTok{c}\NormalTok{(}\DecValTok{1}\NormalTok{, }\DecValTok{1}\NormalTok{)),}
  \AttributeTok{mean.model =} \FunctionTok{list}\NormalTok{(}\AttributeTok{armaOrder =} \FunctionTok{c}\NormalTok{(}\DecValTok{0}\NormalTok{, }\DecValTok{0}\NormalTok{), }\AttributeTok{include.mean =} \ConstantTok{TRUE}\NormalTok{),}
  \AttributeTok{distribution.model =} \StringTok{"norm"}
\NormalTok{)}
\NormalTok{fit }\OtherTok{\textless{}{-}} \FunctionTok{ugarchfit}\NormalTok{(spec, y, }\AttributeTok{solver =} \StringTok{"hybrid"}\NormalTok{)}
\FunctionTok{show}\NormalTok{(fit)}
\end{Highlighting}
\end{Shaded}

\begin{verbatim}
## 
## *---------------------------------*
## *          GARCH Model Fit        *
## *---------------------------------*
## 
## Conditional Variance Dynamics    
## -----------------------------------
## GARCH Model  : sGARCH(1,1)
## Mean Model   : ARFIMA(0,0,0)
## Distribution : norm 
## 
## Optimal Parameters
## ------------------------------------
##         Estimate  Std. Error    t value Pr(>|t|)
## mu      0.029629    0.041371 7.1617e-01  0.47389
## omega   0.001013    0.002640 3.8347e-01  0.70137
## alpha1  0.000000    0.002946 2.0000e-06  1.00000
## beta1   0.999000    0.000057 1.7536e+04  0.00000
## 
## Robust Standard Errors:
##         Estimate  Std. Error    t value Pr(>|t|)
## mu      0.029629    0.032993 8.9804e-01  0.36916
## omega   0.001013    0.002616 3.8709e-01  0.69869
## alpha1  0.000000    0.002869 2.0000e-06  1.00000
## beta1   0.999000    0.000058 1.7184e+04  0.00000
## 
## LogLikelihood : -670.6527 
## 
## Information Criteria
## ------------------------------------
##                    
## Akaike       2.6986
## Bayes        2.7323
## Shibata      2.6985
## Hannan-Quinn 2.7118
## 
## Weighted Ljung-Box Test on Standardized Residuals
## ------------------------------------
##                         statistic p-value
## Lag[1]                      1.791  0.1808
## Lag[2*(p+q)+(p+q)-1][2]     2.561  0.1846
## Lag[4*(p+q)+(p+q)-1][5]     3.855  0.2727
## d.o.f=0
## H0 : No serial correlation
## 
## Weighted Ljung-Box Test on Standardized Squared Residuals
## ------------------------------------
##                         statistic p-value
## Lag[1]                      5.300 0.02132
## Lag[2*(p+q)+(p+q)-1][5]     6.915 0.05439
## Lag[4*(p+q)+(p+q)-1][9]     7.812 0.14019
## d.o.f=2
## 
## Weighted ARCH LM Tests
## ------------------------------------
##             Statistic Shape Scale P-Value
## ARCH Lag[3]     1.251 0.500 2.000  0.2633
## ARCH Lag[5]     2.142 1.440 1.667  0.4407
## ARCH Lag[7]     2.416 2.315 1.543  0.6298
## 
## Nyblom stability test
## ------------------------------------
## Joint Statistic:  1.4952
## Individual Statistics:              
## mu     0.08175
## omega  0.05951
## alpha1 0.05504
## beta1  0.05863
## 
## Asymptotic Critical Values (10% 5% 1%)
## Joint Statistic:          1.07 1.24 1.6
## Individual Statistic:     0.35 0.47 0.75
## 
## Sign Bias Test
## ------------------------------------
##                    t-value    prob sig
## Sign Bias           1.7621 0.07867   *
## Negative Sign Bias  0.7783 0.43678    
## Positive Sign Bias  2.3162 0.02095  **
## Joint Effect        8.2653 0.04084  **
## 
## 
## Adjusted Pearson Goodness-of-Fit Test:
## ------------------------------------
##   group statistic p-value(g-1)
## 1    20     14.56       0.7501
## 2    30     32.92       0.2810
## 3    40     24.80       0.9625
## 4    50     52.20       0.3507
## 
## 
## Elapsed time : 0.324507
\end{verbatim}

\section{状态空间模型与卡尔曼滤波}\label{ux72b6ux6001ux7a7aux95f4ux6a21ux578bux4e0eux5361ux5c14ux66fcux6ee4ux6ce2}

状态空间模型（State Space Model）为时间序列分析提供了一个统一的框架。它将系统分为两部分：

\begin{itemize}
\tightlist
\item
  \textbf{状态方程}：\(\alpha_t = T_t \alpha_{t-1} + R_t \eta_t\)（不可观测的状态如何演化）
\item
  \textbf{观测方程}：\(y_t = Z_t \alpha_t + \varepsilon_t\)（状态如何产生可观测数据）
\end{itemize}

ARIMA、结构性时间序列、时变参数模型都可以写成状态空间形式。

\textbf{卡尔曼滤波}（Kalman Filter）是一组递推公式，用于在线性高斯状态空间模型中，根据截止到当前时刻的观测，最优地估计当前状态及其不确定性。卡尔曼平滑则利用全部观测对历史状态进行回顾性估计。

\section{与本书其他章节的联系}\label{ux4e0eux672cux4e66ux5176ux4ed6ux7ae0ux8282ux7684ux8054ux7cfb}

\begin{itemize}
\tightlist
\item
  ARIMA 模型的估计使用 Ch06 中的最大似然方法。
\item
  GARCH 模型的条件方差方程是 Ch10 回归模型的自然延伸。
\item
  状态空间模型与 Ch09 的贝叶斯动态模型有直接联系。
\item
  时间序列的预测问题与 Ch17 中的序列建模（RNN/LSTM）形成互补视角。
\end{itemize}

\begin{center}\rule{0.5\linewidth}{0.5pt}\end{center}

\textbf{速查卡片}：

\begin{itemize}
\tightlist
\item
  \textbf{核心问题}：如何处理观测值之间存在时间依赖的数据？
\item
  \textbf{关键概念}：平稳性、ACF/PACF、ARMA、差分、GARCH
\item
  \textbf{建模流程}：识别(看ACF/PACF) → 估计(最大似然) → 诊断(残差白噪声检验) → 预测
\item
  \textbf{R命令}：\texttt{arima()}, \texttt{auto.arima()}, \texttt{ugarchfit()}, \texttt{KalmanForecast()}
\item
  \textbf{注意}：伪回归！用非平稳序列直接回归会得到虚假的高 \(R^2\)
\item
  \textbf{参见}：Ch06（MLE）, Ch10（回归分析）, Ch09（贝叶斯计算）
\end{itemize}

\chapter{优化方法}\label{ux4f18ux5316ux65b9ux6cd5}

\begin{quote}
当解析解不可得时，数值优化是统计计算的基石。从最大似然估计到深度学习，优化算法无处不在。
\end{quote}

统计计算让统计推断成为可能。理论告知计算什么，计算告知如何真正计算。本章从最基本的求根方法出发，逐步深入到现代机器学习中广泛使用的随机梯度下降。

\section{优化 Optimisation}\label{ux4f18ux5316-optimisation}

\subsection{牛顿法 Newton Method}\label{ux725bux987fux6cd5-newton-method}

\begin{enumerate}
\def\labelenumi{\arabic{enumi}.}
\item
  二分法是寻找一元函数根的简单方法。

  二分法的基本步骤为：若\(f(a)\)与\(f(b)\)异号，则区间\((a, b)\)内必有根。取中点\(c = (a+b)/2\)，计算\(f(c)\)。若\(f(c) = 0\)，则\(c\)即为根；若\(f(a)f(c) < 0\)，根在\((a, c)\)内；否则根在\((c, b)\)内。每次迭代将区间长度减半，经过\(k\)次迭代后精度达到\((b-a)/2^k\)。二分法保证收敛但收敛速度较慢（线性收敛）。
\item
  牛顿法是一种迭代算法，需要使用原函数和函数的导数，如果有解则会收敛很快。

  迭代步骤： \[
    x_{n+1} = x_n - \frac{f(x_n)}{f'(x_n)}
    \]

  停止规则：牛顿法如果无法到达精确解可能会迭代很久，因此必须设定一个停止规则。我们取一个tolerance level，如果\(|f(x^*)|<\epsilon\)我们认为计算解已经十分接近真实解了

  牛顿法的几何直观为：在点\((x_n, f(x_n))\)处作函数\(f(x)\)的切线，切线与\(x\)轴的交点即为下一次迭代值\(x_{n+1}\)。切线的斜率为\(f'(x_n)\)，由点斜式方程可得\(x_{n+1} = x_n - f(x_n)/f'(x_n)\)。当初始值足够靠近真实根时，牛顿法具有二次收敛速度，远快于二分法。
\item
  比较两种方法的收敛速率

  \begin{enumerate}
  \def\labelenumii{\arabic{enumii}.}
  \item
    线性收敛（比如二分法）

    \[
    \frac{||x_{n+1}-x_{\infty}||}{||x_n-x_{\infty}||}\leq r,r \in (0,1)
    \]
  \item
    超线性收敛

    \[
    lim_{n\rightarrow\infty}\frac{||x_{n+1}-x_{\infty}||}{||x_n-x_{\infty}||}=0
    \]
  \item
    二次收敛（如牛顿法，使用泰勒展开分析）

    \[
    \frac{||x_{n+1}-x_{\infty}||}{||x_n-x_{\infty}||^2}\leq M
    \]
  \end{enumerate}
\item
  使用牛顿法的注意事项

\begin{Shaded}
\begin{Highlighting}[]
\CommentTok{\# Newton–Raphson using deriv() for automatic derivative}
\NormalTok{newton\_raphson }\OtherTok{\textless{}{-}} \ControlFlowTok{function}\NormalTok{(expr, x0, }\AttributeTok{tol =} \FloatTok{1e{-}8}\NormalTok{, }\AttributeTok{max\_iter =} \DecValTok{100}\NormalTok{, }\AttributeTok{verbose =} \ConstantTok{TRUE}\NormalTok{) \{}
  \CommentTok{\# expr should be an expression, e.g. expression(x\^{}2 {-} 5)}
\NormalTok{  var }\OtherTok{\textless{}{-}} \FunctionTok{all.vars}\NormalTok{(expr)}
  \ControlFlowTok{if}\NormalTok{ (}\FunctionTok{length}\NormalTok{(var) }\SpecialCharTok{!=} \DecValTok{1}\NormalTok{) }\FunctionTok{stop}\NormalTok{(}\StringTok{"Only single{-}variable functions supported."}\NormalTok{)}
  \CommentTok{\# Create a function that returns both value and gradient}
\NormalTok{  f\_deriv }\OtherTok{\textless{}{-}} \FunctionTok{deriv}\NormalTok{(expr, var)}
  \CommentTok{\# Initialize}
\NormalTok{  x }\OtherTok{\textless{}{-}}\NormalTok{ x0}
  \ControlFlowTok{for}\NormalTok{ (i }\ControlFlowTok{in} \DecValTok{1}\SpecialCharTok{:}\NormalTok{max\_iter) \{}
    \CommentTok{\# Evaluate f and f\textquotesingle{}}
\NormalTok{    val }\OtherTok{\textless{}{-}} \FunctionTok{eval}\NormalTok{(f\_deriv, }\AttributeTok{envir =} \FunctionTok{list}\NormalTok{(}\AttributeTok{x =}\NormalTok{ x))}
\NormalTok{    fx  }\OtherTok{\textless{}{-}} \FunctionTok{as.numeric}\NormalTok{(val)        }\CommentTok{\# 函数值}
\NormalTok{    fpx }\OtherTok{\textless{}{-}} \FunctionTok{attr}\NormalTok{(val, }\StringTok{"gradient"}\NormalTok{)  }\CommentTok{\# 导数值}
    \CommentTok{\# Newton–Raphson update}
    \ControlFlowTok{if}\NormalTok{ (}\FunctionTok{abs}\NormalTok{(fpx) }\SpecialCharTok{\textless{}} \FloatTok{1e{-}12}\NormalTok{) \{}
      \FunctionTok{warning}\NormalTok{(}\StringTok{"Derivative too close to zero. Stopping."}\NormalTok{)}
      \ControlFlowTok{break}
\NormalTok{    \}}
\NormalTok{    x\_new }\OtherTok{\textless{}{-}}\NormalTok{ x }\SpecialCharTok{{-}}\NormalTok{ fx }\SpecialCharTok{/}\NormalTok{ fpx}
    \ControlFlowTok{if}\NormalTok{ (verbose)}
      \FunctionTok{cat}\NormalTok{(}\FunctionTok{sprintf}\NormalTok{(}\StringTok{"Iter \%d: x = \%.8f, f(x) = \%.8f}\SpecialCharTok{\textbackslash{}n}\StringTok{"}\NormalTok{, i, x\_new, fx))}
    \CommentTok{\# Check convergence}
    \ControlFlowTok{if}\NormalTok{ (}\FunctionTok{abs}\NormalTok{(x\_new }\SpecialCharTok{{-}}\NormalTok{ x) }\SpecialCharTok{\textless{}}\NormalTok{ tol) \{}
      \ControlFlowTok{if}\NormalTok{ (verbose) }\FunctionTok{cat}\NormalTok{(}\StringTok{"Converged!}\SpecialCharTok{\textbackslash{}n}\StringTok{"}\NormalTok{)}
      \FunctionTok{return}\NormalTok{(x\_new)}
\NormalTok{    \}}
\NormalTok{    x }\OtherTok{\textless{}{-}}\NormalTok{ x\_new}
\NormalTok{  \}}
  \FunctionTok{warning}\NormalTok{(}\StringTok{"Did not converge within max\_iter."}\NormalTok{)}
  \FunctionTok{return}\NormalTok{(x)}
\NormalTok{\}}

\CommentTok{\# 调用牛顿法迭代}
\NormalTok{root }\OtherTok{\textless{}{-}} \FunctionTok{newton\_raphson}\NormalTok{(}\FunctionTok{expression}\NormalTok{(x}\SpecialCharTok{\^{}}\DecValTok{2} \SpecialCharTok{{-}} \DecValTok{5}\NormalTok{), }\AttributeTok{x0 =} \DecValTok{1}\NormalTok{)}
\end{Highlighting}
\end{Shaded}

\begin{verbatim}
## Iter 1: x = 3.00000000, f(x) = -4.00000000
## Iter 2: x = 2.33333333, f(x) = 4.00000000
## Iter 3: x = 2.23809524, f(x) = 0.44444444
## Iter 4: x = 2.23606890, f(x) = 0.00907029
## Iter 5: x = 2.23606798, f(x) = 0.00000411
## Iter 6: x = 2.23606798, f(x) = 0.00000000
## Converged!
\end{verbatim}

\begin{Shaded}
\begin{Highlighting}[]
\CommentTok{\#root \textless{}{-} newton\_raphson(expression(x\^{}3 {-} 2*x\^{}2 {-}11*x +12),x0 = 2.35287527)}
\CommentTok{\#root \textless{}{-} newton\_raphson(expression(x\^{}3 {-} 2*x\^{}2 {-}11*x +12),x0 = 2.35284172)}

\CommentTok{\#root \textless{}{-} newton\_raphson(expression(2*x\^{}3+3*x\^{}2+5),x0 = 0.5)}
\CommentTok{\#root \textless{}{-} newton\_raphson(expression(2*x\^{}3+3*x\^{}2+5),x0 = 0)}
\CommentTok{\#注1： Derivative too close to zero. Stopping.}
\CommentTok{\#注2： Did not converge within max\_iter.}
\end{Highlighting}
\end{Shaded}
\item
  小结：牛顿法并不总是能够收敛；初始值的选取很重要，不同初始值可能导致收敛到不同的根；无法计算一阶导数为零的情况；如果函数在根的邻域内是不可微的，牛顿法会无法收敛；
\item
  利用牛顿法求极值

  关注步长的表达式\(-\frac{f'(x)}{f''(x)}\)，符号表示了下一步的迭代方向，大小表示了下一步的迭代步长。二阶导数的绝对值越大，步长会越小。
\item
  多变量优化

  \begin{enumerate}
  \def\labelenumii{\arabic{enumii}.}
  \item
    在确定风险效率最高的投资组合时，回报率是许多权重（每种资产一个权重）的函数。在线性回归模型的最小二乘估计中，平方和是许多系数（每个回归因子一个系数）的函数。
  \item
    当目标函数的自变量是向量的时候，一阶导数计算梯度向量\(\frac{\partial y}{\partial X}\)（对每个变量求偏导数，合并为一个列向量），二阶导数计算海塞矩阵\(\frac{\partial^2y}{\partial X\partial X'}\)，这样牛顿法迭代公式可写为：\[
    x_{n+1} = x_n - (\frac{\partial^2f(x)}{\partial x\partial x'})^{-1}\frac{\partial f(x)}{\partial(x)}
    \]
  \end{enumerate}
\item
  线性模型的极大似然估计

  \[
  y_i = \beta_0 + \beta_1x_i+\epsilon_i,\epsilon_i\sim N(0,1)
  \]

\begin{Shaded}
\begin{Highlighting}[]
\DocumentationTok{\#\# Generate some data}
\NormalTok{beta0 }\OtherTok{\textless{}{-}} \DecValTok{1}
\NormalTok{beta1 }\OtherTok{\textless{}{-}} \DecValTok{3}
\NormalTok{sigma }\OtherTok{\textless{}{-}} \DecValTok{1}
\NormalTok{n }\OtherTok{\textless{}{-}} \DecValTok{1000}
\NormalTok{x }\OtherTok{\textless{}{-}} \FunctionTok{rnorm}\NormalTok{(n, }\DecValTok{3}\NormalTok{, }\DecValTok{1}\NormalTok{)}
\NormalTok{y }\OtherTok{\textless{}{-}}\NormalTok{ beta0 }\SpecialCharTok{+}\NormalTok{ x }\SpecialCharTok{*}\NormalTok{ beta1 }\SpecialCharTok{+} \FunctionTok{rnorm}\NormalTok{(n, }\AttributeTok{mean =} \DecValTok{0}\NormalTok{, }\AttributeTok{sd =}\NormalTok{ sigma)}
\DocumentationTok{\#\# Make the log normal likelihood function}
\NormalTok{func }\OtherTok{=} \ControlFlowTok{function}\NormalTok{(beta)\{}
  \FunctionTok{sum}\NormalTok{((y}\SpecialCharTok{{-}}\NormalTok{beta[}\DecValTok{1}\NormalTok{]}\SpecialCharTok{{-}}\NormalTok{beta[}\DecValTok{2}\NormalTok{]}\SpecialCharTok{*}\NormalTok{x)}\SpecialCharTok{\^{}}\DecValTok{2}\NormalTok{)}
\NormalTok{\}}
\NormalTok{grad }\OtherTok{=} \ControlFlowTok{function}\NormalTok{(beta)\{}
  \FunctionTok{matrix}\NormalTok{(}\FunctionTok{c}\NormalTok{(}\FunctionTok{sum}\NormalTok{(}\SpecialCharTok{{-}}\DecValTok{2}\SpecialCharTok{*}\NormalTok{(y}\SpecialCharTok{{-}}\NormalTok{beta[}\DecValTok{1}\NormalTok{]}\SpecialCharTok{{-}}\NormalTok{beta[}\DecValTok{2}\NormalTok{]}\SpecialCharTok{*}\NormalTok{x)),}\FunctionTok{sum}\NormalTok{(}\SpecialCharTok{{-}}\DecValTok{2}\SpecialCharTok{*}\NormalTok{x}\SpecialCharTok{*}\NormalTok{(y}\SpecialCharTok{{-}}\NormalTok{beta[}\DecValTok{1}\NormalTok{]}\SpecialCharTok{{-}}\NormalTok{beta[}\DecValTok{2}\NormalTok{]}\SpecialCharTok{*}\NormalTok{x))),}\DecValTok{2}\NormalTok{,}\DecValTok{1}\NormalTok{)}
\NormalTok{\}}
\NormalTok{hess }\OtherTok{=} \ControlFlowTok{function}\NormalTok{(beta)\{}
  \FunctionTok{matrix}\NormalTok{(}\FunctionTok{c}\NormalTok{(}\DecValTok{2}\SpecialCharTok{*}\FunctionTok{length}\NormalTok{(x),}\DecValTok{2}\SpecialCharTok{*}\FunctionTok{sum}\NormalTok{(x),}\DecValTok{2}\SpecialCharTok{*}\FunctionTok{sum}\NormalTok{(x),}\DecValTok{2}\SpecialCharTok{*}\FunctionTok{sum}\NormalTok{(x}\SpecialCharTok{\^{}}\DecValTok{2}\NormalTok{)),}\DecValTok{2}\NormalTok{,}\DecValTok{2}\NormalTok{)}
\NormalTok{\}}
\DocumentationTok{\#\# 定义牛顿优化函数}
\NormalTok{newton }\OtherTok{\textless{}{-}} \ControlFlowTok{function}\NormalTok{(f, init, }\AttributeTok{maxit =} \DecValTok{100}\NormalTok{, }\AttributeTok{epsilon =} \FloatTok{1e{-}6}\NormalTok{) \{}
\NormalTok{  beta }\OtherTok{\textless{}{-}}\NormalTok{ init}
  \ControlFlowTok{for}\NormalTok{ (i }\ControlFlowTok{in} \DecValTok{1}\SpecialCharTok{:}\NormalTok{maxit) \{}
\NormalTok{    res }\OtherTok{\textless{}{-}} \FunctionTok{f}\NormalTok{(beta)}
\NormalTok{    g }\OtherTok{\textless{}{-}}\NormalTok{ res}\SpecialCharTok{$}\NormalTok{gradient}
\NormalTok{    H }\OtherTok{\textless{}{-}}\NormalTok{ res}\SpecialCharTok{$}\NormalTok{hessian}
    \ControlFlowTok{if}\NormalTok{ (}\FunctionTok{max}\NormalTok{(}\FunctionTok{abs}\NormalTok{(g)) }\SpecialCharTok{\textless{}}\NormalTok{ epsilon) }\ControlFlowTok{break}
\NormalTok{    H\_inv }\OtherTok{\textless{}{-}} \FunctionTok{solve}\NormalTok{(H)}
\NormalTok{    beta }\OtherTok{\textless{}{-}}\NormalTok{ beta }\SpecialCharTok{{-}} \FunctionTok{as.vector}\NormalTok{(H\_inv }\SpecialCharTok{\%*\%}\NormalTok{ g)}
\NormalTok{  \}}
  \FunctionTok{return}\NormalTok{(beta)}
\NormalTok{\}}
\DocumentationTok{\#\# Newton optimization}
\NormalTok{optimOut }\OtherTok{\textless{}{-}} \FunctionTok{newton}\NormalTok{(}
  \ControlFlowTok{function}\NormalTok{(beta)\{}
  \FunctionTok{list}\NormalTok{(}\AttributeTok{value=}\FunctionTok{func}\NormalTok{(beta), }\AttributeTok{gradient=}\FunctionTok{grad}\NormalTok{(beta), }\AttributeTok{hessian=}\FunctionTok{hess}\NormalTok{(beta))\}}
\NormalTok{  , }\FunctionTok{c}\NormalTok{(}\SpecialCharTok{{-}}\DecValTok{4}\NormalTok{, }\SpecialCharTok{{-}}\DecValTok{5}\NormalTok{))  }

\NormalTok{optimOut}
\end{Highlighting}
\end{Shaded}

\begin{verbatim}
## [1] 0.8166939 3.0599782
\end{verbatim}

\begin{Shaded}
\begin{Highlighting}[]
\NormalTok{beta0Hat }\OtherTok{\textless{}{-}}\NormalTok{ optimOut[}\DecValTok{1}\NormalTok{]}
\NormalTok{beta1Hat }\OtherTok{\textless{}{-}}\NormalTok{ optimOut[}\DecValTok{2}\NormalTok{]}
\NormalTok{myLM }\OtherTok{\textless{}{-}} \FunctionTok{lm}\NormalTok{(y}\SpecialCharTok{\textasciitilde{}}\NormalTok{x)}
\NormalTok{myLMCoef }\OtherTok{\textless{}{-}}\NormalTok{ myLM}\SpecialCharTok{$}\NormalTok{coefficients}
\NormalTok{coefs }\OtherTok{\textless{}{-}} \FunctionTok{data.frame}\NormalTok{(}\AttributeTok{Method =} \FunctionTok{c}\NormalTok{(}\StringTok{\textquotesingle{}Newton\textquotesingle{}}\NormalTok{, }\StringTok{\textquotesingle{}LM\textquotesingle{}}\NormalTok{), }\AttributeTok{intercept =} \FunctionTok{c}\NormalTok{(beta0Hat, myLMCoef[}\DecValTok{1}\NormalTok{]), }\AttributeTok{slope =} \FunctionTok{c}\NormalTok{(beta1Hat,myLMCoef[}\DecValTok{2}\NormalTok{]))}
\NormalTok{coefs}
\end{Highlighting}
\end{Shaded}

\begin{verbatim}
##             Method intercept    slope
##             Newton 0.8166939 3.059978
## (Intercept)     LM 0.8166939 3.059978
\end{verbatim}

\begin{Shaded}
\begin{Highlighting}[]
\FunctionTok{library}\NormalTok{(ggplot2)}
\FunctionTok{qplot}\NormalTok{(x, y) }\SpecialCharTok{+}
  \FunctionTok{geom\_abline}\NormalTok{(}\FunctionTok{aes}\NormalTok{(}\AttributeTok{slope =}\NormalTok{ myLMCoef[}\DecValTok{2}\NormalTok{], }
                  \AttributeTok{intercept =}\NormalTok{ myLMCoef[}\DecValTok{1}\NormalTok{], }\AttributeTok{color =} \StringTok{\textquotesingle{}LM\textquotesingle{}}\NormalTok{), }\AttributeTok{size =} \DecValTok{3}\NormalTok{) }\SpecialCharTok{+}
   \FunctionTok{geom\_abline}\NormalTok{(}\FunctionTok{aes}\NormalTok{(}\AttributeTok{slope =}\NormalTok{ beta1Hat, }\AttributeTok{intercept =}\NormalTok{ beta0Hat, }\AttributeTok{color=}\StringTok{\textquotesingle{}Newton\textquotesingle{}}\NormalTok{),}
               \AttributeTok{linetype=}\StringTok{"dashed"}\NormalTok{, }\AttributeTok{linewidth =} \DecValTok{1}\NormalTok{) }\SpecialCharTok{+} \FunctionTok{labs}\NormalTok{(}\AttributeTok{color =} \StringTok{"Method"}\NormalTok{)}
\end{Highlighting}
\end{Shaded}

  \begin{center}\includegraphics{statistic-book_files/figure-latex/unnamed-chunk-25-1} \end{center}

\begin{Shaded}
\begin{Highlighting}[]
\CommentTok{\# nlm()函数在R中实现了给定初始值，用牛顿法最小化多元函数的方法}
\FunctionTok{nlm}\NormalTok{(func, }\FunctionTok{c}\NormalTok{(}\SpecialCharTok{{-}}\DecValTok{4}\NormalTok{, }\DecValTok{5}\NormalTok{))}\SpecialCharTok{$}\NormalTok{estimate}
\end{Highlighting}
\end{Shaded}

\begin{verbatim}
## [1] 0.8167352 3.0599643
\end{verbatim}
\end{enumerate}

\subsection{单纯形法 Nelder Mead Algorithm}\label{ux5355ux7eafux5f62ux6cd5-nelder-mead-algorithm}

牛顿法用函数一阶导数和二阶导数构建局部二次近似模型，通过迭代求解近似模型的最优解，逐步逼近原函数的最优解。当二阶导数（海塞矩阵）难以计算时，用函数值和一阶导数信息迭代近似海森矩阵，称作拟牛顿法，如BFGS。

当导数不存在时，最常用的是Nelder-Mead单纯形法。

\begin{itemize}
\item
  第一步：选点
\item
  第二步：找重心
\item
  第三步：找反射点
\item
  第四步：分情况得到下一步的迭代点

  具体迭代规则如下：

  \begin{itemize}
  \tightlist
  \item
    若反射点的函数值优于当前最优点：尝试沿反射方向\textbf{扩张（Expansion）}，若扩张点更优则接受扩张点，否则接受反射点。
  \item
    若反射点的函数值介于最好点和最差点之间：\textbf{接受反射点}。
  \item
    若反射点的函数值差于当前所有点：进行\textbf{收缩（Contraction）}，向内收缩或向外收缩，取较优者。
  \item
    若收缩后仍无明显改善：执行\textbf{全局缩小（Shrink）}，所有点向当前最优点靠拢。
  \end{itemize}

  通过以上四种操作，单纯形在参数空间中逐渐缩小并收敛到局部极小值点。
\end{itemize}

\subsection{随机梯度下降 SGD}\label{ux968fux673aux68afux5ea6ux4e0bux964d-sgd}

\[
x_{n+1} = x_n - \alpha \nabla f(x_n),\alpha >0
\]

对比牛顿法和梯度下降法，可以看出牛顿法更快收敛到极值，因为牛顿法用了函数的二阶导数信息，而梯度下降法只用了一阶梯度的信息。

两种方法的比较：
- \textbf{牛顿法}：利用了二阶导数（海塞矩阵）信息，构造局部二次近似模型，收敛速度快（二次收敛），但每步需计算和求逆海塞矩阵，计算成本高。
- \textbf{梯度下降法}：只使用一阶梯度信息，每次迭代计算量小，但收敛慢（线性收敛）。在极值点附近，梯度接近零，步长自动变小；而远离极值点时步长较大。
- \textbf{随机梯度下降（SGD）}：每次只用一个（或小批量）样本的梯度更新参数，适用于大规模数据和在线学习场景，但收敛路径带有随机波动。

什么时候使用随机梯度下降？当训练时间久，样本量大的时候。

\chapter{矩阵计算与数值方法}\label{ux77e9ux9635ux8ba1ux7b97ux4e0eux6570ux503cux65b9ux6cd5}

\begin{quote}
线性代数不仅是数学工具，更是统计计算的引擎。从最小二乘到主成分分析，矩阵分解提供了统一的计算框架。
\end{quote}

\section{Eigenanalysis 特征分析}\label{eigenanalysis-ux7279ux5f81ux5206ux6790}

特征分析可以帮助我们做什么？

特征分析的应用：谷歌的页面排序

\subsection{SVD分解}\label{svdux5206ux89e3}

方阵可以对角化。

非方阵没有特征值，但是可以考虑奇异值。

在 SVD 中，我们不要求输出向量和输入向量在同一个空间：\[Av = \sigma u\]这里，\(v\) 是输入空间的向量，\(u\) 是输出空间的向量，\(\sigma\) 是奇异值。这解决了维度不匹配的问题。

对于任意矩阵 \(A\)，其奇异值 \(\sigma\) 定义为：\[\sigma_i = \sqrt{\lambda_i(A^T A)}\]即：奇异值是 \(A^T A\)（或 \(AA^T\)）特征值的算术平方根。

在公式 \(A = U \Sigma V^T\) 中：

右奇异向量 (\(V\) 的列，记作 \(v_i\))：它们存在于输入空间（原空间）。

\begin{itemize}
\item
  几何意义：它们是单位球面上那些被矩阵 \(A\) 拉伸后，恰好指向椭球各个``半轴''方向的原始向量。
\item
  数学背景：它们是方阵 \(A^T A\) 的特征向量。
\end{itemize}

左奇异向量 (\(U\) 的列，记作 \(u_i\))：它们存在于输出空间（目标空间）。

\begin{itemize}
\item
  几何意义：它们直接定义了变换后椭球的主轴方向。
\item
  数学背景：它们是方阵 \(AA^T\) 的特征向量。
\end{itemize}

\begin{enumerate}
\def\labelenumi{\arabic{enumi}.}
\item
  数学定义：

  对于任何一个 \(m \times n\) 的实数矩阵 \(A\)，SVD 将其分解为三个特定矩阵的乘积：\[A = U \Sigma V^T\]

  \begin{itemize}
  \item
    \(U\) (\(m \times m\))：左奇异向量矩阵。它是正交矩阵，其列向量是 \(AA^T\) 的特征向量。在数据分析中，它通常代表样本的特征。
  \item
    \(\Sigma\) (\(m \times n\))：奇异值矩阵。除对角线外全是 0，对角线上的值 \(\sigma_i\) 称为奇异值（Singular Values），按降序排列。它们衡量了每个维度的重要性。
  \item
    \(V^T\) (\(n \times n\))：右奇异向量矩阵的转置。它也是正交矩阵，其列向量（\(V\) 的列）是 \(A^TA\) 的特征向量。它通常代表变量（特征）的线性组合。
  \end{itemize}
\item
  几何直观意义：拉伸与旋转
\end{enumerate}

SVD 的本质是说：任何线性变换都可以分解为三个简单的步骤：

\begin{itemize}
\item
  旋转：由 \(V^T\) 完成，将原始坐标系旋转到新的方向。
\item
  拉伸：由 \(\Sigma\) 完成，将圆形的单位球在新的轴向上拉伸成一个``超椭圆''。奇异值就是这些半轴的长度。
\item
  再旋转：由 \(U\) 完成，将拉伸后的椭圆旋转到目标空间的方向。
\end{itemize}

\begin{enumerate}
\def\labelenumi{\arabic{enumi}.}
\setcounter{enumi}{2}
\tightlist
\item
  SVD 在现实中有什么用？
\end{enumerate}

A. 降维与 PCA（主成分分析）在处理海量数据时，我们往往只需要保留前几个最大的奇异值及其对应的向量。这可以保留数据中绝大部分的信息，同时滤掉噪声。例子：一个有 100 个特征的数据库，通过 SVD 可能发现前 5 个成分就解释了 90\% 的差异。

B. 图像压缩一张高清图片可以看作一个巨大的矩阵。通过 SVD 分解并只保留前 \(k\) 个较大的奇异值，我们可以用极小的存储空间还原出非常清晰的图像轮廓。

C. 推荐系统（如 Netflix）SVD 可以发现用户和电影之间的``潜在语义''。左奇异向量：代表用户的隐藏偏好。右奇异向量：代表电影的隐藏风格。通过这两个矩阵相乘，算法可以预测你对没看过的电影会打几分。

SVD的性质

对于矩阵 \(M = A^T A\)，其条件数 \(\kappa(M)\) 通常定义为最大奇异值与最小奇异值的比值：\[\kappa(A^T A) = \frac{\sigma_{max}}{\sigma_{min}}\]

如果比值接近1，说明矩阵远离奇异；如果比值接近正无穷，说明矩阵接近奇异。

\subsection{矩阵估计}\label{ux77e9ux9635ux4f30ux8ba1}

Economy SVD

Economy SVD（经济型SVD）：对于\(m \times n\)的矩阵\(A\)（设\(m \geq n\)），标准SVD分解\(A = U_{m \times m} \Sigma_{m \times n} V^T_{n \times n}\)。但\(\Sigma\)中只有前\(n\)个奇异值非零，因此可以只保留\(U\)的前\(n\)列、\(\Sigma\)的前\(n\)行和\(V\)的全部列，得到经济型分解：

\[A = U_{m \times n} \Sigma_{n \times n} V^T_{n \times n}\]

这种紧凑表示显著减少了存储和计算开销，是实践中SVD的默认使用方式。在R中，\texttt{svd(A)\$u}默认返回的就是经济型\(U\)矩阵。

基于奇异值分解的图像压缩（Image Compression）

SVD 分解与低秩近似

对图像矩阵 \(A\)（或彩色图的 \(A_R, A_G, A_B\) 矩阵）进行奇异值分解：\[A = U \Sigma V^T = \sum_{i=1}^{r} \sigma_i u_i v_i^T\]其中 \(r\) 是矩阵的秩（等于 \(\min(m, n)\)）。

奇异值的意义：奇异值 \(\sigma_i\) 衡量了对应分量 \(u_i v_i^T\) 在图像中的重要性或贡献度。通常，大的奇异值对应于图像中的主要轮廓和低频信息；小的奇异值则对应于高频信息，如微小的纹理和噪声。

低秩近似（Low-Rank Approximation）：根据 Eckart--Young--Mirsky 定理，要找到最接近原矩阵 \(A\) 的秩为 \(k\) 的矩阵 \(A_k\)，只需保留前 \(k\) 个最大的奇异值：\[A_k = \sum_{i=1}^{k} \sigma_i u_i v_i^T \quad \text{其中 } k < r\]\(A_k\) 就是原矩阵 \(A\) 的最佳 \(k\) 阶近似。

\subsection{线性回归中关于矩阵不可逆的处理}\label{ux7ebfux6027ux56deux5f52ux4e2dux5173ux4e8eux77e9ux9635ux4e0dux53efux9006ux7684ux5904ux7406}

Pseudo-inverse（伪逆）：为那些没有逆矩阵的矩阵（非方阵或奇异矩阵）提供一个``类似于逆''的矩阵。

\section{最小二乘法 Least Squares}\label{ux6700ux5c0fux4e8cux4e58ux6cd5-least-squares}

\subsection{正规方程与矩阵分解}\label{ux6b63ux89c4ux65b9ux7a0bux4e0eux77e9ux9635ux5206ux89e3}

\begin{longtable}[]{@{}
  >{\raggedright\arraybackslash}p{(\columnwidth - 4\tabcolsep) * \real{0.2533}}
  >{\raggedright\arraybackslash}p{(\columnwidth - 4\tabcolsep) * \real{0.2400}}
  >{\raggedright\arraybackslash}p{(\columnwidth - 4\tabcolsep) * \real{0.5067}}@{}}
\toprule\noalign{}
\endhead
\bottomrule\noalign{}
\endlastfoot
\textbf{分解方法} & \textbf{适用矩阵} & \textbf{特点} \\
\textbf{LU 分解} & 任何方阵 & 通用，但速度较慢。 \\
\textbf{Cholesky 分解} & \textbf{对称正定方阵} & \textbf{最快}，内存占用最少，稳定性最高。 \\
\textbf{SVD 分解} & 任何矩阵 & 最万能，但计算代价最高。 \\
\end{longtable}

\subsection{QR分解}\label{qrux5206ux89e3}

\textbf{QR Decomposition（QR 分解）} 是一种将任意矩阵分解为一个正交矩阵和一个上三角矩阵乘积的技术。

对于\(m \times n\)的设计矩阵\(X\)（\(m \geq n\)），QR分解将其表示为\(X = QR\)，其中：
- \(Q\)为\(m \times n\)的正交矩阵，满足\(Q^T Q = I_n\)（列正交）
- \(R\)为\(n \times n\)的上三角矩阵

利用QR分解求解正规方程\(X'X\beta = X'y\)：代入\(X = QR\)得\(R'Q'QR\beta = R'Q'y\)，即\(R'R\beta = R'Q'y\)。由于\(R\)可逆（满秩时），化简得\(R\beta = Q'y\)。这是一个上三角线性方程组，通过回代（back substitution）即可高效求解\(\beta\)。

QR分解的计算通常采用Householder变换或Gram-Schmidt正交化方法。相比于直接求解正规方程\(X'X\beta = X'y\)，QR分解避免了显式计算\(X'X\)可能带来的数值不稳定性（条件数被平方放大）。在R中，可以使用\texttt{qr(X)}函数进行QR分解，\texttt{qr.solve(X,\ y)}可直接求解最小二乘问题。

\subsection{SVD 解决（更加稳定）}\label{svd-ux89e3ux51b3ux66f4ux52a0ux7a33ux5b9a}

使用SVD求解最小二乘问题：将设计矩阵分解为\(X = U \Sigma V^T\)，则正规方程的解为：

\[\hat{\beta} = V \Sigma^{+} U^T y = \sum_{i=1}^{r} \frac{u_i^T y}{\sigma_i} v_i\]

其中\(\Sigma^{+}\)为\(\Sigma\)的伪逆（将非零奇异值取倒数）。SVD方法的优势在于：
- 即使\(X\)不满秩（奇异值\(\sigma_i\)为零或接近零），仍可通过截断小奇异值给出稳定的解。
- SVD直接揭示了解的方差来源：\(\text{Var}(\hat{\beta}_j) = \sigma^2 \sum_{i=1}^{r} v_{ji}^2 / \sigma_i^2\)。小的奇异值会极大放大估计方差。
- 因此SVD是三种方法中数值稳定性最高的，但计算代价也最大。

\textbf{Updating a Least Squares Solution} 指的是当原始数据集发生变化（例如添加了新样本或新特征）时，如何\textbf{不从零开始重新计算}，而是利用现有的结果快速更新模型参数。

\textbf{Givens Rotations（吉文斯旋转）}：我们使用一系列 \textbf{Givens Rotations}，从上到下逐个消除新行中的元素，直到它重新变回上三角形式。

\subsection{加权最小二乘法}\label{ux52a0ux6743ux6700ux5c0fux4e8cux4e58ux6cd5}

起源于回归模型中误差项的异方差性

在WLS中：\textbf{方差小}的观测值：数据更可靠，赋予\textbf{较大权重}。\textbf{方差大}的观测值：数据噪声多，赋予\textbf{较小权重}。

推导\textbf{加权最小二乘法 (WLS) 的正规方程 (Normal Equations)：}

加权最小二乘的目标函数为最小化加权残差平方和：

\[Q(\beta) = (y - X\beta)' W (y - X\beta)\]

其中\(W = \text{diag}(w_1, w_2, \cdots, w_n)\)，\(w_i = 1/\sigma_i^2\)为第\(i\)个观测的权重（方差的倒数）。

对\(Q(\beta)\)关于\(\beta\)求导并令其为零：

\[\frac{\partial Q}{\partial \beta} = -2 X' W (y - X\beta) = 0\]

得加权正规方程：\(X' W X \hat{\beta} = X' W y\)。

解得\(\hat{\beta}_{WLS} = (X' W X)^{-1} X' W y\)。

转换到同方差系统：

WLS也可以通过数据变换转化为同方差系统（OLS）求解。设\(W = L L'\)为\(W\)的Cholesky分解，令\(\tilde{y} = L' y\)，\(\tilde{X} = L' X\)，则加权残差平方和可写为：

\[Q(\beta) = (y - X\beta)' L L' (y - X\beta) = (\tilde{y} - \tilde{X}\beta)' (\tilde{y} - \tilde{X}\beta)\]

这样WLS问题就等价于在变换后的数据\((\tilde{X}, \tilde{y})\)上进行普通最小二乘。新误差项\(\tilde{\epsilon} = L' \epsilon\)满足\(Var(\tilde{\epsilon}) = L' \sigma^2 W^{-1} L = \sigma^2 I\)，实现了同方差化。这一变换视角也解释了IRLS算法中每一步加权最小二乘的实质。

\chapter{机器学习基础}\label{ux673aux5668ux5b66ux4e60ux57faux7840}

\begin{quote}
机器学习是统计学向高维、非结构化数据延伸的前沿。它不追问''数据服从什么分布''，而是直接问''模型在新数据上能预测多准''。
\end{quote}

与前面各章讨论的基于概率模型的统计推断不同，机器学习更关注如何从数据中自动发现模式并做出预测。本章从机器学习的基本任务和评估方法出发，为后续的分类、集成学习和深度学习奠定基础。

\section{机器学习的基本任务}\label{ux673aux5668ux5b66ux4e60ux7684ux57faux672cux4efbux52a1}

机器学习任务通常分为以下几类：

\begin{itemize}
\item
  \textbf{关联分析}（Association Analysis）：发现数据项之间有趣的关联或相关关系。例如，频繁项集挖掘旨在找出同时出现频率超过给定阈值的项集，并用支持度（Support）和置信度（Confidence）来衡量规则的强度。经典的Apriori算法利用''包含A、B、C项集的支持度不会超过包含A、B项集的支持度''这一单调性进行剪枝，从单元素项集开始逐步扩展。关联分析广泛应用于购物篮分析、电信诈骗检测（团伙作案中多个号码在特定时间段频繁出现）、医疗诊断以及频繁序列挖掘等场景。
\item
  \textbf{分类分析}（Classification）：有监督学习中最常见的任务，目标是基于已有标签的数据构建模型，对新样本的类别进行预测。典型应用包括信用评估与用户画像、垃圾邮件检测、电影票房预测等。
\item
  \textbf{聚类分析}（Clustering）：无监督学习的代表，将相似的数据点归入同一组，使得组内距离最小化、组间距离最大化。聚类主要用于探索性和描述性分析，如文档分组、股票波动模式识别、基因序列模式发现等。
\item
  \textbf{异常检测}（Anomaly Detection）：识别与大多数数据显著不同的观测值，在欺诈检测、工业故障诊断等场景中有重要应用。
\end{itemize}

\section{数据与数据预处理}\label{ux6570ux636eux4e0eux6570ux636eux9884ux5904ux7406}

\subsection{属性类型}\label{ux5c5eux6027ux7c7bux578b}

数据属性的取值可以分为以下四种类型：

\begin{itemize}
\tightlist
\item
  \textbf{名义型}（Nominal）：如学号、邮政编码，取值之间没有大小关系。
\item
  \textbf{序数型}（Ordinal）：如排名、等级，有大小顺序关系但间距不一定相等。
\item
  \textbf{区间型}（Interval）：如日期、温度（摄氏度），可以加减但不可乘除，零点为人为设定。
\item
  \textbf{比例型}（Ratio）：如长度、重量，具有绝对零点，可以比较大小、进行加减乘除运算。
\end{itemize}

对于名义型属性，直接赋予数值会引入不合理的距离关系，通常采用\textbf{独热编码}（One-Hot Encoding）将其转换为多个二值变量。

\subsection{数据集类型}\label{ux6570ux636eux96c6ux7c7bux578b}

常见的数据集结构包括：

\begin{itemize}
\tightlist
\item
  \textbf{表数据}（Tabular Data）：行为样本、列为属性的矩阵形式，是最常见的数据格式。
\item
  \textbf{图数据}（Graph Data）：样本之间存在链接关系，如社交网络、分子结构。
\item
  \textbf{序列数据}（Sequential Data）：数据点按时间或空间顺序排列，如时间序列、DNA序列。
\end{itemize}

\subsection{数据质量问题}\label{ux6570ux636eux8d28ux91cfux95eeux9898}

实际数据往往存在缺失值、重复记录和异常值等问题，需要在建模前进行识别和处理。缺失值的处理策略包括删除含缺失的样本、基于均值/中位数填充，或使用模型预测缺失值。

\subsection{数据预处理的关键技术}\label{ux6570ux636eux9884ux5904ux7406ux7684ux5173ux952eux6280ux672f}

\begin{itemize}
\tightlist
\item
  \textbf{数据聚合}（Aggregation）：将细粒度数据合并为粗粒度汇总，减少数据量同时可能降低噪声。
\item
  \textbf{数据抽样}（Sampling）：当数据量过大或类别不平衡时，通过抽样选取代表性子集。
\item
  \textbf{降维}（Dimensionality Reduction）：通过主成分分析（PCA）等方法减少变量数目，缓解''维数灾难''。
\item
  \textbf{特征筛选}（Feature Selection）：根据特征与目标变量的相关性，选取最有信息量的特征子集。
\end{itemize}

\section{偏差-方差分解}\label{ux504fux5dee-ux65b9ux5deeux5206ux89e3}

偏差-方差分解是理解模型泛化性能的核心理论工具。对于回归问题，在平方损失下，期望预测误差可分解为三个不可约项：

\[
E[(Y - \hat{f}(X))^2] = \underbrace{(f(X) - E[\hat{f}(X)])^2}_{\text{Bias}^2} + \underbrace{E[(\hat{f}(X) - E[\hat{f}(X)])^2]}_{\text{Variance}} + \underbrace{\sigma^2_{\varepsilon}}_{\text{Irreducible Error}}
\]

其中 \(f(X)\) 是真实的（未知的）回归函数，\(\hat{f}(X)\) 是拟合的模型。三项的含义：

\begin{itemize}
\tightlist
\item
  \textbf{偏差²（Bias²）}：模型期望预测与真实值之间的差距。高偏差意味着模型过于简单，无法捕捉数据的真实结构------\textbf{欠拟合}。
\item
  \textbf{方差（Variance）}：模型在不同训练集上的预测波动程度。高方差意味着模型对训练数据的微小变化过于敏感------\textbf{过拟合}。
\item
  \textbf{不可约误差（\(\sigma^2_{\varepsilon}\)）}：数据本身的噪声，任何模型都无法消除。
\end{itemize}

偏差和方差之间存在权衡：简单模型（如线性回归）偏差大、方差小；复杂模型（如深度网络）偏差小、方差大。模型选择本质上是在偏差-方差光谱上寻找最优位置。

\begin{quote}
这一分解贯穿全书：Ch06 中点估计的 MSE = 方差 + 偏差²，正是此公式在参数估计中的对应版本。Ch19 中集成学习的有效性，恰在于 Bagging 降低方差而 Boosting 降低偏差。
\end{quote}

\begin{Shaded}
\begin{Highlighting}[]
\CommentTok{\# 偏差{-}方差权衡示意}
\NormalTok{x }\OtherTok{\textless{}{-}} \FunctionTok{seq}\NormalTok{(}\DecValTok{0}\NormalTok{, }\DecValTok{1}\NormalTok{, }\AttributeTok{length.out =} \DecValTok{50}\NormalTok{)}
\FunctionTok{plot}\NormalTok{(x, }\FunctionTok{sin}\NormalTok{(}\DecValTok{2} \SpecialCharTok{*}\NormalTok{ pi }\SpecialCharTok{*}\NormalTok{ x), }\AttributeTok{type =} \StringTok{"l"}\NormalTok{, }\AttributeTok{lwd =} \DecValTok{2}\NormalTok{, }\AttributeTok{ylim =} \FunctionTok{c}\NormalTok{(}\SpecialCharTok{{-}}\DecValTok{2}\NormalTok{, }\DecValTok{2}\NormalTok{),}
     \AttributeTok{xlab =} \StringTok{"x"}\NormalTok{, }\AttributeTok{ylab =} \StringTok{"y"}\NormalTok{, }\AttributeTok{main =} \StringTok{"偏差{-}方差权衡示意"}\NormalTok{)}
\CommentTok{\# 高偏差（线性拟合）—— 稳定但偏}
\FunctionTok{abline}\NormalTok{(}\FunctionTok{lm}\NormalTok{(}\FunctionTok{sin}\NormalTok{(}\DecValTok{2} \SpecialCharTok{*}\NormalTok{ pi }\SpecialCharTok{*}\NormalTok{ x) }\SpecialCharTok{\textasciitilde{}}\NormalTok{ x), }\AttributeTok{col =} \StringTok{"blue"}\NormalTok{, }\AttributeTok{lwd =} \DecValTok{2}\NormalTok{)}
\CommentTok{\# 高方差（过拟合）—— 贴近训练点但波动大}
\FunctionTok{points}\NormalTok{(x[}\FunctionTok{seq}\NormalTok{(}\DecValTok{1}\NormalTok{, }\DecValTok{50}\NormalTok{, }\DecValTok{5}\NormalTok{)], }\FunctionTok{sin}\NormalTok{(}\DecValTok{2} \SpecialCharTok{*}\NormalTok{ pi }\SpecialCharTok{*}\NormalTok{ x[}\FunctionTok{seq}\NormalTok{(}\DecValTok{1}\NormalTok{, }\DecValTok{50}\NormalTok{, }\DecValTok{5}\NormalTok{)]), }\AttributeTok{pch =} \DecValTok{16}\NormalTok{, }\AttributeTok{col =} \StringTok{"red"}\NormalTok{)}
\FunctionTok{lines}\NormalTok{(}\FunctionTok{smooth.spline}\NormalTok{(x[}\FunctionTok{seq}\NormalTok{(}\DecValTok{1}\NormalTok{, }\DecValTok{50}\NormalTok{, }\DecValTok{5}\NormalTok{)], }\FunctionTok{sin}\NormalTok{(}\DecValTok{2} \SpecialCharTok{*}\NormalTok{ pi }\SpecialCharTok{*}\NormalTok{ x[}\FunctionTok{seq}\NormalTok{(}\DecValTok{1}\NormalTok{, }\DecValTok{50}\NormalTok{, }\DecValTok{5}\NormalTok{)]), }\AttributeTok{spar =} \FloatTok{0.1}\NormalTok{),}
      \AttributeTok{col =} \StringTok{"red"}\NormalTok{, }\AttributeTok{lwd =} \DecValTok{2}\NormalTok{)}
\FunctionTok{legend}\NormalTok{(}\StringTok{"topright"}\NormalTok{, }\FunctionTok{c}\NormalTok{(}\StringTok{"真实函数"}\NormalTok{, }\StringTok{"高偏差 (线性)"}\NormalTok{, }\StringTok{"高方差 (样条)"}\NormalTok{),}
       \AttributeTok{lty =} \FunctionTok{c}\NormalTok{(}\DecValTok{1}\NormalTok{, }\DecValTok{1}\NormalTok{, }\DecValTok{1}\NormalTok{), }\AttributeTok{col =} \FunctionTok{c}\NormalTok{(}\StringTok{"black"}\NormalTok{, }\StringTok{"blue"}\NormalTok{, }\StringTok{"red"}\NormalTok{), }\AttributeTok{cex =} \FloatTok{0.8}\NormalTok{)}
\end{Highlighting}
\end{Shaded}

\begin{center}\includegraphics{statistic-book_files/figure-latex/bias-variance-illustration-1} \end{center}

\section{模型评估与选择}\label{ux6a21ux578bux8bc4ux4f30ux4e0eux9009ux62e9}

\subsection{训练策略}\label{ux8badux7ec3ux7b56ux7565}

\begin{itemize}
\tightlist
\item
  \textbf{留出法}（Holdout）：将数据集按一定比例（如70\%/30\%）划分为训练集和测试集，进行一次性的训练和评估。
\item
  \textbf{交叉验证}（Cross Validation）：进行多次留出法，取平均结果。当折数 \(k\) 等于样本数时，称为留一法交叉验证（LOOCV）。交叉验证是偏差-方差权衡在模型评估中的实践------\(k\) 越小，评估的方差越大但偏差越小。
\item
  \textbf{自助法}（Bootstrap）：有放回地从原始数据中抽样生成多个训练集，未被抽中的样本构成测试集。适用于小样本场景。
\end{itemize}

\subsection{分类性能指标}\label{ux5206ux7c7bux6027ux80fdux6307ux6807}

对于一个二分类问题，混淆矩阵是一个 \(2 \times 2\) 的矩阵：

\begin{longtable}[]{@{}lll@{}}
\toprule\noalign{}
& 预测为正 & 预测为负 \\
\midrule\noalign{}
\endhead
\bottomrule\noalign{}
\endlastfoot
\textbf{实际为正} & TP & FN \\
\textbf{实际为负} & FP & TN \\
\end{longtable}

基于混淆矩阵的关键指标：

\begin{itemize}
\tightlist
\item
  \textbf{准确率}（Accuracy）：\(\text{Accuracy} = \frac{TP + TN}{TP + TN + FP + FN}\)。在类别不平衡时严重失真。
\item
  \textbf{精确率}（Precision）：\(\text{Precision} = \frac{TP}{TP + FP}\)------预测为正的样本中有多少真的为正。
\item
  \textbf{召回率}（Recall）：\(\text{Recall} = \frac{TP}{TP + FN}\)------所有正样本中有多少被成功识别。
\item
  \textbf{F1分数}（F1 Score）：\(\text{F1} = 2 \cdot \frac{\text{Precision} \cdot \text{Recall}}{\text{Precision} + \text{Recall}}\)------精确率与召回率的调和平均数。
\end{itemize}

\subsection{ROC曲线与AUC}\label{rocux66f2ux7ebfux4e0eauc}

大多数分类模型输出的不是类别标签，而是属于正类的概率。设阈值为 \(\tau\)，记真正率 \(TPR = TP/(TP+FN)\)，假正率 \(FPR = FP/(FP+TN)\)。\textbf{ROC曲线}描绘的是 \(\tau\) 从 1.0 变化到 0.0 时 \((FPR, TPR)\) 的轨迹。\textbf{AUC}（Area Under the Curve）为 ROC 曲线下面积，综合衡量模型的排序能力。

\subsection{超参数搜索}\label{ux8d85ux53c2ux6570ux641cux7d22}

\begin{itemize}
\tightlist
\item
  \textbf{网格搜索}（Grid Search）：在预定义的超参数网格上穷举搜索，适合参数空间较小的场景。
\item
  \textbf{随机搜索}（Random Search）：从参数分布中随机采样，在高维空间比网格搜索更高效。
\item
  \textbf{贝叶斯优化}（Bayesian Optimization）：使用高斯过程代理模型，在探索（Exploration）与利用（Exploitation）之间平衡，适用于评估代价高昂的场景（如深度学习）。
\end{itemize}

\section{正则化：从统计学到机器学习}\label{ux6b63ux5219ux5316ux4eceux7edfux8ba1ux5b66ux5230ux673aux5668ux5b66ux4e60}

正则化是防止过拟合的核心技术。它在损失函数上添加对模型复杂度的惩罚项：

\[
\mathcal{L}_{\text{reg}}(\theta) = \mathcal{L}(\theta) + \lambda R(\theta)
\]

\begin{itemize}
\tightlist
\item
  \textbf{L2 正则化（Ridge）}：\(R(\theta) = \|\theta\|_2^2\)。将所有参数向零收缩，但不产生稀疏解。
\item
  \textbf{L1 正则化（Lasso）}：\(R(\theta) = \|\theta\|_1\)。产生稀疏解------部分参数被精确压缩到零，天然实现特征选择。
\item
  \textbf{弹性网络（Elastic Net）}：结合 L1 和 L2 惩罚。
\end{itemize}

从贝叶斯视角看，正则化等价于对参数施加先验分布：L2 正则化等价于高斯先验 \(\theta \sim N(0, \frac{1}{2\lambda})\)，L1 正则化等价于拉普拉斯先验。这直接桥接了 Ch08 中''MAP 估计 = 带正则化的 MLE''的结论。

\begin{quote}
奥卡姆剃刀在统计学中的数学表达，就是正则化。
\end{quote}

\chapter{经典分类方法}\label{ux7ecfux5178ux5206ux7c7bux65b9ux6cd5}

\begin{quote}
分类是监督学习的核心任务。从直觉驱动的决策树到核化的支持向量机，再到基于概率的朴素贝叶斯------每种方法都在偏差-方差光谱上占据不同的位置。理解它们的位置，比记住它们的公式更重要。
\end{quote}

\section{决策树}\label{ux51b3ux7b56ux6811}

设有训练数据 \(\{(x_i, y_i)\}_{i=1}^n\)，其中 \(y_i\) 是类别标签，\(x_i\) 是特征向量。决策树的基本思想是从根节点开始，根据属性的不同取值设置分支条件，希望每次分裂后各子节点内的样本尽可能''纯净''。

\textbf{分裂条件的选择}是决策树的核心问题。两个常用的度量指标为：

\begin{itemize}
\tightlist
\item
  \textbf{基尼系数}（Gini Index）：
\end{itemize}

\[
Gini(D) = 1 - \sum_{k=1}^{K} p_k^2
\]

其中 \(p_k\) 是第 \(k\) 类在节点 \(D\) 中所占的比例。基尼系数越小，节点越纯。决策树遍历所有可能的分裂条件，选择使加权基尼系数下降最多的分裂。

\begin{itemize}
\tightlist
\item
  \textbf{信息熵}（Entropy）：
\end{itemize}

\[
Ent(D) = -\sum_{k=1}^{K} p_k \log_2 p_k
\]

信息增益越大，说明使用该属性进行分裂的效果越好。这一视角与 Ch01 中信息熵的讨论一脉相承------决策树的每次分裂，本质上是在本地做一次''熵减''，逐步将不确定性转化为确定性。

\begin{quote}
用 ID 这样的唯一标识进行分裂，每个节点将只有单个样本，看上去完全纯净，但这样的树没有任何泛化能力------这就是过拟合的直观体现（偏差-方差分解的详细讨论见 Ch17）。
\end{quote}

\textbf{停止条件与剪枝}：为防止过拟合，需要设置停止分裂的条件，如最小基尼系数阈值、最大深度限制、节点最小样本数等。决策树是一种\textbf{贪婪算法}：每次选择当前最优的分裂条件，但不保证得到全局最优的树结构。

\section{支持向量机（SVM）}\label{ux652fux6301ux5411ux91cfux673asvm}

支持向量机的核心目标是寻找一个最大间隔超平面来分离两类样本。其优化问题可形式化为：

\[
\begin{aligned}
\min_{w, b} \quad & \frac{1}{2} \|w\|^2 \\
\text{s.t.} \quad & y_i (w' x_i + b) \geq 1, \quad i = 1, \ldots, n
\end{aligned}
\]

通过拉格朗日乘子法转化为对偶问题，再利用 KKT 条件求解。根据互补松弛条件，只有位于间隔边界上的支持向量对应的乘子不为零------最终的分割超平面仅由少数支持向量决定。

\textbf{序列最小优化算法}（SMO）：每次仅选择两个乘子 \(\alpha_i, \alpha_j\) 进行优化，固定其余所有乘子，大大简化了每次迭代的计算。

\textbf{软间隔}（Soft Margin）：为允许一定的分类错误，引入松弛变量 \(\xi_i\) 和惩罚参数 \(C\)：

\[
\begin{aligned}
\min_{w, b, \xi} \quad & \frac{1}{2} \|w\|^2 + C \sum_{i=1}^{n} \xi_i \\
\text{s.t.} \quad & y_i (w' x_i + b) \geq 1 - \xi_i, \quad \xi_i \geq 0
\end{aligned}
\]

超参数 \(C\) 控制对误分类的惩罚力度：\(C\) 越大，对误分类的容忍度越低。从统计视角看，这一正则化策略等价于在经验风险上添加对模型复杂度的惩罚。

\textbf{核函数}：对于线性不可分的数据，通过映射 \(\phi: \mathcal{X} \to \mathcal{H}\) 将数据升维至高维空间。核函数 \(K(x_i, x_j) = \langle \phi(x_i), \phi(x_j) \rangle\) 相当于在原始空间中直接计算升维后的内积------\textbf{核技巧}避免了显式高维计算的困难。常用核函数：

\begin{longtable}[]{@{}
  >{\raggedright\arraybackslash}p{(\columnwidth - 4\tabcolsep) * \real{0.3636}}
  >{\raggedright\arraybackslash}p{(\columnwidth - 4\tabcolsep) * \real{0.3636}}
  >{\raggedright\arraybackslash}p{(\columnwidth - 4\tabcolsep) * \real{0.2727}}@{}}
\toprule\noalign{}
\begin{minipage}[b]{\linewidth}\raggedright
核名称
\end{minipage} & \begin{minipage}[b]{\linewidth}\raggedright
表达式
\end{minipage} & \begin{minipage}[b]{\linewidth}\raggedright
特点
\end{minipage} \\
\midrule\noalign{}
\endhead
\bottomrule\noalign{}
\endlastfoot
线性核 & \(K(x_i, x_j) = x_i \cdot x_j\) & 无额外超参数，等价于原始空间的线性分类 \\
多项式核 & \(K(x_i, x_j) = (x_i \cdot x_j + c)^d\) & 可控制映射的复杂度（参数 \(d\)） \\
RBF 核 & \(K(x_i, x_j) = \exp(-\gamma\|x_i - x_j\|^2)\) & 最常用，相当于无限维映射 \\
\end{longtable}

\section{朴素贝叶斯分类器}\label{ux6734ux7d20ux8d1dux53f6ux65afux5206ux7c7bux5668}

朴素贝叶斯分类器基于贝叶斯定理：

\[
P(Y = k \mid X = x) = \frac{P(Y = k) P(X = x \mid Y = k)}{P(X = x)}
\]

其''朴素''之处在于假设各特征在给定类别的条件下相互独立：

\[
P(X = x \mid Y = k) = \prod_{j=1}^{p} P(X_j = x_j \mid Y = k)
\]

在这一假设下，分类器的训练本质上是一个计数过程，因此特别适合处理离散型数据。对于文本分类（如垃圾邮件检测）表现尤为出色。对于连续型特征，通常需假设其服从正态分布。

\begin{quote}
若条件独立假设不成立，可以进一步刻画变量之间的依赖关系，即构建\textbf{贝叶斯网络}（Bayesian Network）。
\end{quote}

\section{k-近邻（k-NN）}\label{k-ux8fd1ux90bbk-nn}

k-近邻是最直观的非参数分类方法------它甚至连模型都不显式训练。对于新样本 \(x^*\)，k-NN 在训练集中寻找距离最近的 \(k\) 个点，以多数投票决定 \(x^*\) 的类别：

\[\hat{y} = \arg\max_{c} \sum_{i \in \mathcal{N}_k(x^*)} I(y_i = c)\]

k-NN 天然与 Ch13 的非参数统计理念相通：不做分布假设，让数据自己说话。距离度量的选择（欧氏距离、曼哈顿距离、马氏距离）和 \(k\) 的选择是其两个核心决策。\(k\) 太小则方差大（过拟合），\(k\) 太大则偏差大（欠拟合）------这是偏差-方差权衡（Ch17）的直观展现。

\section{分类器对比}\label{ux5206ux7c7bux5668ux5bf9ux6bd4}

\begin{longtable}[]{@{}llllll@{}}
\toprule\noalign{}
方法 & 偏差 & 方差 & 可解释性 & 计算复杂度 & 适用场景 \\
\midrule\noalign{}
\endhead
\bottomrule\noalign{}
\endlastfoot
决策树 & 中 & 中 & 高 & 低（训练） & 需要解释性规则 \\
SVM (RBF核) & 低 & 低-中 & 低 & 中 & 高维、非线性分类 \\
朴素贝叶斯 & 高 & 低 & 中 & 极低 & 文本分类、小样本 \\
k-NN & 中 & 高 & 中 & 高（预测） & 小数据、非参数场景 \\
\end{longtable}

各方法的集成版本（随机森林、梯度提升）将在 Ch19 中详细讨论。

\chapter{集成学习与模型选择}\label{ux96c6ux6210ux5b66ux4e60ux4e0eux6a21ux578bux9009ux62e9}

\begin{quote}
三个臭皮匠，顶个诸葛亮。集成学习的核心洞见是：将多个弱学习器组合起来，整体的预测性能往往显著优于任何单个学习器。
\end{quote}

\section{为什么集成有效？}\label{ux4e3aux4ec0ux4e48ux96c6ux6210ux6709ux6548}

预测误差可以分解为三个来源（在 Ch04 和 Ch10 中已讨论过平方损失下的偏差-方差分解）：

\[\text{Error} = \text{Bias}^2 + \text{Variance} + \text{Noise}\]

\begin{itemize}
\tightlist
\item
  \textbf{偏差}：模型本身的表达能力不足导致的系统误差。过于简单的模型（如线性回归拟合非线性数据）高偏差。
\item
  \textbf{方差}：模型对训练数据的微扰过于敏感。过于复杂的模型（如深度决策树）高方差。
\item
  \textbf{噪声}：数据本身固有的不可约误差。
\end{itemize}

单个复杂模型的误差往往来自\textbf{方差过大}。集成学习通过组合多个模型来\textbf{降低方差}（Bagging）或\textbf{降低偏差}（Boosting）。

条件：基础学习器之间需要具有\textbf{差异性}（diversity），且每个学习器的错误率不能太高（通常要求 \(< 0.5\)）。

\section{Bagging 与随机森林}\label{bagging-ux4e0eux968fux673aux68eeux6797}

\subsection{Bagging（Bootstrap Aggregating）}\label{baggingbootstrap-aggregating}

\textbf{步骤}：

\begin{enumerate}
\def\labelenumi{\arabic{enumi}.}
\tightlist
\item
  从原始训练集做 \(B\) 次 Bootstrap 抽样（有放回），得到 \(B\) 个自助样本。
\item
  在每个自助样本上训练一个同类型的基学习器。
\item
  对分类问题：用投票法汇总；对回归问题：用平均法汇总。
\end{enumerate}

Bootstrap 抽样的关键性质：每个自助样本大约包含原始数据 \(63.2\%\) 的观测值（因为 \((1 - 1/n)^n \to 1/e \approx 0.368\)），剩余约 \(36.8\%\) 的观测未被抽中------这些就是\textbf{袋外样本}（Out-of-Bag, OOB），可以用来无偏地估计泛化误差。

\subsection{随机森林（Random Forest）}\label{ux968fux673aux68eeux6797random-forest}

随机森林在 Bagging + 决策树的基础上，增加了一层随机性------\textbf{每次分裂时，只从随机选取的 \(m\) 个特征中选择最优分裂特征}（通常 \(m \approx \sqrt{p}\) 用于分类，\(m \approx p/3\) 用于回归）。

这一额外的随机化使得树与树之间的相关性减小，进一步降低了集成后的方差。

\textbf{随机森林的核心超参数}：

\begin{itemize}
\tightlist
\item
  \texttt{ntree}：树的数量，通常越大越好（收益递减）。
\item
  \texttt{mtry}：每次分裂考虑的候选特征数。
\item
  \texttt{max\_depth} 或 \texttt{min\_node\_size}：限制单棵树的复杂度。
\end{itemize}

\textbf{OOB 误差}是随机森林的一个优良特性------不需要单独的验证集即可获得泛化误差的无偏估计。

\begin{Shaded}
\begin{Highlighting}[]
\FunctionTok{library}\NormalTok{(randomForest)}
\FunctionTok{data}\NormalTok{(iris)}
\FunctionTok{set.seed}\NormalTok{(}\DecValTok{123}\NormalTok{)}
\NormalTok{rf }\OtherTok{\textless{}{-}} \FunctionTok{randomForest}\NormalTok{(Species }\SpecialCharTok{\textasciitilde{}}\NormalTok{ ., }\AttributeTok{data =}\NormalTok{ iris, }\AttributeTok{ntree =} \DecValTok{500}\NormalTok{, }\AttributeTok{mtry =} \DecValTok{2}\NormalTok{, }\AttributeTok{importance =} \ConstantTok{TRUE}\NormalTok{)}
\FunctionTok{print}\NormalTok{(rf)                  }\CommentTok{\# OOB 误差}
\end{Highlighting}
\end{Shaded}

\begin{verbatim}
## 
## Call:
##  randomForest(formula = Species ~ ., data = iris, ntree = 500,      mtry = 2, importance = TRUE) 
##                Type of random forest: classification
##                      Number of trees: 500
## No. of variables tried at each split: 2
## 
##         OOB estimate of  error rate: 4.67%
## Confusion matrix:
##            setosa versicolor virginica class.error
## setosa         50          0         0        0.00
## versicolor      0         47         3        0.06
## virginica       0          4        46        0.08
\end{verbatim}

\begin{Shaded}
\begin{Highlighting}[]
\FunctionTok{importance}\NormalTok{(rf)             }\CommentTok{\# 变量重要性}
\end{Highlighting}
\end{Shaded}

\begin{verbatim}
##                 setosa versicolor virginica MeanDecreaseAccuracy
## Sepal.Length  6.495113   7.588736  7.690949            11.098105
## Sepal.Width   4.364415   1.038790  4.546158             5.144762
## Petal.Length 22.157039  34.060864 27.853055            33.535680
## Petal.Width  22.482214  32.571528 30.595022            32.839577
##              MeanDecreaseGini
## Sepal.Length         9.798189
## Sepal.Width          2.236535
## Petal.Length        43.093269
## Petal.Width         44.042372
\end{verbatim}

\begin{Shaded}
\begin{Highlighting}[]
\FunctionTok{varImpPlot}\NormalTok{(rf)             }\CommentTok{\# 变量重要性图}
\end{Highlighting}
\end{Shaded}

\begin{center}\includegraphics{statistic-book_files/figure-latex/unnamed-chunk-27-1} \end{center}

\subsection{变量重要性}\label{ux53d8ux91cfux91cdux8981ux6027}

随机森林提供了两种变量重要性的度量：

\begin{enumerate}
\def\labelenumi{\arabic{enumi}.}
\tightlist
\item
  \textbf{平均精度下降}（Mean Decrease Accuracy）：对某个变量的取值进行随机置换后，OOB 预测精度的下降程度。
\item
  \textbf{平均基尼下降}（Mean Decrease Gini）：在所有树的所有节点分裂中，由该变量带来的基尼不纯度下降之和。
\end{enumerate}

\section{Boosting}\label{boosting}

Boosting 与 Bagging 的哲学不同：Bagging 是\textbf{并行的}------各学习器独立训练然后平均；Boosting 是\textbf{串行的}------每个新学习器关注前一个学习器犯错的样本。

\subsection{AdaBoost}\label{adaboost}

AdaBoost（Adaptive Boosting）的核心机制是\textbf{样本权重的自适应调整}：

\begin{enumerate}
\def\labelenumi{\arabic{enumi}.}
\tightlist
\item
  初始化所有样本权重相同：\(w_i = 1/n\)。
\item
  对 \(t = 1, 2, \ldots, T\)：

  \begin{itemize}
  \tightlist
  \item
    在加权样本上训练基学习器 \(h_t\)。
  \item
    计算加权错误率 \(\epsilon_t\)。
  \item
    计算学习器的权重 \(\alpha_t = \frac{1}{2} \ln \frac{1 - \epsilon_t}{\epsilon_t}\)。
  \item
    增加被误分类样本的权重，减小正确分类样本的权重。
  \end{itemize}
\item
  最终分类器：\(H(x) = \text{sign}\left(\sum_{t=1}^T \alpha_t h_t(x)\right)\)。
\end{enumerate}

\begin{quote}
AdaBoost 的损失函数是指数损失，这恰好等同于用加法模型在前向分步算法下拟合指数损失------这是 AdaBoost 与统计方法的深刻联系。
\end{quote}

\subsection{梯度提升（Gradient Boosting）}\label{ux68afux5ea6ux63d0ux5347gradient-boosting}

梯度提升从\textbf{函数空间的梯度下降}视角理解 Boosting：

\begin{enumerate}
\def\labelenumi{\arabic{enumi}.}
\tightlist
\item
  从初始模型 \(F_0(x) = \arg\min_c \sum L(y_i, c)\) 开始。
\item
  对 \(m = 1, \ldots, M\)：

  \begin{itemize}
  \tightlist
  \item
    计算伪残差：\(r_{im} = -\left[\frac{\partial L(y_i, F(x_i))}{\partial F(x_i)}\right]_{F = F_{m-1}}\)
  \item
    用基学习器 \(h_m\) 拟合伪残差。
  \item
    找到最优步长 \(\rho_m\)，更新 \(F_m = F_{m-1} + \rho_m h_m\)。
  \end{itemize}
\end{enumerate}

不同的损失函数 \(L\) 对应不同的提升方法：平方损失 → L2Boosting；指数损失 → AdaBoost；对数损失 → LogitBoost。

\textbf{XGBoost} 和 \textbf{LightGBM} 是梯度提升的两个工业级实现，在 Kaggle 比赛中统治了表格数据的预测任务。它们在标准梯度提升的基础上加入了正则化项、列抽样、直方图加速等技术。

\begin{Shaded}
\begin{Highlighting}[]
\FunctionTok{library}\NormalTok{(xgboost)}
\CommentTok{\# 准备数据}
\FunctionTok{data}\NormalTok{(iris)}
\NormalTok{X }\OtherTok{\textless{}{-}} \FunctionTok{as.matrix}\NormalTok{(iris[, }\DecValTok{1}\SpecialCharTok{:}\DecValTok{4}\NormalTok{])}
\NormalTok{y }\OtherTok{\textless{}{-}} \FunctionTok{as.numeric}\NormalTok{(iris}\SpecialCharTok{$}\NormalTok{Species) }\SpecialCharTok{{-}} \DecValTok{1}
\CommentTok{\# 训练 XGBoost}
\NormalTok{dtrain }\OtherTok{\textless{}{-}} \FunctionTok{xgb.DMatrix}\NormalTok{(X, }\AttributeTok{label =}\NormalTok{ y)}
\NormalTok{param }\OtherTok{\textless{}{-}} \FunctionTok{list}\NormalTok{(}\AttributeTok{objective =} \StringTok{"multi:softmax"}\NormalTok{, }\AttributeTok{num\_class =} \DecValTok{3}\NormalTok{, }\AttributeTok{max\_depth =} \DecValTok{3}\NormalTok{, }\AttributeTok{eta =} \FloatTok{0.3}\NormalTok{)}
\NormalTok{bst }\OtherTok{\textless{}{-}} \FunctionTok{xgb.train}\NormalTok{(}\AttributeTok{params =}\NormalTok{ param, }\AttributeTok{data =}\NormalTok{ dtrain, }\AttributeTok{nrounds =} \DecValTok{50}\NormalTok{, }\AttributeTok{verbose =} \DecValTok{0}\NormalTok{)}
\NormalTok{pred }\OtherTok{\textless{}{-}} \FunctionTok{predict}\NormalTok{(bst, X)}
\FunctionTok{table}\NormalTok{(pred, y)}
\end{Highlighting}
\end{Shaded}

\begin{verbatim}
##     y
## pred  0  1  2
##    0 50  0  0
##    1  0 50  0
##    2  0  0 50
\end{verbatim}

\section{Stacking（堆叠泛化）}\label{stackingux5806ux53e0ux6cdbux5316}

Stacking 的核心思想是''用模型来组合模型''。不同于 Bagging/Boosting 使用固定的组合规则（投票/加权平均），Stacking 训练一个\textbf{元学习器}（Meta-learner）来学习如何最好地组合基学习器的输出。

\textbf{步骤}：

\begin{enumerate}
\def\labelenumi{\arabic{enumi}.}
\tightlist
\item
  将训练集分为 \(K\) 折。
\item
  对每个基学习器，使用交叉验证生成''元特征''（即每个基学习器对每个训练样本的预测值）。
\item
  用元特征作为输入、原始标签作为输出，训练元学习器（通常是简单的模型如逻辑回归）。
\item
  预测时，先由基学习器生成元特征，再由元学习器做出最终预测。
\end{enumerate}

\section{模型选择与超参数调优}\label{ux6a21ux578bux9009ux62e9ux4e0eux8d85ux53c2ux6570ux8c03ux4f18}

\subsection{偏差-方差分解的实践含义}\label{ux504fux5dee-ux65b9ux5deeux5206ux89e3ux7684ux5b9eux8df5ux542bux4e49}

\begin{longtable}[]{@{}lll@{}}
\toprule\noalign{}
模型 & 偏差倾向 & 方差倾向 \\
\midrule\noalign{}
\endhead
\bottomrule\noalign{}
\endlastfoot
线性回归 & 高 & 低 \\
决策树（深） & 低 & 高 \\
随机森林 & 低 & 中 \\
GBDT & 低 & 中 \\
神经网络（大） & 低 & 高 \\
\end{longtable}

低偏差 + 高方差 → 用 Bagging 降方差。
高偏差 + 低方差 → 用 Boosting 降偏差。

\subsection{超参数搜索策略}\label{ux8d85ux53c2ux6570ux641cux7d22ux7b56ux7565}

\begin{itemize}
\tightlist
\item
  \textbf{网格搜索}（Grid Search）：对所有超参数组合穷举搜索。维度高时计算量爆炸（维数灾难）。
\item
  \textbf{随机搜索}（Random Search）：在超参数空间随机采样。研究证明比网格搜索更高效，因为通常只有少数超参数真正重要。
\item
  \textbf{贝叶斯优化}：用高斯过程建模''超参数 → 验证性能''的映射，每次选择最不确定且可能最优的点进行评估。在评估昂贵的场景下特别有效。
\end{itemize}

\subsection{交叉验证}\label{ux4ea4ux53c9ux9a8cux8bc1}

\(K\) 折交叉验证是模型选择的黄金标准（已在 Ch17 中讨论）：

\[\text{CV}_{(K)} = \frac{1}{K} \sum_{k=1}^K \text{Error}_k\]

\(K = 5\) 或 \(10\) 是常用的折数。\(K = n\) 时退化为留一法交叉验证（LOOCV）。

\begin{quote}
嵌套交叉验证（Nested CV）：外层 CV 估计泛化误差，内层 CV 选择超参数。这是最严格的模型评估方法。
\end{quote}

\begin{center}\rule{0.5\linewidth}{0.5pt}\end{center}

\textbf{速查卡片}：

\begin{itemize}
\tightlist
\item
  \textbf{核心问题}：如何组合多个弱模型得到一个强模型？
\item
  \textbf{Bagging}：并行的，降方差，代表=随机森林
\item
  \textbf{Boosting}：串行的，降偏差，代表=XGBoost/LightGBM
\item
  \textbf{Stacking}：用元学习器组合基学习器
\item
  \textbf{模型选择}：网格搜索 → 随机搜索 → 贝叶斯优化（由简到繁）
\item
  \textbf{R命令}：\texttt{randomForest::randomForest()}, \texttt{xgboost::xgboost()}, \texttt{caret::train()}
\item
  \textbf{参见}：Ch04（偏差-方差不等式）, Ch10（回归的诊断）, Ch18（分类方法）
\end{itemize}

\chapter{深度学习}\label{ux6df1ux5ea6ux5b66ux4e60}

\begin{quote}
如果说经典机器学习（Ch18）是在手工设计的特征空间中进行模式识别，深度学习则是让模型自动学习特征本身。这种从''特征工程''到''特征学习''的范式转换，是深度学习革命的核心。
\end{quote}

\section{从感知机到深度网络}\label{ux4eceux611fux77e5ux673aux5230ux6df1ux5ea6ux7f51ux7edc}

\subsection{感知机：神经网络的原型}\label{ux611fux77e5ux673aux795eux7ecfux7f51ux7edcux7684ux539fux578b}

感知机（Rosenblatt, 1958）是最简单的神经网络：输入加权求和，通过阶跃激活函数输出二分类结果：

\[\hat{y} = \text{sign}(w_1 x_1 + \cdots + w_p x_p + b)\]

单层感知机只能解决\textbf{线性可分}问题。Minsky 和 Papert 在 1969 年证明了感知机甚至无法学习 XOR 函数------这一负面结果导致了第一次''AI 冬天''。

解决方案是：\textbf{堆叠多层}。每一层对上一层的输出做非线性变换，使得网络可以学习复杂的非线性决策边界。

\subsection{前馈神经网络的数学表示}\label{ux524dux9988ux795eux7ecfux7f51ux7edcux7684ux6570ux5b66ux8868ux793a}

一个 \(L\) 层的前馈神经网络（含 \(L-1\) 个隐藏层）可以递归地表示为：

\[\mathbf{h}^{(0)} = \mathbf{x}\]
\[\mathbf{z}^{(l)} = \mathbf{W}^{(l)} \mathbf{h}^{(l-1)} + \mathbf{b}^{(l)}\]
\[\mathbf{h}^{(l)} = g^{(l)}(\mathbf{z}^{(l)}), \quad l = 1, 2, \ldots, L\]
\[\hat{\mathbf{y}} = \mathbf{h}^{(L)}\]

其中 \(g^{(l)}(\cdot)\) 是第 \(l\) 层的激活函数。

\textbf{通用逼近定理}（Universal Approximation Theorem）：一个具有足够多隐藏神经元的单隐藏层神经网络，可以以任意精度逼近任何连续函数。但是，深层网络往往可以用远少于浅层网络的参数达到相同的逼近效果------这是因为深度带来了\textbf{表征的层次化}：底层学习简单特征（边缘），中层学习组合特征（形状），顶层学习抽象概念（物体）。

\subsection{激活函数}\label{ux6fc0ux6d3bux51fdux6570}

激活函数是神经网络非线性能力的来源。如果不使用激活函数（或使用恒等激活），多层网络将退化为单层线性模型。

\begin{longtable}[]{@{}
  >{\raggedright\arraybackslash}p{(\columnwidth - 4\tabcolsep) * \real{0.3333}}
  >{\raggedright\arraybackslash}p{(\columnwidth - 4\tabcolsep) * \real{0.3333}}
  >{\raggedright\arraybackslash}p{(\columnwidth - 4\tabcolsep) * \real{0.3333}}@{}}
\toprule\noalign{}
\begin{minipage}[b]{\linewidth}\raggedright
激活函数
\end{minipage} & \begin{minipage}[b]{\linewidth}\raggedright
表达式
\end{minipage} & \begin{minipage}[b]{\linewidth}\raggedright
特点
\end{minipage} \\
\midrule\noalign{}
\endhead
\bottomrule\noalign{}
\endlastfoot
Sigmoid & \(\sigma(z) = \frac{1}{1 + e^{-z}}\) & 输出 (0,1)，有梯度消失问题 \\
Tanh & \(\tanh(z) = \frac{e^z - e^{-z}}{e^z + e^{-z}}\) & 输出 (-1,1)，零中心化 \\
ReLU & \(\text{ReLU}(z) = \max(0, z)\) & 计算高效，缓解梯度消失，最常用 \\
Leaky ReLU & \(\max(0.01z, z)\) & 解决 ReLU 的''死神经元''问题 \\
\end{longtable}

\begin{quote}
ReLU 的梯度消失缓解：Sigmoid 在 \(z\) 远离 0 时梯度趋于零，深层网络中梯度连乘导致底层几乎无法更新。ReLU 在正半轴的梯度恒为 1，天然避免了这一问题。
\end{quote}

\section{反向传播与训练}\label{ux53cdux5411ux4f20ux64adux4e0eux8badux7ec3}

\subsection{反向传播算法}\label{ux53cdux5411ux4f20ux64adux7b97ux6cd5}

反向传播（Backpropagation）是计算神经网络梯度的算法------它不是学习算法本身，而是高效计算损失函数对每一层权重的偏导数的方法。核心是链式法则的递归应用：

\[\frac{\partial L}{\partial \mathbf{W}^{(l)}} = \frac{\partial L}{\partial \mathbf{h}^{(L)}} \cdot \frac{\partial \mathbf{h}^{(L)}}{\partial \mathbf{h}^{(L-1)}} \cdot \cdots \cdot \frac{\partial \mathbf{h}^{(l+1)}}{\partial \mathbf{h}^{(l)}} \cdot \frac{\partial \mathbf{h}^{(l)}}{\partial \mathbf{W}^{(l)}}\]

前向传播计算损失，反向传播计算梯度------二者交替进行。

\subsection{随机梯度下降（SGD）与变体}\label{ux968fux673aux68afux5ea6ux4e0bux964dsgdux4e0eux53d8ux4f53}

在 Ch15 中我们已经讨论了 SGD 的基本原理。在深度学习实践中，更常用的是带有\textbf{动量}（Momentum）和\textbf{自适应学习率}的变体：

\begin{itemize}
\tightlist
\item
  \textbf{SGD + Momentum}：\(v_{t+1} = \beta v_t + \nabla L(\theta_t)\)，\(\theta_{t+1} = \theta_t - \eta v_{t+1}\)。动量项 \(v\) 累积了历史梯度方向，加速在一致方向上的移动。
\item
  \textbf{Adam}：结合了动量（一阶矩）和自适应学习率（二阶矩）。Adam 是目前最常用的默认优化器。
\end{itemize}

\[\begin{aligned}
m_t &= \beta_1 m_{t-1} + (1 - \beta_1) g_t \\
v_t &= \beta_2 v_{t-1} + (1 - \beta_2) g_t^2 \\
\hat{\theta}_{t+1} &= \theta_t - \eta \frac{\hat{m}_t}{\sqrt{\hat{v}_t} + \epsilon}
\end{aligned}\]

\subsection{损失函数的选择}\label{ux635fux5931ux51fdux6570ux7684ux9009ux62e9}

\begin{longtable}[]{@{}llll@{}}
\toprule\noalign{}
任务 & 损失函数 & 输出层激活 & 统计对应 \\
\midrule\noalign{}
\endhead
\bottomrule\noalign{}
\endlastfoot
二分类 & 二值交叉熵 & Sigmoid & 逻辑回归的对数似然 \\
多分类 & 多类交叉熵 & Softmax & 多项逻辑回归的对数似然 \\
回归 & MSE & 恒等 & OLS 的平方损失 \\
\end{longtable}

\begin{quote}
神经网络的损失函数与统计模型的对数似然存在精确对应：MSE 损失 + 线性输出 = 正态似然；交叉熵 + Sigmoid = 伯努利似然。这一对应关系可以追溯到 Ch01 §分布之间的距离------最小化交叉熵等价于最小化 KL 散度，从而等价于最大化似然。深度学习并非统计学的替代品，而是其自然延伸。
\end{quote}

\section{正则化：防止过拟合}\label{ux6b63ux5219ux5316ux9632ux6b62ux8fc7ux62dfux5408}

深度网络参数远多于样本量，过拟合是最大的风险。以下技术共同形成了现代深度学习的正则化工具箱：

\subsection{Dropout}\label{dropout}

训练时，以概率 \(p\) 随机将神经元输出置零。测试时，所有权重乘以 \(p\)（或输出除以 \(p\)）以补偿训练时的放缩。

\textbf{贝叶斯视角}：Dropout 等价于对网络权重施加特定先验下的变分推断近似。每个 Dropout 掩码对应一个''子网络''，最终预测是大量子网络的集成平均------这本质上是一种集成学习（Ch19）。

\subsection{批归一化（Batch Normalization）}\label{ux6279ux5f52ux4e00ux5316batch-normalization}

对每个 mini-batch 计算均值和方差，将输入归一化到零均值单位方差，然后通过可学习的缩放和平移参数恢复网络的表达能力。批归一化的效果是：允许更大的学习率；减少对初始化的敏感性；自带轻微的正则化效果。

\subsection{早停（Early Stopping）}\label{ux65e9ux505cearly-stopping}

在验证集误差不再下降时停止训练。早停防止了网络在训练集上''死记硬背''而损害泛化能力。

\subsection{L1/L2 正则化}\label{l1l2-ux6b63ux5219ux5316}

对权重施加 \(L_2\) 惩罚（权重衰减）等价于在贝叶斯框架下对权重施加高斯先验。\(L_1\) 惩罚等价于拉普拉斯先验，倾向于产生稀疏权重。

\section{主要网络架构}\label{ux4e3bux8981ux7f51ux7edcux67b6ux6784}

\subsection{卷积神经网络（CNN）}\label{ux5377ux79efux795eux7ecfux7f51ux7edccnn}

CNN 是为处理网格状数据（如图像）设计的架构。其两个核心操作是：

\begin{itemize}
\tightlist
\item
  \textbf{卷积层}：用一组可学习的滤波器（卷积核）在输入上滑动，每个滤波器对局部区域做内积运算。同一个滤波器在整个图像上共享参数------这大大减少了参数量，并且天然地提取了\textbf{平移不变}的特征。
\item
  \textbf{池化层}：对特征图进行降采样（通常是最大值池化），减小空间尺寸的同时保留了主要特征。
\end{itemize}

\begin{quote}
卷积 + 池化的层次化堆叠，使得底层学习边缘和纹理，中层学习形状和部件，顶层学习整体物体------这是一个由局部到全局的表示学习过程。
\end{quote}

\subsection{循环神经网络（RNN）与 LSTM}\label{ux5faaux73afux795eux7ecfux7f51ux7edcrnnux4e0e-lstm}

RNN 处理序列数据 \(x_1, x_2, \ldots, x_T\)，其核心是一个\textbf{隐状态}在时间上传递：

\[\mathbf{h}_t = \tanh(\mathbf{W}_h \mathbf{h}_{t-1} + \mathbf{W}_x \mathbf{x}_t + \mathbf{b})\]

RNN 的问题是\textbf{梯度消失/爆炸}------长序列中梯度在时间上反向传播时容易指数衰减，导致长期依赖无法被学习。

\textbf{LSTM}（长短期记忆）通过引入\textbf{遗忘门}、\textbf{输入门}和\textbf{输出门}三个门控机制，让网络学习''记住什么、忘记什么''，有效缓解了梯度消失问题。LSTM 的出现是序列建模历史上的转折点。

\subsection{Transformer 与注意力机制}\label{transformer-ux4e0eux6ce8ux610fux529bux673aux5236}

Transformer（Vaswani et al., 2017）用\textbf{自注意力}（Self-Attention）机制替代了 RNN 的递归结构，允许序列中任意两个位置直接交互：

\[\text{Attention}(Q, K, V) = \text{softmax}\left(\frac{QK^T}{\sqrt{d_k}}\right) V\]

关键优势：
1. \textbf{并行化}：RNN 必须逐个时间步处理，Transformer 可以一次处理整个序列。
2. \textbf{长距离依赖}：任意两个位置之间的交互路径长度是常数 \(O(1)\)（RNN 是 \(O(n)\)）。
3. \textbf{可扩展性}：Transformer 架构支持大规模预训练（如 GPT 系列），这催生了当前的大语言模型（LLM）时代。

\section{深度学习与统计学的联系}\label{ux6df1ux5ea6ux5b66ux4e60ux4e0eux7edfux8ba1ux5b66ux7684ux8054ux7cfb}

深度学习与前面各章讨论的统计方法并非对立，而是互补：

\begin{itemize}
\tightlist
\item
  神经网络的 MLE 训练等价于极大化条件似然。
\item
  \(L_2\) 正则化等价于高斯先验下的 MAP 估计（Ch08）。
\item
  Dropout 等价于变分推断的特例（Ch09 扩展）。
\item
  Batch Normalization 的''训练-测试不一致''在统计学中对应''估计-预测目标不同''。
\item
  Transformer 的自注意力与核回归有深刻的数学联系。
\item
  生成对抗网络（GAN）的判别器本质上是一个两样本检验的统计量。
\end{itemize}

\begin{quote}
理解了统计学的逻辑，深度学习的各种''技巧''就不再是魔术------它们只是在高维非凸优化场景下，为了保证计算可行性和泛化性能而做的工程妥协。
\end{quote}

\begin{Shaded}
\begin{Highlighting}[]
\CommentTok{\# R 中的简单神经网络示例 (需安装 keras/tensorflow)}
\FunctionTok{library}\NormalTok{(keras)}
\NormalTok{model }\OtherTok{\textless{}{-}} \FunctionTok{keras\_model\_sequential}\NormalTok{() }\SpecialCharTok{\%\textgreater{}\%}
  \FunctionTok{layer\_dense}\NormalTok{(}\AttributeTok{units =} \DecValTok{64}\NormalTok{, }\AttributeTok{activation =} \StringTok{"relu"}\NormalTok{, }\AttributeTok{input\_shape =} \FunctionTok{c}\NormalTok{(}\DecValTok{10}\NormalTok{)) }\SpecialCharTok{\%\textgreater{}\%}
  \FunctionTok{layer\_dropout}\NormalTok{(}\AttributeTok{rate =} \FloatTok{0.3}\NormalTok{) }\SpecialCharTok{\%\textgreater{}\%}
  \FunctionTok{layer\_dense}\NormalTok{(}\AttributeTok{units =} \DecValTok{32}\NormalTok{, }\AttributeTok{activation =} \StringTok{"relu"}\NormalTok{) }\SpecialCharTok{\%\textgreater{}\%}
  \FunctionTok{layer\_dense}\NormalTok{(}\AttributeTok{units =} \DecValTok{1}\NormalTok{, }\AttributeTok{activation =} \StringTok{"sigmoid"}\NormalTok{)}
\NormalTok{model }\SpecialCharTok{\%\textgreater{}\%} \FunctionTok{compile}\NormalTok{(}
  \AttributeTok{optimizer =} \StringTok{"adam"}\NormalTok{,}
  \AttributeTok{loss =} \StringTok{"binary\_crossentropy"}\NormalTok{,}
  \AttributeTok{metrics =} \FunctionTok{c}\NormalTok{(}\StringTok{"accuracy"}\NormalTok{)}
\NormalTok{)}
\end{Highlighting}
\end{Shaded}

\begin{center}\rule{0.5\linewidth}{0.5pt}\end{center}

\textbf{速查卡片}：

\begin{itemize}
\tightlist
\item
  \textbf{核心问题}：如何让模型自动学习层次化的特征表示？
\item
  \textbf{核心机制}：层级非线性变换 → 反向传播求梯度 → SGD 更新权重
\item
  \textbf{三大架构}：CNN（图像）→ RNN/LSTM（序列）→ Transformer（通用序列，当前 SOTA）
\item
  \textbf{关键概念}：激活函数（ReLU）、反向传播、Dropout、批归一化、注意力机制
\item
  \textbf{与统计的联系}：损失函数=负对数似然，正则化=先验，Dropout=集成
\item
  \textbf{R命令}：\texttt{keras::keras\_model\_sequential()}, \texttt{torch::nn\_module()} (R torch)
\item
  \textbf{参见}：Ch08（贝叶斯）, Ch15（优化方法）, Ch19（集成学习）
\end{itemize}

\chapter{多元统计分析}\label{ux591aux5143ux7edfux8ba1ux5206ux6790}

\begin{quote}
当数据从''每个个体一个指标''扩展到''每个个体多个指标''时，变量之间的相关性成为了信息的核心来源。多元统计分析的目标正是揭示和利用这种变量之间的结构。
\end{quote}

前面各章主要关注单个响应变量与若干自变量之间的关系（Ch10, Ch11）。多元统计分析则将所有变量放在同等的地位上，同时研究变量之间的相关结构、降维方式和分组模式。它横跨了经典统计学与机器学习之间的桥梁。

\section{多元数据的描述}\label{ux591aux5143ux6570ux636eux7684ux63cfux8ff0}

\subsection{样本均值向量与协方差矩阵}\label{ux6837ux672cux5747ux503cux5411ux91cfux4e0eux534fux65b9ux5deeux77e9ux9635}

设有 \(n\) 个观测，每个观测有 \(p\) 个变量：\(\mathbf{x}_i = (x_{i1}, x_{i2}, \ldots, x_{ip})'\)。

\textbf{样本均值向量}：

\[\bar{\mathbf{x}} = \frac{1}{n} \sum_{i=1}^n \mathbf{x}_i\]

\textbf{样本协方差矩阵}：

\[\mathbf{S} = \frac{1}{n-1} \sum_{i=1}^n (\mathbf{x}_i - \bar{\mathbf{x}})(\mathbf{x}_i - \bar{\mathbf{x}})'\]

\(\mathbf{S}\) 是一个 \(p \times p\) 的对称正定矩阵。对角线元素是各变量的样本方差，非对角线元素是两两变量之间的样本协方差。

\textbf{样本相关矩阵}：

\[\mathbf{R} = \mathbf{D}^{-1/2} \mathbf{S} \mathbf{D}^{-1/2}\]

其中 \(\mathbf{D} = \text{diag}(s_1^2, \ldots, s_p^2)\)。\(\mathbf{R}\) 的 \((i, j)\) 元即为 \(X_i\) 与 \(X_j\) 的样本相关系数。

\subsection{多元正态分布}\label{ux591aux5143ux6b63ux6001ux5206ux5e03}

\(p\) 维随机向量 \(\mathbf{X}\) 服从多元正态分布 \(N_p(\boldsymbol{\mu}, \boldsymbol{\Sigma})\)，如果其密度函数为：

\[f(\mathbf{x}) = \frac{1}{(2\pi)^{p/2} |\boldsymbol{\Sigma}|^{1/2}} \exp\left(-\frac{1}{2} (\mathbf{x} - \boldsymbol{\mu})' \boldsymbol{\Sigma}^{-1} (\mathbf{x} - \boldsymbol{\mu})\right)\]

\textbf{核心性质}：

\begin{enumerate}
\def\labelenumi{\arabic{enumi}.}
\tightlist
\item
  边缘分布：\(\mathbf{X}\) 的任意子向量也服从正态分布。
\item
  条件分布：给定部分变量后，其余变量的条件分布也是正态分布。
\item
  线性变换：\(\mathbf{Y} = \mathbf{A}\mathbf{X} + \mathbf{b} \sim N_q(\mathbf{A}\boldsymbol{\mu} + \mathbf{b}, \mathbf{A}\boldsymbol{\Sigma}\mathbf{A}')\)。
\item
  独立性与不相关性等价：在多元正态下，\(\text{Cov}(X_i, X_j) = 0 \Leftrightarrow X_i\) 与 \(X_j\) 独立。
\end{enumerate}

\begin{quote}
性质 4 对多元正态成立，但对一般分布不成立------这是 Ch02 中反复强调的''不相关 ≠ 独立''在多元正态中的例外。
\end{quote}

\section{主成分分析（PCA）}\label{ux4e3bux6210ux5206ux5206ux6790pca}

\subsection{基本思想}\label{ux57faux672cux601dux60f3}

主成分分析的目标是找到变量的线性组合，使得这些组合：
1. 方差尽可能大（保留尽可能多的信息）。
2. 彼此之间不相关（去冗余）。

从几何上看，PCA 等价于对数据做\textbf{坐标旋转}，使得新坐标轴的方向与数据最大变异方向对齐。

\subsection{数学推导}\label{ux6570ux5b66ux63a8ux5bfc}

问题：寻找方向 \(\mathbf{a}\)（\(\|\mathbf{a}\| = 1\)），使得投影 \(\mathbf{a}'\mathbf{X}\) 的方差最大化：

\[\max_{\|\mathbf{a}\|=1} \text{Var}(\mathbf{a}'\mathbf{X}) = \max_{\|\mathbf{a}\|=1} \mathbf{a}' \boldsymbol{\Sigma} \mathbf{a}\]

解：\(\mathbf{a}\) 是 \(\boldsymbol{\Sigma}\) 的最大特征值对应的特征向量。这就是\textbf{第一主成分}。

第二主成分在与第一主成分正交的约束下最大化方差，对应第二大的特征向量。以此类推。

\subsection{方差解释比例}\label{ux65b9ux5deeux89e3ux91caux6bd4ux4f8b}

第 \(k\) 个主成分的方差解释比例为：

\[\lambda_k \Big/ \sum_{j=1}^p \lambda_j\]

其中 \(\lambda_1 \geq \lambda_2 \geq \cdots \geq \lambda_p \geq 0\) 是 \(\boldsymbol{\Sigma}\)（或 \(\mathbf{S}\)）的特征值。

\begin{quote}
实践中，通常选取前几个主成分，使得累积方差解释比例达到 80\%-90\%。这既实现了降维，又保留了大部分信息。
\end{quote}

\begin{Shaded}
\begin{Highlighting}[]
\CommentTok{\# PCA 示例：iris 数据集}
\FunctionTok{data}\NormalTok{(iris)}
\NormalTok{pca }\OtherTok{\textless{}{-}} \FunctionTok{prcomp}\NormalTok{(iris[, }\DecValTok{1}\SpecialCharTok{:}\DecValTok{4}\NormalTok{], }\AttributeTok{scale. =} \ConstantTok{TRUE}\NormalTok{)}
\FunctionTok{summary}\NormalTok{(pca)          }\CommentTok{\# 方差解释}
\end{Highlighting}
\end{Shaded}

\begin{verbatim}
## Importance of components:
##                           PC1    PC2     PC3     PC4
## Standard deviation     1.7084 0.9560 0.38309 0.14393
## Proportion of Variance 0.7296 0.2285 0.03669 0.00518
## Cumulative Proportion  0.7296 0.9581 0.99482 1.00000
\end{verbatim}

\begin{Shaded}
\begin{Highlighting}[]
\NormalTok{pca}\SpecialCharTok{$}\NormalTok{rotation           }\CommentTok{\# 载荷（特征向量）}
\end{Highlighting}
\end{Shaded}

\begin{verbatim}
##                     PC1         PC2        PC3        PC4
## Sepal.Length  0.5210659 -0.37741762  0.7195664  0.2612863
## Sepal.Width  -0.2693474 -0.92329566 -0.2443818 -0.1235096
## Petal.Length  0.5804131 -0.02449161 -0.1421264 -0.8014492
## Petal.Width   0.5648565 -0.06694199 -0.6342727  0.5235971
\end{verbatim}

\begin{Shaded}
\begin{Highlighting}[]
\FunctionTok{biplot}\NormalTok{(pca)            }\CommentTok{\# 双标图}
\end{Highlighting}
\end{Shaded}

\begin{center}\includegraphics{statistic-book_files/figure-latex/unnamed-chunk-30-1} \end{center}

\subsection{PCA 与 SVD 的联系}\label{pca-ux4e0e-svd-ux7684ux8054ux7cfb}

对中心化的数据矩阵 \(\tilde{\mathbf{X}}\)（\(n \times p\)），其样本协方差矩阵为 \(\mathbf{S} = \frac{1}{n-1} \tilde{\mathbf{X}}'\tilde{\mathbf{X}}\)。

若对 \(\tilde{\mathbf{X}}\) 做 SVD：\(\tilde{\mathbf{X}} = \mathbf{U} \boldsymbol{\Sigma} \mathbf{V}'\)，则：

\begin{itemize}
\tightlist
\item
  \(\mathbf{V}\) 的列就是主成分方向（\(\mathbf{S}\) 的特征向量）。
\item
  \(\sigma_i^2 / (n-1)\) 就是第 \(i\) 个特征值。
\end{itemize}

这一联系（Ch16 中也有讨论）意味着 PCA 的计算本质上就是一个 SVD 分解。

\section{因子分析}\label{ux56e0ux5b50ux5206ux6790}

因子分析（Factor Analysis）与 PCA 不同：PCA 是纯变量变换，而因子分析假设存在\textbf{隐变量}驱动观测变量之间的相关性。

\[X_i = \lambda_{i1} F_1 + \lambda_{i2} F_2 + \cdots + \lambda_{im} F_m + \varepsilon_i, \quad m < p\]

其中：
- \(F_1, \ldots, F_m\) 是公共因子（不可观测）。
- \(\lambda_{ij}\) 是\textbf{因子载荷}（Factor Loading）。
- \(\varepsilon_i\) 是特殊因子（变量 \(X_i\) 独有的部分）。

\begin{quote}
PCA 和因子分析的区别：PCA 想的是''如何用少数组合作替原始变量''，因子分析想的是''有没有少量隐变量可以解释观测变量间的相关性''。
\end{quote}

\section{聚类分析}\label{ux805aux7c7bux5206ux6790}

聚类分析（Cluster Analysis）将相似的观测归入同一组，使得组内尽可能同质、组间尽可能异质。

\subsection{K-means 聚类}\label{k-means-ux805aux7c7b}

\textbf{算法步骤}：

\begin{enumerate}
\def\labelenumi{\arabic{enumi}.}
\tightlist
\item
  随机初始化 \(K\) 个聚类中心 \(\mathbf{c}_1, \ldots, \mathbf{c}_K\)。
\item
  将每个样本分配到最近的中心。
\item
  重新计算每个聚类的中心（均值）。
\item
  重复步骤 2-3 直至收敛。
\end{enumerate}

K-means 最小化的是组内平方和：\(W(C) = \sum_{k=1}^K \sum_{i \in C_k} \|\mathbf{x}_i - \mathbf{c}_k\|^2\)。

\textbf{注意事项}：K-means 需要预先指定 \(K\)；对初始值敏感；只能发现球形聚类。

\subsection{层次聚类}\label{ux5c42ux6b21ux805aux7c7b}

层次聚类不需要预先指定聚类数，而是构建一个树状结构（系统树图）：

\begin{itemize}
\tightlist
\item
  \textbf{凝聚法}（Agglomerative）：从每个样本各自为一类开始，逐步合并最近的类。
\item
  \textbf{分裂法}（Divisive）：从所有样本为一类开始，逐步分裂。
\end{itemize}

类间距离的定义方式（单连接、全连接、平均连接、Ward 方法）会影响聚类结果。

\begin{Shaded}
\begin{Highlighting}[]
\CommentTok{\# K{-}means 与层次聚类示例}
\FunctionTok{data}\NormalTok{(iris)}
\NormalTok{scaled\_iris }\OtherTok{\textless{}{-}} \FunctionTok{scale}\NormalTok{(iris[, }\DecValTok{1}\SpecialCharTok{:}\DecValTok{4}\NormalTok{])}

\CommentTok{\# K{-}means}
\NormalTok{km }\OtherTok{\textless{}{-}} \FunctionTok{kmeans}\NormalTok{(scaled\_iris, }\AttributeTok{centers =} \DecValTok{3}\NormalTok{)}
\FunctionTok{table}\NormalTok{(km}\SpecialCharTok{$}\NormalTok{cluster, iris}\SpecialCharTok{$}\NormalTok{Species)}
\end{Highlighting}
\end{Shaded}

\begin{verbatim}
##    
##     setosa versicolor virginica
##   1     50          0         0
##   2      0         39        14
##   3      0         11        36
\end{verbatim}

\begin{Shaded}
\begin{Highlighting}[]
\CommentTok{\# 层次聚类}
\NormalTok{hc }\OtherTok{\textless{}{-}} \FunctionTok{hclust}\NormalTok{(}\FunctionTok{dist}\NormalTok{(scaled\_iris), }\AttributeTok{method =} \StringTok{"ward.D2"}\NormalTok{)}
\FunctionTok{plot}\NormalTok{(hc, }\AttributeTok{labels =} \ConstantTok{FALSE}\NormalTok{)}
\end{Highlighting}
\end{Shaded}

\begin{center}\includegraphics{statistic-book_files/figure-latex/unnamed-chunk-31-1} \end{center}

\section{判别分析}\label{ux5224ux522bux5206ux6790}

与聚类分析不同，判别分析是有监督的分类方法------已知每个样本的类别标签，目标是找到能将不同类别最好地分开的方向。

\subsection{线性判别分析（LDA）}\label{ux7ebfux6027ux5224ux522bux5206ux6790lda}

LDA 假设各类别服从多元正态分布，且具有相同的协方差矩阵（但均值不同）。分类规则为：

将 \(\mathbf{x}\) 分配到使后验概率 \(P(Y = k \mid \mathbf{x})\) 最大的类别。

LDA 找到的判别方向恰好是 \(\boldsymbol{\Sigma}^{-1} (\boldsymbol{\mu}_1 - \boldsymbol{\mu}_2)\)（二分类情况），也就是\textbf{最大化类间方差与类内方差之比}的方向。

\begin{quote}
LDA 既可以用于分类，也可以用于降维（找到最具判别力的投影方向）。
\end{quote}

\section{与本书其他章节的联系}\label{ux4e0eux672cux4e66ux5176ux4ed6ux7ae0ux8282ux7684ux8054ux7cfb-1}

\begin{itemize}
\tightlist
\item
  PCA 的计算依赖 Ch16 中的 SVD 分解。
\item
  判别分析与 Ch11（GLM 的二分类）和 Ch18（分类方法）密切相关。
\item
  聚类与 Ch17 中的无监督学习主题重叠。
\item
  协方差矩阵的估计与检验依赖于 Ch07 中的多元检验理论。
\end{itemize}

\begin{center}\rule{0.5\linewidth}{0.5pt}\end{center}

\textbf{速查卡片}：

\begin{itemize}
\tightlist
\item
  \textbf{核心问题}：如何在多变量数据的复杂结构中提取简化的、可解释的模式？
\item
  \textbf{三大任务}：降维（PCA/因子分析）、分组（聚类/判别分析）、相关结构建模
\item
  \textbf{关键方法}：PCA = SVD + 方差最大化；LDA = 最大化类间/类内方差比
\item
  \textbf{R命令}：\texttt{prcomp()}, \texttt{factanal()}, \texttt{kmeans()}, \texttt{hclust()}, \texttt{lda()} (MASS)
\item
  \textbf{参见}：Ch02（多元正态分布）, Ch16（SVD）, Ch17（无监督学习）
\end{itemize}

\chapter{因果推断入门}\label{ux56e0ux679cux63a8ux65adux5165ux95e8}

\begin{quote}
``相关不等于因果''------这是统计学入门第一课就会听到的告诫。但如果相关不等于因果，那我们如何能从数据中推断因果关系？因果推断（Causal Inference）正是为了回答这个问题而发展起来的学科。
\end{quote}

\section{为什么因果推断重要}\label{ux4e3aux4ec0ux4e48ux56e0ux679cux63a8ux65adux91cdux8981}

回顾一下我们学过的内容：Ch06 的假设检验、Ch10 的回归分析、Ch17 的预测建模------它们回答的都是\textbf{关联}问题（``\(X\) 和 \(Y\) 有没有关系？''），而不是\textbf{因果}问题（``改变 \(X\) 会导致 \(Y\) 改变吗？''）。

关联与因果的差距，在现实中意味着巨大的决策差异：

\begin{itemize}
\tightlist
\item
  冰淇淋销量与溺水死亡人数正相关 → 禁止冰淇淋能减少溺水吗？不能。
\item
  受过大学教育的人收入更高 → 扩张大学招生一定会提高人群收入吗？不一定。
\item
  服用某种药物的人康复率更高 → 给所有人用这个药一定有好处吗？不一定------服用药物的人可能本身就更健康。
\end{itemize}

回答因果问题的关键是：\textbf{如果我介入（干预）系统，改变 \(X\)，\(Y\) 会怎么变？}这需要一套不同于传统回归分析的思维框架和数学工具。

\section{潜在结果框架}\label{ux6f5cux5728ux7ed3ux679cux6846ux67b6}

\subsection{Rubin 因果模型（RCM）}\label{rubin-ux56e0ux679cux6a21ux578brcm}

潜在结果框架（Rubin Causal Model）是因果推断最主流的框架。其核心思想是：对于每一个个体，我们定义其在接受处理（\(T=1\)）和不接受处理（\(T=0\)）两种\textbf{反事实}情况下的结果。

对于个体 \(i\)：

\begin{itemize}
\tightlist
\item
  \(Y_i(1)\)：如果个体 \(i\) 接受处理，其潜在结果。
\item
  \(Y_i(0)\)：如果个体 \(i\) 不接受处理，其潜在结果。
\item
  \(T_i \in \{0, 1\}\)：个体 \(i\) 实际接受的处理。
\end{itemize}

个体 \(i\) 的\textbf{因果效应}定义为：

\[\tau_i = Y_i(1) - Y_i(0)\]

\subsection{因果推断的根本问题}\label{ux56e0ux679cux63a8ux65adux7684ux6839ux672cux95eeux9898}

对于每个个体，我们\textbf{只能观测到一个潜在结果}：

\[Y_i^{\text{obs}} = T_i \cdot Y_i(1) + (1 - T_i) \cdot Y_i(0)\]

另一个潜在结果是\textbf{反事实}的------我们永远无法直接观测到同一个人既服药又不服药的结果。这就是因果推断的根本问题（Fundamental Problem of Causal Inference）。

\begin{quote}
既然无法观测个体因果效应，因果推断的目标转向估计\textbf{平均因果效应}（Average Treatment Effect, ATE）：\(\tau = E[Y(1) - Y(0)]\)。
\end{quote}

\subsection{随机化试验}\label{ux968fux673aux5316ux8bd5ux9a8c}

如果处理分配 \(T\) 是\textbf{随机的}（与潜在结果独立），那么：

\[E[Y \mid T=1] - E[Y \mid T=0] = E[Y(1) \mid T=1] - E[Y(0) \mid T=0] = E[Y(1) - Y(0)] = \tau\]

随机化使得处理组和对照组在各方面（可观测和不可观测）都具有可比性，因此简单的均值差就是 ATE 的无偏估计。这就是随机对照试验（RCT）是因果推断''金标准''的原因。

但在大多数现实场景中（经济学、流行病学、社会科学），随机化不可行或不道德------这就是\textbf{观察性研究因果推断}所要解决的问题。

\section{有向无环图（DAG）}\label{ux6709ux5411ux65e0ux73afux56fedag}

\subsection{DAG 的基本语法}\label{dag-ux7684ux57faux672cux8bedux6cd5}

有向无环图（Directed Acyclic Graph, DAG）是用图形语言描述变量间因果关系的工具。

\begin{itemize}
\tightlist
\item
  \textbf{节点}代表变量。
\item
  \textbf{箭头} \(X \rightarrow Y\) 表示 \(X\) 是 \(Y\) 的直接原因。
\item
  \textbf{无环}意味着因果链条不能形成闭环。
\end{itemize}

\subsection{DAG 中的三种基本结构}\label{dag-ux4e2dux7684ux4e09ux79cdux57faux672cux7ed3ux6784}

\begin{enumerate}
\def\labelenumi{\arabic{enumi}.}
\tightlist
\item
  \textbf{链式结构}：\(X \rightarrow Z \rightarrow Y\)

  \begin{itemize}
  \tightlist
  \item
    \(X\) 通过中介变量 \(Z\) 影响 \(Y\)。控制 \(Z\) 会\textbf{阻断}因果路径。
  \end{itemize}
\item
  \textbf{叉式结构（混杂）}：\(X \leftarrow Z \rightarrow Y\)

  \begin{itemize}
  \tightlist
  \item
    \(Z\) 是 \(X\) 和 \(Y\) 的共同原因（混杂因子）。若不控制 \(Z\)，\(X\) 和 \(Y\) 之间会出现虚假的相关。\textbf{控制 \(Z\) 可以消除混杂}。
  \end{itemize}
\item
  \textbf{对撞结构}：\(X \rightarrow Z \leftarrow Y\)

  \begin{itemize}
  \tightlist
  \item
    \(Z\) 是 \(X\) 和 \(Y\) 的共同结果（对撞变量）。\(X\) 和 \(Y\) 本身独立，但\textbf{控制 \(Z\) 会打开一条虚假的关联路径}。这正是''选择偏差''的来源。
  \end{itemize}
\end{enumerate}

\begin{quote}
三大结构的区分是画出 DAG、识别混杂因子和选择控制变量的基础。
\end{quote}

\subsection{do-算子与干预}\label{do-ux7b97ux5b50ux4e0eux5e72ux9884}

Pearl 引入了 \textbf{do-算子} 来形式化干预的概念：

\begin{itemize}
\tightlist
\item
  \(P(Y \mid X = x)\)：在\textbf{观测到} \(X = x\) 的条件下，\(Y\) 的分布（关联）。
\item
  \(P(Y \mid do(X = x))\)：在\textbf{强制设定} \(X = x\)（切断所有指向 \(X\) 的箭头）后，\(Y\) 的分布（因果）。
\end{itemize}

两者的区别正是因果推断的核心。在 DAG 的框架下，\(P(Y \mid do(X = x))\) 可以通过\textbf{后门准则}（Backdoor Criterion）从观测数据中识别。

\section{观察性研究中的因果推断方法}\label{ux89c2ux5bdfux6027ux7814ux7a76ux4e2dux7684ux56e0ux679cux63a8ux65adux65b9ux6cd5}

\subsection{倾向性得分（Propensity Score）}\label{ux503eux5411ux6027ux5f97ux5206propensity-score}

倾向性得分是个体根据其协变量被分配到处理组的概率：

\[e(\mathbf{X}) = P(T = 1 \mid \mathbf{X})\]

Rosenbaum 和 Rubin 证明：\textbf{在给定倾向性得分的条件下，处理分配与协变量独立}。这意味着我们可以用一维的''得分''替代多维的协变量来控制混杂。

常用的基于倾向性得分的方法：

\begin{enumerate}
\def\labelenumi{\arabic{enumi}.}
\tightlist
\item
  \textbf{倾向性得分匹配（PSM）}：为每个处理组个体匹配一个倾向性得分相似的对照组个体。
\item
  \textbf{逆概率加权（IPW）}：用倾向性得分的倒数对观测值加权，构造一个''伪随机化''样本。
\item
  \textbf{分层}：按倾向性得分分层，在每层内估计处理效应，然后加权合并。
\end{enumerate}

\begin{Shaded}
\begin{Highlighting}[]
\CommentTok{\# 倾向性得分匹配示例}
\FunctionTok{library}\NormalTok{(MatchIt)}
\FunctionTok{data}\NormalTok{(lalonde, }\AttributeTok{package =} \StringTok{"MatchIt"}\NormalTok{)}
\CommentTok{\# 用逻辑回归估计倾向性得分}
\NormalTok{m.out }\OtherTok{\textless{}{-}} \FunctionTok{matchit}\NormalTok{(treat }\SpecialCharTok{\textasciitilde{}}\NormalTok{ age }\SpecialCharTok{+}\NormalTok{ educ }\SpecialCharTok{+}\NormalTok{ race }\SpecialCharTok{+}\NormalTok{ married }\SpecialCharTok{+}\NormalTok{ re74 }\SpecialCharTok{+}\NormalTok{ re75,}
                 \AttributeTok{data =}\NormalTok{ lalonde, }\AttributeTok{method =} \StringTok{"nearest"}\NormalTok{, }\AttributeTok{ratio =} \DecValTok{1}\NormalTok{)}
\FunctionTok{summary}\NormalTok{(m.out)}
\end{Highlighting}
\end{Shaded}

\begin{verbatim}
## 
## Call:
## matchit(formula = treat ~ age + educ + race + married + re74 + 
##     re75, data = lalonde, method = "nearest", ratio = 1)
## 
## Summary of Balance for All Data:
##            Means Treated Means Control Std. Mean Diff. Var. Ratio eCDF Mean
## distance          0.5722        0.1845          1.8020     0.8651    0.3763
## age              25.8162       28.0303         -0.3094     0.4400    0.0813
## educ             10.3459       10.2354          0.0550     0.4959    0.0347
## raceblack         0.8432        0.2028          1.7615          .    0.6404
## racehispan        0.0595        0.1422         -0.3498          .    0.0827
## racewhite         0.0973        0.6550         -1.8819          .    0.5577
## married           0.1892        0.5128         -0.8263          .    0.3236
## re74           2095.5737     5619.2365         -0.7211     0.5181    0.2248
## re75           1532.0553     2466.4844         -0.2903     0.9563    0.1342
##            eCDF Max
## distance     0.6428
## age          0.1577
## educ         0.1114
## raceblack    0.6404
## racehispan   0.0827
## racewhite    0.5577
## married      0.3236
## re74         0.4470
## re75         0.2876
## 
## Summary of Balance for Matched Data:
##            Means Treated Means Control Std. Mean Diff. Var. Ratio eCDF Mean
## distance          0.5722        0.3664          0.9563     0.7101    0.1280
## age              25.8162       25.5135          0.0423     0.4700    0.0803
## educ             10.3459       10.6757         -0.1640     0.5207    0.0350
## raceblack         0.8432        0.4703          1.0259          .    0.3730
## racehispan        0.0595        0.2216         -0.6857          .    0.1622
## racewhite         0.0973        0.3081         -0.7113          .    0.2108
## married           0.1892        0.2000         -0.0276          .    0.0108
## re74           2095.5737     2353.7678         -0.0528     1.3642    0.0520
## re75           1532.0553     1696.8096         -0.0512     1.4223    0.0516
##            eCDF Max Std. Pair Dist.
## distance     0.4162          0.9564
## age          0.2432          1.4112
## educ         0.1405          1.1910
## raceblack    0.3730          1.0259
## racehispan   0.1622          1.1429
## racewhite    0.2108          0.7478
## married      0.0108          0.8005
## re74         0.3027          0.7552
## re75         0.1946          0.7298
## 
## Sample Sizes:
##           Control Treated
## All           429     185
## Matched       185     185
## Unmatched     244       0
## Discarded       0       0
\end{verbatim}

\begin{Shaded}
\begin{Highlighting}[]
\CommentTok{\# 匹配后的因果效应估计}
\NormalTok{m.data }\OtherTok{\textless{}{-}} \FunctionTok{match.data}\NormalTok{(m.out)}
\FunctionTok{t.test}\NormalTok{(re78 }\SpecialCharTok{\textasciitilde{}}\NormalTok{ treat, }\AttributeTok{data =}\NormalTok{ m.data)}
\end{Highlighting}
\end{Shaded}

\begin{verbatim}
## 
##  Welch Two Sample t-test
## 
## data:  re78 by treat
## t = -1.1787, df = 344.45, p-value = 0.2393
## alternative hypothesis: true difference in means between group 0 and group 1 is not equal to 0
## 95 percent confidence interval:
##  -2290.9583   574.0328
## sample estimates:
## mean in group 0 mean in group 1 
##        5490.681        6349.144
\end{verbatim}

\subsection{工具变量（Instrumental Variable）}\label{ux5de5ux5177ux53d8ux91cfinstrumental-variable}

当存在未观测的混杂因子时，倾向性得分方法失效。这时可以考虑使用\textbf{工具变量}（IV）：

一个有效的工具变量 \(Z\) 必须满足三个条件：

\begin{enumerate}
\def\labelenumi{\arabic{enumi}.}
\tightlist
\item
  \textbf{相关性}：\(Z\) 与处理 \(T\) 相关（影响处理概率）。
\item
  \textbf{外生性}：\(Z\) 与未观测的混杂因子 \(U\) 独立。
\item
  \textbf{排他性}：\(Z\) 仅通过 \(T\) 影响结果 \(Y\)（无直接效应）。
\end{enumerate}

\begin{quote}
经典例子：Angrist 和 Krueger 用''出生季度''作为''受教育年限''的工具变量，估计教育对收入的因果效应。前提是：出生季度影响受教育年限（因入学截止日），但与能力等未观测因素无关，且只通过教育影响收入。
\end{quote}

两阶段最小二乘（2SLS）是 IV 估计的标准方法：

\begin{enumerate}
\def\labelenumi{\arabic{enumi}.}
\tightlist
\item
  第一阶段：\(T = \alpha_0 + \alpha_1 Z + \mathbf{X}\boldsymbol{\gamma} + \varepsilon\)
\item
  第二阶段：\(Y = \beta_0 + \beta_1 \hat{T} + \mathbf{X}\boldsymbol{\delta} + \eta\)
\end{enumerate}

\(\beta_1\) 即是 ATE 的一致估计。

\subsection{双重差分（Difference-in-Differences, DiD）}\label{ux53ccux91cdux5deeux5206difference-in-differences-did}

当无法随机化但有面板数据时，DiD 利用处理组和对照组在时间上的变化来识别因果效应：

\[\hat{\tau}_{\text{DiD}} = (\bar{Y}_{\text{treated, post}} - \bar{Y}_{\text{treated, pre}}) - (\bar{Y}_{\text{control, post}} - \bar{Y}_{\text{control, pre}})\]

DiD 的核心假设是\textbf{平行趋势}（Parallel Trends）：如果没有处理，处理组和对照组的时间趋势应该是相同的。

\subsection{断点回归（Regression Discontinuity, RD）}\label{ux65adux70b9ux56deux5f52regression-discontinuity-rd}

当处理分配由一个阈值（Cutoff）决定时，RD 利用阈值附近的''准随机化''来估计因果效应。例如：考试成绩高于分数线才能入学------比较分数线上下微小邻域内的学生，可以近似认为他们是随机分配的。

\section{因果推断与统计建模的关系}\label{ux56e0ux679cux63a8ux65adux4e0eux7edfux8ba1ux5efaux6a21ux7684ux5173ux7cfb}

因果推断不是对传统统计学的否定，而是它的扩展：

\begin{longtable}[]{@{}lll@{}}
\toprule\noalign{}
问题类型 & 统计工具 & 因果工具 \\
\midrule\noalign{}
\endhead
\bottomrule\noalign{}
\endlastfoot
变量间是否有关联？ & 相关性、回归系数 & --- \\
改变 \(X\) 会使 \(Y\) 变多少？ & --- & ATE、ATT \\
效应通过什么路径传导？ & --- & 中介分析 \\
效应对谁起效、谁不起效？ & --- & 异质性处理效应 \\
\end{longtable}

\begin{quote}
一本统计学教材讨论因果推断，是想提醒读者：统计模型本身不附带因果解释。因果性来自于研究设计、识别策略和领域知识------统计只是检验的工具。
\end{quote}

\section{与本书其他章节的联系}\label{ux4e0eux672cux4e66ux5176ux4ed6ux7ae0ux8282ux7684ux8054ux7cfb-2}

\begin{itemize}
\tightlist
\item
  倾向性得分的估计使用 Ch11 中的逻辑回归。
\item
  工具变量法使用 Ch10 中的两阶段最小二乘。
\item
  因果效应在随机化试验中的估计与 Ch07 中的两样本 t 检验等价。
\item
  贝叶斯因果推断结合了 Ch08 的先验设定和 Ch09 的 MCMC 计算。
\item
  DAG 与 Ch20（深度学习）中的贝叶斯网络在数学上同源。
\end{itemize}

\begin{center}\rule{0.5\linewidth}{0.5pt}\end{center}

\textbf{速查卡片}：

\begin{itemize}
\tightlist
\item
  \textbf{核心问题}：如何从数据中推断''改变 \(X\) 会导致 \(Y\) 改变多少''？
\item
  \textbf{核心框架}：潜在结果（Rubin） + 因果图（Pearl）
\item
  \textbf{关键概念}：反事实、混杂、中介、对撞、do-算子
\item
  \textbf{观察性研究方法}：匹配(PSM) → 加权(IPW) → 工具变量(IV) → 双重差分(DiD) → 断点回归(RD)
\item
  \textbf{R命令}：\texttt{MatchIt::matchit()}, \texttt{AER::ivreg()}, \texttt{rdrobust::rdrobust()}
\item
  \textbf{参见}：Ch07（假设检验）, Ch10（回归分析）, Ch11（GLM）
\end{itemize}

\end{document}
